% !TEX program = xelatex
\documentclass[12pt,a4paper,fontset=fandol]{ctexbook}

% ==================== 页面 ====================
\usepackage[top=2.5cm,bottom=2.5cm,left=2.5cm,right=2.5cm]{geometry}

% ==================== 字体 ====================
\usepackage{fontspec}
\usepackage{iffont}
% CJK 钉 fandol（两侧都取自 TeX 树里的同一份 FandolSong-Regular.otf）；
% mono 钉 DejaVu Sans Mono：CI 由 fontconfig 命中同名，本机没有该系统字体，
% 就用 MiKTeX dejavu 包里的同一版本 TTF（kpathsea 按文件名命中）。
\IfFontExistsTF{DejaVu Sans Mono}
  {\setmonofont{DejaVu Sans Mono}[Scale=MatchLowercase]}
  {\setmonofont[Extension=.ttf,
      UprightFont=DejaVuSansMono,
      BoldFont=DejaVuSansMono-Bold,
      ItalicFont=DejaVuSansMono-Oblique,
      BoldItalicFont=DejaVuSansMono-BoldOblique]
     {DejaVuSansMono}[Scale=MatchLowercase]}

% ==================== 颜色 ====================
\usepackage[dvipsnames,svgnames,x11names]{xcolor}
\definecolor{codebg}{HTML}{F5F5F5}
\definecolor{codeframe}{HTML}{E0E0E0}
\definecolor{cppkeyword}{HTML}{1A1AFF}
\definecolor{cppcomment}{HTML}{2E8B2E}
\definecolor{cppstring}{HTML}{B22222}
\definecolor{linenumgray}{HTML}{999999}
\definecolor{toolheadbg}{HTML}{2D3E50}
\definecolor{yamlkey}{HTML}{7A3E9D}
\definecolor{diffadd}{HTML}{1A7F37}
\definecolor{diffdel}{HTML}{B62324}

% ==================== listings ====================
\usepackage{listings}

% ---- 所有代码块共用的排版（先定义，各语言 style 组合它）----
\lstdefinestyle{codebase}{
  basicstyle=\small\ttfamily,
  keywordstyle=\color{cppkeyword}\bfseries,
  commentstyle=\color{cppcomment}\itshape,
  stringstyle=\color{cppstring},
  numberstyle=\tiny\color{linenumgray},
  numbers=left,
  frame=single,
  rulecolor=\color{codeframe},
  framesep=6pt,
  breaklines=true,
  tabsize=4,
  columns=flexible,
  keepspaces=true,
  backgroundcolor=\color{codebg},
  showstringspaces=false,
  escapeinside={(*@}{@*)},
}

\lstdefinestyle{cppstyle}{
  style=codebase,
  language=C++,
  morekeywords={constexpr,consteval,constinit,decltype,auto,
    concept,requires,co_await,co_yield,co_return,
    noexcept,override,final,nullptr,static_assert,
    template,typename,namespace,using,mutable,volatile,
    explicit,operator,friend,inline,virtual,
    thread_local,alignas,alignof},
  deletekeywords={int,char,bool,float,double,long,short,unsigned,signed,void,wchar_t},
  emph={size_t,int32_t,uint32_t,int64_t,uint64_t,ssize_t,
    string,vector,map,set,unique_ptr,shared_ptr,optional,variant,any,
    span,string_view,FILE,wchar_t},
  emphstyle=\color{teal!80!black},
}
\lstset{style=cppstyle}

% ---- C：编解码器、设备与滤镜的 C ABI 边界 ----
\lstdefinestyle{cstyle}{
  style=codebase,
  language=C,
  morekeywords={bool,FILE,size_t,restrict,inline},
  emph={size_t,int32_t,uint32_t,int64_t,uint64_t,ssize_t,
    AVFrame,AVPacket,AVCodec,AVCodecContext,AVDictionary,AVFormatContext,
    SwsContext,SwrContext,OpusEncoder,x264_t,x265_analysis_data,v4l2_format},
  emphstyle=\color{teal!80!black},
}

% ---- YAML：CI 描述文件、转码与推流流水线配置 ----
\lstdefinestyle{yamlstyle}{
  style=codebase,
  morecomment=[l]{\#},
  morestring=[b]",
  morestring=[d]',
  alsoletter={-,:},
  morekeywords=[1]{Version,Checks,CheckOptions,HeaderFilterRegex,FormatStyle,
    User,WarningsAsErrors,CompilationDatabase,ExtraArgs,Add,Remove,
    Diagnostics,ClangTidy,UnusedIncludes,MissingIncludes,index,background-index,
    clang-tidy,InlayHints,Style,BasedOnStyle,IndentWidth,TabWidth,UseTab,
    ColumnLimit,BreakBeforeBraces,AlignAfterOpenBracket,AllowShortIfStatementsOnASingleLine,
    IncludeBlocks,SortIncludes,SeparateDefinitionBlocks,SpacesInAngles,
    Language,If,Else,Disable,Key,Name,files,exclude,repo,rev,hooks,id,entry,
    language_types,on,push_pull_request,steps,uses,runs-on,defaults},
  keywordstyle=[1]\color{yamlkey}\bfseries,
}

% ---- JSON：ffprobe 的机读输出、预设与工具配置 ----
\lstdefinestyle{jsonstyle}{
  style=codebase,
  morestring=[b]",
  alsoletter={-},
  morekeywords={directory,file,output,arguments,command,version,
    cmakeMinimumRequired,name,generator,binaryDir,sourceDir,caches,
    toolchainFile,C_CXX_FLAGS,defineConstants,environments,
    editorMode,rulerColumn,tabSize,insertSpaces,files,settings,
    codec_name,codec_type,width,height,pix_fmt,r_frame_rate,avg_frame_rate,
    duration,bit_rate,sample_rate,sample_fmt,channel_layout,nb_frames,pkt_size},
  keywordstyle=\color{yamlkey}\bfseries,
}

% ---- Shell：命令行与工具调用 ----
\lstdefinelanguage{ShellLD}{
  morekeywords={if,then,else,elif,fi,for,in,do,done,while,case,esac,
    function,return,local,export,source,echo,cd,set,exec,trap,read,
    printf,test,exit,shift,eval,alias},
  sensitive=true,
  morecomment=[l]{\#},
  morestring=[b]",
  morestring=[d]',
}
\lstdefinestyle{shellstyle}{
  style=codebase,
  language=ShellLD,
  keywordstyle=\color{cppkeyword}\bfseries,
  emph={ffmpeg,ffprobe,ffplay,x264,x265,aomenc,SvtAv1EncApp,opusenc,opusdec,
    mediainfo,mp4box,packager,srt-live-transmit,werepeater,
    v4l2-ctl,pactl,arecord,aplay,alsauff,amixer,jackd,qjackctl,
    gst-launch-1.0,gst-inspect-1.0,gstreamer,nvidia-smi,vaapi-info,dxva2,
    grep,sed,awk,xargs,jq,python3,pip,pip3,pwsh,bash,git,
    cmake,ninja,make,ccache,lldb,gdb,code,nvim,tar,mkdir,touch,ls,cat,
    pkg-config,ldconfig,strace,ltrace,perf,valgrind,gh,yamllint,wc,diff,find,sort,uniq,tee,rm},
  emphstyle=\color{teal!80!black}\bfseries,
}

% ---- CMake ----
\lstdefinelanguage{CMakeLD}{
  morekeywords={cmake_minimum_required,project,add_executable,
    add_library,target_link_libraries,target_include_directories,
    target_compile_definitions,target_compile_options,set,if,elseif,
    else,endif,find_package,option,message,install,file,
    target_sources,add_subdirectory,enable_language,set_target_properties,
    FetchContent_Declare,FetchContent_MakeAvailable,list,string},
  sensitive=false,
  morecomment=[l]{\#},
  morestring=[b]",
}
\lstdefinestyle{cmakestyle}{
  style=codebase,
  language=CMakeLD,
  keywordstyle=\color{cppkeyword}\bfseries,
}

% ---- Python：PyAV/ffmpeg 工具脚本与批量转码流水线 ----
\lstdefinestyle{pythonstyle}{
  style=codebase,
  language=Python,
  morekeywords={self,lambda,async,await,yield,with,as,pass,del,
    global,nonlocal,is,in,not,and,or,True,False,None},
  keywordstyle=\color{cppkeyword}\bfseries,
  emph={av,subprocess,np,numpy,plt,json,os,sys,pathlib,struct,numpy},
  emphstyle=\color{teal!80!black},
}

% ---- 统一 diff：补丁式改动对比 ----
\lstdefinestyle{diffstyle}{
  style=codebase,
  numbers=none,
  frame=none,
  backgroundcolor=\color{white},
  morecomment=[l][\color{diffdel}]{-},
  morecomment=[l][\color{diffadd}]{+},
  morecomment=[l][\color{blue!70!black}\bfseries]{@@},
}

% ==================== tcolorbox ====================
\usepackage[most]{tcolorbox}

\newtcolorbox{guideline}[1][]{
  enhanced,breakable,
  colback=blue!5!white,colframe=blue!75!black,
  fonttitle=\bfseries,title=要点,#1,
  boxed title style={colback=blue!75!black},
  attach boxed title to top left={yshift=-2mm,xshift=2mm},
  arc=2mm,boxrule=0.8mm,
}
\newtcolorbox{compare}[1][]{
  enhanced,breakable,
  colback=green!3!white,colframe=green!50!black,
  fonttitle=\bfseries,title=对比,#1,
  boxed title style={colback=green!50!black},
  attach boxed title to top left={yshift=-2mm,xshift=2mm},
  arc=2mm,boxrule=0.8mm,
}
\newtcolorbox{warning}[1][]{
  enhanced,breakable,
  colback=red!5!white,colframe=red!75!black,
  fonttitle=\bfseries,title=常见陷阱,#1,
  boxed title style={colback=red!75!black},
  attach boxed title to top left={yshift=-2mm,xshift=2mm},
  arc=2mm,boxrule=0.8mm,
}
\newtcolorbox{note}[1][]{
  enhanced,breakable,
  colback=orange!5!white,colframe=orange!80!black,
  fonttitle=\bfseries,title=注意,#1,
  boxed title style={colback=orange!80!black},
  attach boxed title to top left={yshift=-2mm,xshift=2mm},
  arc=2mm,boxrule=0.8mm,
}
\newtcolorbox{bestpractice}[1][]{
  enhanced,breakable,
  colback=purple!5!white,colframe=purple!60!black,
  fonttitle=\bfseries,title=最佳实践,#1,
  boxed title style={colback=purple!60!black},
  attach boxed title to top left={yshift=-2mm,xshift=2mm},
  arc=2mm,boxrule=0.8mm,
}

% ==================== 章节标题 ====================
\usepackage{titlesec}
\usepackage{fancyhdr}
\usepackage{graphicx}
\usepackage{amssymb}
\usepackage{amsmath}
\usepackage{tabularx}
\usepackage{booktabs}
\usepackage{longtable}
\usepackage{array}
\usepackage{multirow}
\usepackage{enumitem}
\usepackage{seqsplit}
\usepackage{float}

\titleformat{\chapter}[display]
  {\normalfont\huge\bfseries\color{toolheadbg}}
  {\CTEXthechapter}{20pt}{\Huge}
\titleformat{\section}
  {\normalfont\Large\bfseries\color{blue!70!black}}
  {\thesection}{1em}{}
\titleformat{\subsection}
  {\normalfont\large\bfseries\color{blue!60!black}}
  {\thesubsection}{1em}{}

\pagestyle{fancy}
\fancyhf{}
\fancyhead[LE,RO]{\thepage}
\fancyhead[RE]{\leftmark}
\fancyhead[LO]{\rightmark}
\setlength{\headheight}{14.5pt}

\setlist[itemize]{leftmargin=2em,itemsep=2pt}
\setlist[enumerate]{leftmargin=2em,itemsep=2pt}
\emergencystretch=6em

% ==================== 超链接 ====================
\usepackage{hyperref}
\hypersetup{
  colorlinks=true,
  linkcolor=blue!70!black,
  urlcolor=blue!60!black,
  bookmarks=true,
  pdfauthor={C++ Reference},
  pdftitle={C/C++ 音视频开发参考手册：设备栈、编解码与实时流媒体管线},
}

% ==================== 自定义命令 ====================
\newcommand{\cpp}[1]{C++#1}
\newcommand{\Fn}[1]{\texttt{#1()}}
\newcommand{\Tp}[1]{\texttt{#1}}
\newcommand{\Opt}[1]{\texttt{-\/-#1}}
\newcommand{\Bk}{\discretionary{}{}{}}
\newcommand{\Lib}[1]{\textbf{\texttt{#1}}}
\newcommand{\Lang}[1]{\textsf{#1}}
\newcommand{\CL}{\Lang{C}}
\newcommand{\CPP}{\Lang{C++}}
\newcommand{\PY}{\Lang{Python}}
% 域标签：本书讲"两侧对比"时统一用这四个写法，避免同义多形
% \DomAud/{\Bk}\DomVid 标音/视频各自的链路，\DomRT/{\Bk}\DomBatch 标实时与离线批处理
\newcommand{\DomAud}{\Lang{Audio}}
\newcommand{\DomVid}{\Lang{Video}}
\newcommand{\DomRT}{\Lang{Realtime}}
\newcommand{\DomBatch}{\Lang{Offline}}

% ====================================================================
% book 双页模式默认 \flushbottom，会靠拉伸行距把每页撑满；本仓库框体与长表多，
% 撑页产生的 Underfull \vbox 是排版噪声而非缺陷，这里统一改用小底部对齐。
\raggedbottom

\begin{document}

\frontmatter

\begin{titlepage}
\centering
\vspace*{4cm}
{\Huge\bfseries C/C++ 音视频开发参考手册\par}
\vspace{1cm}
{\LARGE 设备栈、编解码内核与实时流媒体管线\par}
\vspace{0.4cm}
{\Large ALSA、WASAPI、FFmpeg、x264、Opus 与 GStreamer、WebRTC 的 C/C++ 接口\par}
\vfill
{\large \today\par}
\end{titlepage}

\tableofcontents

\mainmatter

\cleardoublepage

% ====================================================================
% 章节正文由 frag_chNN.tex 合并插入到下面这行之前（merge_frags.py）

% ============================================================
% ---- frag_ch01.tex ----
% ---------------------------------------------------------------------
% 第 1 章 分界与成因：为什么音视频内核落在 C/C++
% ---------------------------------------------------------------------
\part{音视频的两条时间轴：离线批处理与硬实时}

\chapter{分界与成因：为什么音视频内核落在 C/C++}
\label{ch:ch01-boundary}

同样一段 \DomAud{} 代码，跑在转码命令行里和跑在通话软件里，是两个不同的
工程问题。前者可以把一秒钟的音频花三秒钟去处理，只要总产量够高；后者必须
在声音真正被扬声器放出来之前把这一帧填好，晚一点就是可听见的断裂。本册
把这条分界当作全书的地基：先说清两条时间轴各自的验收标准
（第~\ref{sec:ch01-two-timelines} 节），
再量一量实时那一侧的约束到底有多硬
（第~\ref{sec:ch01-hard-budget} 节），然后解释带宽量级为什么会把设计逼向零拷贝
（第~\ref{sec:ch01-bandwidth} 节），以及这一域里 \CL/{\Bk}\CPP{} 与上层语言的实际分工
（第~\ref{sec:ch01-where-cpp} 节）。

\section{两条时间轴：验收标准先于代码}
\label{sec:ch01-two-timelines}

\DomBatch{}（转码、剪辑渲染、离线降噪、字幕与指纹批量提取）的货币是
\emph{吞吐}：单位墙钟时间内产出多少秒素材。它的失败是可以撤销的——一次
崩溃就重跑一次，一段参数不对就再跑一遍，甚至可以把同一个输入同时喂给
两组参数做 A/B 对比。因此批处理允许缓存整个文件、允许为提升压缩率把
单帧耗时提高十倍、允许在第一遍里统计直方图再在第二遍里分配比特。

\DomRT{}（语音通话、会议、直播连麦、乐器监听、VR 音频）的货币是
\emph{延迟与连续性}：端到端延迟要压在某个预算内，且每一段都得按时交付。
它的失败是不可撤销的——错过一个周期，用户听到的就是``咔''的一声，看到的
就是花屏或音画错位，没有任何重跑的机会。也正因为如此，实时侧的性能
指标不能只看平均值：一次 30 ms 的 GC 停顿藏在平均 8 ms 的后面，仍然会让
那一路会议彻底不可用。

两类系统对同一段代码提出的要求经常是互相冲突的，表~\ref{tab:ch01-timeline}
把它列成可对照的八个维度。工程上最有效的做法不是折中，而是把链路拆成
``热路径''与``编排路径''两段：热路径只负责按时交付缓冲，编排路径负责决策、
重试与统计。

\begin{longtable}{>{\raggedright\arraybackslash}p{0.20\textwidth}
                  >{\raggedright\arraybackslash}p{0.31\textwidth}
                  >{\raggedright\arraybackslash}p{0.31\textwidth}}
\caption{\DomRT{} 与 \DomBatch{} 的验收维度对照}\label{tab:ch01-timeline} \\
\toprule
\textbf{维度} & \textbf{\DomRT{}：通话、直播、监听} & \textbf{\DomBatch{}：转码、渲染} \\
\midrule
\endfirsthead
\toprule
\textbf{维度} & \textbf{\DomRT{}} & \textbf{\DomBatch{}} \\
\midrule
\endhead
\midrule
\bottomrule
\endlastfoot
目标函数 & 端到端延迟上界与连续丢帧率 & 吞吐（每秒产出多少秒素材）与码率效率 \\
失败的代价 & 不可撤销：用户当场听到、看到 & 可撤销：重跑、换参数再跑一遍 \\
时间预算 & 固定，由采样率与帧率反推出来 & 弹性，可以慢于素材时长 \\
分配策略 & 预先分配、池化，回调里不分配 & 按需分配，允许大缓冲与临时文件 \\
并发模型 & 优先级固定，避免无界锁与优先级反转 & 线程池、任务图，允许阻塞等待 \\
可观测手段 & 只能采样：统计本身也进热路径 & 全量日志、逐帧指标、离线分析 \\
优化方向 & 提前算、查表、把决策挪出回调 & 多遍编码、更贵的率失真搜索 \\
典型验收 & 95/99 分位延迟、xrun 次数、MOS & 完成时间、CPU/GPU 占用、VMAF \\
\end{longtable}

\section{硬实时那一侧的量级：一个回调到底有多少微秒}
\label{sec:ch01-hard-budget}

把量级算出来，才知道``实时''两个字意味着什么。

\begin{itemize}
  \item \DomAud{}：48 kHz 意味着 1 秒 48000 帧。一个 128 帧的周期是
        $128/48000$ 秒，约 2.67 ms；10 ms 的帧则是 480 帧。立体声 S16
        格式下 128 帧只有 512 字节，填它的时间比一次堆分配的系统调用
        开销大不了多少。
  \item \DomVid{}：60 fps 只有 16.7 ms 一整帧的预算，还要减去显示与
        合成的时间；30 fps 的 33.3 ms 看起来宽松，但一条链路里
        采集、编码、网络、解码、渲染五段平分下来，每段只剩几个毫秒。
  \item 抖动：两次回调的间隔不是恒定的，被中断、调度与内存回收拉扯。
        预算必须按最坏情况而不是平均值留。
\end{itemize}

于是热路径上有四类事情是明确不该出现的：\emph{垃圾回收或任何不可预测的
停顿式内存管理}；\emph{页错误}（包括首次触碰新分配页、换页、内存映射
文件缺页）；\emph{无界锁}（等待时间取决于别的线程什么时候放手，还可能
撞上优先级反转）；\emph{阻塞式 IO 与任何进入内核后可能睡眠的调用}。
列表里的每一条在批处理里都稀松平常，在回调里都是故障源。

代码~\ref{lst:ch01-alsa-loop} 是一段真实的 ALSA 写循环：它把这两个约束
体现在错误处理上——正常路径里只有算术与一次 \Fn{snd\_pcm\_writei}，
只有出错的分支才允许睡眠，因为那时已经不再是稳态。

\begin{lstlisting}[style=cstyle,
  caption={ALSA 写循环：正常路径极简，xrun 与 suspend 走恢复分支},
  label={lst:ch01-alsa-loop}]
#include <alsa/asoundlib.h>
#include <stdio.h>
#include <unistd.h>

/* 48 kHz / 2 ch / S16_LE：128 帧的 period 约 2.67 ms。
 * 正常路径里没有 malloc、没有锁、没有磁盘 IO。 */
static int write_frames(snd_pcm_t *pcm, const int16_t *src, size_t frames)
{
    size_t sent = 0;

    while (sent < frames) {
        snd_pcm_sframes_t wr = snd_pcm_writei(pcm,
                                             src + sent * 2,   /* 交织：每帧 2 采样 */
                                             frames - sent);
        if (wr == -EPIPE) {          /* xrun：追不上 DMA，状态机进 XRUN */
            if (snd_pcm_recover(pcm, (int) wr, 0) < 0)
                return -1;           /* recover 内部做 prepare/setparams */
            continue;                /* 同一批数据重投一次 */
        }
        if (wr == -ESTRPIPE) {       /* 设备被挂起：只能等它回来 */
            while (snd_pcm_resume(pcm) == -EAGAIN)
                usleep(1000);        /* 只有这条分支可以睡 */
            if (snd_pcm_recover(pcm, -ESTRPIPE, 0) < 0)
                return -1;
            continue;
        }
        if (wr < 0)
            return -1;               /* 其他 errno：交给上层重建整条流 */
        sent += (size_t) wr;         /* 实际写入的帧数可能少于请求 */
    }
    return 0;
}
\end{lstlisting}

\begin{warning}[title={把恢复逻辑写成``再睡一会儿''}]
第~\ref{sec:ch01-hard-budget} 节里的 \Fn{usleep} 只出现在 \Tp{-ESTRPIPE}
分支，因为设备挂起本来就不是稳态；一旦有人在正常分支里加
``睡 1 ms 等等看''，这条流就变成了周期性的 xrun 发生器。判断依据很简单：
任何以``时间''为条件的等待，都必须以``状态''为条件的重试来替换。
状态机恢复的正确顺序是 \Fn{snd\_pcm\_prepare}（或直接用
\Fn{snd\_pcm\_recover}）之后重新投递同一批帧，而不是丢掉这批帧再睡。
\end{warning}

\section{带宽的量级：为什么必须零拷贝与引用计数}
\label{sec:ch01-bandwidth}

\DomVid{} 的另一个约束是数据量。3840 $\times$ 2160、60 fps 每秒约
5 亿个像素。按 4:2:0 计算，8 bit 每像素合 1.5 字节，量级是每秒
0.7 GB；10 bit 若装进 16 bit 容器（几乎所有管线与显卡都这样存），
每像素合 3 字节，量级翻到每秒 1.5 GB，即十 Gbit 以上——一条 10 GbE
链路装不下它。即使按更常见的 1080p 计算，未压缩数据也是每秒数百 MB
的量级，而 CPU 单线程的 \Fn{memcpy} 吞吐不过个位数到十几 GB/s，
还要占用整条内存带宽。

结论是：热路径上不允许``传一帧''等于``拷一帧''。业界的统一答案是把缓冲
变成带引用计数的句柄，帧在管线里以指针流动：\Lib{libavcodec} 用
\Tp{AVFrame} 与 \Tp{AVBufferRef}，\Fn{av\_frame\_ref} 只把引用计数加一；
GStreamer 用 \Tp{GstBuffer} 加 \Tp{GstMemory}；PipeWire 用
\Tp{spa\_buffer} 与共享内存；\DomAud{} 侧 JACK 与 CoreAudio 直接交
客户端一个指向环形缓冲的指针。表~\ref{tab:ch01-copy} 给出同一帧在不同
策略下的量级差别（1080p 8 bit 4:2:0，约 3.1 MB，60 fps）。

\begin{longtable}{>{\raggedright\arraybackslash}p{0.30\textwidth}
                  >{\raggedright\arraybackslash}p{0.26\textwidth}
                  >{\raggedright\arraybackslash}p{0.26\textwidth}}
\caption{一帧数据在三种传递策略下的代价（量级估算，具体依版本与硬件而异）}\label{tab:ch01-copy} \\
\toprule
\textbf{传递方式} & \textbf{每帧额外带宽} & \textbf{能否放进回调} \\
\midrule
\endfirsthead
\toprule
\textbf{传递方式} & \textbf{每帧额外带宽} & \textbf{能否放进回调} \\
\midrule
\endhead
\midrule
\bottomrule
\endlastfoot
深拷贝一遍 & 读加写约 6 MB，缓存被冲掉 & 勉强可以，但会挤掉后续环节 \\
引用计数句柄 & 一次原子自增，数十字节 & 可以，这是各框架的默认做法 \\
映射与零拷贝 & 只有页表与同步开销 & 可以，前提是缓冲生命周期受控 \\
\end{longtable}

零拷贝换来的第二笔账是内存占用：没有引用计数时，每个环节都留着
一份私有副本；有引用计数时，同一份像素可以同时在采集、缩放、编码
与预览四只手上传，最后一个持有者释放。代价是生命周期变复杂——
你必须保证被引用的缓冲在使用期间不被原始拥有者归还给系统，
这也是各框架在 API 里区分``借用''与``持有''两种语义的原因。

\section{\DomBatch{} 的时间可以怎么花：一个编排层的正确与错误用法}
\label{sec:ch01-python-role}

\PY{} 与 \Lang{JavaScript}{} 在音视频工程里的位置是编排：拼命令行、
解析 ffprobe 的 JSON、调度任务、做统计与报表。代码~\ref{lst:ch01-py-batch}
的上半段是这批语言的舒适区——一次调用跑几秒到几分钟，失败就重试；
下半段则是同一批语言被误放进热路径的样子。

\begin{lstlisting}[style=pythonstyle,
  caption={同一个文件里的两种命运：上半段是合格的批处理编排，下半段是不能上生产的反例},
  label={lst:ch01-py-batch}]
import subprocess
import time

import numpy as np
import soundfile as sf


def transcode_worker(in_path: str, out_path: str) -> None:
    ``''``批处理：慢没关系，可重跑，失败可以换个参数再来。''``''
    subprocess.run([``ffmpeg'', ``-y'', ``-i'', in_path,
                    ``-c:a'', ``libopus'', ``-b:a'', ``64k'', out_path],
                   check=True)                    # 阻塞等待进程结束


def bad_realtime_tick(frames: int) -> None:
    ``''``反例：这一段如果被接进音频回调，故障是必然的。''``''
    buf = np.zeros((frames, 2), dtype=np.int16)   # 分配：新页首次触碰会缺页
    sf.write(``/tmp/tick.wav'', buf, 48000)         # 磁盘 IO：时长完全不可控
    time.sleep(0.002)                             # 一个 period 只有约 2.67 ms
    global_counter += 1                           # 全局状态，跨线程还要加锁
\end{lstlisting}

三条理由足以把下半段请出回调：其一，解释器逐条派发字节码，函数调用
与对象属性访问的开销是固定的几百纳秒量级，累积起来吃掉预算；其二，
内存回收（\PY{} 的引用计数加循环检测，以及其他托管语言的移动式回收）
的停顿时刻不由你决定；其三，\Fn{subprocess.run}、\Fn{sf.write}、
\Fn{time.sleep} 都会进入内核并可能睡眠，而睡眠一次的时长由整机负载
决定，与你的 2.67 ms 无关。

需要强调的是：这三条并不是``\PY{} 慢''，而是``这三类操作的时间分布
与硬预算不兼容''。同一台机器上，用 \PY{} 起 200 个 \Fn{ffmpeg}
进程把一天的素材压完是很好的方案，因为批处理只关心动望。

\section{这一域里 \CL/\CPP{} 的实际位置}
\label{sec:ch01-where-cpp}

打开任何一个音视频工程的依赖树，会发现一个稳定的分层：\emph{能被
调用的那一层几乎全是 \CL{} 的 ABI}。具体地：
\begin{itemize}
  \item \DomVid{}/{\Bk}\DomAud{} 编解码内核：\Lib{libavcodec}、\Lib{libavformat} 与
        \Lib{libavutil}，加上 \Lib{libswscale}、\Lib{libswresample}，全 \CL{}
        接口；\Lib{x264}、\Lib{x265}、\Lib{libopus}、\Lib{libvpx}
        对外也是 \CL{} 头文件，即使 \Lib{x265} 内部是 \CPP{} 实现。
  \item 设备层：Linux 的 ALSA（\Fn{snd\_{\Bk}pcm\_{\Bk}writei} 一族的 \CL{} 函数）
        与 PipeWire（\CL{} 对象与回调表），Windows 的 WASAPI
        （以 \CL{} 风格调用的 COM 接口），Apple 的 CoreAudio 与 AudioUnit
        （\CL{} 函数加 \CL{} 结构体），跨平台的 JACK（也是 \CL{}）。
  \item 框架层：GStreamer 的框架与滤镜都是 \CL{} 的 GObject 体系；
        WebRTC 的媒体内核是 \CPP{}，但对外交出的仍是
        \Tp{rtc::scoped\_refptr} 这类引用计数句柄；
        FFmpeg 的滤镜图是 \CL{}，Vulkan 是 \CL{} API，
        D3D11 的媒体扩展则是 COM 风格接口。
\end{itemize}

上层语言通过 FFI 接进这个 ABI：\PY{} 的 ctypes/cffi、\Lang{Node.js}{}
的 N-API 与 WebAssembly、\Lang{Rust}{} 的 bindgen、\Lang{C\#}{} 的
P/Invoke、Android 的 JNI。它们的共同特征是\emph{调用点必须稀疏}：
一次跨语言边界的开销在微秒级，对批处理无所谓，对 2.67 ms 的周期
就是几个百分点的预算。因此本册的立场很明确——\DomBatch{} 的胶水
随便写，\DomRT{} 的热路径只用 \CL/{\Bk}\CPP{} 把缓冲按时交付，
决策留在别处。

\section{本册路线}
\label{sec:ch01-roadmap}

后面几章按``先时间、再设备、再编解码、最后链路''的顺序展开，与
表~\ref{tab:ch01-timeline} 的两条时间轴一一对应：

\begin{itemize}
  \item 第~\ref{sec:ch01-hard-budget} 节的量化在
        第~\ref{ch:ch02-clock-buffer} 章里变成时钟漂移与缓冲层级的具体公式，
        并引出 jitter buffer 的自适应目标延迟。
  \item 设备侧：第~\ref{ch:ch03-audio-device} 章逐个拆 ALSA、PipeWire、
        WASAPI、ASIO、CoreAudio 与 JACK 的 C API 形态；
        第~\ref{ch:ch04-capture-display} 章处理采集与显示的时钟对齐。
  \item 编解码内核：第~\ref{ch:ch06-codec-api} 章的
        \Fn{avcodec\_send\_packet} 与 \Fn{avcodec\_receive\_frame}
        两阶段模型是本册的中枢，配合第~\ref{ch:ch07-rate-control} 章
        的码控与第~\ref{ch:ch08-decode-hw} 章的硬件解码。
  \item 链路侧：第~\ref{ch:ch09-filter-resample} 章的滤镜与重采样、
        第~\ref{ch:ch10-mux-sync} 章的封装与同步、
        第~\ref{ch:ch11-streaming} 章的实时传输，以及
        第~\ref{ch:ch14-congestion-abr} 章的拥塞控制与自适应码率。
  \item 工程化：第~\ref{ch:ch12-realtime-thread} 章讲线程与优先级，
        第~\ref{ch:ch13-framework-integration} 章讲与 GStreamer、WebRTC
        这类框架的集成；选型速查见第~\ref{ap:apC-selection} 附录。
\end{itemize}

\begin{bestpractice}[title={先写下验收标准，再碰第一行热路径代码}]
\begin{itemize}
  \item 立项时把延迟预算写成三个数字：目标分位（例如 99 分位）、
        允许的最大值、以及超预算时的降级动作。没有这三个数字，
        ``实时''只是形容词。
  \item 在代码里明确划出热路径边界：边界内只用预分配内存、无锁或
        有界锁、非阻塞调用；所有决策、日志与重试挪到边界外。
  \item 用真实量级检验设计：把周期时长、一帧字节数、每秒帧数写在
        注释或文档里（例如 128 帧 $@$ 48 kHz 约 2.67 ms），
        让每次评审都能一眼看出某次拷贝是否付得起。
  \item 上层语言只做编排与统计，不进热路径；跨语言调用点要稀疏，
        并把``一次调用做一批事''作为默认形态。
\end{itemize}
\end{bestpractice}

% ============================================================
% ---- frag_ch02.tex ----
% ---------------------------------------------------------------------
% 第 2 章 采样、时钟与缓冲：从 PCM 到 jitter buffer
% ---------------------------------------------------------------------
\chapter{采样、时钟与缓冲：从 PCM 到 jitter buffer}
\label{ch:ch02-clock-buffer}

第~\ref{ch:ch01-boundary} 章给出预算的数字，本章给出预算的来源。音频链路
的每个环节都在和时间换算：多少采样等于多少毫秒，多少毫秒的缓冲能吸收
多大的抖动。把这些换算写清楚，\DomAud{} 与 \DomVid{} 两侧的设计就不再是
两套术语，而是同一套账本的两种记法。

\section{PCM 的三个参数与一个布局}
\label{sec:ch02-pcm}

一段 \DomAud{} 数据由三个参数完全描述：

\begin{itemize}
  \item \emph{采样率}：每秒多少帧。人声链路常用 8 kHz 与 16 kHz，通话与
        音乐用 48 kHz，母带级用 96 kHz 或 192 kHz；\Lang{CD}{} 用 44.1 kHz。
        奈奎斯特定则要求采样率高于被采样信号最高频率的两倍，这决定了
        上限而不是下限——用 48 kHz 录 20 kHz 的声音可以，反过来不行。
  \item \emph{位深与存储格式}：16 bit 的量化噪声理论动态范围约 96 dB，
        24 bit 约 144 dB，远超实际话筒与本底噪声的差值，因此传输侧
        几乎都是 S16，只有录制与混音中间环节保留 S32 或 \Tp{float}。
        浮点在 \Lib{FFmpeg}{} 里是 \Tp{AV\_SAMPLE\_FMT\_FLT}，
        在 ALSA 里是 \Tp{SND\_PCM\_FORMAT\_FLOAT\_LE}，
        两者都表示 32 bit 浮点。
  \item \emph{声道布局}：单声道、立体声、四声道、5.1 与 7.1。
        \Lib{libavutil}{} 用 \Tp{AVChannelLayout} 描述，
        旧写法是 \Tp{channel\_layout} 位掩码，例如 5.1 对应
        \Tp{AV\_CH\_LAYOUT\_5POINT1} 的六个声道；ALSA 只协商声道数，
        哪个位置是哪个声道由驱动与规范约定，不参与参数协商。
\end{itemize}

关键概念是\emph{帧}（frame）：同一时刻每个声道各一个采样，帧数是所有
声道共享的时间刻度。一帧立体声 S16 是 4 字节，一帧 5.1（六个声道）是
12 字节。第~\ref{ch:ch01-boundary} 章的 ALSA 写循环里用
\Tp{snd\_pcm\_uframes\_t} 而不是字节数说话，正是这个原因。

交织与平面是同一批数据的两种排法：

\begin{itemize}
  \item \emph{交织}（interleaved）：L R L R L R，第 $i$ 帧第 $c$ 个声道
        落在第 $i$ 乘声道数加 $c$ 个采样处。设备与网络格式几乎都是
        交织的，因为一次 DMA 搬一段连续内存。
  \item \emph{平面}（planar）：先写完所有 L，再写所有 R，偏移是
        第 $c$ 乘帧数加 $i$。DSP 处理（增益、混音、滤波）偏爱
        平面，因为可以按声道做向量化。
\end{itemize}

\Lib{libavutil}{} 把两者编进采样格式枚举：\Tp{AV\_{\Bk}SAMPLE\_{\Bk}FMT\_{\Bk}S16}
是交织，\Tp{AV\_{\Bk}SAMPLE\_{\Bk}FMT\_{\Bk}S16P} 是平面；
\Tp{AV\_{\Bk}SAMPLE\_{\Bk}FMT\_{\Bk}S32} 与它的 P 结尾版本、
\Tp{AV\_{\Bk}SAMPLE\_{\Bk}FMT\_{\Bk}FLT} 与 \Tp{AV\_{\Bk}SAMPLE\_{\Bk}FMT\_{\Bk}FLTP} 同理
（最后这一对正是 AAC 编码器要的排法）。\Tp{AVFrame::data}
是指针数组，平面格式下每个声道占一个元素，交织格式下只有第 0 个有效；
\Tp{AVFrame::nb\_{\Bk}planes}（旧的判法是 \Fn{av\_{\Bk}get\_{\Bk}planar\_{\Bk}sample\_{\Bk}fmt}）
回答该用哪个循环。ALSA 侧对应的是访问模式，取值由
\Tp{SND\_{\Bk}PCM\_{\Bk}ACCESS\_{\Bk}} 前缀加四种组合：读写交织、读写非交织、
内存映射交织与内存映射非交织。

\begin{warning}[title={格式假设的两种失败方式}]
把平面数据当交织读，得到的不是噪声而是\emph{整段偏一个声道的错位}：
左声道全是前一半采样，右声道全是后一半，听感上像``人声跑到一边''，
但波形没有任何异常，因此很难在日志里发现。反过来把交织当平面读，
只会听到刺耳的嘶声。更隐蔽的是采样格式与位序的混淆：
\Tp{SND\_PCM\_FORMAT\_S16\_LE} 与 \Tp{SND\_PCM\_FORMAT\_S16\_BE}
只差一个枚举值，在少数大端架构上就会踩到。测试时分别送左右
不同的正弦，是最便宜的防护。
\end{warning}

\section{两台独立晶振：漂移不是意外而是必然}
\label{sec:ch02-drift}

采集卡与声卡各有各的晶振。消费级时钟的精度在几十 ppm 量级
（1 ppm 是百万分之一，常见的 25 ppm 温漂与老化叠加后，
把正负 50 ppm 当作设计假设是合理的），而正负 50 ppm 在 48 kHz 上意味着
每秒最多相差 2.4 个采样，一分钟约 144 个采样、也就是 3 ms，
十分钟就累积到 30 ms 量级
——且方向随时钟温度变化而反转。任何``两端时钟不同源''的
链路（独立 USB 声卡、蓝牙耳机、另一台机器上的对端）都会遇到它。

于是产生一个绕不开的事实：\emph{不处理漂移的长链路必然 xrun}。
播放端每小时需要 172.8 M 帧（48000 $\times$ 3600），只要生产者按自己的
时钟走，缓冲区就会缓慢地朝空或满的一侧漂，漂到头就是断音或积压。
工程上有三类对策，代价依次升高：

\begin{enumerate}
  \item \emph{丢帧或插静音}：最简单，但每次数毫秒的可闻断裂，只适合
        漂移极慢且缓冲足够大的场合。
  \item \emph{时间拉伸或连续重采样}：把比例偏差摊到很长一段里，
        每次只修正千分之几，听不见。\Lib{libsamplerate}{} 这类重采样器
        （\Tp{src\_simple}、\Fn{src\_process}）适合与采样率换算合并做；
        要保持音高则用 WSOLA 一类的拉伸算法。代价是每采样多做几次乘加。
  \item \emph{拉时钟}：让生产者服从播放端节拍（\DomRT{} 音频引擎的常规做法，
        见第~\ref{ch:ch12-realtime-thread} 章），或直接让两端共用一个
        时钟源：专业声卡的 word clock、网络的 PTP（IEEE 1588），
        以及 \Lang{WebRTC}{} 用 RTCP 发送方报告把 RTP 时间戳
        与墙上时间对齐。
\end{enumerate}

命令行里最常见的补救是 \Lib{FFmpeg}{} 的 \Fn{aresample} 滤镜：代码~\ref{lst:ch02-shell}
的第 4 条命令让它按时间戳补或删采样，把已经对不齐的两条流重新拉回同一
节拍。

\begin{lstlisting}[style=shellstyle,
  caption={与设备能力、漂移补救有关的四条命令},
  label={lst:ch02-shell}]
# 1) 问设备：hw:0,0 接受哪些采样率、格式、period 与 buffer 区间
arecord -D hw:0,0 --dump-hw-params \
        -f S16_LE -r 48000 -c 2 -d 1 /dev/null

# 2) 绕开一切插件，直接采裸 PCM 再原样播回去（验证链路只有 DMA）
arecord -D hw:0,0 -f S16_LE -r 48000 -c 2 -t raw -d 5 mic.raw
aplay -D hw:0,0 -f S16_LE -r 48000 -c 2 -t raw mic.raw

# 3) 同一批数据交给插件层：plughw 会替你完成格式转换与重采样
arecord -D plughw:0,0 -f S32_LE -r 96000 -c 2 -d 5 mic32.raw

# 4) 时钟对不齐：让 aresample 依据时间戳补/删采样（阈值单位依版本而异）
ffmpeg -f s16le -ar 48000 -ac 2 -i mic.raw \
       -af aresample=async=1:first_pts=0 -f wav aligned.wav
\end{lstlisting}

\begin{note}[title={\Tp{plughw} 与 \Tp{hw} 的分工}]
\Tp{hw:0,0} 直通内核驱动：参数越界就直接拒绝，延迟最低、约束最硬，
适合本册说的热路径。插件层 \Tp{plughw:0,0} 允许你交一个设备并不支持的
格式，它内部插入重采样与格式转换，代价是多一级缓冲与一段不可控的 CPU
工作。调试时先确认自己走的是哪一层，否则所有延迟数字都对不上。
\end{note}

\section{缓冲的层次与延迟的合成}
\label{sec:ch02-buffer-layers}

一次播放至少穿过四级缓冲，每级都为平滑付出一点延迟：

\begin{enumerate}
  \item \emph{应用缓冲}：你自己的环形队列或无锁 SPSC 队列，决定生产
        与消费解耦到什么程度。
  \item \emph{驱动与硬件环}：ALSA 的 buffer 是被 DMA 循环读写的那段
        连续内存，period 是每填满或读完一段触发一次中断的粒度。
        约定是 \Tp{buffer\_size} 等于 \Tp{period\_size} 的整数倍，
        常见 2 到 16 个 period。
  \item \emph{服务器层}：PipeWire 或 PulseAudio 的 sink 缓冲、WASAPI
        共享模式下的系统混音器，通常以毫秒而不是采样说话。
  \item \emph{网络层}：jitter buffer（第~\ref{sec:ch02-jitter} 节），
        把不可预测的到达时刻整理成可预测的读取时刻。
\end{enumerate}

代码~\ref{lst:ch02-hwparams} 是这一段的核心：它把``约束协商''写完整，
并显式暴露了延迟的合成方式——起播阈值越大越不怕卡顿，但也越晚出声。

\begin{lstlisting}[style=cstyle,
  caption={ALSA 参数协商：hw\_params 定形状，sw\_params 定节拍与唤醒时机},
  label={lst:ch02-hwparams}]
#include <alsa/asoundlib.h>

/* 期望：48 kHz / 2 ch / period 128 帧 / buffer 512 帧（4 个 period）。 */
static int negotiate(snd_pcm_t **pcm, const char *dev, unsigned rate,
                     unsigned channels, snd_pcm_uframes_t *period,
                     snd_pcm_uframes_t *buffer)
{
    snd_pcm_hw_params_t *hw;
    snd_pcm_sw_params_t *sw;
    int err, dir = 0;

    err = snd_pcm_open(pcm, dev, SND_PCM_STREAM_PLAYBACK, 0);
    if (err < 0)
        return err;                     /* 打不开设备：换卡或换会话 */
    snd_pcm_hw_params_alloca(&hw);      /* alloca 族：栈上分配，无需 free */
    snd_pcm_hw_params_any(*pcm, hw);    /* 从设备能力全集起步 */
    snd_pcm_hw_params_set_access(*pcm, hw,
                    SND_PCM_ACCESS_RW_INTERLEAVED);
    snd_pcm_hw_params_set_format(*pcm, hw, SND_PCM_FORMAT_S16_LE);
    snd_pcm_hw_params_set_channels(*pcm, hw, channels);
    /* near 系列把请求值改成设备实际接受值，dir 指示偏大还是偏小 */
    snd_pcm_hw_params_set_rate_near(*pcm, hw, &rate, &dir);
    snd_pcm_hw_params_set_period_size_near(*pcm, hw, period, &dir);
    snd_pcm_hw_params_set_buffer_size_near(*pcm, hw, buffer);
    if ((err = snd_pcm_hw_params(*pcm, hw)) < 0)
        return err;                     /* 约束互相矛盾：整组失败 */

    /* 回读实际值：near 只保证接近，不保证等于请求 */
    snd_pcm_hw_params_get_period_size(hw, period, &dir);
    snd_pcm_hw_params_get_buffer_size(hw, buffer);

    snd_pcm_sw_params_alloca(&sw);
    snd_pcm_sw_params_current(*pcm, sw);
    /* avail 攒够一个 period 才叫醒上层：一次唤醒 = 一次交付 */
    snd_pcm_sw_params_set_avail_min(*pcm, sw, *period);
    /* 攒满整个 buffer 才起播：最抗卡顿，起播延迟也最大 */
    snd_pcm_sw_params_set_start_threshold(*pcm, sw, *buffer);
    return snd_pcm_sw_params(*pcm, sw); /* sw 参数同样要一次提交 */
}
\end{lstlisting}

延迟合成的读法就藏在里面：播放方向的排队延迟约等于
\Tp{start\_threshold} 对应的时间，唤醒粒度是 \Tp{period\_size}，
而可容忍的积压上限是 \Tp{buffer\_size} 减去 \Tp{period\_size}。
表~\ref{tab:ch02-latency} 把常见组合换算成毫秒，用于反推预算。

\begin{longtable}{>{\raggedright\arraybackslash}p{0.15\textwidth}
                  >{\raggedright\arraybackslash}p{0.22\textwidth}
                  >{\raggedright\arraybackslash}p{0.15\textwidth}
                  >{\raggedright\arraybackslash}p{0.28\textwidth}}
\caption{采样率、单帧大小与缓冲帧数决定的延迟（只算缓冲本身，未计处理与网络）}\label{tab:ch02-latency} \\
\toprule
\textbf{采样率} & \textbf{单帧大小（时长）} & \textbf{缓冲帧数} & \textbf{缓冲延迟} \\
\midrule
\endfirsthead
\toprule
\textbf{采样率} & \textbf{单帧大小（时长）} & \textbf{缓冲帧数} & \textbf{缓冲延迟} \\
\midrule
\endhead
\midrule
\bottomrule
\endlastfoot
48 kHz & 128 帧（2.67 ms） & 4 & 约 10.7 ms，乐器监听的量级 \\
48 kHz & 128 帧（2.67 ms） & 16 & 约 42.7 ms，通话可接受 \\
48 kHz & 480 帧（10 ms） & 4 & 40 ms，会议软件常见起点 \\
44.1 kHz & 512 帧（11.6 ms） & 4 & 约 46.4 ms \\
48 kHz & 960 帧（20 ms） & 4 & 80 ms，转码或直播上行 \\
16 kHz & 160 帧（10 ms） & 3 & 30 ms，窄带通话引擎 \\
96 kHz & 256 帧（2.67 ms） & 8 & 约 21.3 ms \\
\end{longtable}

三点提醒：其一，表里的数字要乘上环节数，采集加播放两套缓冲就是两倍；
其二，关掉 \Tp{period\_event} 之后，\Fn{snd\_pcm\_wait} 的唤醒时机
不再是 period 边界，靠它定节拍的代码要跟着改；
其三，多出一级服务器（PipeWire 的 quantum 或 WASAPI 共享混音器）
就多出一份等价的缓冲，这就是共享模式延迟高于独占模式的根本原因，
详见第~\ref{ch:ch03-audio-device} 章。

\section{\DomVid{} 侧的对应物：帧队列、vsync 与显示时间戳}
\label{sec:ch02-video-layers}

\DomVid{} 把同一套账本换成帧做单位。双缓冲让绘制与显示交替，
但一帧被改一半就送上门会导致撕裂；三缓冲再加一级排队，代价是
一到两个刷新周期的额外延迟。刷新周期由显示模式决定：60 Hz 是 16.7 ms，
120 Hz 是 8.3 ms，而 23.976 fps 的素材周期是约 41.7 ms
（$1000/23.976$），两者不成整数倍，这就是播放 24 fps 素材需要
3:2 pulldown 或时间戳重排的根源——与第~\ref{sec:ch02-drift} 节的晶振
漂移在数学上是同一件事：两个不同源的节拍。

\DomVid{} 的时间尺度是 presentation timestamp，即 PTS。它挂在
\Tp{AVFrame} 的 \Tp{pts} 上（解码后更稳的读法是
\Tp{best\_effort\_timestamp}），单位由 \Tp{AVFrame::time\_base}
或流级 \Tp{AVStream::time\_base} 定义，
于是 \Fn{av\_rescale\_q} 成了全书最常见的换算：把 \Tp{int64\_t} 的时间
从一种时基搬到另一种。两种时基值得记住：\Lib{MPEG-TS}{} 与 RTP 用
90 kHz，因为它能被 24、25、30、50、60 整除；\DomAud{} 封装常用
采样率当时基，此时 1 个 tick 就是 1 个采样。

\begin{compare}[title={\DomAud{} 与 \DomVid{} 的对应概念}]
48 kHz 的 480 帧对应 10 ms，60 fps 的 1 帧对应 16.7 ms：音频以毫秒为单位
说话，视频以帧为单位说话，两者靠 PTS 对齐。音频缓冲越深越不怕抖动、
越怕延迟；视频缓冲越深越不怕丢包，越怕操作反馈变慢。音频的失败是
咔哒声，视频的失败是花屏或冻结。音频必须逐帧交付，视频可以丢帧、
可以跳到一个关键帧重来——这是第~\ref{ch:ch10-mux-sync} 章同步策略
里``丢视频保音频''的依据。
\end{compare}

\section{jitter buffer：把不规则的到达换成稳定的读取}
\label{sec:ch02-jitter}

网络上一个 RTP 包的到达时刻由排队、重传与拥塞决定，抖动可达几十毫秒。
jitter buffer 的作用就一句话：\emph{用固定延迟把随机到达整形成固定速率
的读取}。代价是把延迟预算换成抗抖能力，收益是把丢包与乱序吸收掉。

自适应 jitter buffer 的控制目标是当前缓冲水位：它统计到达间隔的分位数，
把目标延迟推到能覆盖近期绝大多数抖动的水平，并在网络转好后缓慢回落。
水位上下限一般写成毫秒的区间，实际取值随实现与场景而异，
但形状一致：下挡保证不会因欠载而卡顿，上挡保证延迟不失控。
工程约束有三条：

\begin{itemize}
  \item 调整要平滑。每步只改一个包时长（20 ms 的编码器帧就是一个一个
        包地加减），并用时间拉伸吸收由此产生的采样数变化，否则
        每次扩容都是一声咔哒。
  \item 缓冲深度以包计而非以毫秒计，因为丢包隐藏（PLC）与重排都以
        包为最小单位；目标水位通常落在几十毫秒的量级，并随抖动统计
        在几秒内缓慢移动。
  \item 前置的处理顺序是：先按序号重排与去重，再解码，最后才把
        空洞交给 PLC。\Lib{Opus}{} 解码时把数据指针传 \Tp{NULL}
        就是请求对上一帧做 concealment；它的 inband FEC 让下一个包
        携带上一个包的冗余，因此恢复丢帧需要多等一个包的到达
        （能力开关依编码参数而异）。
\end{itemize}

\begin{bestpractice}[title={把延迟预算写成一张表再动手}]
\begin{enumerate}
  \item 先定总额：通话建议端到端 150 ms 以内（约 80 ms 可被多数用户
        接受，150 ms 是常见的上限），乐器监听压到 20 ms 以内，
        直播单向可以到 500 ms 以上。
  \item 再切分：把总额分给采集缓冲、处理帧长、网络与 jitter buffer、
        播放缓冲四段，允许每段有 20\% 的浮动，写进配置而不是记在脑子里。
  \item 最后核对漂移：按几十 ppm 估算最长通话时长下的累积采样差，
        检查选定的重采样或拉时钟方案能否吸收它。
  \item 表~\ref{tab:ch02-latency} 是你唯一的换算工具：任何一处改动
        都先换算成毫秒再评审，不换算就不合并。
\end{enumerate}
\end{bestpractice}

% ============================================================
% ---- frag_ch03.tex ----
% ---------------------------------------------------------------------
% 第 3 章 音频设备栈与它们的 C API
% ---------------------------------------------------------------------
\chapter{音频设备栈与它们的 C API}
\label{ch:ch03-audio-device}

第~\ref{ch:ch02-clock-buffer} 章把延迟拆成了几级缓冲，本章去看这些缓冲
各自由谁提供。每个平台都不止一套音频栈，它们的名字经常出现在同一份
依赖表里：ALSA 与 PipeWire、共享模式 WASAPI 与 ASIO、CoreAudio 与
AudioUnit。搞清楚每一层的 C API 长什么样、缓冲模型是什么、你能把延迟
压到多少，才知道该在哪一层写代码。

\section{Linux：内核里的 ALSA，进程里的 PipeWire}
\label{sec:ch03-alsa-pipewire}

ALSA 分两半：内核里按声卡驱动的 \Tp{pcm} 设备节点，和用户态的
\Lib{libasound}{}。设备名决定你走哪条路：\Tp{hw:card,device} 直通内核，
\Tp{plughw:} 加上格式转换插件，\Tp{default} 走 ALSA 配置文件里的路由
（可能是软件混音 \Tp{dmix}）。协商过程就是第~\ref{ch:ch02-clock-buffer}
章那一套 \Tp{hw\_params} 与 \Tp{sw\_params}；协商完成后
\Fn{snd\_pcm\_state} 报出的状态机只有六个取值需要记住：
\Tp{SND\_PCM\_STATE\_SETUP}（还没参数化）、\Tp{PREPARED}、
\Tp{RUNNING}、\Tp{XRUN}、\Tp{SUSPENDED}、\Tp{DRAINED}，
外加一个 \Tp{DISCONNECTED}（设备被拔）。恢复动作的差别很实际：
\Tp{XRUN} 用 \Fn{snd\_pcm\_recover} 拉回 \Tp{PREPARED}；
\Tp{SUSPENDED} 要先 \Fn{snd\_pcm\_resume}；\Tp{DISCONNECTED}
只能重建整条流。

取数方式决定你能不能做到零拷贝：
\begin{itemize}
  \item \Tp{SND\_PCM\_ACCESS\_RW\_INTERLEAVED} 是最常见的写法：
        \Fn{snd\_pcm\_writei} 一次提交若干帧，库把数据拷进环形缓冲。
  \item 内存映射（mmap）访问模式配合 \Fn{snd\_pcm\_mmap\_begin}
        与 \Fn{snd\_pcm\_mmap\_commit}：前者返回一个
        \Tp{snd\_pcm\_mmap\_altsbuffer\_t}，其中的指针就指着那段
        DMA 缓冲，你直接在它上面混音，省掉一次拷贝。注意它给出的
        可用帧数可能跨过一次回绕，必须按返回的长度处理，
        不能自己假设内存连续。
  \item 非交织访问让每个声道一段指针，配合平面数据很方便，
        但硬件支持得少，常常要靠插件层转换。
\end{itemize}

PipeWire 是这几年的实际赢家：同一个守护进程里同时提供 PulseAudio 与
JACK 的兼容接口，发行版默认已经把它摆在这两者的位置上。它把音频与视频、
播放与采集、共享与低延迟收进同一张 graph：节点在同一进程内被调度，
节拍由图上的``量子''（quantum）驱动，默认在 1024 帧量级（约 21 ms，
具体值依版本与配置而异），需要低延迟的客户端可以把它调小。客户端 API 是
\Tp{pw\_stream}（流式播放与采集）与 \Tp{pw\_filter}（处理节点），
形态都是``注册一组回调，\Tp{process} 被叫到时读写 \Tp{spa\_buffer}''；
缓冲底层是共享内存，于是服务器与客户端之间只有一次映射的成本。

老栈并没有消失：\Lib{PulseAudio}{} 的简易接口 \Fn{pa\_simple\_new}
配 \Fn{pa\_simple\_write} 仍是很多工具的默认输出，它走异步
\Tp{pa\_operation} 加共享内存块，延迟与抖动都比直连 ALSA 大一级；
JACK 则留在音乐制作圈（第~\ref{sec:ch03-jack} 节）。选型时真正要问的是
``数据要穿过几级进程''，因为每多一级就多一次调度与一份缓冲。

\section{Windows：WASAPI 的两种共享模式与两种取数方式}
\label{sec:ch03-wasapi}

WASAPI 不是独立的守护进程，而是系统音频服务暴露的一组 COM 接口，
以 \CL{} 风格调用：每个方法的第一个参数是接口指针，写成
\Tp{IAudioClient\_Start(pClient)} 这样的宏式包装。入口链固定：
用设备枚举器取默认渲染设备，\Fn{IMMDevice\_Activate} 拿到
\Tp{IAudioClient}，然后 \Fn{IAudioClient\_Initialize} 一次性把模式、
缓冲时长与格式定下来。

两个维度决定行为：

\begin{enumerate}
  \item \emph{共享还是独占}。\Tp{AUDCLNT\_SHAREMODE\_SHARED}
        把数据交给系统混音器，格式必须问
        \Fn{IAudioClient\_GetMixFormat}（常见是 48 kHz、32 bit 浮点，
        声道数随扬声器配置而异），别人可以同时用设备。
        \Tp{AUDCLNT\_SHAREMODE\_EXCLUSIVE} 直接对接硬件缓冲，
        绕过混音器，格式必须是设备原生格式（用
        \Fn{IAudioClient\_IsFormatSupported} 问，必要时接受它给出的
        最接近值），而且独占期间别的应用发不出声。
  \item \emph{事件驱动还是轮询}。Initialize 时带上
        \Tp{AUDCLNT\_STREAMFLAGS\_EVENTCALLBACK}，再用
        \Fn{IAudioClient\_SetEventDevice} 交一个自动重置事件，
        系统每到周期边界就把它置起。不带这个标志就是轮询，
        自己按时间读 \Fn{IAudioClient\_GetCurrentPadding} 推算。
        轮询实现简单，但唤醒时机与设备节拍无关，容易多垫一层缓冲。
\end{enumerate}

代码~\ref{lst:ch03-wasapi} 是事件驱动循环的骨架。注意两件事：每次要填
的帧数是``缓冲总帧数减当前填充量''算出来的，不是常量；取到的指针
指向系统缓冲，就地写，别再开一块内存拷过去。

\begin{lstlisting}[style=cstyle,
  caption={WASAPI 共享模式的事件驱动渲染循环骨架（省略错误处理与混音实现）},
  label={lst:ch03-wasapi}]
#include <mmdeviceapi.h>
#include <audioclient.h>
#include <avrt.h>

static IAudioClient       *g_client;
static IAudioRenderClient *g_render;
static HANDLE              g_evt;          /* 自动重置事件 */
static UINT32              g_bufFrames;    /* 缓冲总帧数 */
static UINT32              g_blockAlign;   /* 每帧字节数 */

/* 事件被置起 = 一个周期边界到了，该补多少帧就补多少。 */
static void render_loop(void)
{
    for (;;) {
        UINT32 pad = 0, next;
        BYTE  *data = NULL;

        if (WaitForSingleObject(g_evt, INFINITE) != WAIT_OBJECT_0)
            break;                    /* 句柄失效：线程该退出了 */
        IAudioClient_GetCurrentPadding(g_client, &pad);
        next = g_bufFrames - pad;     /* 本周期要补的帧数 */
        if (next == 0)
            continue;

        if (FAILED(IAudioRenderClient_GetBuffer(g_render,
                                               next, &data)))
            break;
        mix_into(data, next, g_blockAlign);   /* 就地填，不分配 */
        IAudioRenderClient_ReleaseBuffer(g_render, next, 0);
    }
}

/* 一次性定下模式、时长、事件与格式；顺序不能颠倒。 */
static HRESULT start_shared(IAudioClient *client)
{
    WAVEFORMATEX *fmt = NULL;        /* GetMixFormat 会分配它 */
    HRESULT hr = IAudioClient_GetMixFormat(client, &fmt);
    if (FAILED(hr))
        return hr;
    g_blockAlign = fmt->nBlockAlign;

    hr = IAudioClient_Initialize(client, AUDCLNT_SHAREMODE_SHARED,
            AUDCLNT_STREAMFLAGS_EVENTCALLBACK,
            100000,                  /* 缓冲时长 10 ms，单位 100 ns */
            0,                       /* 共享模式此处必须传 0 */
            fmt, NULL);
    CoTaskMemFree(fmt);              /* 混音格式是系统分配的，要还回去 */
    if (FAILED(hr))
        return hr;

    g_evt = CreateEventEx(NULL, NULL, 0, EVENT_ALL_ACCESS);
    IAudioClient_SetEventDevice(client, g_evt);
    IAudioClient_GetService(client, &IID_IAudioRenderClient,
                            (void **) &g_render);
    IAudioClient_GetBufferSize(client, &g_bufFrames);
    return IAudioClient_Start(client);
}
\end{lstlisting}

独占模式下取数换成 \Fn{IAudioRenderClient\_Fill}，每次交付的帧数由
Initialize 的 \Tp{Periodicity} 参数决定（同样以 100 ns 计）。共享模式的
周期长度由系统定，可以用 \Fn{IAudioClient\_GetDefaultDevicePeriod}
和 \Fn{IAudioClient\_GetStreamLatency} 问出建议周期与最小延迟——
两者都是量级参考，不是你能承诺的实际延迟。

线程优先级要交给 MMCSS：\Fn{Av{\Bk}Set{\Bk}Mm{\Bk}Thread{\Bk}Characteristics} 传入
``Pro Audio''（也可用 ``Playback'' 或 ``Capture''）拿到一个任务句柄，
退出时 \Fn{Av{\Bk}Revert{\Bk}Mm{\Bk}Thread{\Bk}Characteristics} 归还；再配合
\Fn{Av{\Bk}Set{\Bk}Mm{\Bk}Thread{\Bk}Priority} 明确要求低延迟档。这不是装饰：
没注册 MMCSS 的音频线程在系统繁忙时会被普通负载挤掉，表现就是偶发
一次断裂——共享模式下往往是一段静音（提交时给
\Tp{AUDCLNT\_{\Bk}BUFFERFLAGS\_{\Bk}SILENT}），独占模式下则直接是丢周期。

\begin{note}[title={ASIO 与 WASAPI 独占不是同一件事}]
ASIO 是厂商定义的驱动接口契约，宿主进程直接和用户态驱动交换数据，
本来就绕开系统混音器；WASAPI 独占则是操作系统提供的直通通道，仍然走
系统音频服务。在 Windows 上做低延迟，常见次序是：先试 WASAPI 独占，
不满足再要求用户装 ASIO 驱动，或者用桥接层把 ASIO 请求转成 WASAPI。
\end{note}

\section{ASIO：把驱动写进宿主进程}
\label{sec:ch03-asio}

ASIO 的特别之处是没有守护进程、也没有自己的内核线程：宿主
（数字音频工作站）与驱动之间是一组双向约定的函数指针，数据搬运发生在
宿主进程的地址空间里。核心契约有四点：

\begin{itemize}
  \item 缓冲协商：\Fn{ASIOGet{\Bk}Buffer{\Bk}Size} 返回最小值、最大值、推荐值与
        粒度，宿主在这个区间里按 2 的幂挑一个；声道与格式由
        \Tp{ASIOBuffer{\Bk}Info} 数组描述，\Fn{ASIOCreate{\Bk}Buffers} 一次建好，
        \Fn{ASIOStart} 之后才开始流动。
  \item 乒乓双缓冲：驱动与宿主交替读写两块缓冲，每次切换通过回调通知，
        所以一次交付的单位就是缓冲时长——这也是 ASIO 延迟能做低的
        结构性原因。
  \item 回调表：注册时交一个 \Tp{ASIOCallbacks}，其中的
        \Tp{bufferSwitchTimeInfo}（或更早的 \Tp{bufferSwitch}）
        在每次切换时被调用，附带 \Tp{ASIOTime} 与当前缓冲区索引；
        采样率变化走 \Tp{sampleRateChanged}，其余能力协商走
        \Tp{asioMessage} 与 \Fn{ASIOFuture}，例如请求 64 bit
        原生格式或查询设备能力。
  \item 延迟上报：\Fn{ASIOGetLatency} 分别给出输入与输出延迟，宿主用它
        做录音补偿；少了这一步，即使缓冲很小，返听依然是错位的。
\end{itemize}

代价是 ASIO 不做任何``友好''的事：不重采样、不混音、不降级，设备被占用
就直接失败，参数不对也不解释。它在生产环境里的定位因此很清楚——
专业音频，而不是通用播放。

\section{Apple：CoreAudio HAL 与 AudioUnit 的渲染回调}
\label{sec:ch03-coreaudio}

Apple 侧分两层。底层 CoreAudio 的硬件抽象层用一套对象属性接口描述
设备：\Fn{AudioObjectGetPropertyData} 以 \Tp{kAudioDeviceProperty}
为前缀的属性名寻址，读写缓冲帧数、标称采样率与通道数；数据以
\Tp{AudioBufferList} 传递，它是``若干段等长缓冲''的数组，比 ALSA
的交织约定更宽。

上层 AudioUnit 才是应用真正接的地方。输出单元在桌面上是默认输出
子类型，移动端是 RemoteIO 子类型；应用注册一个渲染回调，
系统的高优先级音频线程按周期调用它，让你把 \Tp{inNumberFrames}
帧写进 \Tp{ioData}。代码~\ref{lst:ch03-coreaudio} 是这个契约的最小实现。

\begin{lstlisting}[style=cppstyle,
  caption={AudioUnit 渲染回调：按系统给的缓冲列表就地填浮点帧},
  label={lst:ch03-coreaudio}]
#include <AudioToolbox/AudioToolbox.h>
#include <algorithm>

class AudioEngine;      /* 内部是一条预分配的环形队列 */

/* 被系统的高优先级线程调用：不分配、不阻塞、不打日志。 */
static OSStatus OnRender(void *refCon,
                         AudioUnitRenderActionFlags *ioActionFlags,
                         const AudioTimeStamp *inTimeStamp,
                         UInt32 inBusNumber,
                         UInt32 inNumberFrames,
                         AudioBufferList *ioData)
{
    auto *self = static_cast<AudioEngine *>(refCon);

    for (UInt32 b = 0; b < ioData->mNumberBuffers; ++b) {
        AudioBuffer &buf = ioData->mBuffers[b];
        auto *dst = static_cast<Float32 *>(buf.mData);
        size_t want = buf.mDataByteSize / sizeof(Float32);
        size_t n = std::min(want, static_cast<size_t>(inNumberFrames));

        self->pullFrames(dst, n);    /* 只从预分配队列取 n 帧 */
        buf.mDataByteSize = static_cast<UInt32>(n * sizeof(Float32));
    }
    return noErr;                    /* 出错只能返回状态码，不能抛异常 */
}

/* 注册回调：target 与 proc 成对给出，bus 0 是输出总线。 */
static OSStatus InstallCallback(AudioUnit unit, AudioEngine *engine)
{
    AURenderCallbackStruct cb{};
    cb.inputProc = &OnRender;
    cb.inputProcRefCon = engine;
    return AudioUnitSetProperty(unit, kAudioUnitProperty_SetRenderCallback,
                                kAudioUnitScope_Input, 0, &cb, sizeof(cb));
}
\end{lstlisting}

三点约定值得单独记住。其一，回调参数里的 \Tp{ioData} 是系统提供的
缓冲，你写它而不是换成自己的指针，所以一次渲染的分配成本必须为零。
其二，AudioUnit 的内部格式是 32 bit 浮点小端，描述它用的是
\Tp{AudioStreamBasicDescription}：格式标志给线性 PCM 加浮点位，
交织还是平面由声道数和一个``非交织''标志决定——此时
\Tp{AudioBufferList} 的``一个声道一段'' 才真正派上用场。
其三，采集与播放可以在同一个回调里完成：在输出总线的回调里再调
\Fn{AudioUnitRender} 向输入总线取数据，这是做即时返听与回声消除的
标准姿势，两端的延迟就是第~\ref{ch:ch02-clock-buffer} 章那两级缓冲之和。

\section{JACK：process 回调的硬实时约定}
\label{sec:ch03-jack}

JACK 的客户端 API 是纯 \CL{}，约定比别家都直白：
\Fn{jack\_set\_process\_callback} 注册的那个函数跑在实时线程上，
签名是``一个帧数加一个用户指针''，返回非零就等于请求断开。

\begin{itemize}
  \item 端口与缓冲：\Fn{jack\_port\_register} 建端口时给的缓冲长度只是
        提示，实际数据长度每次由 \Fn{jack\_port\_get\_buffer}
        按本次 \Tp{nframes} 返回。音频端口的采样类型是 32 bit 浮点，
        每端口一个声道、交织无意义，所以回调里通常只做乘加，
        不做格式转换。
  \item 帧数会变小：重新连线之后、或者服务器改了缓冲大小之后，
        \Tp{nframes} 就换了量级（缓冲大小变化另有通知回调
        \Fn{jack\_set\allowbreak\_buffer\_size\allowbreak\_callback}）。
        回调里必须按参数算，绝不能用启动时缓存的值，否则不是越界就是断音。
  \item 禁止清单：不分配、不睡眠、不做可能阻塞的系统调用、不打日志、
        不拿可能与普通线程对争的锁。要跨线程传状态就用无锁队列，
        或者把决策挪到 \Fn{jack\_cycle\_timer} 注册的函数里——它在每个
        周期末尾被调用，仍在实时线程内，但已经不在音频回调里。
  \item 拓扑变化不在回调里处理：连线、断开与客户端注册都走图回调，
        音频回调只读已经定形的端口指针。
\end{itemize}

这份约定正是 PipeWire 走的路：兼容层把 JACK 客户端接进同一张图的调度里，
让多个客户端共享一个实时线程，比经典 JACK``一个进程一个线程''
更省上下文切换。代价是你的回调被并入别人的周期，最坏情况要等一个
quantum。

\section{六个栈的横向对比与选型}
\label{sec:ch03-compare}

表~\ref{tab:ch03-stacks} 按五个维度横向比较。``可达延迟'' 一列指稳态下
单向缓冲的量级，实际取决于设备、采样率与是否独占，不要当成承诺；
``可移植性'' 说的是同一套代码能带走的程度。

\begin{longtable}{>{\raggedright\arraybackslash}p{0.14\textwidth}
                  >{\raggedright\arraybackslash}p{0.15\textwidth}
                  >{\raggedright\arraybackslash}p{0.15\textwidth}
                  >{\raggedright\arraybackslash}p{0.12\textwidth}
                  >{\raggedright\arraybackslash}p{0.11\textwidth}
                  >{\raggedright\arraybackslash}p{0.11\textwidth}}
\caption{六套音频栈的 C API 形态、缓冲模型与延迟量级对比}\label{tab:ch03-stacks} \\
\toprule
\textbf{栈} & \textbf{C API 形态} & \textbf{缓冲模型} & \textbf{可达延迟} & \textbf{实时优先级} & \textbf{可移植性} \\
\midrule
\endfirsthead
\toprule
\textbf{栈} & \textbf{C API 形态} & \textbf{缓冲模型} & \textbf{可达延迟} & \textbf{实时优先级} & \textbf{可移植性} \\
\midrule
\endhead
\midrule
\bottomrule
\endlastfoot
ALSA & 函数族加句柄，状态机显式 & 驱动内环形区，分 period 与 buffer，可选内存映射 & 数毫秒级 & 需自行配置线程与调度 & 仅 Linux，接口最稳定 \\
PipeWire & 回调表加共享内存缓冲 & 图上按 quantum 调度，缓冲按引用传递 & 十毫秒级，可下调 & 服务器统一实时线程 & 兼顾两种兼容接口 \\
WASAPI & COM 接口，按 \CL{} 风格调用 & 系统混音器或独占设备环 & 独占可到数毫秒 & MMCSS 类调度 & 仅 Windows \\
ASIO & 宿主与驱动双向回调表 & 乒乓双缓冲，长度可协商 & 数毫秒级 & 由驱动实现决定 & 跨系统与平台，需装驱动 \\
CoreAudio & 属性读写加单元回调 & 系统给的缓冲列表，按帧就地填 & 十毫秒级 & 系统实时线程 & 仅 Apple 平台 \\
JACK & 纯 \CL{} 回调，进程内直连 & 服务器分配的统一周期缓冲 & 十毫秒级 & 实时线程，约定最严 & 跨平台，制作圈常用 \\
\end{longtable}

代码上的结论很朴素：把``取一帧、填一帧''的逻辑抽成一个与栈无关的函数，
每个栈只留一层薄适配器；状态机、恢复与重协商各平台自己实现，因为
它们的状态名和触发条件都不一样。线程与优先级见
第~\ref{ch:ch12-realtime-thread} 章，框架对接见
第~\ref{ch:ch13-framework-integration} 章。

\begin{guideline}[title={选栈的三条判断}]
\begin{enumerate}
  \item 先问``要不要和别人共用设备''。共用就得接受混音层（ALSA 的
        默认路由、PipeWire、WASAPI 共享、CoreAudio），延迟预算按
        十毫秒量级起步；不共用才谈独占或 ASIO。
  \item 再问``失败模式能否接受''。独占与 ASIO 在设备被占用时直接失败，
        共享模式永远能出声但延迟不可控。播放类应用通常两级都做：
        先试独占，失败退回共享，并把这件事显示给用户。
  \item 最后问``代码能带走多少''。需要跨平台就别指望单一栈：把接口
        收窄到``打开、协商、回调、恢复、关闭''五个动作，每个平台一份
        实现，公共部分只处理采样格式与缓冲引用。
        各 API 的签名速查见第~\ref{ap:apB-api-quickref} 附录。
\end{enumerate}
\end{guideline}

% ============================================================
% ---- frag_ch04.tex ----
\chapter{视频采集与显示路径}
\label{ch:ch04-capture-display}

\DomVid{} 管线上游是摄像头与采集卡，下游是屏幕。这两段是全册里平台差异最
集中的地方：同一路 1080p30 信号，在 Linux 上是一个 \Tp{/dev/video0} 节点，
在 Windows 上是一个 Media Foundation 激活对象，在转播车里是一段 SDI 同轴
电缆上的裸帧；显示侧能不能零拷贝、多久出画面，决定权在显示驱动与窗口系统，
而不在用户态库。本章按状态机的粒度走通每条采集与显示路径，最后把从曝光到
编码的延迟拆成可以逐项压缩的预算。

\section{V4L2：Linux 采集接口的状态机}

V4L2 把摄像头暴露成设备节点，全部操作靠 \Fn{ioctl} 推进，状态机的次序是
固定的：

\begin{itemize}
\item \Tp{VIDIOC\_QUERYCAP}：读 \Tp{v4l2\_capability\Bk}，确认
      \Tp{V4L2\_CAP\_VIDEO\_CAPTURE}（或其 \Tp{\_MPLANE} 变体）能力位，
      顺带拿到驱动名与总线信息用于多设备识别；
\item \Tp{VIDIOC\_G\_FMT} / \Tp{VIDIOC\_S\_FMT}：读或写分辨率、帧率与
      \Tp{pixelformat}（fourcc）。驱动有权就地改写你提交的参数，设置后
      必须从同一个结构体读回实际谈定的格式；\Tp{VIDIOC\_TRY\_FMT} 可以
      在不落地的前提下先试探结果；
\item \Tp{VIDIOC\_{\Bk}REQBUFS}：申请缓冲队列。内存类型在
      \Tp{V4L2\_{\Bk}MEMORY\_{\Bk}MMAP} 与 \Tp{V4L2\_{\Bk}MEMORY\_{\Bk}DMABUF} 之间选择，
      驱动会把请求数向下收敛到它允许的范围；
\item \Tp{VIDIOC\_{\Bk}QUERYBUF}：逐块取回缓冲区尺寸与映射偏移，\Fn{mmap}
      进进程后用 \Tp{VIDIOC\_QBUF} 全部挂到入账队列；
\item \Tp{VIDIOC\_{\Bk}STREAMON} 起流，进入 ``DQBUF $\to$ 消费 $\to$ QBUF''
      循环，收尾用 \Tp{VIDIOC\_{\Bk}STREAMOFF}。
\end{itemize}

队列纪律只有一条：每个 \Tp{VIDIOC\_DQBUF} 取走的 buffer 都必须归还
\Tp{VIDIOC\_QBUF}。入账队列一旦枯竭，帧率不是下降而是断流，且症状常被
误判成相机故障。清单 \ref{lst:ch04-v4l2buf} 给出从打开节点到取一帧的最小
骨架，错误处理按位省略。

\begin{lstlisting}[style=cstyle,caption={V4L2 采集：从 QUERYCAP 到 DQBUF 的缓冲队列骨架（MMAP 路径）},label={lst:ch04-v4l2buf}]
/* V4L2 采集状态机骨架：省略错误处理，正式代码每步都要查返回值 */
#include <errno.h>
#include <fcntl.h>
#include <linux/videodev2.h>
#include <sys/ioctl.h>
#include <sys/mman.h>
#include <unistd.h>

enum { NBUFS = 4 };                       /* 队列深度：2~4 帧 */
static void  *maps[NBUFS];
static size_t lens[NBUFS];

int v4l2_capture_open(const char *dev)
{
    int fd = open(dev, O_RDWR | O_NONBLOCK, 0);
    struct v4l2_capability cap;           /* 状态机第一步 */
    ioctl(fd, VIDIOC_QUERYCAP, &cap);
    if (!(cap.capabilities & V4L2_CAP_VIDEO_CAPTURE))
        return -1;                        /* 不是采集节点 */

    struct v4l2_format fmt;               /* 协商格式 */
    memset(&fmt, 0, sizeof(fmt));
    fmt.type = V4L2_BUF_TYPE_VIDEO_CAPTURE;
    fmt.fmt.pix.width       = 1280;
    fmt.fmt.pix.height      = 720;
    fmt.fmt.pix.pixelformat = V4L2_PIX_FMT_YUYV;
    fmt.fmt.pix.field       = V4L2_FIELD_NONE;
    ioctl(fd, VIDIOC_S_FMT, &fmt);        /* 读回 fmt 里的实际值！ */

    struct v4l2_requestbuffers req;       /* 建队列 */
    memset(&req, 0, sizeof(req));
    req.count  = NBUFS;
    req.type   = V4L2_BUF_TYPE_VIDEO_CAPTURE;
    req.memory = V4L2_MEMORY_MMAP;
    ioctl(fd, VIDIOC_REQBUFS, &req);

    for (unsigned i = 0; i < req.count; i++) {
        struct v4l2_buffer buf;           /* 逐块问尺寸再映射 */
        memset(&buf, 0, sizeof(buf));
        buf.type   = V4L2_BUF_TYPE_VIDEO_CAPTURE;
        buf.memory = V4L2_MEMORY_MMAP;
        buf.index  = i;
        ioctl(fd, VIDIOC_QUERYBUF, &buf);
        maps[i] = mmap(NULL, buf.length, PROT_READ | PROT_WRITE,
                       MAP_SHARED, fd, buf.m.offset);
        lens[i] = buf.length;
        ioctl(fd, VIDIOC_QBUF, &buf);     /* 空 buffer 先挂进去 */
    }

    enum v4l2_buf_type type = V4L2_BUF_TYPE_VIDEO_CAPTURE;
    ioctl(fd, VIDIOC_STREAMON, &type);    /* 起流 */
    return fd;
}

int v4l2_grab(int fd, unsigned *idx, size_t *used)
{
    struct v4l2_buffer buf;
    memset(&buf, 0, sizeof(buf));
    buf.type   = V4L2_BUF_TYPE_VIDEO_CAPTURE;
    buf.memory = V4L2_MEMORY_MMAP;
    if (ioctl(fd, VIDIOC_DQBUF, &buf) < 0)
        return (errno == EAGAIN) ? -2 : -1;
    *idx  = buf.index;   /* 数据在 maps[idx] */
    *used = buf.bytesused;
    ioctl(fd, VIDIOC_QBUF, &buf);         /* 用完立刻归还 */
    return 0;
}
\end{lstlisting}

MMAP 路径的 ``零拷贝'' 指的是内核 DMA 缓冲直接映射进用户地址页表，省掉驱动
到应用的一次搬运；若要帧字节完全不经过 CPU，走 DMABUF：采集端用
\Tp{VIDIOC\_EXPBUF} 把某块 buffer 导出为 fd，消费端各自 import——V4L2 M2M
硬件编码器在 \Tp{VIDIOC\_REQBUFS} 时声明
\Tp{V4L2\_MEMORY\Bk\_\_DMABUF} 并把 capture 端的 fd 填进
\Tp{v4l2\_buffer.m.fd}；DRM 侧用 \Fn{drmPrimeFDToHandle} 把同一个 fd 变成
framebuffer；OpenGL 侧经 \Tp{EGL\_EXT\_image\_dma\_buffer} 绑成纹理。这条
链能否端到端打通，取决于各驱动是否愿意导出，嵌入式 SoC 支持度高，PCIe 采集
卡则普遍不提供。

\begin{warning}[title={S\_FMT 谈成的格式未必是你请求的格式}]
\Tp{VIDIOC\_S\_FMT} 调用成功不代表驱动接受了参数：很多 UVC 相机在请求的
分辨率没有原生模式时，会静默降到最近的档位再返回。所有后续计算（stride、
buffer 尺寸、帧率上限）必须以读回的结构体为准，而不是以你写进去的为准。
\end{warning}

\section{命令行枚举与直通采集：v4l2-ctl 与 FFmpeg}

排查采集问题的第一步永远是确认设备 ``自称支持什么''。\Lib{v4l2-ctl} 直接
转储驱动能力，FFmpeg 的 \Opt{f} \Tp{v4l2} 输入设备则给出管线侧视角；两者
对不上时，通常是用户态版本或内核驱动太旧。清单 \ref{lst:ch04-v4l2ctl} 是
一组常用的枚举与采集命令。

\begin{lstlisting}[style=shellstyle,caption={用 v4l2-ctl 枚举能力，用 ffmpeg/ffplay 直通采集},label={lst:ch04-v4l2ctl}]
# 枚举每个像素格式下支持的分辨率与帧率（先看菜单再点菜）
v4l2-ctl -d /dev/video0 --list-formats-ext

# 一次性打印全部能力、当前格式与控件
v4l2-ctl -d /dev/video0 --all

# 让 FFmpeg 报告设备端格式清单（v4l2 输入设备实现）
ffmpeg -f v4l2 -list_formats all -i /dev/video0

# 指定输入编码与尺寸采 10 秒：MJPG 源直接落盘走重编码
ffmpeg -f v4l2 -input_format mjpeg -video_size 1920x1080 \
       -framerate 30 -i /dev/video0 -t 10 cap.mp4

# 实时预览：显式 -f v4l2 跳过探测，避免被音频设备截胡
ffplay -f v4l2 -video_size 1280x720 -framerate 30 -i /dev/video0

# 只测采集与编码的真实帧率，丢弃输出
ffmpeg -f v4l2 -video_size 1280x720 -i /dev/video0 \
       -t 30 -c:v libx264 -preset ultrafast -f null -
\end{lstlisting}

\section{Media Foundation 采集与 UVC 格式协商}

Windows 侧的枚举入口是 ``激活对象''。\Fn{MFEnumDeviceSources} 接受一组
属性——键 \Tp{MF\_DEVSOURCE\_ATTRIBUTE\Bk\_SOURCE\_TYPE}，值
\Tp{...VIDCAP\_GUID\Bk}——返回 \Tp{IMFActivate} 数组；从激活对象属性里
读显示名与符号链接即可区分多台相机，\Fn{ActivateObject} 之后拿到
\Tp{IMFMediaSource}。相机 App、会议软件与直播工具的采集图都建在这套对象
之上；\Tp{IMFCaptureEngine} 是官方提供的更高层封装，把预览、拍照与录制
收进同一个引擎回调。

格式协商沿 \Tp{IMFPresentation{\Bk}Descriptor} 的 0 号流取
\Tp{IMFStream{\Bk}Descriptor}，再问出 \Tp{IMFMedia{\Bk}Type{\Bk}Handler}，用
\Fn{GetCount} 与 \Fn{Media{\Bk}Type{\Bk}By{\Bk}Index} 枚举每个媒体类型（分辨率、帧率、
media subtype 三元组），选定后 \Fn{SetType} 提交。subtype 常见的是
\Tp{MFVideo{\Bk}Format\_{\Bk}YUY2}、\Tp{MFVideo{\Bk}Format\_{\Bk}MJPG}、
\Tp{MFVideo{\Bk}Format\_{\Bk}NV12} 与桌面源的 \Tp{MFVideo{\Bk}Format\_{\Bk}RGB32}。

UVC 相机的原生格式只有三类，带宽是硬约束：YUY2 与 NV12 是裸流，按像素率
计费，1080p30 的 YUY2 约需 1\,Gbps，远超 USB 2.0 的有效带宽，所以绝大多数
相机会把 YUY2 限到 1080p 五帧以下、只让 MJPG 跑满 30 fps；NV12 常见于
UVC 1.5 之后的新设备与更高带宽的接口。MJPG 是相机内部做了运动 JPEG 编码
再由主机解码，画质上有块噪声、延迟上多一拍编解码。带宽谈不拢时系统往往不
报错，只是枚举结果里对应档位 ``恰好不存在''，或者插到不同 hub 上行为悄然
变化——排查时以 \Fn{GetCount} 的实际枚举为准。

\section{专业采集：SDI 路径与帧同步}

广电与专业视频不走 USB，走 SDI 同轴：1.5G（SMPTE ST 292-1）、3G
（ST 424）、12G（ST 213）分别承载不同档位的未压缩 10-bit 4:2:2 信号。
采集卡 SDK（Blackmagic DeckLink、AJA 等）返回的帧常见为 \Tp{v210}——
一种每 3 个 32-bit 字塞 4 个 10-bit 分量的打包格式——或 P210 之类的平面
格式；时间码从 ancillary 数据（ST 12-1）读出，与帧号绑定。

SDI 路径的卖点不是带宽大，而是确定性：帧是整帧，不裁剪、不压缩、不排队；
genlock 参考信号（black burst 或 tri-level sync，多机时基用 ST 2059 的
PTP）让多台机器在同一相位出帧；采集卡上的 ``through'' 口是硬件直通监看，
不占软件延迟。需要逐样精确的多相机采集、外场回显或和 broadcast 链路对齐时，
这条路径仍是唯一选择，代价是线缆、采集卡与同步发生器的整套基建。

\section{显示侧：DRM/KMS plane、vsync 反馈与窗口系统}

显示侧的 ``采集'' 是反过来的：把帧交给扫描引擎，并问清它何时亮。

\textbf{DRM/KMS。} Linux 内核显示栈由 CRTC（时序发生器）、encoder、
connector 和 plane 组成；plane 分 primary、overlay、cursor 三类，硬件
overlay 允许把 YUV 帧直接在扫描时缩放与色彩转换，避开 GPU。传统接口用
\Fn{drmModeAddFB} 建 framebuffer，只认 RGB 类 fourcc；\Fn{drmModeAddFB2}
接受任意 fourcc，每平面给 handle、pitch、offset 三个数组，于是 DMABUF
导入的 handle 可以直接成为可扫描的 FB（清单 \ref{lst:ch04-v4l2buf} 导出的
fd 在这里闭环）。提交用 atomic：给 plane 的 \Tp{FB\_ID}、\Tp{CRTC\_ID}、
\Tp{SRC\_W} 等属性填值后 \Fn{drmModeAtomicCommit}，先用
\Tp{DRM\_MODE\_ATOMIC\_TEST\_ONLY} 验证模式不花屏。vsync 事件经
\Fn{drmWaitVBlank}（legacy）或 atomic commit 携带的事件对象回报，时间戳
取自内核，精度足以做管线标定。

\textbf{Wayland。} 合成器代管上屏时机，客户端需要 ``到底何时亮'' 的反馈。
presentation 扩展提供 \Tp{wpresentation} 与每帧一个
\Tp{wpresentation\_feedback} 对象：\Fn{presented} 事件回送同步标记、
完成时刻的时钟读数（客户端先声明要 monotonic 还是 realtime）、刷新周期与
序号，并以 flags 区分 ``本帧被复制合成'' 与 ``被更新帧取代''。帧调度器
据此把提交提前量压到刚好赶上上一个 vsync。fast path 下合成器把 buffer
以 DMABUF 直接交给 KMS，无拷贝；一旦需要合成（多窗口、缩放、受保护内容）
就会退化为 GPU 搬运。

\textbf{X11。} 传统路径里 swap 之后画面何时亮起，GLX 不给答案。
\Tp{OML\_sync\_control} 扩展补上了三块表：\Fn{glXGetSyncValuesOML} 读
ust（微秒时钟）、msc（vblank 计数）、sbc（swap 计数），
\Fn{glXGetMscRateOML} 报刷新率，\Fn{glXSwapBuffersMscOML} 可以把 swap
排程到目标 msc。它与 SGI\_video\_sync 一起构成 Mesa 的
\Tp{video\_swap\_chain} 同步算法的基础。但窗口若被合成器接管，帧会多走
一跳，反馈的 msc 就不再等于真实上屏——这也是低延迟显示近年整体迁向
Wayland 的原因之一。

\section{从曝光到编码：延迟预算分解}

端到端 ``玻璃到玻璃'' 延迟要逐项记账才能压缩，典型量级如下：

\begin{itemize}
\item 曝光：快门周期的下限就是 1 帧（30fps 约 33\,ms），高光低噪场景无法
      靠软件绕开；
\item 传感器读出与相机内 ISP：行读出加降噪/锐化，数 ms 到一帧；
\item 相机内部缓冲：UVC MJPG 相机常缓存 1--2 帧凑带宽，这是 USB 摄像头
      ``摸起来慢'' 的主因；
\item 传输与主机栈：USB 等时传输约 1 帧周期对齐，采集卡 DMA 接近零；
\item 入账队列深度：DQBUF 侧每多留一拍缓冲，端到端就多一拍延迟；
\item 色彩转换与缩放：CPU 转换 1080p 约 1--3\,ms，硬件 overlay 可摊平；
\item 编码器前瞻：lookahead 与 B 帧重排会引入多帧延迟（第
      \ref{ch:ch07-rate-control} 章细账）；
\item 网络与抖动缓冲：实时链路里最大的一块，见第
      \ref{ch:ch11-streaming} 部与第 \ref{ch:ch12-realtime-thread} 章。
\end{itemize}

\begin{longtable}{>{\raggedright\arraybackslash}p{0.24\textwidth} >{\raggedright\arraybackslash}p{0.17\textwidth} >{\raggedright\arraybackslash}p{0.26\textwidth} >{\raggedright\arraybackslash}p{0.13\textwidth}}
\caption{常见采集与显示路径的延迟量级、零拷贝能力与平台分布}\label{tab:ch04-paths} \\
\toprule
路径 & 典型延迟 & 零拷贝能力 & 平台 \\
\midrule
\endfirsthead
\toprule
路径 & 典型延迟 & 零拷贝能力 & 平台 \\
\midrule
\endhead
\bottomrule
\endfoot
V4L2 MMAP（UVC 相机） & 1--3 帧 & 映射即共享内核 DMA 页，免用户态拷贝；格式转换仍过 CPU & Linux \\
V4L2 MPLANE + DMABUF & 同上再省搬运 & \Fn{VIDIOC\_EXPBUF} 的 fd 直通编码器与显示，端到端零拷贝 & Linux（SoC 驱动为主） \\
MF 采集 UVC（MJPG） & 2--4 帧 & MF sample 是 CPU 缓冲，进 GPU 需一次上传 & Windows \\
SDI 采集卡（v210） & 约 1 帧 & SDK 给设备侧缓冲，多数支持显存互操作 & 跨平台专业 \\
DRM/KMS atomic + DMABUF FB & 等 vsync，1--2 帧 & prime fd 组成 FB，硬件 plane 直接扫描 & Linux \\
Wayland + presentation & 1--2 帧 & fast path 零拷贝，需合成时退化为 GPU 复制 & Linux 新栈 \\
X11 + GLX（OML 同步） & 2--3 帧 & 依赖驱动 direct scanout，历史路径多一跳 & 跨平台遗留 \\
\end{longtable}

压缩延迟的操作顺序也从这张表读出：先砍相机侧（选裸流格式、少留队列、
关相机内缓冲），再砍搬运侧（让帧以 DMABUF 或硬件 plane 的身份走完全程），
最后才轮到编码器前瞻与网络缓冲——后两者的收益每 10\,ms 都要拿画质或
鲁棒性来换。

\begin{bestpractice}[title={先量再调：用事件时间戳给采集链路记账}]
改任何参数之前，先把每个环节的 ``帧到达时刻'' 打点：V4L2 的
\Tp{v4l2\_{\Bk}buffer.{\Bk}timestamp} 给采集侧基准，DRM 事件与 Wayland
\Fn{presented} 给显示侧基准。有了两端时间戳，曝光到上屏的预算才能分解；
否则所有 ``感觉变快了'' 的调优都无法复现，也无法在回归时守住。
\end{bestpractice}

% ============================================================
% ---- frag_ch05.tex ----
\chapter{像素与平面格式：YUV、色度次采样与对齐}
\label{ch:ch05-pixel-format}

像素格式是 \DomVid{} 管线里最小的货币单位：采集端产出它，滤波器之间流转
它，编码器只接受它允许的子集，显示端再把它展开回光线。看似只是 ``一帧
图像怎么排字节''，实际上四套独立约定——色彩矩阵、量化范围、次采样与采样
点位置、平面布局与对齐——交织出了全册最容易 ``颜色不对'' 的一段。本章把
四套约定分别讲清，再给出跨库命名的对照与拷贝代码的正确写法。

\section{从 RGB 到 Y\hspace{0.1em}'CbCr：带宽与兼容性}

RGB 是显示域的坐标：三个通道各自逼近一种原色的刺激量。但传输域不必照搬：
人眼对亮度细节的分辨力远高于色度细节，于是把 RGB 线性组合成一路亮度
\(Y'\) 与两路色差 \(C_b = k_b (B - Y')\)、\(C_r = k_r (R - Y')\)，色度
通道的采样率就可以放心地砍半而不被察觉。这条路线还有历史动因：NTSC 与
PAL 彩色电视要求彩色信号能被黑白电视机直接接收——黑白机只看 \(Y'\)，
色差分量对它无意义，兼容性由构造保证。

亮度权重随 primaries 变化，三个常见矩阵是：

\begin{itemize}
\item BT.601（SD 时代）：\(Y' = 0.299R + 0.587G + 0.114B\)；
\item BT.709（HD 时代）：\(Y' = 0.2126R + 0.7152G + 0.0722B\)；
\item BT.2020（UHD 广角 primaries）：\(Y' = 0.2627R + 0.6780G + 0.0593B\)，
      HDR 的 BT.2100 沿用这组 primaries 与其亮度权重，另定 transfer
      特性（PQ 与 HLG 两条曲线）。
\end{itemize}

符号上的撇不是装饰：矩阵作用在已做传递函数编码的信号 \(R'G'B'\) 上，得到
的是 ``亮度信号''（luma）而非物理亮度（luminance）——对伽马值做线性加权
并不等于对光能量做加权，\(Y'\) 因此只是感知近似的亮度通道。任何 ``先在
伽马域平均、再解伽马'' 之类的操作在数学上都是错的，几何滤波色度可以、
滤波 \(Y'\) 的均值则要三思。

拿 709 矩阵解 601 素材（或反之）不会崩溃，只会让饱和度悄悄偏移——红色
变暗或变艳，肉眼熟悉后甚至察觉不出，这是跨源拼接项目里潜伏最深的错误。

\section{矩阵、范围与标准族：601 / 709 / 2100 的三层面具}

``BT.601 视频'' 一词至少混指三件事，FFmpeg 流层面把它们拆成独立字段：
\Tp{color\_{\Bk}primaries}（三原色与白点）、\Tp{color\_trc}（传递函数，
gamma、PQ、HLG）、\Tp{color\_{\Bk}space}（\(Y'C_bC_r\) 换算用矩阵，
\Tp{AVCOL\_{\Bk}SPC\_{\Bk}SMPTE170M} / \Tp{AVCOL\_{\Bk}SPC\_{\Bk}BT709} /
\Tp{AVCOL\_{\Bk}SPC\_{\Bk}BT2020\_{\Bk}NCL} 等）。转码改分辨率不改色彩语义，但字段要
原样透传；缩放时真正做矩阵运算的是 swscale（第 \ref{ch:ch09-filter-resample}
章）。

第四件事是量化范围（\Tp{color\_range}）：
\begin{itemize}
\item limited（``视频范围''，legacy 称为 video range/legit）：8-bit 下
      \(Y'\) 取 16--235，色差取 16--240（中性 128）。留出的上下沿是给
      模拟过冲的 headroom，数字时代沿用成广播默认；
\item full（``JPEG 范围''）：0--255 满量程，桌面抓取与游戏画面的默认。
\end{itemize}
10-bit 下 limited 对应 \(Y'\) 64--940、色差 64--960。范围标记错了的典型
症状：黑场发灰或暗部糊死、白色过曝——修复只需要重打 tag 不需要重新
量化；但 ``内容真实范围与 tag 相反'' 时就必须做 \((v-16)\times 255/219\)
一类的映射转手，swscale 打开时的色彩空间参数就是在矩阵与范围之间做
换算的开关。

这些字段要在一次转码里活下来，得靠搬运链条：容器层有 MP4 的 colr 盒与
Matroska 的 Colour 元素，比特流层有 H.264/H.265 的 VUI 语义，解封装与
封装各自负责映射；``容器 tag 与比特流 VUI 说法不一'' 的文件，播放器与
编辑软件可能给出两副颜色。工程上以比特流内的声明为准、容器 tag 作缓存，
发现两者矛盾就重打标签而不是重新量化。

\section{4:4:4、4:2:2、4:2:0 与采样点位置}

次采样写法 \(J:A:B\) 描述以 4 个水平亮度采样为参照的色度密度：4:4:4
每像素一组色差；4:2:2 色度水平减半、垂直全保留；4:2:0 横竖都减半（只有
前者一半的色度样本数，不是 ``四分之一亮度''）。4:2:0 是编码器与硬件的
绝对主流，4:2:2 常见于采集与专业接口，4:4:4 留给屏幕内容与中间处理。

同样是 4:2:0，色度样本 ``钉在哪'' 有两族约定。水平方向：MPEG-2 与 DV 的
\emph{cosited}（色度与同行的左亮度同位）对比 MPEG-4 系（H.26x）的
色度居中西式；垂直方向：MPEG-420（色度在两行亮度正中）对比
DV-420（与上一行同位），JPEG 则同时是 cosited 且顶部对齐。换算错
半像素，上采样后边缘出现彩色镶边、彩色细线偏移，而静态截图可能完全
看不出——要靠运动帧的边缘才能暴露。库间搬运时优先询问接口是否处理
siting：swscale 与编解码器两侧各有默认假定，V4L2 与硬件编码器一般
按 MPEG-420 居中的约定处理。

\begin{compare}[title={三档次采样的取舍}]
4:2:0 把色度样本数压到亮度的四分之一，编码器色度残差的熵编码开销同步下降，
是所有交付格式的默认；代价是色度细节损失，抠像与红蓝屏幕合成时边缘发毛。
4:2:2 保留垂直色度分辨率，采集卡、专业摄像机与调色链路用它在带宽和后期
空间之间折中；编码侧只有部分档位的硬件编码器与专业中间格式接受它。
4:4:4 无损色度，留给屏幕内容、合成中间层与母版归档，带宽是 4:2:0 的四倍。
同一管线内换档要经过重采样，位置约定（siting）随之改变，不是免费的。
\end{compare}

\section{布局三态：planar、semi-planar、packed 与 fourcc 命名网}

同一份 4:2:0 数据有三种排法家族：
\begin{itemize}
\item planar（如 FFmpeg 的 \Tp{yuv420p}）：\Tp{data[0..2]} 分别指向
      独立的 \(Y' / C_b / C_r\) 平面，各带各的 linesize；
\item semi-planar（\Tp{nv12}、\Tp{p010}）：亮度一个平面，色度把两个
      分量交错成单平面（UVUV 逐字节对），硬件最爱，因为色度只需一次
      访存；
\item packed（\Tp{uyvy422}、\Tp{v210}、\Tp{y410}）：单平面内按扫描线
      交错各分量，UYVY 每两个像素以 U、Y0、V、Y1 四字节为一组，
      v210 把 4 个 10-bit 分量塞进 3 个 32-bit 字，y410 则是 10-bit
      YUVA 4:4:4:4 打包，常见于 D3D11/VA 硬管线。
\end{itemize}

RGB 一族同样有这三种排法：\Tp{rgb24} 与 \Tp{bgra} 按像素打包三或四个
通道，10-bit 桌面内容用 \Tp{x2rgb10le}；滤镜中间层偏爱 planar 的
\Tp{gbrp}（三通道各一个平面，便于逐通道 SIMD 处理曲线与卷积）。显示
扫描输出几乎只接受 packed RGB，而编码器几乎不接受它，两头之间永远隔着
一次转换。

清单 \ref{lst:ch05-formats} 展示查询格式清单与读回文件色彩三元组的
命令行。

\begin{lstlisting}[style=shellstyle,caption={用 ffmpeg -pix\_fmts 与 mediainfo/ffprobe 读格式与色彩语义},label={lst:ch05-formats}]
# FFmpeg 认识的全部像素格式（含标志位、组件数与位深列）
ffmpeg -hide_banner -pix_fmts | less

# 只看某档：名字、别名与位数
ffmpeg -hide_banner -pix_fmts | grep -E "yuv420p|nv12|p010"

# 读文件的色彩三元组与位深（封装层标签原样取回）
ffprobe -hide_banner -v quiet -select_streams v:0 \
  -show_entries stream=pix_fmt,color_space,color_primaries,\
color_trc,color_range -of default clip.mp4

# MediaInfo 的字段式输出：色彩空间/次采样/位深一行读完
mediainfo --Inform="Video;%ColorSpace% %ChromaSubsampling% \
%BitDepth% %ColorRange%" clip.mkv
\end{lstlisting}

命名是跨栈最烦人的部分：FFmpeg 用小写助记符（\Tp{yuv420p}，经
\Fn{av\_{\Bk}get\_{\Bk}pix\_{\Bk}fmt\_{\Bk}name} 取得），一部分格式直接由 fourcc 宏定义，
且大端小端的字母序相反（如 \Tp{AV\_{\Bk}PIX\_{\Bk}FMT\_{\Bk}UYVY}）；V4L2 用
\Tp{V4L2\_{\Bk}PIX\_{\Bk}FMT\_{\Bk}} 前缀加自己的 fourcc，同一布局可能与
FFmpeg 名字不同字符（planar 4:2:0 在 V4L2 是 \Tp{V4L2\_{\Bk}PIX\_{\Bk}FMT\_{\Bk}YUV420}，
即 fourcc \Tp{YU12}，而 GStreamer 与容器侧更多写 \Tp{I420}）。同一语义在
三个生态有三种拼法，换栈时永远查表对照，不要按字面猜。表
\ref{tab:ch05-pixfmt} 给出常见格式的工程属性。

\begin{longtable}{>{\raggedright\arraybackslash}p{0.21\textwidth} >{\raggedright\arraybackslash}p{0.08\textwidth} >{\raggedright\arraybackslash}p{0.19\textwidth} >{\raggedright\arraybackslash}p{0.12\textwidth} >{\raggedright\arraybackslash}p{0.20\textwidth}}
\caption{常见像素格式的布局、位深与后端支持对照}\label{tab:ch05-pixfmt} \\
\toprule
格式（FFmpeg 名） & 平面 & 采样/位深 & 布局 & 常见编码器与硬件后端 \\
\midrule
\endfirsthead
\toprule
格式（FFmpeg 名） & 平面 & 采样/位深 & 布局 & 常见编码器与硬件后端 \\
\midrule
\endhead
\bottomrule
\endfoot
\Tp{yuv420p} & 3 & 4:2:0，8-bit & planar & x264/x265 与几乎全部硬件编码器默认 \\
\Tp{yuv420p10le} & 3 & 4:2:0，10-bit & planar & x265 10-bit、专业中间件 \\
\Tp{nv12} & 2 & 4:2:0，8-bit & semi-planar & QSV、NVENC、VAAPI、AMF 通用输入 \\
\Tp{p010le} & 2 & 4:2:0，10-bit & semi-planar & Intel/高通硬编 10-bit 档 \\
\Tp{uyvy422} & 1 & 4:2:2，8-bit & packed & UVC 相机原生、SDI 采集卡 \\
\Tp{v210} & 1 & 4:2:2，10-bit & packed & 广播采集卡与剪辑中间格式 \\
\Tp{y410} & 1 & 4:4:4:4，10-bit & packed & D3D11/VA 硬解与合成管线 \\
\Tp{bgra} & 1 & RGB，8-bit & packed & 桌面抓取、显示扫描输出 \\
\end{longtable}

\section{linesize 与对齐：为什么逐行 memcpy 必须写成循环}

\Tp{AVFrame->linesize[p]} 是平面 \(p\) 一行占用的字节数，它 ``通常大于
可见像素占的字节数''，原因有三类：向量化代码要求行起点按 32/64 字节对齐
以走满宽载入；硬件 surface（D3D11、V4L2 M2M、视频解码器输出）常按 128
或 256 字节向上取整 pitch；次采样平面在半宽分辨率时还要把奇数宽度进位
（宽 17 的 4:2:0 色度行是 \((17+1)/2 = 9\) 字节可见，再加 padding）。

于是 ``整平面一次 memcpy'' 只在极窄的条件下成立：源与目标的
\Tp{linesize} 恰好都等于可见行字节数。工程上正确的写法是逐行循环，行数
用高度、每行字节用可见宽度乘单位格字节数，指针按各自的 linesize 前进。
清单 \ref{lst:ch05-copy} 以 8-bit planar 为例；NV12 的色度是单平面、
每行 \(cw \times 2\) 字节、行高仍是 \(ch\)，packed 格式则只有一个平面
但行字节为整宏像素数。16-bit 格式（\Tp{yuv420p10le}、\Tp{p010le}）
每样本占 2 字节，行字节数再乘 2。

帧池是第二个陷阱：\Fn{av\_buffer\_ref} 回收缓冲时，同分辨率但步长不同的
两块内存可能混进同一条管线——一块由 \Fn{av\_frame\_get\_buffer} 按默认
对齐产出，另一块由硬件接口包装进来。把 (width, height, linesize) 当作
一组键值缓存的代码，在切换分辨率后读到错误步长，就会产出整体错位的帧。
每帧都从帧自身的 \Tp{linesize} 读取，若步长变化就重建队列，不要试图
替它心算。

同一条约定贯穿缩放接口：\Fn{sws\_scale} 以 ``平面指针数组加步长数组''
为输入，按切片（起始行与行数）推进处理，但它从不怀疑你给的步长——
把可见宽度冒充 \Tp{linesize} 传进去，输出就会整体倾斜。手工铺帧没有
收益，交给 \Fn{av\_frame\_get\_buffer} 永远是更短的正确路径。

\begin{lstlisting}[style=cstyle,caption={按 AVFrame->linesize 逐平面、逐行拷贝：只有步长恰好等于可见行宽时才允许合并成一次 memcpy},label={lst:ch05-copy}]
#include <libavutil/frame.h>
#include <string.h>

/* 逐行搬运一个平面：src_ls/dst_ls 是各自的 linesize */
static void copy_rows(uint8_t *dst, int dst_ls,
                      const uint8_t *src, int src_ls,
                      int row_bytes, int rows)
{
    if (row_bytes == src_ls && src_ls == dst_ls) {
        /* 无 padding 才允许合并成一次 memcpy */
        memcpy(dst, src, (size_t)row_bytes * rows);
        return;
    }
    for (int y = 0; y < rows; y++) {
        memcpy(dst + (size_t)y * dst_ls,
               src + (size_t)y * src_ls, row_bytes);
    }
}

/* 8-bit planar yuv420p：3 平面，色度横竖都减半 */
void copy_yuv420p(AVFrame *d, const AVFrame *s)
{
    int cw = (s->width  + 1) >> 1;      /* 色度可见宽度 */
    int ch = (s->height + 1) >> 1;      /* 色度可见行数 */

    copy_rows(d->data[0], d->linesize[0],
              s->data[0], s->linesize[0], s->width, s->height);
    copy_rows(d->data[1], d->linesize[1],
              s->data[1], s->linesize[1], cw, ch);
    copy_rows(d->data[2], d->linesize[2],
              s->data[2], s->linesize[2], cw, ch);
}
\end{lstlisting}

\begin{note}[title={手写 linesize 是最常见的绿条来源}]
自己拼 \Tp{AVFrame} 时把 \Tp{linesize} 填成 ``宽度''而不是 ``分配的行
步长''，解码或缩放输出就会按短步长写、按长步长读（或反之），表现为画面
斜切或绿色竖条。让 \Fn{av\_frame\_get\_buffer} 或 \Fn{av\_image\_alloc}
替你算步长；确实要手工铺指针时，用 \Fn{av\_image\_fill\_linesizes} 与
\Fn{av\_image\_fill\_pointers}，它们保证平面尺寸与对齐约定一致。
\end{note}

\section{把四套约定钉在帧上}

一帧 ``正确'' 的表示需要四个维度同时落位：\(Y'C_bC_r\) 与 RGB 的选择
（传输域与显示域）、矩阵与范围（601/709/2020 乘 full/limited）、次采样
与 siting（4:4:4/4:2:2/4:2:0 乘两族对齐）、布局与步长（平面划分乘
linesize）。任何一格标错都不会报错，只会让颜色 ``差一点''。

这四个维度都可以实证核对：用 \Tp{rawvideo} 输出一帧再按平面做十六进制
dump，平面数、步长与样本次序一目了然；颜色不对时先拿一段已知语义的
彩条片做基准，逐字段比较矩阵、范围与 siting 哪个 tag 错了。自我断言的
成本低，远比发布后靠用户截图返工便宜。

\begin{guideline}[title={格式选择的最小充分集}]
项目里维护一份显式的 ``色彩四元组''（pix\_fmt + primaries/trc +
matrix + range + siting 默认值），采集、滤波、编码三端各自声明能输入与
输出的组合，落在交集外的帧一律在边界处显式转换并留日志。转换本身不花钱，
``隐式转换'' 才是返工账单的来源。
\end{guideline}

% ============================================================
% ---- frag_ch06.tex ----
\part{编解码内核：从帧到码流}

\chapter{编码器与解码器的 C API 形态}
\label{ch:ch06-codec-api}

前面几章处理的都是``还没压缩''的数据。本章起进入压缩内核，而压缩内核的第一
件事不是算法，是接口形态：一帧裸数据怎么交出去、换回来的字节什么时候才算属于
你、为什么五个主流编码库给出五种互不相同的答案。参数怎么填留给第
\ref{ch:ch07-rate-control} 章，解码与硬件通路见第 \ref{ch:ch08-decode-hw} 章。

\section{libavcodec 的两阶段收发模型}

旧接口 \Fn{avcodec\_encode\_video2} 把输入帧与输出包塞进同一次调用，返回值同时
承担``编码了几字节''与``还要更多输入吗''两种语义。这与真实编码器的状态机冲突：
B 帧重排与 lookahead 让输入输出的比例不是 $1:1$，一帧可能吐出零个包；带预滚
延迟的音频编码器要攒好几帧才出货。3.1 版把它拆成 send 与 receive 两阶段，4.0
版删掉旧接口，收益是每个调用只回答一个问题：送帧 \Fn{avcodec\_send\_frame} 只管
输入是否被接受，收包 \Fn{avcodec\_receive\_packet} 只管有没有输出，两侧的``暂时
没有''都用 \Tp{AVERROR(EAGAIN)} 表达，取值矩阵见表 \ref{tab:ch06-ret}。一次
receive 至多一个包，而一个包可能含多个 NAL（H.264 的 AUD、SPS、PPS 与 slice 常在
同一个包里），容器侧分帧是第 \ref{ch:ch10-mux-sync} 章的事。

\subsection{完整流程：打开、送一帧、收一包、排空}

\begin{lstlisting}[style=cstyle,caption={编码侧最小闭环：打开编码器、送一帧、收一包、发 NULL 排空},label={lst:ch06-sendrecv}]
#include <libavcodec/avcodec.h>
/* 三个对象都由库分配: 字段随小版本追加, 自己 malloc 必越界 */
static AVCodecContext *enc;  static AVFrame *frame;  static AVPacket *pkt;

static int enc_open(int w, int h, int64_t bitrate)
{
    const AVCodec *c = avcodec_find_encoder(AV_CODEC_ID_H264);
    if (!c) return AVERROR(ENOSYS);
    enc = avcodec_alloc_context3(c);  frame = av_frame_alloc();
    pkt = av_packet_alloc();
    if (!enc || !frame || !pkt) return AVERROR(ENOMEM);
    enc->width = w;  enc->height = h;  enc->pix_fmt = AV_PIX_FMT_YUV420P;
    enc->time_base = (AVRational){ 1, 30 };  enc->framerate = (AVRational){ 30, 1 };
    enc->bit_rate = bitrate;  enc->gop_size = 60;  enc->max_b_frames = 2;
    enc->thread_count = 4;  enc->flags |= AV_CODEC_FLAG_GLOBAL_HEADER;
    if (avcodec_open2(enc, c, NULL) < 0) return AVERROR(EINVAL);
    frame->format = enc->pix_fmt;   /* 以打开后的协商结果为准 */
    frame->width = w;  frame->height = h;
    return av_frame_get_buffer(frame, 0);
}
/* 取干当前可取的包: cb 返回后立刻 unref, 再进下一轮 */
static int enc_drain(int (*cb)(void *, AVPacket *), void *ctx)
{
    int ret;
    while ((ret = avcodec_receive_packet(enc, pkt)) == 0) {
        ret = cb(ctx, pkt);  av_packet_unref(pkt);
        if (ret < 0) return ret;
    }
    return (ret == AVERROR(EAGAIN) || ret == AVERROR_EOF) ? 0 : ret;
}
static int enc_send(int64_t pts, int (*cb)(void *, AVPacket *), void *ctx)
{
    int ret;
    frame->pts = pts;
    ret = avcodec_send_frame(enc, frame);
    if (ret == AVERROR(EAGAIN)) {   /* 输出满: 先干一次, 再重试同一帧 */
        if ((ret = enc_drain(cb, ctx)) < 0) return ret;
        ret = avcodec_send_frame(enc, frame);
    }
    return ret < 0 ? ret : enc_drain(cb, ctx);
}
static int enc_finish(int (*cb)(void *, AVPacket *), void *ctx)
{
    int ret = avcodec_send_frame(enc, NULL);  /* 排空: NULL 只送一次 */
    if (ret < 0 && ret != AVERROR_EOF) return ret;
    return enc_drain(cb, ctx);
}
\end{lstlisting}

\begin{warning}[title={排空不等于关闭}]
只有 \Fn{avcodec\_send\_frame(enc, NULL)} 被接受之后 receive 才返回
\Tp{AVERROR\_EOF}。正常编码途中循环 receive 拿到的是 \Tp{AVERROR(EAGAIN)}，它的
意思是``再多喂一帧''而不是``结束了''，把它当 EOF 是延迟累积类故障的常见成因，附录
\ref{ap:apC-selection} 的排障表也列了它。想中止本轮、丢掉挂起帧，调
\Fn{avcodec\_flush\_buffers}。
\end{warning}

\begin{longtable}{@{}>{\raggedright\arraybackslash\footnotesize}p{0.25\textwidth} >{\raggedright\arraybackslash\footnotesize}p{0.26\textwidth} >{\raggedright\arraybackslash\footnotesize}p{0.29\textwidth}@{}}
\caption{两阶段收发的返回值契约}\label{tab:ch06-ret} \\
\toprule
返回值 & 含义 & 正确反应 \\
\midrule
\endfirsthead
\Tp{0} & 调用被接受 & receive 侧包有效，用完 unref \\
\Tp{AVERROR\Bk\_(EAGAIN)} & 输入满或输出空 & 换到另一侧循环，\textbf{原帧}重发 \\
\Tp{AVERROR\_EOF} & 已排空且无残留 & 停止本轮，再送帧须先 flush \\
\Tp{AVERROR(ENOMEM)} & 内部缓冲分配失败 & 视为致命，不要当 EAGAIN 重试 \\
\Tp{AVERROR\Bk\_(INVALID\Bk\_DATA)} & 解码侧包本身坏 & 丢包继续，参考帧可能已污染 \\
\Tp{AVERROR(ENOSYS)} & 打开时该编码器无此能力 & 多为像素格式或参数组合非法 \\
\bottomrule
\end{longtable}

\section{引用计数缓冲与帧包所有权}

\subsection{所有权的三条规则}

这是本章最容易写错的地方，而且写错之后不一定立刻崩，所以必须背下来。

\begin{itemize}
\item \Tp{AVFrame} 侧是\textbf{借}。\Fn{avcodec\_send\_frame} 不接管帧的引用计数
  缓冲：库在调用期间读像素，若要跨调用保留（B 帧重排、lookahead、环路参考帧），
  库自己做一次内部的 \Fn{av\_frame\_ref}。调用返回后你可以立刻 unref 或复用这帧；
  唯一硬要求是调用期间像素必须有效——零拷贝复用同一块 DMA 循环缓冲时，必须先确认
  上一帧已经返回。
\item \Tp{AVPacket} 侧是\textbf{给}。\Fn{avcodec\_receive\_packet} 在包里放一份新
  分配（或来自内部池）的引用计数数据，用完必须 \Fn{av\_packet\_unref}；复用同一个
  \Tp{AVPacket} 之前也要先 unref，否则是引用泄漏而非双自由。
\item 解码侧的 \Fn{avcodec\_send\_packet} 与帧侧相反：成功后包数据的归属转给库，
  而且\textbf{即使返回失败也算已消费}。想留着这个包，先 \Fn{av\_packet\_ref}。
\end{itemize}

\subsection{缓冲、池与写时复制}

\Tp{AVFrame} 的像素不在结构体里：\Tp{buf} 数组存 \Tp{AVBufferRef} 指针（引用
计数），\Tp{data} 数组只是带偏移的裸指针。所以 \Fn{av\_frame\_ref} 的代价是浅拷贝
加几次引用递增，与分辨率无关；两个帧共享同一块内存时，任一方写像素另一方立刻看到。
要保证独占写就调 \Fn{av\_frame\_make\_writable}：引用数大于一才复制，等于一时零
成本。而 \Tp{pts}、\Tp{linesize}、\Tp{flags}（\Tp{AV\_FRAME\_FLAG\_KEY} 这一类）
是\textbf{逐帧复制}的标量，不随引用计数共享。分辨率变化时先 \Fn{av\_frame\_unref}
再换几何参数，否则旧的 \Tp{linesize} 会套在新尺寸上，症状是斜切的花屏。跨线程传帧
的正确做法是每跳一次 \Fn{av\_frame\_ref}，收端负责 unref。把分配权收回应用侧的手段
是给 \Tp{get\_buffer2} 挂回调，内部用 \Tp{AVBufferPool} 供块。

\begin{lstlisting}[style=cstyle,caption={用缓冲池让解码器输出帧复用同一组像素缓冲，避开每帧分配},label={lst:ch06-pool}]
static AVBufferPool *pool;  static size_t pool_size;

static int alloc_frame(AVCodecContext *avctx, AVFrame *frame, int flags)
{
    size_t size = av_image_get_buffer_size(avctx->pix_fmt, avctx->width,
                                           avctx->height, 32);
    AVBufferRef *buf;
    if (size != pool_size) {  /* 几何变了: 关旧池, 借出的块归还后自行释放 */
        av_buffer_pool_uninit(&pool);
        if (!(pool = av_buffer_pool_init(size, NULL))) return AVERROR(ENOMEM);
        pool_size = size;
    }
    if (!(buf = av_buffer_pool_get(pool))) return AVERROR(ENOMEM);
    frame->buf[0] = buf;      /* 引用计数交给 frame 管; 已归还的块不可再写 */
    return av_image_fill_arrays(frame->data, frame->linesize, buf->data,
                                avctx->pix_fmt, avctx->width, avctx->height, 32);
}
\end{lstlisting}

\section{AVCodecContext 的参数分层与 ABI 稳定性}

\subsection{三层参数}

同一个数值常常有三条写入路径，选错会静默无效。第一层是\textbf{公共字段}：
\Tp{width}、\Tp{pix\_fmt}、\Tp{time\_base} 这类尺寸与时钟字段，\Tp{gop\_size}、\Tp{bit\_rate}、
\Tp{rc\_{\Bk}max\_{\Bk}rate}、\Tp{rc\_{\Bk}buffer\_{\Bk}size}、\Tp{qcompress} 这类码率控制字段，
以及 \Tp{profile}、\Tp{level}、\Tp{thread\_{\Bk}count}、\Tp{extradata}——类型明确、编译期可查，但只覆盖
与容器和时钟相关的最小集合。第二层是 \Tp{AVOption} 通用面：几乎所有公共字段同时是
选项，\Fn{av\_{\Bk}opt\_{\Bk}set\_{\Bk}int} 用字符串名 ``\Tp{bit\_rate}'' 写同一个字段，区别是
名字要运行期查找，写错得到 \Tp{AVERROR(ENOENT)} 而不是编译错误。第三层是\textbf{私有
选项表}，属于具体编码器而不是 libavcodec：\Tp{crf}、\Tp{qp}、\Tp{preset}、
\Tp{aq\_mode} 对 \Lib{libx264} 成立，对 \Lib{libx265} 只有部分同名，
\Tp{cpu-used} 只在 \Lib{libsvtav1} 与 \Lib{libaom} 的包装里有；而
\Tp{x264-params} 这类``原始串透传''是逃生舱，把整串交给库自己的解析器。查这张表可用
\Fn{av\_{\Bk}opt\_{\Bk}next} 遍历，命令行下 \texttt{ffmpeg -h encoder=libx264} 打印的就是它。

\subsection{选项注入与剩余项检查}

\Fn{avcodec\_open2} 的第三个参数是 \Tp{AVDictionary} 的二重指针，被当成
``待消费的选项队列''：打开时被库内解析器命中的键会移走。所以拼错的键不会报错，
它只是留在字典里没人管，想抓到就必须查余量。

\begin{lstlisting}[style=cppstyle,caption={用 AVDictionary 注入选项，并把未消费的键视为配置错误},label={lst:ch06-dict}]
int open_x264(AVCodecContext *enc, const AVCodec *codec, const char *preset,
              std::vector<std::string> &rejected)
{
    AVDictionary *opts = nullptr;
    av_dict_set(&opts, "preset", preset, 0);
    av_dict_set(&opts, "tune", "zerolatency", 0);
    av_dict_set(&opts, "crf", "23", 0);
    av_dict_set(&opts, "x264-params",
                "keyint=60:min-keyint=15:scenecut=40:aq-mode=3", 0);
    int ret = avcodec_open2(enc, codec, &opts);
    for (const AVDictionaryEntry *e = av_dict_get(opts, "", nullptr, 0);
         e; e = av_dict_get(opts, "", e, 0))
        rejected.push_back(std::string(e->key) + "=" + e->value);
    av_dict_free(&opts);
    if (!rejected.empty()) return AVERROR(EINVAL);  /* 拼错的选项别静默 */
    return ret;
}
\end{lstlisting}

\begin{compare}[title={三条写入路径的取舍}]
公共字段：类型安全、无查找开销，但改字段名就是 ABI 事件。选项字符串：跨版本稳，
同一套代码可复用到不同包装，代价是拼错只在运行期暴露。私有透传串（\Tp{x264-params}、
\Tp{x265-params}、\Tp{svtav1-params}）表达力最强，等于把官方 CLI 的解析器搬进进程，
但它不受 libavcodec 的校验与文档约束，升级库时要单独回归。
\end{compare}

\subsection{ABI 稳定性靠什么维持}

FFmpeg 的公开结构体（\Tp{AVCodecContext}、\Tp{AVFrame}、\Tp{AVPacket}）不是
``字段布局冻结''式的稳定，而是三条约定的合力。

\begin{enumerate}
\item \textbf{只加不改}：新字段一律追加到结构体尾部，已有字段的类型与语义在小
  版本内不变；废弃字段先标 \Tp{deprecated}，保留到下一个大版本才删。
\item \textbf{库分配}：结构体只能由 \Fn{avcodec\_{\Bk}alloc\_{\Bk}context3}、
  \Fn{av\_{\Bk}frame\_{\Bk}alloc}、\Fn{av\_{\Bk}packet\_{\Bk}alloc} 创建；尺寸随版本增长，调用方按
  \texttt{sizeof} 造对象传给新库就越界读，所以这条是上一条成立的前提。内部状态一律
  藏在库里私有的结构中，应用不许自己算偏移。
\item \textbf{版本三元组}：\texttt{libavcodec/{\Bk}version.h} 的
  \Tp{LIBAVCODEC\_VERSION\_MAJOR} 与 MINOR、MICRO 对应运行期的
  \Fn{avcodec\_version}；追加字段与新增选项升 MINOR，破坏性变更升 MAJOR，只修 bug
  升 MICRO。\Fn{avcodec\_configuration} 与 \Fn{avcodec\_license} 说明这份库是用什么
  configure 参数编出来的。
\end{enumerate}

由此得到实用的降级技巧：要写一个新字段、又要兼容旧发行版时，先试
\Fn{av\_{\Bk}opt\_{\Bk}set\_{\Bk}int}，返回 \Tp{AVERROR(ENOENT)} 说明这版库没有该选项，再退回旧
写法。选项名不是 ABI，但它比 \texttt{sizeof} 安全得多。

\section{x264 与 x265：原生 API 与 FFmpeg 包装的关系}

\Lib{x264} 也是 C ABI，但形态与 libavcodec 完全不同：句柄是不透明的
\Tp{x264\_t} 指针，参数是\textbf{一个巨大的结构体} \Tp{x264\_{\Bk}param\_{\Bk}t}
（打开时被读一次，生效值用 \Fn{x264\_{\Bk}encoder\_{\Bk}parameters} 回读），
输出单位是 \textbf{NAL} 而不是包。

\begin{itemize}
\item \Fn{x264\_param\_default} 填保守默认；\Fn{x264\_param\_default\_preset}
  等价于 CLI 的 \Opt{preset} 加 \Opt{tune}。\Fn{x264\_param\_parse} 接受与 CLI
  长选项同名（不带前缀横线）的键，返回值能区分``名字不认识''与``值不合法''，
  所以它可以当校验器用。
\item \Fn{x264\_picture\_init}(pic) 只做一件事：把属于调用方的元数据
  （\Tp{i\_pts}、\Tp{opaque}、\Tp{i\_type}）复位到``自动''。像素内存仍由你给；
  想让库给就调 \Fn{x264\_picture\_alloc}，之后 \Fn{x264\_picture\_clean}。
\item \Fn{x264\_encoder\_encode} 的返回值是本次输出的\textbf{总字节数}，出参给
  NAL 个数；每个 \Tp{x264\_nal\_t} 带 \Tp{payload}、\Tp{i\_size}、\Tp{i\_type}
  与是否长起始码的标志。\Tp{pic\_out} 传非空才拿到重建帧。
\item 排空：\Tp{pic\_in} 传 NULL 反复调用直到返回 0，
  \Fn{x264\_encoder\_delayed\_frames} 先告诉你挂着几帧。\Fn{x264\_encoder\_reconfig}
  允许中途换参数、下一帧生效，但不是所有参数都可换，且不得与 encode 并发。
\end{itemize}

\Lib{libx264} 包装维持``一个 \Tp{AVCodecContext} 对一个 \Tp{x264\_t}''：open
时把公共字段单向映射进 \Tp{x264\_param\_t}（\Tp{bit\_rate} 到 \Tp{rc.i\_bitrate}
且单位换成 kbit/s、\Tp{gop\_size} 到 \Tp{keyint}、\Tp{max\_b\_frames} 到
\Tp{bframes}、\Tp{qcompress} 到 \Tp{qcomp}），再按
\Tp{AV\_CODEC\_FLAG\_GLOBAL\_HEADER} 决定 SPS/PPS 是收进 \Tp{extradata}
（AVCC 形态）还是留在带起始码的码流里（Annex B）；send 阶段用帧的 \Tp{data}
与 \Tp{linesize} 填平面、调 encode、把返回的 NAL 数组\textbf{拼成一个包}，
并用 \Tp{opaque} 带回 pts 与关键帧标志。

代价是表达力的单向性：\Tp{AVCodecContext} 模型里没有 open GOP、sliced threads、
psy-rd、zones 这些概念，它们只能从透传串进去；反过来 x264 的 \Tp{reconfig} 在
libavcodec 的收发契约里也没有对应动作，想运行中改码率到底生不生效必须实测
（第 \ref{ch:ch07-rate-control} 章给了两条可验证写法）。\Lib{x265} 类似，
但有两个额外现实：位深变体（8、10、12 位是不同符号，官方推荐用
\Fn{x265\_api\_get} 取函数表而不是直链符号）、排空延迟可先问
\Fn{x265\_encoder\_get\_frame\_delay}。强制 IDR 在 x265 是 \Tp{x265\_picture}
的 \Tp{forceKeyFrame}，在 x264 是 \Tp{i\_type} 设成 \Tp{X264\_TYPE\_IDR}；
x265 的字节数要逐块从 NAL 长度字段累加，字段命名是驼峰风格，
移植代码时这是最容易漏的一处。

\section{aom、vpx 与 EbApi：三种不同的句柄模型}

\Lib{libvpx} 的 \Tp{vpx\_{\Bk}codec\_{\Bk}ctx\_{\Bk}t} 由\textbf{调用方在栈上}给出（传地址），
配置来自 \Fn{vpx\_{\Bk}codec\_{\Bk}enc\_{\Bk}config\_{\Bk}default} 填好的
\Tp{vpx\_{\Bk}codec\_{\Bk}enc\_{\Bk}cfg\_{\Bk}t}，接口对象是 \Fn{vpx\_{\Bk}codec\_{\Bk}vp9\_{\Bk}cx}；微调走变长
参数的 \Fn{vpx\_{\Bk}codec\_{\Bk}control}（如 \Tp{VP9E\_{\Bk}SET\_{\Bk}TILE\_{\Bk}COLUMNS}），两遍编码靠
把 \Tp{g\_pass} 设成 first pass 与 last pass。输出不是返回值而是\textbf{迭代器}：
\Fn{vpx\_{\Bk}codec\_{\Bk}get\_{\Bk}cx\_{\Bk}data} 每次给一个 \Tp{vpx\_{\Bk}codec\_{\Bk}cx\_{\Bk}pkt\_{\Bk}t}，
类型由 \Tp{kind} 区分，帧字节在 \Tp{data.frame}。

\Lib{libaom} 沿用同一骨架（\Tp{aom\_codec\_iter\_t}、\Fn{aom\_codec\_get\_cx\_data}），
但把 deadline 从每次 encode 的参数改成配置字段 \Tp{g\_deadline}，补了
\Fn{aom\_codec\_set\_option} 做字符串选项（end-usage、cq-level、cpu-used、
lag-in-frames），ctl 族按 \Tp{AOME} 与 \Tp{AV1E} 分工：前者沿袭通用项
（\Tp{AOME\_SET\_TUNING}），后者是 AV1 专属（\Tp{AV1E\_SET\_CPUUSED}、
\Tp{AV1E\_SET\_ROW\_MT}）。解码侧对称：接口换成 \Fn{aom\_codec\_av1\_dx}，压缩数据用
\Fn{aom\_codec\_decode} 送入，重建帧同样靠 \Fn{aom\_codec\_get\_frame\_data} 迭代取出。

\Lib{svt-av1} 的 \Tp{EbApi} 又不一样：进程级初始化单独成一步，句柄用二级指针由
库分配，参数集中在一个配置对象里一次灌入，缓冲用统一的结构化描述块（含通道号、
标志位与用户数据指针），用其中的结束标志表达排空；取回的码流块必须由调用方显式
归还，否则内部缓冲池耗尽，症状是跑几分钟后不再出包而不是崩溃。它在 FFmpeg 侧只
暴露 \Tp{preset}（p0 最慢、p13 最快）与 \Tp{svtav1-params} 透传串，这是目前最
省事的用法（第 \ref{ch:ch13-framework-integration} 章的 \Lib{GStreamer} 通路同理）。

\begin{longtable}{@{}>{\raggedright\arraybackslash\footnotesize}p{0.11\textwidth} >{\raggedright\arraybackslash\footnotesize}p{0.145\textwidth} >{\raggedright\arraybackslash\footnotesize}p{0.16\textwidth} >{\raggedright\arraybackslash\footnotesize}p{0.13\textwidth} >{\raggedright\arraybackslash\footnotesize}p{0.145\textwidth} >{\raggedright\arraybackslash\footnotesize}p{0.115\textwidth}@{}}
\caption{六个编码器句柄模型的 API 形态对比}\label{tab:ch06-handles} \\
\toprule
模型 & 创建方式 & 参数注入 & 线程设置 & 排空语义 & 许可证 \\
\midrule
\endfirsthead
\toprule
模型 & 创建方式 & 参数注入 & 线程设置 & 排空语义 & 许可证 \\
\midrule
\endhead
\midrule
\bottomrule
\endlastfoot
libavcodec & 库分配的 \Tp{AV\Bk\_Codec\Bk\_Context} & 公共字段与选项字典 & \Tp{thread\Bk\_count} 与类型 & 送 NULL 帧至 \Tp{AVERROR\Bk\_EOF} & LGPL 2.1\Bk+；开 GPL 后 GPL 2\Bk+ \\
x264 原生 & \Fn{x264\Bk\_encoder\Bk\_open} 给句柄 & \Fn{x264\Bk\_param\Bk\_parse} 或直写 param & 线程数与 lookahead 线程 & \Tp{pic\_in} 传 NULL 至返回 0 & GPL 2.0\Bk+，另有商用授权 \\
libvpx & \Tp{vpx\Bk\_codec\Bk\_ctx\_t} 在栈上 & 配置结构体与 \Fn{vpx\Bk\_codec\Bk\_control} & \Tp{g\_threads} 或多句柄 & img 传 NULL，迭代器取空 & BSD 风格，附免专利费授权 \\
libaom & \Tp{aom\Bk\_codec\Bk\_ctx\_t} 在栈上 & 配置结构体、ctl 族与选项串 & \Tp{g\_threads} & img 传 NULL；\Tp{g\_lag\Bk\_in\Bk\_frames} 越大挂越多 & BSD 2 条款加专利承诺 \\
SVT-AV1 & 二级指针由库分配 & 配置对象一次灌入加透传串 & 配置项与 \Tp{preset} & 置结束标志，包须归还 & BSD 2 条款加专利承诺 \\
libopus & \Fn{opus\Bk\_encoder\Bk\_create} 给堆对象 & 变长参数 ctl & 无，实现单线程 & 无排空延迟；\Tp{OPUS\Bk\_RESET\Bk\_STATE} & BSD 3 条款 \\
\end{longtable}

\section{音频侧：Opus 与 FDK-AAC 的生命周期}

音频编码器普遍没有重排延迟（CELT 只有很小的算法延迟），所以接口形态比视频简单：
一次调用出结果。\Lib{Opus} 用堆上的不透明 \Tp{OpusEncoder}，
\Fn{opus\_encoder\_create} 带一个错误码出参；参数只有\textbf{变长参数的 ctl 族}，
没有配置结构体。\Fn{opus\_encoder\_ctl} 送 \Tp{OPUS\_SET\_BITRATE} 时接受的是
\textbf{bit/s 而不是 kbit/s}，这个单位差是踩坑最多的地方。\Fn{opus\_encode}
自己写好 TOC 字节与长度，返回的就是可以直接塞进 RTP 或 Ogg 的一个完整包；
输出缓冲必须不小于 \Tp{OPUS\_MAX\_PACKET\_SIZE}，否则库拒绝而不是截断。
码率、带宽、复杂度与 \Tp{OPUS\_SET\_INBAND\_FEC} 都能运行中随时改、下一包生效，
这正是 \DomRT{} 链路（第 \ref{ch:ch11-streaming} 与第 \ref{ch:ch14-congestion-abr}
章）需要的性质。解码侧对称：\Fn{opus\_decoder\_create} 配 \Fn{opus\_decode}，输出
\Tp{int16\_t} 交错样本（每声道样本数由调用方显式给），丢包时传 NULL 即触发隐式丢包
补偿。

\Lib{fdk-aac} 是另一代设计：句柄是 \Tp{HANDLE\_AACENCODER}，\Fn{aacEncOpen}
收二级指针；配置通过 \Fn{aacEncSpecialCommand} 的``命令号加参数指针''三元组下发
（选编码器配置、选码率档、开 SBR、指定带宽），没有字符串选项也没有 ctl。
它的收发同步但带\textbf{分数延迟}：输入与输出各用一个``缓冲加标识符''的描述块，
编码器可以只吞掉一部分样本，所以要循环喂到吃完，输出侧同样可能为空。它报出的编码器
延迟与前置填充要在容器里写成 edit list 才能对上时间轴（第 \ref{ch:ch10-mux-sync} 章）。
两者的差异不止技术：\Lib{fdk-aac} 是
Fraunhofer IIS 的专有许可，编进可再分发的二进制需单独授权，FFmpeg 只有加
\Opt{enable-nonfree} 才允许用它，而这种构建既不是 LGPL 也不是 GPL；\Lib{Opus}
则是干净的 BSD 加免版税格式许可。同一段代码换库，法务成本差一个量级，
这在附录 \ref{ap:apC-selection} 里展开。

\begin{bestpractice}[title={写编解码封装层的四条纪律}]
\begin{enumerate}
\item 把所有权写进类型：帧用带析构的句柄封装、离开 send 之后释放，包在回调边界
  上明确由谁 unref，别让两个模块都以为对方负责。
\item 永远检查 \Fn{avcodec\_open2} 之后字典的余量，把未消费的键当成配置错误抛出，
  而不是静默降级。
\item 能用包装层就用包装层；需要 \Tp{reconfig}、zones、open GOP 这类原生特性时
  直接调原生 API，并把 NAL 到包的转换写成可单测的纯函数。
\item 许可证随构建配置变化：把 \Fn{avcodec\_configuration} 与
  \Fn{avcodec\_license} 打进日志与``关于''面板，出货前对照附录
  \ref{ap:apC-selection} 的许可矩阵复核。
\end{enumerate}
\end{bestpractice}

% ============================================================
% ---- frag_ch07.tex ----
\chapter{码率控制与参数面：QP、CRF、VBV 与延迟档}
\label{ch:ch07-rate-control}

第 \ref{ch:ch06-codec-api} 章解决的是``把帧交出去、把包拿回来''的接口形态；本章解决
同一批字节里\textbf{每一帧该花多少}的问题。这两件事在代码里挨得很近——都发生在
\Tp{AVCodecContext} 的字段赋值和 \Fn{avcodec\_open2} 之前——但它们的失败模式完全不
同：形态写错会当场返回错误码，额度写错会安静地跑很多天，直到某个场景切换把带宽顶
穿，或者某个低码率片段出现整片色带。解码侧与硬件通路见第
\ref{ch:ch08-decode-hw} 章，滤镜链在帧进入本章之前做了什么见第
\ref{ch:ch09-filter-resample} 章。

\section{量化是唯一真正被决定的量}

编码器里能被``额度''直接操纵的只有量化：变换系数除以量化步长、四舍五入、丢掉余数。
QP 与量化步长的对应关系是几何的——\textbf{每加 6 个 QP，步长翻一倍}，于是同一批系数
留下的比特大致减半；H.264 与 H.265 的 QP 都在 0 到 51 之间，AV1 用 0 到 63 的另一把
尺。除量化之外的所有手段——模式决策、运动搜索、环路滤波、熵编码——都只是在决定
``哪些系数值得留''，最终仍归结为一个每块 QP。

这条链给出本章全部讨论的地基：\textbf{码控不是分配比特，而是分配 QP}。因此

\begin{itemize}
\item 给定内容，QP 与码率之间是强关系；给定\emph{编码器}，QP 与\emph{主观质量}之间
  才是强关系。跨码流比 QP 数字、跨编解码器比 CRF 数字，都只是在比两把不同的尺。
\item 给定码率目标，质量取决于内容复杂度，而复杂度在时间轴上差一个数量级：静态会议
  画面与手持运动镜头之间不是``稍贵一点''，而是十倍量级。任何``固定每帧字节''的策略
  都会在这种起伏上失败。
\item 所以额度契约必须同时规定三件事：平均值往哪里靠、瞬时峰值允许多高、以及允许
  质量下探到哪一档（\Tp{qmin}/{\Bk}\Tp{qmax}，或各家等价的量化范围钳位）。
\end{itemize}

FFmpeg 侧的通用入口是 \Tp{AVCodecContext} 的四个字段：\Tp{bit\_rate}（平均额度）、
\Tp{rc\_max\_rate}（峰值额度）、\Tp{rc\_buffer\_size}（允许峰值持续多久，\textbf{单位
是比特}）、以及 \Tp{global\_quality} 与 \Tp{qmin}/{\Bk}\Tp{qmax} 这组质量钳位。注意
\Tp{global\_quality} 的量纲随编码器而变，不要把它当跨库可比的数字用。

\begin{warning}[title={QP、CRF 与全局质量是三把尺，不能拿数字对话}]
同一段素材，\Lib{x264} 的 \Opt{crf} 23 大致是``视觉近无损的起点''，\Lib{x265} 的
\Opt{crf} 23 明显更粗，\Lib{libsvtav1} 的 \Opt{crf} 23 则接近透明。各家把 CRF 定标在
自己的率失真模型上，没有公共基准。可比的做法只有两种：把主观质量交给同一个人评审，
或者固定质量指标（VMAF/SSIM 的目标分）反推码率。写自动化流程时，把
\Fn{avcodec\_configuration} 与版本号一起记进产出清单，否则``同样的 CRF 为什么换了
构建就变了''无法归因。
\end{warning}

\section{三种额度契约：定QP、恒定质量与缓冲约束}

契约的名字在各家不同，但只有三种形状。

\textbf{定量化（CQP、constant-QP）。}每帧用同一个 QP（或按帧类型偏移一个小量）。
质量稳定、行为最好预测，代价是码率完全不封顶：一个场景切换就是一次尖峰。它适合
``产物体积不重要、可复现性重要''的场合，比如离线存档、回归测试基准、需要逐字节比对
的编码器开发。

\textbf{恒定质量（CRF、one-pass quality-based VBR）。}按目标质量出块，额度随内容
起伏。压缩效率是三者里最好的——复杂处多花、简单处少花——但平均码率与峰值码率都
只是结果，不是约束。它是本地转码与归档的默认选择，也常常是直播误配的起点：CRF 模式
下``我已经设了 4000 kbps''这类叙述往往是假的，因为该字段根本没参与决策。

\textbf{缓冲约束（CBR/VBV、two-pass ABR）。}引入解码端缓冲模型：平均额度是
\Tp{bit\_rate}，瞬时斜率上限是 \Tp{rc\_max\_rate}，允许的突发由
\Tp{rc\_buffer\_size} 决定。编码器在每个帧上做一次``缓冲会不会被抽干''的预测，一旦
预测越界就立刻粗化 QP。这是直播与RTC的唯一选择，也是唯一需要你把带宽预算、缓冲深度
和可容忍延迟三件事同时说清楚的契约。

\begin{longtable}{@{}>{\raggedright\arraybackslash\footnotesize}p{0.17\textwidth} >{\raggedright\arraybackslash\footnotesize}p{0.13\textwidth} >{\raggedright\arraybackslash\footnotesize}p{0.15\textwidth} >{\raggedright\arraybackslash\footnotesize}p{0.14\textwidth} >{\raggedright\arraybackslash\footnotesize}p{0.23\textwidth}@{}}
\caption{三种额度契约的可控维度}\label{tab:ch07-contract} \\
\toprule
契约 & 峰值可控 & 质量稳定 & 额外延迟 & 典型用途 \\
\midrule
\endfirsthead
定QP（CQP） & 否 & 高（QP 固定） & 无 & 存档、回归基准、编码器对拍 \\
恒定质量（CRF） & 否 & 中（随复杂度） & 取决于 \Tp{rc\Bk\_\_lookahead} & 本地转码、离线归档 \\
单遍 CBR/VBV & 是 & 低（复杂度顶不住时粗化） &  lookahead 置零时可到帧级 & 直播上行、RTC \\
两遍目标码率 & 平均值可控 & 中 & 不适用（离线） & 交付固定体积的成品 \\
低延迟档 & 是 & 低到中 & 一帧到两帧 & 会议、云游戏、硬件在环 \\
\bottomrule
\end{longtable}

表 \ref{tab:ch07-contract} 的``额外延迟''一列是本册反复强调的账：额度契约不只是质量
问题，它直接决定编码器被允许看多远、被允许重排多远。

\section{VBV 的账本：缓冲充实度与允许的突发}

VBV 模型在编码器里维护一只虚拟缓冲：每个包的字节数从缓冲里取出，缓冲以
\Tp{rc\_max\_rate} 的速率回填，容量是 \Tp{rc\_buffer\_size}。把它换算成时间更直观：缓冲深度（秒）等于 \Tp{rc\_buffer\_size} 除以 \Tp{rc\_max\_rate}，
每帧额度（比特）等于 \Tp{rc\_max\_rate} 除以帧率。

于是``缓冲''的真实含义是\textbf{允许的突发时长}：深度 0.5 秒意味着一段连续高复杂度的
画面可以把瞬时速率顶在 \Tp{rc\_max\_rate} 上约半秒再被压回去；深度 4 秒则宽松得多，
但解码端要多等 4 秒才拿到完整的缓冲充实度——这正是``同样的码率，直播平台的 VBV 设置
比本地 ffmpeg 更糊''的原因。

三种越界方式各有症状：

\begin{enumerate}
\item \textbf{上溢}（缓冲被填满）：编码器必须把这一帧做大，代价是延迟或丢帧；多数
  实现选择粗化 QP 而不是丢帧，于是症状是``画面一变复杂就糊''。
\item \textbf{下溢}（缓冲被抽干）：违反合规性，真实解码器可能缓冲不足而卡顿。编码器
  的对策是把后续帧的额度压低——症状是一段``过度平滑''的低质量窗口，通常在场景切换后
  一到两秒出现。
\item \textbf{首帧超支}：IDR/I 帧的体积常常是平均帧的十倍量级，如果缓冲初始充实度
  设得太小，第一个关键帧就违反约束。做法是给关键帧单独放宽（各家的
  \Opt{qcomp}、\Opt{ipratio}、关键帧量化偏移一族），并把 GOP 长度与网络侧的丢包恢复
  需求一起定，见第 \ref{ch:ch14-congestion-abr} 章。
\end{enumerate}

\begin{lstlisting}[style=cstyle,caption={额度契约的完整一次表达：通用字段加私有选项，并验证哪些键真的被消费},label={lst:ch07-vbv}]
#include <libavcodec/avcodec.h>
#include <libavutil/opt.h>

static int enc_setup_cbr(AVCodecContext *enc, int fps, int kbps_avg, int kbps_peak)
{
    /* 通用层：三个字段的单位都是 bit/s，buffer 是 bit（不是秒！）          */
    enc->bit_rate        = (int64_t)kbps_avg  * 1000;
    enc->rc_max_rate     = (int64_t)kbps_peak * 1000;
    enc->rc_buffer_size  = kbps_peak * 1000 / 2;   /* 0.5 秒突发额度       */
    enc->qmin            = 18;                      /* 别用无限额度换坏质量 */
    enc->qmax            = 48;
    enc->gop_size        = 2 * fps;                 /* 关键帧间隔按秒计     */
    enc->max_b_frames    = 0;                       /* RTC 侧不做重排       */

    /* 私有层：字符串选项，名字与单位由具体编码器决定                       */
    av_opt_set(enc->priv_data, "preset",    "veryfast", 0);
    av_opt_set(enc->priv_data, "tune",      "zerolatency", 0);
    av_opt_set(enc->priv_data, "vbv_maxrate", "4400", 0);  /* x264: kbit/s */
    av_opt_set(enc->priv_data, "vbv_bufsize", "2200", 0);

    /* 可验证写法之一：每个字符串选项都单独看返回值，非 0 就是名字或单位错  */
    int r = av_opt_set(enc->priv_data, "vbv_bufsize", "2200", 0);
    if (r == AVERROR_OPTION_NOT_FOUND) return AVERROR(EINVAL);

    /* 可验证写法之二：open 之后字典里剩下的键必须为空，否则配置根本没进去 */
    AVDictionary *opts = NULL;
    av_dict_set(&opts, "x264-params", "keyint=60:min-keyint=30:scenecut=40", 0);
    int ret = avcodec_open2(enc, avcodec_find_encoder(enc->codec_id), &opts);
    AVDictionaryEntry *left = av_dict_get(opts, "", NULL, AV_DICT_IGNORE_SUFFIX);
    av_dict_free(&opts);
    if (ret < 0) return ret;
    if (left) return AVERROR(EINVAL);   /* 有键没被消费：当配置错误抛出 */
    return 0;
}
\end{lstlisting}

第 \ref{ch:ch06-codec-api} 章``把未消费的键当成配置错误抛出''那条纪律，具体到码控上
就是上面这两处检查。它的价值在于抓住一类非常安静的故障：选项名随小版本改动、单位从
kbit/s 变成 bit/s、或者你其实走的是硬件编码器而它不支持该参数——三种情况都不会让
\Fn{avcodec\_open2} 失败，只会让实际码率与你以为的不一样。

\begin{note}[title={验证一条额度设置是否生效，只看两样东西}]
不要看日志里打印的参数表，看产物本身：一是逐帧字节序列（\Fn{ffprobe} 的
\Tp{pkt\_size} 列，见第 \ref{ch:ch10-mux-sync} 章），把它对时间求和，检查斜率是否
始终不超过 \Tp{rc\_max\_rate}、滑动均值是否落在 \Tp{bit\_rate} 附近；二是缓冲充实度
曲线，多数编码器能从调试输出里取到，越界点应当与画面发糊的时刻对齐。两处都对上，契约
才算成立。
\end{note}

\section{先分析再分配：lookahead、自适应量化与场景切换}

单遍 CBR 之所以能用，是因为编码器被允许先看未来几帧再决定当前帧的额度。这一族参数的
总称是 lookahead：

\begin{itemize}
\item \textbf{分析窗口}（\Opt{rc-lookahead} 一族）决定``能看到多远''。窗口越大，
  场景切换前的额度可以提前攒、切换后的粗化可以避免；代价是每帧要缓存未来的原始帧，
  内存按窗口帧数线性涨，延迟同线性涨。
\item \textbf{自适应量化}（\Opt{aq-mode}、\Opt{aq-strength} 一族）在\emph{同一帧内}
  重新分配：对暗部、纹理密集区、人眼不敏感区给不同 QP。它改善主观质量，但也让``帧的
  字节数与其像素方差成正比''这条粗略直觉失效——排查码率异常时不能只看画面复杂度。
\item \textbf{时域传播加权}（宏块级率失真优化这一族）把``这一帧被后续多少帧引用''计入
  代价，被引用多的区域给更小的 QP。它的收益在多参考帧、长 GOP 时最明显。
\item \textbf{场景切换检测}（\Opt{scenecut} 阈值一族）在帧间统计突变处决定是否插入
  帧内编码。插得太少，切换后长时间``拖影''；插得太多，一堆昂贵的 IDR 把缓冲顶穿。
  工程上常见的是``阈值 + 最小间隔''成对设置，最小间隔至少覆盖 GOP 的一小段。
\end{itemize}

这里有一条与采集侧直接相关的算术（第 \ref{ch:ch04-capture-display} 章留下的细账）。
裸码率是宽乘高乘帧率乘每像素比特，$1920 \times 1080$、30 fps、每像素 1.5 字节（
YUV420P）约 94 Mbit/s；4:2:2 10 bit 则是它的 1.78 倍。压缩到 4 Mbit/s 意味着约
24:1 的总压缩比，而这个额度要由``每帧约 133 kbit''的分配去实现——静态帧用不完、
切换帧严重不够。\textbf{lookahead 与场景切换的本质，是把这条不可行的均匀分配改成
可行的不均匀分配}；额度契约越紧，这族参数的收益越大，这也是为什么低码率直播比高码率
存档更依赖参数面。

\begin{compare}[title={参数在哪一层被暴露，决定了它能不能被验证}]
FFmpeg 路径有三层：\Tp{AVCodecContext} 的通用字段（跨编码器可比，但只覆盖平均值、
峰值、缓冲、GOP、B 帧这些公约数）；\Tp{priv\_data} 上的字符串选项（\Fn{av\_opt\_set}
或 \Opt{x264-params} 透传，覆盖面最广，但名字与单位随实现走，拼错不报错）；原生 API
的参数结构体（\Lib{x264} 的 \Fn{x264\_param\_parse} 一路、\Lib{libaom} 的控制项、
\Lib{svt} 的 \Fn{EbApiSet} 一族，能碰到包装层没暴露的特性，如分区策略与环路滤波细节）。
可验证性沿这三层递减：通用字段可以用产物字节数反证，字符串选项要靠``未消费键必须为
空''，原生参数则常常要打开编码器自己的调试输出才能确认。选封装层时要按``这个参数出了
错我会不会知道''来判断。
\end{compare}

\section{两遍编码：把复杂度统计变成额度表}

离线交付要的是``体积一定、质量尽量好''。单遍做不到，因为它不知道后面有什么。两遍
编码的分工是：第一遍用轻量分析产出逐帧（或逐宏块）复杂度统计，第二遍按统计把总预算
切成一张时间轴上的额度表，再真正编码。

它成立的前提是\textbf{两遍看到的是同一件事}，因此有五条一致性要求：分辨率与帧率、
像素格式、GOP 结构、preset/分析深度、以及素材本身必须完全一致；换素材或改尺寸后旧
统计一律作废。多数实现对统计文件做了版本或参数指纹校验，也有不校验的——不校验时
症状是``输出体积对，但质量分布诡异''，比报错更难查。

\begin{lstlisting}[style=shellstyle,caption={三组两遍写法对照：x264 手工两遍、FFmpeg 透传、AV1 一次调用内两遍},label={lst:ch07-twopass}]
# 1) x264 直接调：第一遍只为统计，输出丢弃
x264 --fps 30 --preset slow --crf 23 --pass 1 --stats show.log \
     --bitrate 4000 --vbv-maxrate 4400 --vbv-bufsize 2200 in.y4m -o /dev/null
x264 --fps 30 --preset slow --crf 23 --pass 2 --stats show.log \
     --bitrate 4000 --vbv-maxrate 4400 --vbv-bufsize 2200 in.y4m -o out.264

# 2) 走 FFmpeg：-pass 与 -passlogfile 是通用层，编码器私有项透传
ffmpeg -y -i in.mp4 -c:v libx264 -preset slow -b:v 4000k -maxrate 4400k \
       -bufsize 2200k -pass 1 -passlogfile show -an -f mp4 /dev/null
ffmpeg -y -i in.mp4 -c:v libx264 -preset slow -b:v 4000k -maxrate 4400k \
       -bufsize 2200k -pass 2 -passlogfile show -c:a copy out.mp4

# 3) AV1 参考实现：--passes=2 在同一次调用里跑两遍，--fpf 是统计文件
aomenc --codec=av1 --passes=2 --fps-num=30 --fps-denom=1 \
       --end-usage=cbr --bitrate=4000 --buf-sz=2200 --buf-optimal-sz=1100 \
       --fpf=stats.json --tile-columns=2 in.ivf -o out.ivf

# 4) 音频侧没有第二遍：Opus 的额度是逐帧决策，只做平均/约束选择
opusenc --bitrate 24 --frame-size 20 --speech in.wav out.opus
\end{lstlisting}

三组命令的对照里有两条要点。\Opt{end-usage}（\Lib{libaom}/{\Bk}\Lib{SVT-AV1} 一族）与
\Lib{x264} 的 \Opt{crf}/{\Bk}\Opt{bitrate} 不是同一词汇表：前者是``用什么模式分配''，后者
是``分配到什么值''，各家把这两件事打包进不同的选项名里，跨库脚本必须逐库重写。
\Tp{buf-optimal-sz} 一类是``缓冲的目标充实度''，与 \Tp{buf-sz}（容量）共同决定突发
行为——只设容量不设目标充实度，第一遍之后的填充策略就交给默认值。

两遍什么时候不值得：直播与 RTC（素材还没发生）、时长长于预算的离线批次（第一遍本身
超预算）、以及内容复杂度可用轻量分析估出的场景（此时单遍 VBV 加足够 lookahead 即可）。

\section{延迟档：把码控参数换成毫秒}

额度契约对延迟的影响是可算的。端到端延迟里编码器负责的这部分，由三项构成：

\begin{enumerate}
\item \textbf{重排延迟}：有 B 帧就要把参考顺序与显示顺序分开，最小是``连续两个 B''
  的间距，通常一到两帧。关掉 B 帧（\Tp{max\_b\_frames} 置 0）是云游戏与硬件在环的
  第一刀。
\item \textbf{前瞻延迟}：\Opt{rc-lookahead} 与\emph{帧级}决策要看的未来帧数，直接
  等于编码器``必须攒够才敢出手''的帧数。置 0 后额度决策退化为纯贪心，质量代价在
  场景切换处最明显。
\item \textbf{缓冲延迟}：\Tp{rc\_buffer\_size} 换算出的秒数不会出现在编码器输出
  延迟里，但会出现在解码端合规性上——接收缓冲至少要这么大才不抽干，见第
  \ref{ch:ch02-clock-buffer} 章的 jitter buffer 账。
\end{enumerate}

``低延迟档''在不同层是不同东西：\Lib{x264} 的 \Opt{tune=zerolatency} 是一组开关的
打包（关 lookahead、关重排、限制关键帧开销、改变码控的短视程度）；FFmpeg 通用层的
\Tp{flags} 里有 \Tp{AV\_CODEC\_FLAG\_LOW\_DELAY}；硬件编码器（第
\ref{ch:ch08-decode-hw} 章）通常在创建时会话时就选定 GOP 与 B 帧策略，运行期改不了。
三者的共同结果是：低延迟档下同样的平均码率，质量必然低于离线档，差距集中在切换与
高运动段。

\begin{warning}[title={被静默降级的三个参数，低延迟配置里最常见}]
其一，\Opt{preset} 类``分析深度''选项在硬件编码器上没有对应物，透传后不进决策；其二，
音频侧的复杂度/带宽类控制在某些包装层被固定为默认，``改了就有效果''只对本机那一版
成立；其三，帧率与时间基不一致时（第 \ref{ch:ch10-mux-sync} 章），编码器以为的
``每帧额度''与实际封装速率差一个倍数，症状是码率恒定但画面按固定周期抖一下——这通常是
时间基问题而不是码控问题。三者共同点是都不报错。
\end{warning}

\section{参数面与跨版本稳定性}

码控参数是编码器行为差异的主要来源，也是升级后最需要复核的部分：新增选项、改名、
默认值变化、以及同一选项换实现（包装层把某参数映射到另一个内部字段）。工程上把这种
风险压到可管理的做法有三条：把``生效参数''而不是``提交参数''落盘（从编码器自己的
统计/调试输出里取），把构建配置字符串与库版本号打进每个产物元数据，以及为额度契约留
一条回归指标——同一素材、同一脚本，逐帧字节曲线的形状应与历史一致，斜率上限与滑动
均值是最便宜的对拍量。

\begin{guideline}[title={把额度当成契约来写，而不是当成一个数}]
\begin{enumerate}
\item 先写三件东西再填参数：平均额度、峰值额度、允许的突发时长。三者缺一，``4000 kbps''
  这种单一数字无法判断是否满足下游。
\item 恒定质量与缓冲约束不要混用意图：设了 CRF 又指望峰值不超，必须同时显式给出
  VBV 两项，且要验证它们是否被消费。
\item 每个配置项配一条``没生效我就知道''的断言：未消费字典为空、逐帧字节曲线斜率
  不越界、关键帧间隔落在预期范围内。
\item 低延迟预算要先换算成帧数再选参数：重排、前瞻、缓冲三项各自的毫秒账写在配置
  注释里，避免下一次``为了画质''顺手把 lookahead 调大。
\end{enumerate}
\end{guideline}

\begin{bestpractice}[title={码控参数面的四条可执行纪律}]
\begin{enumerate}
\item 额度写成三元组（平均、峰值、突发时长），并且始终用产物的逐帧字节数反证，
  不要相信参数打印。
\item 跨编码器不比较参数数值；把目标定在质量指标或主观评审上，让码控去解出码率。
\item 每帧额度 $=$ 峰值速率除以帧率，这一行算术要出现在代码注释或配置文档里——
  它是判断``这个 GOP 在这个带宽下是否现实''的唯一快速手段。
\item 两遍编码只在素材与参数完全冻结时有意义；任何一次尺寸、帧率或 preset 变化都要
  重新跑第一遍，并把统计文件与产物一起归档。
\end{enumerate}
\end{bestpractice}

% ============================================================
% ---- frag_ch08.tex ----
\chapter{解码与硬件加速：把帧留在设备侧}
\label{ch:ch08-decode-hw}

这一章走 \Lib{FFmpeg} 的硬件解码路径：先弄清两种硬件上下文各自的分工，
再对比 DXVA2、QSV、 VA-API、NVDEC、VideoToolbox 五个后端的 API 形状与
支持边界，最后算清楚''什么时候必须回主内存''这笔账，以及解码输出与
显示渲染侧的接口在哪里交接。纹理互操作、DMABUF 送合成的细节属于渲染
侧，本章只把边界画清楚。

\section{为什么解码器的输出不一定在主内存里}

软件解码路径下，\Tp{AVFrame::data} 指向的是内存中的 YUV 平面，拿到
帧等于拿到字节。硬件解码改变了这一点：解码器把重建帧写进 GPU（或
SoC）的显存表面，\Tp{AVFrame::data} 里装的只是句柄或描述符，CPU 直接
解引用会段错误。为了在类型系统里区分这两种帧，libavutil 给每个后端
设了一个''硬件像素格式''：

\begin{itemize}
  \item \Tp{AV\_PIX\_FMT\_VAAPI}、\Tp{AV\_PIX\_FMT\_QSV}：Linux 侧
        Intel/AMD/NVIDIA 驱动共享的两个成员；
  \item \Tp{AV\_PIX\_FMT\_DXVA2\_VLD}、\Tp{AV\_PIX\_FMT\_D3D11}：
        Windows 两代图形接口的输出；
  \item \Tp{AV\_PIX\_FMT\_CUDA}：NVIDIA 解码与计算共享内存的格式；
  \item \Tp{AV\_PIX\_FMT\_VIDEOTOOLBOX}：Apple 系统帧，data 指针里
        携带的就是系统像素缓冲；
  \item \Tp{AV\_PIX\_FMT\_VDPAU}：NVIDIA Linux 的老接口，逐步退场，
        但仍能在存量机器上碰到。
\end{itemize}

这些格式不描述内存布局，只声明''这帧归哪个后端管''。帧的实际字节
形态记录在同池的 \Tp{sw\_format} 字段里——下载回主内存后按它展开。

解码器能不能输出硬件格式，由能力枚举决定。现代 API 是
\Fn{avcodec\_{\Bk}get\Bk\_{\Bk}hw\_{\Bk}config}：按序号迭代，返回一组
\Tp{AVCodecHWConfig}，每项带 \Tp{pix\_fmt}、\Tp{device\_{\Bk}type} 与
\Tp{methods} 位集。四个位共享前缀
\Tp{AV\_{\Bk}CODEC\_{\Bk}HW\_{\Bk}CONFIG\_{\Bk}METHOD\_{\Bk}}，去掉前缀看语义：
\Tp{HW\_{\Bk}DEVICE\_{\Bk}CTX} 表示挂上设备上下文即可；
\Tp{HW\_{\Bk}FRAMES\_{\Bk}CTX} 表示必须自备帧池；\Tp{INTERNAL} 是解码器内置
设备、无需外挂；\Tp{AD\_HOC} 则是''能跑但没有正规协商通道''的旧式
路径。打开编码器或解码器前先迭代这张表，是硬件路径唯一可靠的探测
手段。

\begin{guideline}[title={硬件解码路径的三个判断点}]
  一，解码器是否声明了目标硬件格式（迭代 \Fn{avcodec\_get\Bk\_hw\_config}）；
  二，接入方式是设备上下文还是必须显式帧池（看 methods 位）；
  三，输出格式能否满足下游——若下游要求下载，硬件解码省下的 CPU
  会被 PCIe 搬运和驱动开销吃掉一半以上。三个点都确认，才谈得上加速。
\end{guideline}

\section{设备上下文与帧池：av\_hwdevice\_ctx\_create 的分工}

设备上下文代表''与驱动握过手''：\Fn{av\_{\Bk}hwdevice\_{\Bk}ctx\_{\Bk}create} 按
\Tp{AV\_{\Bk}HWDEVICE\_{\Bk}TYPE\_{\Bk}VAAPI} 这样的类型值打开后端，成功返回一个
\Tp{AVBufferRef}，内部是 \Tp{AVHWDevice{\Bk}Context}。参数里的 location
随后端而变——VA-API 与 QSV（Linux）填渲染节点如
\Tp{/dev/{\Bk}dri/{\Bk}renderD128}，CUDA 与 NVDEC 可以填 GPU 序号，DXVA2 与
D3D11VA 填适配器编号或留空用首选卡，VideoToolbox 通常传 \Tp{NULL}。
把引用挂到 \Tp{AVCodec{\Bk}Context::hw\_{\Bk}device\_{\Bk}ctx}，解码器就能在设备
上分配表面并自行推导一个隐式池：池的硬件格式、软件格式、宽高取自
码流参数。

帧池上下文 \Tp{AVHWFrames\Bk Context} 是这层隐式行为的显式化：先
\Fn{av\_{\Bk}hwframe\Bk\_{\Bk}ctx\_{\Bk}alloc}（传入设备引用），填 \Tp{format}、
\Tp{sw\_format}、\Tp{width}、\Tp{height}、\Tp{initial\_{\Bk}pool\_{\Bk}size}，
再 \Fn{av\_{\Bk}hwframe\_{\Bk}ctx\_{\Bk}init}，最后挂到 \Tp{hw\_{\Bk}frames\_{\Bk}ctx}。
两种上下文的分工可以概括为：设备上下文回答''用哪块卡、哪个驱动''，
帧池回答''表面长什么样、有多少个''。滤镜图与硬件编码器之间传帧时，
池子同时是背压机制——表面耗尽，产出端就得等，命令行里
\Tp{-extra\_{\Bk}hw\_{\Bk}frames} 预留的正是这点余量。走
\Tp{AV\_{\Bk}CODEC\_{\Bk}HW\_{\Bk}CONFIG\_{\Bk}METHOD\_{\Bk}HW\_{\Bk}FRAMES\_{\Bk}CTX} 路线的解码器，
帧池是硬性要求；其余后端在分辨率中途变化或需要指定 \Tp{sw\_format} 时，
也应当显式建池。

完整流程如下：打开 VA-API 设备，选带硬件后缀的解码器，挂设备，
收帧，只在必须看字节时才下载。

\begin{lstlisting}[style=cstyle,language=C,
  caption={打开硬件解码器并取回帧的完整流程},
  label={lst:ch08-hwdecode}]
#include <libavcodec/avcodec.h>
#include <libavformat/avformat.h>
#include <libavutil/hwcontext.h>
#include <errno.h>

static AVBufferRef *hw_device_ref = NULL;

/* 以 H.264 为例；其他码流换掉解码器名字即可 */
int open_hw_decoder(const char *url, AVFormatContext **out_fmt,
                    AVCodecContext **out_ctx, int *out_idx)
{
    AVFormatContext *fmt = NULL;
    int err = av_hwdevice_ctx_create(&hw_device_ref,
              AV_HWDEVICE_TYPE_VAAPI, "/dev/dri/renderD128",
              NULL, 0);
    if (err < 0)
        return err;          /* 渲染节点不存在或驱动异常 */

    if (avformat_open_input(&fmt, url, NULL, NULL) < 0 ||
        avformat_find_stream_info(fmt, NULL) < 0)
        return -1;

    int idx = av_find_best_stream(fmt, AVMEDIA_TYPE_VIDEO,
                                  -1, -1, NULL, 0);
    const AVCodec *dec = avcodec_find_decoder_by_name("h264_vaapi");
    if (!dec) {
        /* 没有硬件变体：退回软解加 hw_device_ctx 的旧式路线，
           或干脆报错。这里演示显式选择。 */
        return AVERROR(ENOSYS);
    }

    AVCodecContext *ctx = avcodec_alloc_context3(dec);
    avcodec_parameters_to_context(ctx, fmt->streams[idx]->codecpar);
    ctx->hw_device_ctx = av_buffer_ref(hw_device_ref);
    ctx->pix_fmt = AV_PIX_FMT_VAAPI;   /* 显式声明要硬件帧 */
    if (avcodec_open2(ctx, dec, NULL) < 0)
        return -1;

    *out_fmt = fmt; *out_ctx = ctx; *out_idx = idx;
    return 0;
}

/* 送一包并取帧：hw_frame 可能始终留在 GPU，sw_frame 只在下载后有效 */
int pump_frames(AVCodecContext *ctx, AVPacket *pkt, AVFrame *hw)
{
    AVFrame *sw = av_frame_alloc();
    int ret = avcodec_send_packet(ctx, pkt);
    if (ret == AVERROR(EAGAIN))
        return ret;        /* 队列满：先收帧再送包 */
    while (ret >= 0) {
        ret = avcodec_receive_frame(ctx, hw);
        if (ret == AVERROR(EAGAIN) || ret == AVERROR_EOF)
            break;
        if (hw->format == AV_PIX_FMT_VAAPI &&
            need_cpu_pixels(hw)) {
            /* 唯一的下载点：GPU 表面拷回主内存 */
            if (av_hwframe_transfer_data(sw, hw, 0) < 0)
                return AVERROR(ENOMEM);
            /* 此刻 sw->data[] 才是 CPU 可见的 NV12 平面 */
            sw->pts = hw->pts;
        }
    }
    av_frame_free(&sw);
    return ret;
}
\end{lstlisting}

\subsection*{协商代码的形状}

生产代码不应硬编码''h264\_vaapi 一定存在''。稳妥做法是迭代解码器
能力，挑一个本机可用的设备类型：

\begin{lstlisting}[style=cstyle,language=C,
  caption={迭代 \#avcodec-get-hw-config 挑选可用后端},
  label={lst:ch08-negotiate}]
const AVCodecHWConfig *cfg;
enum AVPixelFormat want = AV_PIX_FMT_NONE;
for (int i = 0; (cfg = avcodec_get_hw_config(dec, i)); i++) {
    if (!(cfg->methods &
          AV_CODEC_HW_CONFIG_METHOD_HW_DEVICE_CTX))
        continue;
    /* 逐个尝试打开设备：能握手成功的才算数 */
    if (av_hwdevice_ctx_create(&ref, cfg->device_type,
                               NULL, NULL, 0) >= 0) {
        want = cfg->pix_fmt;
        break;
    }
}
if (want == AV_PIX_FMT_NONE)
    return AVERROR(ENOSYS);      /* 退回软件解码路径 */
\end{lstlisting}

listing 里的 \verb|need_cpu_pixels| 是示意函数：代表''下游是否要求
字节''的判断，真实项目里往往是像素格式白名单。

\section{五个后端：一张支持矩阵}

\small
\begin{longtable}{>{\raggedright\arraybackslash}p{0.12\textwidth}
                  >{\raggedright\arraybackslash}p{0.12\textwidth}
                  >{\raggedright\arraybackslash}p{0.19\textwidth}
                  >{\raggedright\arraybackslash}p{0.15\textwidth}
                  >{\raggedright\arraybackslash}p{0.12\textwidth}
                  >{\raggedright\arraybackslash}p{0.10\textwidth}}
\caption{五个硬件后端的编解码支持矩阵（FFmpeg 接口，截至 2026）}\label{tab:ch08-hwmatrix} \\
\toprule
后端 & 平台 & 解码格式 & 编码能力 & 位深 & 4:4:4 \\
\midrule
\endfirsthead
\toprule
后端 & 平台 & 解码格式 & 编码能力 & 位深 & 4:4:4 \\
\midrule
\endhead
\midrule
\endfoot
\bottomrule
\endlastfoot
DXVA2 & Windows 7+ & h264\Bk/hevc\Bk/vc1\Bk/mpeg2\Bk/vp9（驱动而定） &
  无（FFmpeg 侧解码专用） & 8\Bk/10 & 否 \\
D3D11VA（同族后继） & Win\Bk/ 8+ & 接近 DXVA2\Bk 再加 av1（新驱动） &
  无 & 8\Bk/10 & 否 \\
QSV & Win\Bk/ Linux（Intel） & h264\Bk/hevc\Bk/vp8\Bk/vp9\Bk/mpeg2\Bk/jpeg\Bk/av1（新平台） &
  h264\Bk/hevc\Bk/vp9\Bk/av1\Bk/mjpeg（代际而定） & 8\Bk/10 & 基本无 \\
VA-API & Linux & h264\Bk/hevc\Bk/vp9\Bk/av1\Bk/jpeg\Bk/mpeg2（看驱动） &
  h264\Bk/hevc\Bk/vp9\Bk/av1\Bk/mjpeg\Bk/vp8（看驱动） & 8\Bk/10 & 视驱动而定 \\
NVDEC & Win\Bk/ Linux（NVIDIA） & h264\Bk/hevc\Bk/vp9\Bk/av1（Turing\Bk+）\Bk/mpeg2\Bk/vc1\Bk/jpeg &
  无（编码走 NVENC 另一条线） & 8\Bk/10\Bk，hevc 部分 12 & 否 \\
Video\Bk Toolbox & macOS\Bk/ iOS & h264\Bk/hevc 为主\Bk，av1 等新系统才有条目 &
  h264\Bk/hevc\Bk 成熟\Bk，prores 等 & 8\Bk/10 & 基本无 \\
\end{longtable}
\normalsize

矩阵之外记住三条口径。其一，''格式列''写的是 FFmpeg 侧真实存在的
解码器条目（\Tp{h264\_vaapi}、\Tp{hevc\_nvdec}、\Tp{av1\_qsv} 等），
实际可用性仍受 GPU 代际与驱动版本钳制，AV1 尤其如此。其二，DXVA2
与 D3D11VA 是同一台 Windows 机器上的两代接口，能力表相似，D3D11VA
的帧还能直接被系统纹理共享，新代码默认选后者。其三，VideoToolbox
是系统级调度，应用不指定硬件；它输出的帧携带系统像素缓冲，与显示
侧几乎零成本衔接，这也是 Apple 平台播放器偏爱它的原因。

\section{回到 CPU 的那一刻：下载与接口边界}

\Fn{av\_hwframe\_transfer\_data} 是唯一正规的''搬运工''：按源帧的
后端查池子，把 GPU 表面（如 VA 表面、D3D 纹理、CUDA 内存）解码成
主内存平面，或反向上传。参数顺序是目的帧在前。一次下载的代价是
一帧的 PCIe 往返加驱动同步：1080p NV12 约 3\,MB，4K 约 12\,MB，
几十路并发时这本身就是瓶颈。\textbf{一旦下载，这条链路就不再是
零拷贝}——硬件解码剩下的收益只有''把重建的算力挪出 CPU''这一条，
必须按这个口径记账。

典型的下载理由有四类：滤镜链里有后端没有对应硬件变体的成员（大量
音频滤镜与部分绘制类）；需要把像素喂给 OpenCV、推理引擎或自研算法；
目标像素格式（如 4:4:4 或 RGB）超出硬件池的 sw\_format 范围；调试
或校验需要逐字节可读。若只是显示，不该下载——接口边界上还有更省
的通道：

\begin{itemize}
  \item Apple：硬件帧自带系统像素缓冲，直接交 Metal 层或显示原语；
  \item Windows：D3D11 纹理句柄可从 \Tp{data} 指针取出，
        \Fn{av\_hwframe\_map} 可把一帧映射成另一后端的视图；
  \item Linux：VA-API 表面可以经驱动接口导出为 DMABUF 文件描述符，
        FFmpeg 里也有以 \Tp{AV\_PIX\_FMT\_DRM\_PRIME} 描述帧的通道，
        句柄背后是 \Tp{AVDRMFrameDescriptor}。
\end{itemize}

导出与互操作的消费端属于渲染侧的课题，本章到交接面为止；但设计
管线时''帧在哪一侧落地''必须在画图上先回答，而不是在编码时补。

\begin{warning}[title={硬件解码的三个高频误区}]
  一是把''用了硬件解码''当成''零拷贝''：中间夹一次 transfer 就已经
  不是；二是池子尺寸想当然：硬件编码器与显示队列各占表面，池耗尽
  时整条链停在等待上，需要显式建池并用 \Tp{initial\_pool\_size} 与
  \Tp{-extra\_hw\_frames} 预留；三是假设能力表等于本机现实：同一
  型号 GPU 在不同驱动下支持列表不同，上线前必须用协商迭代实测。
\end{warning}

\section{命令行全链路与解码错误恢复}

命令行是 API 的镜子，也是最快捷的验证工具。一组真实拼写：

\begin{lstlisting}[style=shellstyle,
  caption={VAAPI/QSV/NVDEC 全链路命令与错误检测开关},
  label={lst:ch08-cli}]
# VAAPI 全链路：设备、硬件解码、硬件缩放、硬件编码
ffmpeg -vaapi_device /dev/dri/renderD128 \
       -hwaccel vaapi -hwaccel_output_format vaapi \
       -i input.mp4 \
       -vf 'scale_vaapi=w=1280:h=720:format=nv12' \
       -c:v hevc_vaapi -b:v 2000k out_hevc.mp4

# 同一设备只做解码，帧留在 GPU，测吞吐
ffmpeg -vaapi_device /dev/dri/renderD128 \
       -hwaccel vaapi -hwaccel_output_format vaapi \
       -i input.mp4 -f null -

# QSV 路线：按名选硬件解码器，vpp_qsv 缩放，hevc_qsv 编码
ffmpeg -hwaccel qsv -c:v h264_qsv -i input.mp4 \
       -vf 'vpp_qsv=w=1280:h=720' \
       -c:v hevc_qsv -global_quality 32 out_qsv.mp4

# NVDEC 路线：帧留在 CUDA 内存，不解回主内存
ffmpeg -hwaccel cuda -hwaccel_output_format cuda \
       -c:v hevc_nvdec -i input.mp4 -f null -

# 错误检测：校验和、语法、越界任一命中即报错，不静默掩盖
ffmpeg -err_detect crccheck+bitstream+buffer+explode \
       -i broken.ts -f null -
\end{lstlisting}

\Tp{-hwaccel\_output\_format} 是链路''留在设备侧''的开关：不设它，
解码帧会被立刻下载；设成 \Tp{vaapi} 或 \Tp{cuda} 后，滤镜必须用
对应硬件变体接住，否则图配置阶段就报格式不支持。

硬件路径的错误恢复与软解同源但有额外一层。\Tp{-err\_detect} 的
字符串位（crccheck、bitstream、buffer、careful、compliant、
aggressive、explode）控制解码器的错误判定强度：careful/compliant/
aggressive 是语法与参数的严格度递增，explode 把''报告错误''升级成
''返回错误''，批处理转码用它防止静默产出坏帧。错误掩盖的强度由
\Tp{-ec} 调节。丢包或时间戳跳变后重建参考链的正规动作是
\Fn{avcodec\_flush\_buffers}：清空帧线程队列，从下一个关键帧重新
起链——帧线程数越大，flush 后需要喂的关键帧前置包越多。硬件后端
还有自己的一层：驱动返回的表面级错误（DXVA2 的失败位图、
VideoToolbox 的 per-frame status）在 FFmpeg 内映射回普通错误码，
但行为随驱动而定，压测时要把''单帧解码失败''当作常态路径测试。
时间戳损坏的更多容错手段在章 \ref{ch:ch10-mux-sync} 展开。

\begin{bestpractice}[title={硬件加速路径选型清单}]
  先迭代能力表确认解码器声明了目标硬件格式，再决定接入方式是设备
  上下文还是显式帧池；整条链路（解码、缩放、编码）尽量用同一后端，
  跨后端就是一次映射或下载；下游只要显示与再编码，就不要碰
  \Fn{av\_hwframe\_transfer\_data}；最后用 \Tp{-hwaccel\_output\_format}
  把链路钉死在设备侧，跑一遍吞吐，确认收益来自算力而不是错觉。
\end{bestpractice}

% ============================================================
% ---- frag_ch09.tex ----
\chapter{缩放、重采样与滤镜图}
\label{ch:ch09-filter-resample}

这一章覆盖 \DomVid\ 与 \DomAud\ 两条链路的''变形''环节：像素域的尺寸
与格式转换（\Lib{libswscale}）、采样率与声道布局转换
（\Lib{libswresample}），以及把任意处理串成管线的 \Lib{libavfilter}
图模型。三个库都在帧边界上做引用计数与格式协商，理解它们的代价
结构，才能在 \DomRT\ 与 \DomBatch\ 两种上下文里做出正确取舍。

\section{libswscale：算法族与色彩语义}

\Fn{sws\_get\Bk\_Context} 按源与目的的宽高、像素格式创建一个缩放
上下文，\Fn{sws\_scale} 逐切片执行；两个函数只处理''几何 + 格式''，
色彩语义要靠第三个参数——算法标志位。算法族从便宜的
\Tp{SWS\_POINT}（最近邻）、\Tp{SWS\_BILINEAR}、\Tp{SWS\_AREA}
（缩小时的面积平均）到高质量的 \Tp{SWS\_BICUBIC}、\Tp{SWS\_GAUSSIAN}、
\Tp{SWS\_LANCZOS}、\Tp{SWS\_SPLINE} 排开；\Tp{SWS\_ACCURATE\_RND}
强制精确舍入、\Tp{SWS\_BITEXACT} 要求跨平台逐位一致，两者都是拿
吞吐换确定性。矩阵系数常量（\Tp{SWS\_CS\_} 前缀一族，对应 601/709/2020
等）决定 YUV 与 RGB 互转时用哪组权重。

同一对格式、同一个算法，结果仍可能和命令行不一致——原因在色彩
语义的取法。裸调 swscale 时，范围与矩阵按传入常量与像素格式默认值
处理，它不会去读码流里的标记；而 \Tp{-vf scale=...} 滤镜会读取帧的
\Tp{color\_space}、\Tp{color\_range} 属性做协商，并暴露
\Tp{in\_range}/{\Bk}\Tp{out\_range} 与 \Tp{in\_color\_matrix}/
\Tp{out\_color\_matrix} 选项显式覆盖。BT.709 的素材按 601 系数转
RGB，色偏肉眼可辨——''结果不同''多半就是这一层。

\begin{compare}[title={scale 滤镜与直调 sws 的差别}]
  直调 \Fn{sws\_get\Bk\_Context} 只认你传进去的参数，适合管线已
  固定、语义自己保证的场合；scale 滤镜则参与图的格式协商、能读帧
  属性、支持表达式（如 \Tp{scale=iw/2:ih}）、有切片多线程，还能
  挂硬件变体（\Tp{scale\_vaapi}、\Tp{scale\_cuda}）。离线转码优先
  用滤镜，嵌入式里直调库才划算。
\end{compare}

\section{libswresample：引擎、声道布局与延迟补偿}

音频侧的对应物是 \Lib{libswresample}：采样率、样本格式、声道布局
三个维度的任意组合转换。现代入口是 \Fn{swr\_alloc\_set\Bk\_opts2}
（取代旧的两级指针前的便利函数，旧版已在长期弃用名单上），传入输出
与输入的 \Tp{AVChannelLayout}、\Tp{AV\_SAMPLE\_FMT\_} 前缀格式与采样率，
拿到 \Tp{SwrContext} 再 \Fn{swr\_init}。逐帧转换用
\Fn{swr\_convert}，或直接 \Fn{swr\_convert\_frame} 让布局与格式从
\Tp{AVFrame} 上自取。

声道布局是这一层的重点：down-mix 有定义好的系数（5.1 到立体声、
立体声到单声道取加权平均），up-mix 则是按布局补声道；两者的矩阵都
可以用 \Fn{swr\_{\Bk}set\_{\Bk}matrix}/{\Bk}\Fn{swr\_{\Bk}get\_{\Bk}matrix} 读改。重采样
引擎可选——多相滤波器是默认，此外还有线性插值等低成本档，通过
context 的 AVOptions 传字典调窗口长度、截断点与精度。异步场景里
\Fn{swr\_{\Bk}get\_{\Bk}delay} 报出内部缓冲折算成输入率还欠多少样本，
\Fn{swr\_{\Bk}set\_{\Bk}compensation} 可以主动吞掉或插入样本来追时钟——
声卡漂移的补偿正是靠这一对接口。

\section{sws 与 swr 的组合调用}

视频帧与音频帧各走各的管线，但工程上常要同一函数成对处理。下面
一次给出两条腿的最小可用形态：

\begin{lstlisting}[style=cstyle,language=C,
  caption={sws 缩放一帧与 swr 重采样一帧的组合用法},
  label={lst:ch09-swsswr}]
#include <libswscale/swscale.h>
#include <libswresample/swresample.h>
#include <libavutil/channel_layout.h>
#include <libavutil/opt.h>

/* 视频：1920x1080 yuv420p -> 1280x720 yuv420p，lanczos */
int scale_frame(AVFrame *in, AVFrame *out)
{
    struct SwsContext *sws = sws_getContext(
        in->width, in->height, (enum AVPixelFormat)in->format,
        out->width, out->height, (enum AVPixelFormat)out->format,
        SWS_LANCZOS | SWS_ACCURATE_RND, NULL, NULL, NULL);
    if (!sws)
        return -1;
    out->format = AV_PIX_FMT_YUV420P;
    out->width = 1280; out->height = 720;
    if (av_frame_get_buffer(out, 32) < 0) {
        sws_freeContext(sws);
        return -1;
    }
    int done = sws_scale(sws, (const uint8_t * const *)in->data,
                         in->linesize, 0, in->height,
                         out->data, out->linesize);
    sws_freeContext(sws);
    return done == in->height ? 0 : -1;
}

/* 音频：立体声 fltp 44100 -> 单声道 s16 48000 */
int resample_frame(AVFrame *in, AVFrame *out, SwrContext **pswr)
{
    if (!*pswr) {
        AVChannelLayout in_lay = AV_CHANNEL_LAYOUT_STEREO;
        AVChannelLayout out_lay = AV_CHANNEL_LAYOUT_MONO;
        if (swr_alloc_set_opts2(pswr, &out_lay, AV_SAMPLE_FMT_S16,
                                48000, &in_lay, AV_SAMPLE_FMT_FLTP,
                                44100, 0, NULL) < 0)
            return -1;
        /* 引擎与截断点等走 AVOptions 再调 */
        av_opt_set(*pswr, "cutoff", "0.95", 0);
        if (swr_init(*pswr) < 0)
            return -1;
    }
    /* 欠了多少输入样本：异步补偿的读数口 */
    int64_t lat = swr_get_delay(*pswr, in->sample_rate);
    (void)lat;
    out->sample_rate = 48000;
    out->format = AV_SAMPLE_FMT_S16;
    av_channel_layout_default(&out->ch_layout, 1);
    out->nb_samples = (int)(av_rescale_q(in->nb_samples,
            (AVRational){1, in->sample_rate},
            (AVRational){1, 48000}) + 16);
    if (av_frame_get_buffer(out, 0) < 0)
        return -1;
    int n = swr_convert(*pswr, out->data, out->nb_samples,
                        (const uint8_t **)in->data, in->nb_samples);
    if (n < 0)
        return n;
    out->nb_samples = n;
    out->pts = in->pts;   /* 时基同为 1/采样率，帧号按输出率重算 */
    return 0;
}
\end{lstlisting}

两处易错点：\Fn{sws\_scale} 返回的是''已处理的源行数''而非错误码，
凑不满总行数就是部分失败；重采样后的 pts 若沿用输入帧号，时间轴会
被拉长或压缩——严格做法是按两套采样率换算。

\section{libavfilter：图模型与帧引用}

滤镜图的数据结构是''节点 + 带名字的 pad + link''：每个滤镜实例
（\Tp{AVFilter{\Bk}Context}）暴露若干输入/输出 pad，pad 之间建立
\Tp{AVFilter{\Bk}Link}，链两端在配置期跑格式协商（各滤镜声明支持的
格式集合，交集决定链上实际格式）。图的两端由系统滤镜钉死：入口
\Tp{buffer}/{\Bk}\Tp{abuffer} 用参数字符串声明尺寸、像素格式、时基或
采样率与布局，出口 \Tp{buffersink}/{\Bk}\Tp{abuffersink} 收帧。

代码上最常见的是''手工建端点 + 字符串描述中段''：

\begin{lstlisting}[style=cstyle,language=C,
  caption={buffersrc 到 buffersink 的图装配与逐帧送取},
  label={lst:ch09-graph}]
#include <libavfilter/avfilter.h>
#include <libavfilter/buffersink.h>
#include <libavfilter/buffersrc.h>

static AVFilterGraph *graph;
static AVFilterContext *buffersrc_ctx, *buffersink_ctx;

int build_graph(const char *desc, AVFrame *probe, AVRational tb)
{
    char args[256];
    snprintf(args, sizeof(args),
        "video_size=%dx%d:pix_fmt=%d:time_base=%d/%d:pixel_aspect=1/1",
        probe->width, probe->height, probe->format, tb.num, tb.den);

    graph = avfilter_graph_alloc();
    int err = avfilter_graph_create_filter(&buffersrc_ctx,
              avfilter_get_by_name("buffer"), "in",
              args, NULL, graph);
    if (err < 0)
        return err;
    err = avfilter_graph_create_filter(&buffersink_ctx,
              avfilter_get_by_name("buffersink"), "out",
              NULL, NULL, graph);
    if (err < 0)
        return err;

    AVFilterInOut *outputs = avfilter_inout_alloc();
    AVFilterInOut *inputs  = avfilter_inout_alloc();
    outputs->name = av_strdup("in");
    outputs->filter_ctx = buffersrc_ctx;
    outputs->pad_idx = 0;
    inputs->name = av_strdup("out");
    inputs->filter_ctx = buffersink_ctx;
    inputs->pad_idx = 0;
    /* 命名 pad 的复杂描述用 parse_ptr；无命名链可用 parse2 */
    err = avfilter_graph_parse_ptr(graph, desc, &inputs,
                                   outputs, NULL);
    if (err >= 0)
        err = avfilter_graph_config(graph, NULL);
    return err;
}

int push_pull(AVFrame *f, AVFrame *out)
{
    /* 默认会吃掉引用：想保住 f 就带 KEEP_REF */
    int err = av_buffersrc_add_frame_flags(buffersrc_ctx, f,
              AV_BUFFERSRC_FLAG_KEEP_REF);
    if (err < 0)
        return err;
    return av_buffersink_get_frame(buffersink_ctx, out);
}
\end{lstlisting}

帧在图内以引用计数传递：buffersrc 收到帧后按链上格式决定是否需要
转换，转换即换一次引用；下游每级滤镜按需克隆或改写，\Fn{av\_frame\_unref}
在消费完成处释放。图级并发上，命令行有 \Tp{-filter\_threads} 与
\Tp{-filter\_complex\_threads} 控制线程数，切片式滤镜（如 scale、
noise 一族）在单帧内多线程，而帧间流水由调用方的送帧节奏决定——
图的''线程模型''更多体现在单滤镜切片，而不是自动并行整条链。

\section{音频 DSP 链的常见件}

\Tp{loudnorm} 实现 EBU R128 双口径目标（\Tp{I} 积分响度、\Tp{TP}
真峰值、\Tp{LRA} 范围）：单遍动态模式逐帧调整，可流式但有泵声
风险；两遍线性模式第一遍输出 JSON 测量值，第二遍以
\Tp{linear=true} 加 \Tp{measured\_} 前缀的一组测量值做一次全局增益，要求把输出
提到 192\,kHz 才能保住线性相位——离线交付用它。降噪一族里，
\Tp{arnndn} 加载 RNNoise 类的预训练模型文件（\Tp{model=xxx.rnnn}），
\Tp{afftdn} 走频谱减法门道，\Tp{anlmdn} 是非局部均值；它们的
C 接口形态没有特殊之处——都注册成滤镜，靠图描述串或滤镜参数 API
接入，代价模型见下面的表。\Tp{atempo} 变速不变调，单级历史范围
0.5--2.0，超界靠两级串联实现。

\begin{lstlisting}[style=shellstyle,
  caption={filter\_complex 图：scale、setpts 与 amix 的混排},
  label={lst:ch09-fcomplex}]
# 缩放 + 时间戳归零 + 把背景音垫在解说下面
ffmpeg -i host.mp4 -i bed.wav -filter_complex \
  "[0:v]scale=1280:720:flags=lanczos,setpts=PTS-STARTPTS[v]; \
   [0:a]asetpts=PTS-STARTPTS[host]; \
   [1:a]atempo=1.2,afade=t=in:d=2[bed]; \
   [host][bed]amix=inputs=2:duration=first:dropout_transition=2[a]" \
  -map "[v]" -map "[a]" -c:v libx264 -c:a aac out.mp4

# 两遍 loudnorm：第一遍测量，第二遍线性校正（输出须过采样）
ffmpeg -i in.wav -af loudnorm=I=-16:TP=-1.5:LRA=11:print_format=json \
       -f null -
ffmpeg -i in.wav -af 'loudnorm=I=-16:TP=-1.5:LRA=11:linear=true: \
       measured_I=-16.23:measured_TP=-1.98:measured_LRA=11.4: \
       measured_thresh=-26.5' -ar 192000 out.wav

# 语音降噪：高通滤掉低频隆隆声，再进 RNNoise 类模型
ffmpeg -i noisy.wav -af 'highpass=f=200,arnndn=model=denoise.rnnn' \
       clean.wav
\end{lstlisting}

\section{实时与离线：两种代价账本}

\small
\begin{longtable}{>{\raggedright\arraybackslash}p{0.13\textwidth}
                  >{\raggedright\arraybackslash}p{0.24\textwidth}
                  >{\raggedright\arraybackslash}p{0.12\textwidth}
                  >{\raggedright\arraybackslash}p{0.18\textwidth}
                  >{\raggedright\arraybackslash}p{0.13\textwidth}}
\caption{常用滤镜的计算代价、状态内存与可流式性}\label{tab:ch09-filters} \\
\toprule
滤镜 & 用途 & CPU & 内存与延迟 & 可逐帧流式 \\
\midrule
\endfirsthead
\toprule
滤镜 & 用途 & CPU & 内存与延迟 & 可逐帧流式 \\
\midrule
\endhead
\midrule
\endfoot
\bottomrule
\endlastfoot
scale\Bk/bilinear & 常规缩放 & 低 & 无状态，单帧 & 是 \\
scale\Bk/lanczos & 高质量缩放 & 中高 & 卷积窗口行缓存 & 是 \\
scale\_vaapi 等 & GPU 缩放 & 卸载到 GPU & 占硬件表面 & 是 \\
fps & 帧率转换 & 低--中 & 缓冲不超一帧间隔 & 是 \\
setpts & 改写时间戳 & 极低 & 零状态 & 是 \\
overlay & 画中画/水印 & 中 & 持有两路输入帧 & 是 \\
drawtext & 文字渲染 & 中 & 字形缓存常驻 & 是 \\
aresample & 采样率转换 & 低--中 & 滤波器拖尾 & 是 \\
atempo & 变速不变调 & 中 & 秒级窗口，固定延迟 & 是 \\
loudnorm 动态 & 响度归一 & 中 & 前瞻窗（默认秒级） & 是 \\
loudnorm 线性 & 全局响度校正 & 低 & 需整段测量值 & 否 \\
arnndn & RNN 降噪 & 中高 & 模型状态，帧级延迟 & 是 \\
afftdn & 频谱降噪 & 中 & FFT 窗口（数十毫秒） & 是 \\
amix & 多路混音 & 低 & 取最长输入缓冲 & 是 \\
concat & 音画拼接 & 低 & 近乎零 & 是 \\
\end{longtable}
\normalsize

\DomRT\ 的读法：只选''可逐帧流式''的行，并把每帧成本压进预算——
缩放用 bilinear 或 GPU 变体，响度用动态模式，降噪选 arnndn 这类
帧级延迟的实现；任何需要''先看全片''的件（线性 loudnorm、两遍
统计）都不允许出现在直播链路里。\DomBatch\ 反过来：两遍流程、
lanczos/spline、精确舍入都可以上，唯一要记账的是墙钟时间与内存
峰值。同一张表在两种上下文里的''贵''完全不同——实时怕的是方差，
离线怕的是总量。

\begin{bestpractice}[title={管线选型速查}]
  几何变换优先问硬件后端有没有对应变体；色彩语义从帧属性走，别在
  调用点硬编码矩阵；音频链里把无状态件（重采样、高通）放最前、
  有状态件（atempo、loudnorm）放后，方便单测；实时链路用表
  \ref{tab:ch09-filters} 里''可逐帧流式=是''的子集做白名单，离线
  链路才允许引入整段测量类滤镜。
\end{bestpractice}

% ============================================================
% ---- frag_ch10.tex ----
\chapter{封装、时间基与容错}
\label{ch:ch10-mux-sync}

这一章把镜头从帧拉回到包：容器格式在文件与网络上''装''了什么——
索引结构、时间戳体系、参数集的安放位置，以及在丢包、截断、缺索引
这三类事故里，哪些损失是可恢复的。所有换算都围绕
\Lib{FFmpeg} 的 \Fn{av\_rescale\_q} 展开，示例拼写来自
\Tp{avformat}、\Tp{ffprobe} 的真实接口。

\section{容器骨架：box 树、EBML 与 188 字节}

MP4（ISO BMFF 家族）是一棵自述的盒树：\Tp{ftyp} 报品牌，
\Tp{mdat} 装裸数据，\Tp{moov} 装一切索引——\Tp{mvhd} 给全片时长与
timescale，每个 \Tp{trak} 内 \Tp{mdhd} 声明流时间刻度，\Tp{stbl}
样本表是核心：\Tp{stsd} 记录编解码描述（含参数集盒，如
\Tp{avcC}/{\Bk}\Tp{hvcC}），\Tp{stts} 把样本数映射到时长，\Tp{stss} 列
同步样本（可跳读的位置），\Tp{stsc}/{\Bk}\Tp{stsz}/{\Bk}\Tp{stco} 把样本
换算到文件字节偏移。这棵树写完才能播放，所以默认 \Tp{moov} 在
\Tp{mdat} 之后；\Tp{-movflags +faststart} 把它搬到头部换取边下边播。

分片 MP4（fMP4）换了思路：开头放一个''空 \Tp{moov}''（只声明轨与
默认值，不带样本表），此后每个分片是 \Tp{moof} 加 \Tp{mdat}：
\Tp{moof} 里的 \Tp{traf} 用 \Tp{tfhd}（默认值）、\Tp{tfdt}（解码
时间基准）与 \Tp{trun}（逐样本的时长、偏移、标志）自描述；
\Tp{sidx} 提供分片级索引。写入端不需要回填任何旧字节，天然适合
追加式发布。

Matroska（MKV/WebM）是 EBML：变长前缀整数 + 分层元素 ID。骨架是
\Tp{Segment} 下的 \Tp{SegmentInfo}、\Tp{Tracks} 与一串
\Tp{Cluster}；每个簇带自己的时间码，块（\Tp{Block}/
\Tp{SimpleBlock}）的时标相对簇起算。\Tp{TimecodeScale} 默认一百万
纳秒，即容器时标以毫秒计——这就是 \Tp{mkv} 流的
\Tp{time\_base} 恒为 1/1000000 的原因。检索用的 \Tp{Cues} 在文件
末尾，与 \Tp{moov} 同病：不写完不能快速跳。

MPEG-TS 反其道：没有树，只有等长 188 字节包。包头 4 字节：同步字节
\Tp{0x47}、13 位 \Tp{PID}、适配域标志与 4 位连续计数器；适配域可带
填充与 \Tp{PCR}（基数字段 90 kHz，扩展字段细到 27 MHz），负载是
 elementary 流切片。节目-信息表构成最小索引：\Tp{PAT}（PID 恒为 0）
指到各路 \Tp{PMT}，\Tp{PMT} 声明流类型、 elementary PID 与 PCR PID。
连续计数器逐 PID 独立递增，接收端靠它发现丢包——这是 TS 唯一的
内置完整性检查。

\section{时间基：time\_base 与 av\_rescale\_q}

\Tp{AVRational} 是（分子，分母）对；\Tp{AVStream::time\_base} 声明
该流 \Tp{AVPacket::pts}/{\Bk}\Tp{dts} 的单位，\Tp{AVFrame::pts} 则按
解码器的时基走（打开时用 \Tp{AVCodecContext} 的时间基参数对齐）。
换算只有一个正规口：\Fn{av\_rescale\_q} 按三元组
\Tp{x}、\Tp{src\_tb}、\Tp{dst\_tb} 做有理数换算，涉及取整方向时用
\Fn{av\_rescale\_q\_rnd}。常见刻度：MP4 视频轨默认由帧率推导
（\Tp{-video\_track\_timescale 90000} 可显式钉死），MKV 恒为
1/1000000，TS 的 PTS/DTS 固定 90 kHz。同一帧在不同容器里的
\Tp{pts} 数字因此毫无可比性，比的是 \Tp{pts\_time} 这类换算值。

\begin{guideline}[title={时间基三条军规}]
  一，永远不要在包上带着源流时基的数字直接塞给目的流——先过
  \Fn{av\_{\Bk}rescale\_{\Bk}q}；二，跨''取整歧义''用 \Fn{av\_{\Bk}rescale\_{\Bk}q\_{\Bk}rnd}
  显式定方向，变帧率素材尤其敏感；三，调试时以 \Tp{pts\_time} 的
  秒值为准对齐两套读数，不要对着裸整數吵。
\end{guideline}

\section{B 帧、DTS 与 PTS：两套顺序的账}

显示顺序与解码顺序在 B 帧下分道：包必须按解码顺序写入容器，所以
\Tp{dts} 单调不减是硬要求；\Tp{pts} 可以乱序，两者之差就是重排
延迟。开启双向预测后，流开头的重排缓冲让第一包的 \Tp{dts} 变成
负数——ffprobe 直接可见：

\begin{lstlisting}[style=jsonstyle,
  caption={ffprobe 包视图：\#pts 与 \#dts 的负偏移},
  label={lst:ch10-packets}]
{
  "packets": [
    {
      "codec_type": "video",
      "pts": 0,
      "pts_time": "0.000000",
      "dts": -512,
      "dts_time": "-0.005689",
      "duration": 3000,
      "size": 81234,
      "flags": " B"
    },
    {
      "codec_type": "video",
      "pts": 3000,
      "pts_time": "0.033333",
      "dts": -212,
      "duration": 3000,
      "size": 5411,
      "flags": " P"
    }
  ]
}
\end{lstlisting}

MP4 的 \Tp{stsd} 时代只能表达''非负时间''，于是负的重排偏移被拆成
两半：\Tp{ctts} 盒记录每个样本显示时刻相对解码时刻的正偏移；开头
的负起点由编辑列表 \Tp{elst} 的负 \Tp{media\_time} 表达——
\Tp{-use\_editlist 0} 会禁用这条通道，改用单纯加偏移，代价是负
PTS 的素材在部分播放器上出现首帧跳动。MKV 与 TS 没有编辑列表
概念，负时标由写入端整体平移解决。TS 还有一笔特殊的账：PTS 只有
33 位，90 kHz 下不到 27 小时回绕一次，接收端必须按单调外推处理
跳变，而非当成损坏。

命令行上直接读两套时间基与包序列：

\begin{lstlisting}[style=shellstyle,
  caption={ffprobe 读流级 time\_base 与包级 pts/dts},
  label={lst:ch10-ffprobe}]
# 流级：时间基、平均帧率、起点
ffprobe -v error -show_entries \
  stream=index,codec_name,time_base,avg_frame_rate,start_time \
  -of default=noprint_wrappers=1 in.mp4

# 包级：pts/dts 与其秒值、包类型标志
ffprobe -v error -select_streams v:0 -show_packets \
  -show_entries packet=pts,pts_time,dts,dts_time,duration,flags \
  -of csv=p=0 in.mp4 | head -40

# 整体机读输出，交给脚本核对换算
ffprobe -v error -print_format json \
  -show_format -show_streams in.mp4 > probe.json
\end{lstlisting}

\section{extradata 与参数集：带外一份，带内一份}

解码器要拿到 SPS/PPS（HEVC 还加 VPS）才能起链。存放位置分两派：
MP4 把参数集收进样本描述盒（\Tp{avcC}/{\Bk}\Tp{hvcC}），包体内只有
长度前缀的 NAL 数据；TS/裸流走 Annex-B，起始码 \Tp{00 00 01} 分隔
NAL，参数集随码流带内重复。FFmpeg 内部用
\Tp{AVCodecContext::extradata} 承载''带外''那一份，封入 MP4 时由
muxer 转成盒，转出 TS 时经比特流滤镜换成起始码内嵌——
\Tp{h264\_mp4toannexb}、\Tp{hevc\_mp4toannexb} 做正向，
\Tp{extract\_extradata} 做反向收割；带 \Tp{AVFMT\_FLAG\_AUTO\_BSF}
时这些桥由复用器自动挂上。

带内重复不是冗余：直播接收者从中途入会，等下一个参数集才能起链，
所以广播侧按秒级周期重发 SPS/PPS；转码管线则应把参数集固定放带外
一份、按需带内点缀，否则每份副本都在白占带宽。

\section{可 seek、可流式与残缺文件的账}

写分片 MP4 的最小完整流程——建输出上下文、设 movflags、换算时基、
交错写入：

\begin{lstlisting}[style=cstyle,language=C,
  caption={\#avformat 写出分片 MP4（含时间基换算）},
  label={lst:ch10-fmp4}]
#include <libavformat/avformat.h>
#include <libavutil/mathematics.h>

int write_fmp4(AVFormatContext *in, int idx, const char *path)
{
    AVFormatContext *oc = NULL;
    if (avformat_alloc_output_context2(&oc, NULL, "mp4", path) < 0)
        return -1;

    AVStream *st_in = in->streams[idx];
    AVStream *st_out = avformat_new_stream(oc, NULL);
    if (!st_out)
        return AVERROR(ENOMEM);
    if (avcodec_parameters_copy(st_out->codecpar,
                                st_in->codecpar) < 0)
        return -1;
    st_out->time_base = (AVRational){1, 90000};

    /* 空 moov 开头 + 关键帧处切分片 + trun 用绝对基 */
    av_opt_set(oc->priv_data, "movflags",
        "frag_keyframe+empty_moov+default_base_moof", 0);
    av_opt_set_int(oc->priv_data, "frag_size",
        (int64_t)1 << 20, 0);

    if (!(oc->oformat->flags & AVFMT_NOFILE))
        if (avio_open(&oc->pb, path, AVIO_FLAG_WRITE) < 0)
            return -1;
    if (avformat_write_header(oc, NULL) < 0)
        return -1;

    for (;;) {
        AVPacket *pkt = av_packet_alloc();
        if (av_read_frame(in, pkt) < 0) {
            av_packet_free(&pkt);
            break;
        }
        if (pkt->stream_index != idx) {
            av_packet_free(&pkt);
            continue;
        }
        pkt->stream_index = 0;
        /* 源流时基 -> 输出流时基：dts/pts/duration 一起换算 */
        av_packet_rescale_q(pkt, st_in->time_base,
                            st_out->time_base);
        /* interleaved 写入负责 dts 单调化与重排缓冲 */
        if (av_interleaved_write_frame(oc, pkt) < 0) {
            av_packet_free(&pkt);
            return -1;
        }
        av_packet_free(&pkt);
    }
    av_write_trailer(oc);
    if (oc->pb && !(oc->oformat->flags & AVFMT_NOFILE))
        avio_closep(&oc->pb);
    avformat_free_context(oc);
    return 0;
}
\end{lstlisting}

\Tp{frag\_{\Bk}keyframe} 决定切片只落在关键帧边界，这正是''每片独立
可解''的前提；\Tp{av\_{\Bk}packet\_{\Bk}rescale\_{\Bk}q} 内部走
\Fn{av\_{\Bk}rescale\_{\Bk}q}，顺带把 \Tp{duration} 也换掉；
\Fn{av\_{\Bk}interleaved\_{\Bk}write\_{\Bk}frame} 维护各流的交错顺序并保证
\Tp{dts} 不减，手动 \Fn{av\_{\Bk}write\_{\Bk}frame} 就得自己兜底。

\small
\begin{longtable}{>{\raggedright\arraybackslash}p{0.11\textwidth}
                  >{\raggedright\arraybackslash}p{0.14\textwidth}
                  >{\raggedright\arraybackslash}p{0.15\textwidth}
                  >{\raggedright\arraybackslash}p{0.18\textwidth}
                  >{\raggedright\arraybackslash}p{0.21\textwidth}}
\caption{容器横向对比：低延迟、可 seek、流式友好与生态}\label{tab:ch10-containers} \\
\toprule
容器 & 低延迟可播 & 可 seek & 流式写入友好 & 生态与兼容性 \\
\midrule
\endfirsthead
\toprule
容器 & 低延迟可播 & 可 seek & 流式写入友好 & 生态与兼容性 \\
\midrule
\endhead
\midrule
\endfoot
\bottomrule
\endlastfoot
MP4 整文件 & 差（等 \Tp{moov} 落位） & 优（\Tp{stss}/{\Bk}\Tp{stco} 全表） &
  差（需回填索引） & 通用：点播、编辑、归档首选 \\
fMP4 & 良（分片即出即解） & 片内靠 \Tp{sidx}，跨片靠清单 &
  优（纯追加） & HLS/DASH 事实标准，浏览器 MSE 可喂 \\
MKV & 中（簇粒度约 1 s） & 有 \Tp{Cues} 才优 &
  中（可追加但索引补写） & 本地库与剪辑工具友好，网页弱 \\
WebM & 同 MKV & 同 MKV & 同 MKV & 浏览器原生；WebRTC 录制 \\
MPEG\Bk-TS & 优（188 字节粒度） & 弱（仅 \Tp{PID} 级定位） &
  优（为流而生） & 广播、卫星、RTP 封装底座 \\
HLS 切片 & 取决于切片长度 & 按切片跳 & 优（清单+分片） &
  HTTP 缓存与 CDN 分发 \\
\end{longtable}
\normalsize

容错三本账。丢包：TS 靠连续计数器发现，MP4/fMP4 的分片边界丢失
意味着该 \Tp{trun} 样本数与真实字节不符，此后整片作废——直播端要
按''分片为最小可信单元''设计重发；MKV 簇内块损坏只能整簇丢。
截断：未写完的 MP4 缺尾部 \Tp{moov}，文件近乎报废，抢救靠扫描
\Tp{mdat} 里的帧边界重建索引（专业工具或重封装修复）；MKV 缺
\Tp{Cues} 不伤播放只伤跳读，重新封装即可补全。时间戳损坏：
\Tp{-fflags +genpts} 允许读取端补造单调时标，
\Tp{-fflags +discardcorrupt} 把损坏包直接丢干净，跳变与回绕的定位
在 \Tp{AVFMT\_TS\_DISCONT} 标志下由调用方裁决；解码侧掩盖强度见
章 \ref{ch:ch08-decode-hw} 的 \Tp{-err\_detect} 一节。

\begin{warning}[title={三个常见翻车点}]
  一是拿 MP4 直接当直播容器：不设分片标志时，样本索引始终在等
  回填，延迟无上界；二是跨流复制包忘了换 \Tp{time\_base}，
  数字全对、节奏全错；三是把 TS 的连续计数器当校验和——它只证明
  ''包没丢''，不证明''内容没错''，内容级校验要 PCR 与解码侧
  错误检测配合。
\end{warning}

\begin{bestpractice}[title={封装层设计检查单}]
  先按消费端定容器：浏览器与 CDN 分发选 fMP4，广播链路选 TS，
  本地档案选 MKV，点播文件加 \Tp{+faststart}；再按时延预算定切片
  或簇的粒度；所有跨容器拷贝统一走 \Fn{av\_packet\_rescale\_q}；
  最后给每路输入预设''缺索引''的降级路径——重封装、补
  \Tp{Cues}、或改按分片消费。
\end{bestpractice}

% ============================================================
% ---- frag_ch11.tex ----
\part{实时链路与工程化}

\chapter{流媒体协议栈：从 RTMP 到 WebRTC}
\label{ch:ch11-streaming}

前面十章把一帧画面送到编码器入口，从这一章开始要回答另一个问题：码流离开
本机之后，怎么走完剩下的几十到几千公里。第~\ref{ch:ch10-mux-sync}~章已经处理
过封装与时间戳，本章把这些包交到网络协议手里；而协议一旦选定，端到端延迟的
量级就基本定死了，后面所有工程手段都只能在这个量级内部做微调。

本章按延迟预算把协议排在一条轴上，再逐个看它们的丢包恢复方式、NAT 穿越能力
和 \CL{} 侧接入形态。需要提醒的是，实时链路的复杂度不会因为选型而消失：
它在传输层、应用层与线程模型之间来回转移，第~\ref{ch:ch12-realtime-thread}~章
会看到同样的规律。

\section{把协议排在延迟预算这条轴上}

选型的第一步是问``我允许观众看到几秒之前的画面''。这个问题之所以是第一步，
是因为协议不是功能的多寡，而是延迟预算的分配方案。TCP 的可靠重传会把一个
丢包变成整条连接的等待，队头阻塞的代价在第~\ref{ch:ch02-clock-buffer}~章
缓冲一节已经算过；基于 UDP 的协议则要在``等多久重传''与``直接放弃这一帧''
之间显式取一个数。

表~\ref{tab:ch11-protocol-map} 按延迟从低到高排开，最后一列是部署成本的
粗略判断，只考虑并发规模与中转跳数，不考虑内容分发的商业条款。

\begin{longtable}{>{\raggedright\arraybackslash}p{0.15\textwidth}
                 >{\raggedright\arraybackslash}p{0.15\textwidth}
                 >{\raggedright\arraybackslash}p{0.20\textwidth}
                 >{\raggedright\arraybackslash}p{0.16\textwidth}
                 >{\raggedright\arraybackslash}p{0.13\textwidth}}
\caption{流媒体协议横向对比：延迟、恢复、穿越与成本}\label{tab:ch11-protocol-map} \\
\toprule
协议 & 端到端延迟 & 丢包恢复 & NAT 穿越 & 部署成本 \\ \midrule
\endfirsthead
\toprule
协议 & 端到端延迟 & 丢包恢复 & NAT 穿越 & 部署成本 \\ \midrule
\endhead
\bottomrule
\endfoot
WebRTC & 约 0.2 到 1.5 秒 & NACK 与 PLI，另带 FEC 与前向冗余 & ICE 内置，STUN 与 TURN 兜底 & 最高 \\
SRT & 约 1 到 3 秒 & 选择性 ARQ，窗口即延迟 & 需中继或端口映射 & 中 \\
RIST & 约 1 到 3 秒 & ARQ 与发送端 FEC & 点对点可直连 & 中 \\
LL-HLS 与 LL-DASH & 约 2 到 6 秒 & 靠 HTTP 分块与分片重取 & 走 443 端口，天然穿透 & 低 \\
RTSP 加 RTP & 约 0.3 到 2 秒 & 原生只丢，靠 RTCP 上报 & 需 ALG 或转 TCP & 中 \\
RTMP & 约 1 到 3 秒 & TCP 重传，队头阻塞 & 走 TCP，穿透容易 & 低 \\
\end{longtable}

这张表的排法值得说明：延迟最低的 WebRTC 同时是部署成本最高的，因为它把
拥塞控制、抖动缓冲、NAT 穿越与密钥协商全部塞进了端点；而成本最低的
LL-HLS 复用了既有 CDN，代价是它的每一个分片都必须完整落地才能被请求。
``用延迟换工程量''与``用工程量换延迟''是同一笔交易的两侧。

\begin{compare}[title={三种典型链路的延迟构成}]
\DomRT{} 监控回传选 SRT：延迟 1 到 2 秒可以接受，但要跑在移动网络上，
丢包率高，需要 ARQ。互联网直播选 LL-HLS：延迟 3 到 5 秒，走 CDN 边缘缓存，
并发不受源站限制。互动连麦选 WebRTC：延迟 400 毫秒以内，代价是必须维护
一整套信令与 TURN 基础设施。三者的编码参数可以完全一致，差别全在传输层。
\end{compare}

\section{RTMP 与 RTSP：两种正在退场的设计}

RTMP 是 Adobe 为 Flash Player 设计的协议，它把一切压在 TCP 上。握手用
\Tp{C0C1C2} 与 \Tp{S0S1S2} 三段交换完成，随后双方各持一条 TCP 连接，把消息
切成固定长度的块（chunk）来发。默认块大小是 128 字节，块头有四种编码形式，
其中 type 3 的延续块只带最少的复用信息，用来续写上一个未发完的消息。消息
时间戳是 24 位毫秒，超过 \Tp{0xFFFFFF} 时改用扩展时间戳字段，这个设计在
长时间推流上出过不少兼容性问题。

块小是为了降低首包延迟，但在 TCP 上它只增加了头部开销：每个块都要一次
块头，块大小不协商就会让带宽利用率下降，而真正的延迟来源是 TCP 重传与
拥塞窗口，块切得再小也躲不过队头阻塞。RTMP 被淘汰的原因不是性能，而是
承载它的 Flash 在 2020 年底退出浏览器；播放侧随即转向 HTTP 分片，推流侧
反而保留了很久，直到今天仍有大量 CDN 用 RTMP 做 ingest，因为它的服务端
实现（\Lib{nginx-rtmp}、\Lib{SRS}、\Lib{MediaMTX}）足够成熟稳定。

RTSP 走的是相反的路。它是一个纯信令协议，负责 DESCRIBE、SETUP、PLAY 与
TEARDOWN，媒体本身交给 RTP 走 UDP，RTCP 走同一地址的下一个偶数端口。
SETUP 的 \Tp{Transport} 头里协商端口、SSRC 与是否交织（interleaved）到 TCP
之内，这一条逃生通道让 RTSP 在严格 NAT 环境里仍可工作。RTSP 1.0 见
RFC~2326，2.0 见 RFC~7826，后者补齐了传输安全与规范文本。

问题在于：RTSP 与 RTP 的组合里，丢包恢复完全不在协议义务之内。RTP 只保证
时间戳与序号的语义，接收端要靠 RTCP 接收报告（RR）把丢包统计反馈回去，
但发送端要不要重传、重传多久，规范留白。于是监控行业里常见``能播就行''的
实现，一旦网络抖动就花屏。要可靠就得上 SRT 或 RIST，要互动就得上 WebRTC，
这两者都把 RTSP 的信令角色换掉了。

\begin{warning}[title={RTSP 的 TCP 交织不等于可靠传输}]
把 RTP 交织进 TCP（\Tp{Transport} 头里带 \Tp{interleaved}）只解决了穿越问题，
不改变 RTP 不重传的事实。此时丢包变成 TCP 的字节流重排，接收端解码器看到的
仍是损坏的帧。误把``走 TCP''读成``有 ARQ''是这条链路上最常见的设计错误。
\end{warning}

\section{SRT 与 RIST：用延迟买重传机会}

SRT（Secure Reliable Transport）由 Haivision 主导，血统来自 \Lib{UDT}：
一条 UDP 双向会话，发送端维护发送缓冲，接收端在 \Tp{latency} 给定的窗口内
用 NACK 请求重传。这个窗口是整套设计的核心旋钮：它规定``一个包最多可以被
迟到多久''，接收端的 TSBPD（时间戳基打包交付）会把包缓存到各自的呈现时刻
再交给解码器，于是重传来得及，呈现也不抖。窗口设 120 毫秒还是 500 毫秒，
决定了同一网络下你能承受的丢包率数量级。

清单~\ref{lst:ch11-srt-sender} 给出 \CL{} 侧建立监听端的最小代码。选项层参数
传 0 表示 SRT 自己的选项空间；加密口令用 \Tp{SRTO\_PASSPHRASE}，其缓冲区
长度是 79 字节。会合（Rendezvous）模式需要设 \Tp{SRTO\_MODE} 为
\Tp{SRT\_MODE\_RENDEZVOUS}，两端都``主动''连即可穿透 NAT，无需公网中继。

\begin{lstlisting}[style=cstyle,caption={SRT 发送端：建立监听会话并设置重传窗口},label={lst:ch11-srt-sender}]
/* SRT sender: listener with an explicit retransmit window. */
#include <srt/srt.h>
#include <arpa/inet.h>
#include <string.h>

int srt_open_listener(int port, int32_t latency_ms) {
  int opt = 0;   /* 0 = SRT socket option level */
  int yes = 1;
  struct sockaddr_in sa;
  SRT_SOCKET ss = srt_create_socket();
  if (ss == SRT_INVALID_SOCK) {
    return -1;
  }
  memset(&sa, 0, sizeof sa);
  sa.sin_family = AF_INET;
  sa.sin_addr.s_addr = INADDR_ANY;
  sa.sin_port = htons((uint16_t)port);

  srt_setsockopt(ss, opt, SRTO_LATENCY, &latency_ms,
                 sizeof latency_ms);
  srt_setsockopt(ss, opt, SRTO_TSBPDMODE, &yes, sizeof yes);

  if (srt_bind(ss, (struct sockaddr *)&sa, sizeof sa)
      == SRT_ERROR) {
    srt_close(ss);
    return -2;
  }
  if (srt_listen(ss, 2) == SRT_ERROR) {
    srt_close(ss);
    return -3;
  }
  int len = sizeof sa;
  SRT_SOCKET cs = srt_accept(ss, (struct sockaddr *)&sa, &len);
  if (cs == SRT_INVALID_SOCK) {
    return -4;
  }
  return (int)cs;   /* caller owns this socket */
}
\end{lstlisting}

命令行与 GStreamer 的接入形态如清单~\ref{lst:ch11-srt-cli}。
\Tp{pkt\_size=1316} 不是随意数字：一个 MPEG-TS 包 188 字节，
$188\times 7=1316$，正好凑成一个 UDP 数据报，避免分片。
\Lib{srt-live-transmit} 用来做协议搬运与中继，参数是输入 URI 加输出 URI。

\begin{lstlisting}[style=shellstyle,caption={SRT 链路的推流、取流与中继},label={lst:ch11-srt-cli}]
# push MPEG-TS over SRT, 200 ms retransmit window
ffmpeg -re -i cam.mp4 -c copy -f mpegts \
  "srt://relay:9000?mode=caller&latency=200&pkt_size=1316"

# GStreamer sender: listener on every interface
gst-launch-1.0 -v videotestsrc is-live=true \
  ! x264enc tune=zerolatency ! mpegtsmux \
  ! srtsink uri="srt://:9000?mode=listener&latency=200"

# protocol bridging: multicast in, SRT out
srt-live-transmit udp://@239.0.0.1:1234 "srt://:9002"
\end{lstlisting}

GStreamer 侧要注意元素代际：1.18 起 \Tp{srtsrc} 与 \Tp{srtsink} 被标为
废弃，改用 \Tp{srtstream} 与 \Tp{srtsocketsrc} 一族，因为它们支持一个会话
服务多个对端。用旧元素的工程在新版本上会看到弃用警告。

RIST（Reliable Internet Stream Transport）由 RIST Forum 制定，站在 RTP 之上
做增量扩展而非重写。它分两档：Simple Profile 提供 ARQ 与连接校验，行为与
SRT 接近；Advanced Profile 补上 FEC、VNF（virtual frame numbering，虚拟帧
编号）与多路复用，其中虚拟帧编号是关键差异：允许对单个流声明``允许迟到
多少帧''，比按字节窗口更贴近视频语义。参考实现是 \Lib{librist}，中继工具
是 \Tp{werepeater}；\CPP{} 工程可直接链 \Lib{librist}，\DomAud{} 与
\DomVid{} 共用一条 RIST 会话。

两者的取舍可以简化成一句话：SRT 生态更成熟（编码器、切换台、单机工具几乎
都带），RIST 的规范表达力更强（多流复用、每流独立的容忍度、FEC 与 ARQ
并存）。跨机构互联且需要多路复用时看 RIST，与消费级设备对接时看 SRT。

\section{HLS 与 LL-HLS：把文件切成可缓存的小片}

HLS 与 DASH 是``用 HTTP 传媒体''的两套方言。HLS 的清单是文本播放列表，
DASH 的是 XML 的 MPD，两者都把内容切成独立的分片对象交给 CDN 缓存，
客户端自己按时间线拼回去。这条路的优点是可以水平扩展：边缘节点缓存命中
之后，源站几乎不承担并发压力。

传统 HLS 的延迟来自三个地方叠加：分片时长（\Tp{\#EXT-X-TARGETDURATION}
决定的 4 到 6 秒一片）、客户端至少缓冲 3 个分片、以及清单轮询周期。
LL-HLS（Low-latency HLS，2019 年进入 HLS 规范修订）逐个拆掉它们。

清单~\ref{lst:ch11-hls-tags} 是 LL-HLS 片段的骨架。要点四条：分片内部再切
partial segment，用 \Tp{\#EXT-X-PART} 描述；\Tp{PART-TARGET} 声明 partial
时长；服务器用 HTTP 分块传输（chunked transfer encoding）在分片写完之前
就把已完成的 partial 推给客户端；客户端在请求上带 \Tp{\_HLS\_msn} 与
\Tp{\_HLS\_part} 让服务器阻塞到该位置才有响应，这就是 blocking playlist
reload。还有 \Tp{\#EXT-X-PRELOAD-HINT}，让客户端提前订阅服务器正在写的那
一个 partial。

\begin{lstlisting}[style=shellstyle,caption={LL-HLS 主清单与媒体清单中的关键标签},label={lst:ch11-hls-tags}]
# master.m3u8
#EXT-X-VERSION:9
#EXT-X-SERVER-CONTROL:CAN-BLOCK-RELOAD=YES,PART-TARGET=0.500
#EXT-X-STREAM-INF:BANDWIDTH=4800000,RESOLUTION=1920x1080
media/720p.m3u8

# media.m3u8
#EXT-X-TARGETDURATION:4
#EXT-X-PART-INF:PART-TARGET=0.333
#EXTINF:4.000,
seg-1024.m4s
#EXT-X-PART:DURATION=0.333,URI="part-1025.0.m4s",INDEPENDENT=YES
#EXT-X-PART:DURATION=0.333,URI="part-1025.1.m4s"
#EXT-X-PRELOAD-HINT:TYPE=PART,URI="part-1025.2.m4s"
\end{lstlisting}

DASH 侧的对应机制写在 MPD 的属性里：\Tp{availabilityTimeOffset} 声明
``分片在时间线上比它落地时刻早多少可用''，客户端据此发起提前请求；
\Tp{ServiceDescription} 里的 \Tp{Latency@target} 给出目标延迟，供播放器
自适应；低延迟还要求客户端与服务器的墙钟对齐，MPD 用 \Tp{UTCTiming} 元素
指定校时源。CMAF 则是承载层的收敛：HLS 与 DASH 共用同一批 fMP4 分片，
一个分片由若干 \Tp{moof} 加 \Tp{mdat} 的成对结构串成，这些对既是分片也是
低延迟语境下的 chunk，所以打包一次即可两边分发。

HTTP 分发必须处理跨源取回：清单与分片走 HTTPS，媒体清单常放在另一个域，
于是响应头 \Tp{Access-Control-Allow-Origin} 缺失会让播放器卡在取清单一步；
带查询参数的 blocking reload 也不能被 CDN 当成不同对象反复回源。

\begin{note}[title={CDN 与低延迟分片缓存并不总是相容}]
LL-HLS 的价值在于``不等整片落地''，但很多 CDN 默认只在完整响应上缓存，
分块响应会被整体缓冲，延迟立刻退回整片模式。上线前必须确认边缘节点允许
分块透传，也要确认 \Tp{\_HLS\_msn} 这类查询参数不会被忽略，否则阻塞请求会
拿到一份陈旧清单。
\end{note}

\section{WebRTC：完整分层与数据通道}

WebRTC 的分层是把上面所有协议各取一段拼起来的：ICE 负责连通性（RFC~8445，
配套 STUN 见 RFC~8489、TURN 见 RFC~5766），DTLS 负责密钥协商，SRTP 负责
媒体加密（RFC~3711），SCTP 承载数据通道（RFC~8831 与 RFC~8832），RTP 与
RTCP 仍是最底层的媒体语法。会话描述用 SDP，通过应用自带的信令通道交换，
协议本身不含信令。

这一套分层解释了它为什么延迟低、为什么贵：ICE 让端点自己找直连路径，
打洞失败才回落 TURN 中继；拥塞控制是端点内建的 GCC，用传输宽度的 RTCP
反馈逐包定位丢包，再用 NACK 精确补包、PLI 请求关键帧。
延迟敏感的场合把 \Tp{maxRetransmits} 设成 0 或 1，宁可花屏也不等重传。
数据通道提供两种可靠性：完全有序可靠，或设 \Tp{ordered:false} 与重传上限
做``过期即弃''的通道，适合控制消息与字幕。

\begin{lstlisting}[style=shellstyle,caption={GStreamer 经 webrtcbin 组一条连麦推流},label={lst:ch11-webrtcbin}]
# webrtcbin handles ICE/DTLS/SRTP; SDP needs own signaling
gst-launch-1.0 -v webrtcbin name=w bundle-policy=max-bundle \
  stun-server=stun://stun.example:19302 \
  videotestsrc is-live=true \
  ! video/x-raw,framerate=30/1 ! vp8enc error-resilient=1 \
  ! rtpvp8pay ! w.sink_0
\end{lstlisting}

\CPP{} 工程里落地一般直接取 \Lib{libwebrtc}：它的 C++ 接口提供
\Tp{Peer{\Bk}Connection{\Bk}Interface}、轨道与转发的抽象，同时把 jitterbuffer、
AEC/AGC/ANS 的 \DomAud{} 处理链和视频软件编码管线一并带进来。代价是构建
体量与迭代节奏，多数团队会选预编译发行版或封装层，而不是自己啃上游。
GStreamer 的 \Tp{webrtcbin} 是另一条路：它把 ICE 交给 \Lib{libnice}，
自己组装 RTP、DTLS-SRTP 与 SCTP，SDP 通过 \Tp{create-{\Bk}offer} 与
\Tp{set-{\Bk}remote-{\Bk}description} 这些 action 驱动，信令仍需自备。

工程上容易低估的三点：TURN 是有状态的中继，带宽成本按峰值计费；多端
一致性依赖 SFU 转发，转发者要理解 simulcast 与关键帧请求的语义；
数据通道默认走 DTLS 应用协议，与服务端非 WebRTC 组件对接时往往是障碍，
此时更常见的组合是媒体走 WebRTC、控制消息走 WebSocket。

\section{QUIC 上的媒体传输与 MoQ}

HTTP/2 已经暴露了把多路复用架在单一 TCP 上的后果：一个包丢，全部流停住。
QUIC（RFC~9000）把这层搬回 UDP，每条流独立重传、独立拥塞控制，握手与
TLS 1.3 交织完成，连接迁移支持换网不断线。对媒体而言，最直接的收益是
``一条连接上并行拉多个分片，慢的那个不再拖累别人''。

在此之上，IETF 的 MoQ 工作组正在标准化 Media over QUIC：以发布/订阅为中
心的对象模型，把一帧切成若干 object 归入 group，由 relay 做缓存与转发，
目标是让延迟与缓存命中率第一次同时可调。截至 2026 年，MoQ 的传输规范仍
处于草案阶段，多厂商互操作尚未收敛，落地时建议把它当架构参考而非现成的
部署选项。相对成熟的用法是 CMAF 分片通过 HTTP/3 分发，收益是消除队头
阻塞、缩短首屏，而不需要改动清单格式。

\begin{lstlisting}[style=cstyle,caption={用 RTCP 发送报告把 RTP 时间戳钉到墙钟},label={lst:ch11-clock}]
#include <libavutil/time.h>

/* Map a wall clock reading onto the NTP epoch (1900, not 1970). */
static void ntp_from_wallclock(int64_t wall_us, uint32_t *sec,
                               uint32_t *frac) {
  int64_t secs = wall_us / 1000000;
  int64_t f = (wall_us % 1000000) * 4294LL;   /* 2^32 / 10^6 */
  /* NTP epoch 1900 vs Unix epoch 1970: 2208988800 seconds. */
  *sec = (uint32_t)(secs + 2208988800LL);
  *frac = (uint32_t)(f / 1000000);
}

/* RTCP SR carries both halves together with the RTP timestamp. */
int stamp_sender_report(uint32_t rtp_ts, uint64_t out[2]) {
  uint32_t ntp_sec, ntp_frac;
  ntp_from_wallclock(av_gettime(), &ntp_sec, &ntp_frac);
  out[0] = ((uint64_t)ntp_sec << 32) | ntp_frac;
  out[1] = rtp_ts;   /* the pair is the media-to-wall mapping */
  return 0;
}
\end{lstlisting}

跨层对齐的通用规则可以写成：发送端把媒体时间戳与墙钟成对上报，接收端据此
把自己的墙钟折算成媒体时间轴，再与本地单调钟比较，得到漂移速率并平滑修正
呈现节奏。RTCP 发送报告用两个 32 位字段带出完整的 NTP 时间戳，与它同段的
RTP 时间戳就是映射的另一端，这正是清单~\ref{lst:ch11-clock} 的形状；
WebRTC 侧还有 \Tp{abs-send-time} 扩展，把同一时刻的秒小数截成 24 位塞进每个
RTP 包，供发送端拥塞控制逐包测延迟。HLS 用 \Tp{\#EXT-X-PROGRAM-DATE-TIME}
表达同样的映射；DASH 靠 \Tp{availabilityStartTime} 与 \Tp{UTCTiming} 完成。
客户端的墙钟来自 NTP、PTP（Linux 上可由 \Tp{/dev/ptp0} 配合
\Tp{SO\_TIMESTAMPING} 拿到网卡硬件时间戳），墙钟与单调钟的分工见
第~\ref{ch:ch02-clock-buffer}~章。

\begin{guideline}[title={传输层选型的四条判据}]
\begin{itemize}
  \item 延迟预算先定数量级：亚秒级只剩 WebRTC 与调小窗口的 SRT 两条路；
        三秒级 SRT 与 RIST 更省心；五秒以上交给 LL-HLS 或 LL-DASH。
  \item 丢包恢复机制要与链路匹配：公网高丢包选 ARQ，弱网且可容忍花屏选
        FEC 加 ``过期即弃''，专线低丢包时原生 RTP 已经够用。
  \item NAT 形态决定拓扑：有公网 IPv4 直连即可；对称 NAT 后必须准备 TURN
        或 SRT 的 Rendezvous 模式，别指望 ``UDP 一定能打洞''。
  \item 接入形态与既有设施对齐：自研 socket 拿最大控制权，用库
        （\Lib{libsrt}、\Lib{librist}、\Lib{libwebrtc}）省时间，用 GStreamer
        元素（\Tp{srtsink}、\Tp{webrtcbin}）省整条管线的胶水代码。
\end{itemize}
\end{guideline}

% ============================================================
% ---- frag_ch12.tex ----
\chapter{实时线程模型与无锁回调}
\label{ch:ch12-realtime-thread}

前面十章里，缓冲、时钟与设备都是``数据''的问题；本章换成``时间''的问题：
谁在什么时候被叫醒，被叫醒之后必须在多长时间内干完。这一章的结论在很大程度上
否定了 \DomBatch{} 链路养成的习惯（第~\ref{ch:ch10-mux-sync}~章的封装与同步可以
慢慢重试，这里不行），也解释了为什么实时音频的代码总在跟语言本身较劲。

\section{硬实时回调的契约}

先看这条契约的具体形状。以 \DomAud{} 为例：驱动把回调登记为 ``在硬件周期
到点时被调用，必须在下一个周期到点前返回''。这不是一句性能建议，而是一串
可判定的禁止事项。

第一条是禁止堆分配。分配器内部有锁与线程局部状态，撞上别的线程正在释放时会
自旋或阻塞，最长等待无上界。第二条是禁止可缺页的内存访问。第一次触碰未映射
的页会触发次缺页，内核要清零页框，时间不确定；这也是为什么``只做个 memcpy''
仍可能出事，如果目标缓冲是刚分配的、还没被写过。第三条是禁止任何无上界等待
的锁，包括看似无害的引用计数与单例懒初始化。第四条是禁止阻塞式 IO：写管道、
打印、落盘都在出局之列，只读的无锁查询与不会睡眠的时间戳接口可以保留。
第五条是禁止会睡眠的同步原语，条件变量、读写锁、信号量在这一层全部出局。

违反的代价不是``慢''，而是可闻的咔哒声与可见的画面跳，因为它把``延迟''变成
了``延迟不确定''。这也解释了实现策略：一次性预分配所有帧、把所有跨线程交接
换成无锁结构、把所有慢活推到别的线程，回调线程只剩下搬字节与移动游标。

\begin{guideline}[title={回调线程里允许存在的东西}]
预分配缓冲上的 memcpy 与算术、原子读写（显式 memory order）、只读的
\Fn{snd\_pcm\_avail\_update}、以及自己维护的下标与状态机。其余一律先证明
无界等待不存在，再考虑放进来。调试期的断言宏（例如把违反契约的调用直接
断言失败）比事后 profiling 有效得多。
\end{guideline}

\section{SPSC 环形缓冲：游标、内存序与缓存行}

无锁交接的最小可用形态是单生产者单消费者环形缓冲（SPSC）。它成立的前提是
``恰好两个线程''：一个只会写 \Tp{head}，一个只会读 \Tp{tail}，两者都是单调
递增的下标，取模靠 2 的幂容量与掩码完成。单调性让``谁覆盖了谁''变成不可能，
于是无需 CAS，也无需回退重试。

内存序是关键。裸 \CPP{} 的 \Tp{std::atomic} 默认是 \Tp{seq\_cxx}，安全但
在弱序架构上偏保守；SPSC 的正确配方是：写方先写数据、再以 release 发布游标，
读方以 acquire 观察游标、再读数据。这样 release 之前的写入对看到该值的线程
一律可见，而无需全局串行化。消费者判断可用字节数用 \Tp{head - tail}，无符号
回绕恰好给出正确差值，这是把 ``游标`` 写成单调计数而不是环内下标的原因。

第二件事是假共享。头尾游标若落在同一缓存行，两个核会不停地把这一行互相夺回，
吞吐可以掉一个数量级。\cpp{17} 给出的缓存行大小上界常量
\Tp{hardware\_destructive\Bk\_interference\Bk\_size} 可以用，
但工程上更普遍的做法是直接按 64 或 128 字节对齐两个游标，让每个热字段独占一行。
清单~\ref{lst:ch12-ring} 是一个可用的骨架。

\begin{lstlisting}[style=cppstyle,caption={带显式内存序与缓存行填充的 SPSC 环形缓冲},label={lst:ch12-ring}]
// Lock-free single-producer / single-consumer ring buffer.
class Spsc {
public:
  explicit Spsc(size_t cap_pow2)
    : buf_(cap_pow2), mask_(cap_pow2 - 1) {
    // Touch every page once; never fault inside the callback.
    memset(buf_.data(), 0, cap_pow2);
  }

  // Producer side: runs on the callback thread.
  bool try_push(const uint8_t *src, size_t n) noexcept {
    if (n > free_bytes()) {
      return false;   // full: caller drops or retries
    }
    const size_t w = head_.load(std::memory_order_relaxed);
    memcpy(&buf_[w & mask_], src, n);  // payload first
    // Publish the cursor only after the payload is in place.
    head_.store(w + n, std::memory_order_release);
    return true;
  }

  // Consumer side: returns the bytes actually copied out.
  size_t pop(uint8_t *dst, size_t n) noexcept {
    const size_t r = tail_.load(std::memory_order_relaxed);
    const size_t avail = used_bytes();
    const size_t k = n < avail ? n : avail;
    if (k) {
      memcpy(dst, &buf_[r & mask_], k);
      tail_.store(r + k, std::memory_order_release);
    }
    return k;
  }

  size_t used_bytes() const noexcept {
    const size_t h = head_.load(std::memory_order_acquire);
    const size_t t = tail_.load(std::memory_order_relaxed);
    return h - t;
  }
  size_t free_bytes() const noexcept {
    return buf_.size() - used_bytes();
  }

private:
  alignas(64) std::atomic<size_t> head_;  // producer writes
  alignas(64) std::atomic<size_t> tail_;  // consumer writes
  std::vector<uint8_t> buf_;
  size_t mask_;
};
\end{lstlisting}

三点提醒：容量取 2 的幂，掩码才成立；\Fn{try\_push} 满了就返回 false，绝不
等待，等待会把无锁的意义抵消掉；跨界写入（一次 memcpy 跨过环尾）在这份骨架里
简化掉了，真实实现要拆成两次 memcpy。另外，生产者与消费者各自的 ``观察对方
游标'' 用 acquire，``修改自己的游标'' 用 release，读自己的游标 relaxed 足够。

\begin{warning}[title={无锁不等于无等待，也不等于可重入}]
SPSC 只保证没有锁竞争，不保证在硬中断或更高优先级抢占下仍能按时完成：如果
另一个实时线程优先级更高，回调照样会被推迟。同时环形缓冲不是信号安全的，
不要指望在信号处理函数里重入它；驱动侧的唤醒机制
（ALSA 的 \Tp{poll} 描述符、WASAPI 的事件句柄）才是合法的触发路径。
\end{warning}

\section{抖动的来源与两侧补偿}

生产侧的抖动多数来自``活不在固定时刻来''：解码线程被调度推迟、滤镜耗时随内容变化、网络侧把一帧的数据分成两次到达。消费侧的抖动来自硬件周期本身与驱动缓冲管理，量级通常更小但更规律。两者的合流点就是环形缓冲的深度曲线。

补偿策略必须区分``谁更快''。消费方持续比生产方快（欠载）时，正确做法是补
静音而不是让设备停住，因为停一次就要重新走一遍预启动；消费方持续比生产方慢
（积压）时，必须主动丢弃，否则延迟单调增长。第~\ref{ch:ch02-clock-buffer}~章
讨论过重采样的频率漂移，这里的丢帧与补静音是同一件事在离散侧的表现。

\DomAud{} 在 Linux 上的再同步入口是 \Fn{snd\_pcm\_avail\_update}：它返回驱动当前
可写的帧数，回调据此决定这一周期搬多少。缓冲区耗尽后写入返回 \Tp{-EPIPE}，此时
不要慌，标准动作是 \Fn{snd\_pcm\_prepare} 或 \Tp{snd\_pcm\_recover}，恢复之后
再重新对齐阈值（\Fn{snd\_pcm\_set\_start\_threshold}），并顺手用
\Fn{snd\_pcm\_delay} 与硬件时间戳校正呈现起点。Windows 的 \DomAud{} 侧对应
\Tp{IAudioRenderClient} 的 \Fn{GetBuffer} 与 \Fn{ReleaseBuffer} 配对，事件驱动
模式下由 \Fn{GetNextPacketSize} 决定取多少；macOS 与 iOS 的 \Tp{AudioUnit} 回调
里只能碰 \Tp{AudioBufferList} 给的那块内存。三条路的共同点是：容量与阈值都在
打开设备时谈妥，回调只做搬运。

清单~\ref{lst:ch12-alsa} 给出 \CL{} 侧的回调形态。注意 \Fn{calloc} 与
\Fn{snd\_pcm\_open} 都发生在打开阶段，回调函数体内只剩 \Fn{memset}、\Fn{memcpy}、
游标推进与 ALSA 的无锁查询。

\begin{lstlisting}[style=cstyle,caption={ALSA 播放回调：只做搬运、补静音与游标推进},label={lst:ch12-alsa}]
/* Playback drain: no allocation, no lock, no logging. */
#include <alsa/asoundlib.h>
#include <string.h>

size_t ring_pop(struct Spsc *r, uint8_t *dst, size_t n);

typedef struct {
  snd_pcm_t *pcm;
  struct Spsc *ring;
  uint8_t *stage;    /* one period, allocated at open */
  size_t period_bytes;
} Sink;

/* Device setup: never runs on the RT thread. */
static int sink_open(Sink *s, const char *dev, size_t bytes) {
  s->stage = calloc(1, bytes);
  if (!s->stage) {
    return -1;
  }
  s->period_bytes = bytes;
  if (snd_pcm_open(&s->pcm, dev, SND_PCM_STREAM_PLAYBACK, 0)
      < 0) {
    return -2;
  }
  snd_pcm_set_params(s->pcm, SND_PCM_FORMAT_S16_LE,
                     SND_PCM_ACCESS_RW_INTERLEAVED,
                     2, 48000, 0, 500000);
  return snd_pcm_prepare(s->pcm) < 0 ? -3 : 0;
}

/* The RT body. 2 channels, 16 bit: 4 bytes per frame. */
static void drain_one_period(Sink *s) {
  snd_pcm_sframes_t want = snd_pcm_avail_update(s->pcm);
  const size_t per = s->period_bytes / 4;
  while (want >= (snd_pcm_sframes_t)per) {
    if (ring_pop(s->ring, s->stage, s->period_bytes)
        < s->period_bytes) {
      memset(s->stage, 0, s->period_bytes);  /* silence */
    }
    if (snd_pcm_writei(s->pcm, s->stage, per) < 0) {
      return;      /* recover on a control thread */
    }
    want = snd_pcm_avail_update(s->pcm);
  }
}
\end{lstlisting}

\section{实时调度与内存锁定}

\DomRT{} 线程需要比普通线程更高的调度地位。Linux 侧是
\Tp{SCHED\_FIFO}：同优先级之间先到先跑得完为止，不同优先级严格抢占。
它需要权限，判据是 \Tp{RLIMIT\_RTPRIO}（即 \Tp{ulimit -r}）与能力位
\Tp{CAP\_SYS\_NICE}，容器里还要放开对应的 seccomp 与 cgroup 限制。
优先级取值由 \Fn{sched\_get\_priority\_min} 与 \Fn{sched\_get\_priority\_max}
给出，1 到 99，工程上取中段留出上调空间。

真正的坑是 RT throttling：内核默认每个 1 秒的周期里只允许实时任务烧掉
950000 微秒，剩下 50000 微秒留给普通任务，参数是
\Tp{kernel.sched\_rt\_runtime\_us} 与 \Tp{kernel.sched\_rt\_period\_us}。
一条``永不调用阻塞接口''的实时线程很容易让人忘记这个开关存在，一旦多个实时
线程叠加，回调会被整体掐住接近一个周期，症状是每秒一次的规律性卡顿。
需要更严格的可调度性保证时，\Tp{SCHED\_DEADLINE} 用运行时间、截止期与周期
三元组声明需求，内核做准入控制，比 FIFO 更可控。

内存侧要 \Fn{mlockall} 把当前与将来的页都锁住，配合启动时遍历一遍所有缓冲
（清单~\ref{lst:ch12-ring} 构造函数的 \Fn{memset} 就是这件事），次缺页才
真的消失。锁页上限是 \Tp{RLIMIT\_MEMLOCK}，容器与老内核上默认很小。
CPU 亲和性把回调钉在隔离核上，隔离靠引导参数与中断重排；否则一次网卡软中断
就能吃掉整个周期。

清单~\ref{lst:ch12-rtsetup} 是这一节的三种手段合流后的形态。Windows 侧
不用手工设优先级，而是向 MMCSS 声明任务类别：
\Fn{AvSetMmThreadCharacteristicsW} 传 \Tp{Pro Audio} 这样的任务名换取
一个句柄与优先级，之后可用 \Fn{AvSetMmThreadPriority} 调整；MMCSS 会统一
管理优先级、IRQ 与电源策略，比自提 \Tp{REALTIME\_PRIORITY} 更适合声卡共用
场景。macOS 与 iOS 上应用一般拿不到系统级实时策略，\Tp{qos\_class\_t}
（\Tp{QOS\_CLASS\_USER\_INTERACTIVE} 一族）是普通线程的优先级分层，
音频实时性由 CoreAudio 托管的 render 线程提供，工程侧要做的是不在回调里
拖慢它，而不是设法提权。

\begin{lstlisting}[style=cstyle,caption={锁页、提权与配额自检},label={lst:ch12-rtsetup}]
/* One-time setup: pin memory, then raise priority. */
#include <sched.h>
#include <sys/mman.h>
#include <pthread.h>
#include <string.h>

int lock_process_memory(void) {
  /* Never let the callback take a page fault. */
  if (mlockall(MCL_CURRENT | MCL_FUTURE) != 0) {
    return -1;
  }
  return 0;
}

int make_thread_realtime(pthread_t th, int prio) {
  struct sched_param sp;
  memset(&sp, 0, sizeof sp);
  sp.sched_priority = prio;   /* 1..99 for SCHED_FIFO */
  /* pthread_setaffinity_np can pin this to one core. */
  return pthread_setschedparam(th, SCHED_FIFO, &sp);
}

int rt_setup_supported(void) {
  int lo = sched_get_priority_min(SCHED_FIFO);
  int hi = sched_get_priority_max(SCHED_FIFO);
  /* RT throttling: runtime_us taken out of period_us.
   *   the default pair is 950000 out of 1000000.
   */
  return (lo >= 1 && hi <= 99) ? 0 : -1;
}
\end{lstlisting}

运维与自检的命令面如下。

\begin{lstlisting}[style=shellstyle,caption={实时配置的确认与抢占观测},label={lst:ch12-rtshell}]
# can we take SCHED_FIFO? check limits before promoting
ulimit -r
grep CapEff /proc/self/status
# how much CPU may realtime burn per scheduling period
cat /proc/sys/kernel/sched_rt_runtime_us
cat /proc/sys/kernel/sched_rt_period_us
# promote and pin the mixing thread, then watch preemptions
sudo chrt -f -p 80 $(pgrep mymixer)
sudo perf record -e sched:sched_switch -a -- sleep 5
# a stray malloc() on the RT thread is loud under these tools
valgrind --tool=helgrind ./mymixer
\end{lstlisting}

\begin{note}[title={周期、缓冲与延迟不是同一件事}]
\DomAud{} 的端到端延迟下界约等于一个 period 加上启动阈值对应的时长，
跟缓冲区总大小关系不大：缓冲区加大只是给了更多抗抖动余量，并不自动降低
延迟，反而让一次积压更难被丢掉。调低 period 会提高回调频率与单位时间开销，
在弱 CPU 上反而制造丢帧。这两个方向要一起测，不能只调一个。
\end{note}

\section{视频侧的对应物：呈现线程、背压与丢弃策略}

\DomVid{} 很少有一个``驱动强制返回''的硬实时线程，但有同一类问题的软实时版：
呈现必须踩在垂直同步上，而解码与滤镜的耗时是内容相关的。标准拆分是三个
线程：接收或读取、解码、呈现，前两个跑``尽力而为''，第三个由显示后端驱动。
第~\ref{ch:ch04-capture-display}~章谈过采集与显示的时序，这里补上中间的
队列策略。呈现线程要短小：只做等待同步、提交与取下一帧，把色彩转换与缩放
推到工作线程，否则垂直同步一到手上却没有可提交的帧，等价于丢帧。

背压用有界队列表达。队列深度是延迟的直接来源，``满了''这件事必须给出
选择：阻塞上游（延迟涨、不丢内容）、丢最旧（延迟低、画面跳）、丢最新
（画面顺、延迟涨）。延迟敏感的链路里答案通常是丢最旧，因为旧帧已经不代表当前
时刻；而交互式渲染更倾向丢最新，宁可延迟也不要闪烁。丢弃判定还要看参考结构：
丢一个被后续帧参考的 \DomVid{} 帧，代价是到下一个帧内参考点为止全部解不出来，
所以真正的做法是丢一个 GOP 的尾巴或者请求关键帧，而不是随机丢单帧。

框架里这些旋钮是有名字的。GStreamer 的 \Tp{appsrc} 与 \Tp{appsink} 提供
\Tp{max-by-buffers}、\Tp{emit-signals} 与 \Tp{drop}，其中 \Tp{drop=TRUE}
正是``满了丢最旧''；播放器的同步由 \Tp{GstBaseSink} 的 preroll 与 render
两步完成，\Tp{sync} 属性关掉就变成把节奏判断交给上游。FFmpeg 的播放示例
\Tp{ffplay} 用互斥与条件变量加超时代替无锁，因为它的线程不是硬实时；
Media Foundation 的源读取器与事件队列走异步回调，呈现时刻由媒体时钟决定。
这些机制可以混用，但同一个缓冲只能有一套丢弃策略，否则两处同时丢会造成倒退。

\begin{compare}[title={音频与视频实时契约的差异}]
\DomAud{} 侧：周期固定、时长可预期，违约的表现是咔哒与爆音，因此可以严格
禁止分配与阻塞。
\DomVid{} 侧：帧时长随内容变化、解码耗时不确定，违约的表现是丢帧或跳帧，
几乎无人能察觉，因此允许慢路径与重试。
共同点：两侧的缓冲都必须有界，都必须定义``满了怎么办''，都必须把墙钟与
单调钟分开用（见第~\ref{ch:ch02-clock-buffer}~章）。
\end{compare}

\section{常见反模式与检测}

表~\ref{tab:ch12-antipattern} 把这些年反复出现的模式列全。``检测方法''一列
只写手段，不写具体工具版本；核心思路都一样：把回调的进入与退出时间戳记下来
做直方图，任何长尾都能现形，比读代码可靠。

\begin{longtable}{>{\raggedright\arraybackslash}p{0.19\textwidth}
                 >{\raggedright\arraybackslash}p{0.21\textwidth}
                 >{\raggedright\arraybackslash}p{0.20\textwidth}
                 >{\raggedright\arraybackslash}p{0.19\textwidth}}
\caption{实时回调的反模式、症状、检测与修法}\label{tab:ch12-antipattern} \\
\toprule
反模式 & 典型症状 & 检测手段 & 修法 \\ \midrule
\endfirsthead
\toprule
反模式 & 典型症状 & 检测手段 & 修法 \\ \midrule
\endhead
\bottomrule
\endfoot
回调里分配堆内存 & 周期性爆音，偶发几毫秒尖峰 & 回调耗时直方图，或拦截分配器 & 预分配帧池，运行期不分配 \\
回调里打日志 & 平均延迟没变但抖动明显 & 对比关掉日志后的抖动曲线 & 回调写无锁日志缓冲，另线程落盘 \\
回调里取互斥锁 & 卡顿与系统负载正相关 & 进入与退出各打一次时间戳并断言超时 & 换成单生产者单消费者的原子交接 \\
运行期反复创建小对象 & 长时间运行后延迟单调爬升 & 统计分配次数与常驻内存 & 对象池加放置构造 \\
没有消除次缺页 & 冷启动头几秒规律咔哒 & 观察缺页计数在运行期是否增长 & 启动时锁页并预先触碰所有缓冲 \\
忽略实时调度配额 & 每秒一次的整周期停顿 & 查两个配额开关的取值 & 提高配额，或改用截止期调度类 \\
无界队列加阻塞等待 & 延迟一路上涨且不回落 & 队列深度与最大滞留时间指标 & 队列封顶，满了丢最旧或就地补静音 \\
忙等共享标志 & 该核占用满格，帧率没变好 & 观察核占用与缓存失效率 & 用事件加超时等待，或让出该核 \\
\end{longtable}

\begin{bestpractice}[title={实时侧的落地顺序}]
\begin{enumerate}
  \item 先定接口：把回调与慢线程的职责切干净，跨线程只留一条 SPSC
        通道，参数（容量、周期、通道数、格式）在打开阶段全部谈妥。
  \item 再做预分配：帧池、暂存区、日志缓冲、编解码器上下文，一次分配
        并在构造期触碰一遍页，运行期只剩游标。
  \item 然后提权与锁页：\Tp{SCHED\_FIFO} 或 MMCSS 二选一按平台定，
        同时确认配额、能力位与亲和性，别只改线程不改内核参数。
  \item 最后压测：把链路跑在``比预期更坏''的网络与负载下，看回调耗时
        直方图的尾分位与队列深度曲线，而不是只看平均值。
\end{enumerate}
\end{bestpractice}

% ============================================================
% ---- frag_ch13.tex ----
\chapter{框架与引擎集成：把解码器接进别人的进程}
\label{ch:ch13-framework-integration}

前面十二章讨论的都是``自己进程里的一条管线''：设备取帧、时钟对齐、编解码、滤镜、封装、推流。
本章换一个视角。现实中绝大多数音视频内核并不是终端产品，而是别人产品里的一个依赖——
接进剪辑器、接进直播客户端、接进游戏引擎、接进浏览器。
一旦位置变成``依赖''，决定成败的就不再是 \Fn{avcodec\_{\Bk}send\_{\Bk}packet} 的返回值，
而是链接形态、许可义务、崩溃半径与升级成本这四件事。

\section{四档嵌入形态：链接、管道、服务与沙箱}

把解码器交给宿主进程，只有四种基本姿势，代价各自完全不同。

\begin{longtable}{>{\raggedright\arraybackslash}p{0.13\textwidth}
                  >{\raggedright\arraybackslash}p{0.15\textwidth}
                  >{\raggedright\arraybackslash}p{0.14\textwidth}
                  >{\raggedright\arraybackslash}p{0.13\textwidth}
                  >{\raggedright\arraybackslash}p{0.16\textwidth}}
\caption{嵌入形态对比：延迟、许可、维护与硬件解码可达性}\label{tab:ch13-embed-forms} \\
\toprule
形态 & 额外延迟 & 许可风险 & 可维护性 & 硬件解码可达性 \\
\midrule
\endfirsthead
\toprule
形态 & 额外延迟 & 许可风险 & 可维护性 & 硬件解码可达性 \\
\midrule
\endhead
\midrule
\endfoot
\bottomrule
\endlastfoot
进程内链接 \Bk libav*、\Bk GStreamer、\Bk Media Foundation &
毫秒级，只在帧拷贝与回调上花时间 &
最高：整份二进制的许可被依赖决定 &
符号级可见，崩溃半径是整个宿主 &
最高，但要自己走厂商 API 或 \Bk GStreamer 的 vaapi\Bk/ d3d11 元素 \\
fork 子进程喂管道 &
一帧到一帧的调度开销，几十毫秒 &
低：子进程可独立许可、独立替换 &
只共享字节，日志与错误要自己串 &
可达，靠命令行 \Tp{-{\Bk}hwaccel}，失败要能退回软解 \\
独立转码服务 &
网络往返加排队，秒级起步 &
最低，服务端内部使用最自由 &
多一层协议、部署与容量管理 &
由服务端统一决定，客户端不参与 \\
浏览器 WASM &
大：容器解封装、拷贝、线性内存边界 &
低：加载哪个解码器由页面决定 &
调试与性能归因困难 &
几乎不可达，WebCodecs 才另开一条路 \\
\end{longtable}

判据可以简化成一句：需要逐帧像素、需要在帧里改数据、需要亚帧延迟，就进进程内；
只需要端到端的文件进、文件出，就用子进程或服务。
中间态（例如把 GStreamer 塞进一个不欢迎 GObject 的 \CPP{} 引擎）通常比两端都贵。

\subsection*{故障隔离的真实边界}
进程内链接时，解码器的一次越界写会污染宿主的分配器，崩溃栈落在第三方 \Tp{.{\Bk}so} 里，
复现依赖具体码流，符号还常被 strip 掉；子进程与服务则把这件事变成``任务失败''，
可以重启、可以换节点，代价是状态与缓冲都要跨边界重建。
对 \DomRT{} 链路，隔离与延迟是直接冲突的两项，必须显式选一个。

\section{GStreamer：caps 协商替你省下的转换代码}

\subsection*{对象模型：先接受它，再使用它}
GStreamer 的类型系统是 GObject 的 \Tp{GType} 动态类体系，不是 \CPP{} 模板。
三件事必须内化：第一，引用计数，跨线程传递要 \Fn{g\_{\Bk}object\_{\Bk}ref}、用完 \Fn{g\_{\Bk}object\_{\Bk}unref}；
第二，属性，运行时按字符串名读写，\Fn{g\_{\Bk}object\_{\Bk}set} 传 ``sync'' 这类键，写错键只在控制台告警；
第三，四态状态机 NULL、READY、PAUSED、PLAYING——sink 在 PAUSED 就要拿到首帧做预滚动，
所以 \DomAud{} 的``先填一帧再播放''在 GStreamer 里是由状态机负责，而不是由你写。

\subsection*{caps 协商决定了你几乎不写格式转换}
每个 pad 挂一个 \Tp{GstCaps} 模板：最一般的形式在查询中冒泡，具体形式由下游固定。
一旦某个 sink 需要 1280$\times$720、YUY2、30 fps，链上的 \Tp{capsfilt{\Bk}er} 与自动插件会让
\Tp{decodebin} 一族在候选解码器里试到能产出兼容 caps 为止；像素格式不合就插 `videoconvert`，
速率不合就插 `videorate`，音频就插 `audioconvert` 与 `audioresample`。
于是你的转换代码缩成一行 caps 字符串——这不是魔法，是把第 5、9 章的格式与滤镜问题
交还给了注册表里已经存在的元素。
代价是行为依赖运行时注册表：容器环境里必须校验元素与插件是否存在，否则协商失败发生在现场。

\subsection*
{手搭管线与 pad probe 取帧}
\begin{lstlisting}[style=cstyle,caption={GStreamer：手搭管线、在 sink pad 上挂 probe 取帧并读取时钟},label={lst:ch13-gst-probe}]
/* 目标：uridecodebin 解码 -> appsink，probe 只读统计不改数据 */
static GstPadProbeReturn
on_buf(GstPad *pad, GstPadProbeInfo *info, gpointer user_data)
{
    FrameStats *st = (FrameStats *)user_data;
    GstBuffer *buf = GST_PAD_PROBE_INFO_BUFFER(info);
    GstMapInfo map;

    if (gst_buffer_map(buf, &map, GST_MAP_READ)) {
        st->bytes += map.size;
        st->pts_us = (gint64)GST_BUFFER_PTS(buf) / 1000;   /* 纳秒转微秒 */
        gst_buffer_unmap(buf, &map);
    }
    /* 只要元数据时返回 OK 即可；DROP 会把这帧从链上摘掉 */
    return GST_PAD_PROBE_OK;
}

gboolean attach_probe(GstElement *sink, gpointer user_data)
{
    GstPad *pad = gst_element_get_static_pad(sink, "sink");
    if (!pad)
        return FALSE;
    gulong id = gst_pad_add_probe(pad, GST_PAD_PROBE_TYPE_BUFFER,
                                  on_buf, user_data, NULL);
    gst_object_unref(pad);
    return id != 0;
}

GstElement *build(const char *uri, GError **err)
{
    GError *local = NULL;
    char *desc = g_strdup_printf("uridecodebin uri=%s ! "
        "video/x-raw,format=I420 ! appsink name=sink sync=false drop=true", uri);
    GstElement *p = gst_parse_launch(desc, &local);
    if (local) { g_propagate_error(err, local); g_free(desc); return NULL; }
    g_free(desc);

    GstElement *sink = gst_bin_get_by_name(GST_BIN(p), "sink");
    GstCaps *want = gst_caps_from_string("video/x-raw,width=1920,height=1080");
    g_object_set(sink, "caps", want, NULL);   /* 具体 caps 由下游向上传 */
    gst_caps_unref(want);
    attach_probe(sink, g_new0(FrameStats, 1));
    gst_object_unref(sink);
    return p;   /* 调用方负责 set_state 与 unref */
}
\end{lstlisting}

probe 挂在链路上有四个维度：类型（BUFFER、EVENT、QUERY、REFCOUNT）、方向（upstream、downstream、
block 与非 block）、返回语义（\Tp{GST\_{\Bk}PAD\_{\Bk}PROBE\_{\Bk}OK}、\Tp{GST\_{\Bk}PAD\_{\Bk}PROBE\_{\Bk}DROP}、
\Tp{GST\_{\Bk}PAD\_{\Bk}PROBE\_{\Bk}REMOVE}）与阻塞性。
关键约束：pipeline 的锁在 probe 里不可重入，想在 probe 中改结构必须返回 \Tp{GST\_{\Bk}PAD\_{\Bk}PROBE\_{\Bk}BLOCK}
并用 \Fn{gst\_{\Bk}element\_{\Bk}send\_{\Bk}event} 一类的异步手段，或干脆只做入队、把工作丢给别的线程。
GStreamer 的 WebRTC 栈（\Tp{webrtcbin}）把 RTCP、DTLS、SRTP 拆成独立元素，靠 `internal` 池 pad 串起，
第 11、12 章讲的抖动与线程边界在这里变成元素拓扑问题。

\section{把 libav* 当依赖：构建与许可的坑}

\subsection*{configure 开关如何改变整份产品的许可}
\Lib{libav*} 本体按 LGPL 构建：默认即如此，前提是不把 GPL-only 的部分编进去。
难点在于``能力''与``许可''被同一个开关绑在一起：\Opt{enable-{\Bk}libx264} 与
\Opt{enable-{\Bk}libx265} 只有同时给出 \Opt{enable-{\Bk}gpl} 才会被 configure 接受
（x264 与 x265 自身是 GPL），于是这一条开关把整份 \Tp{avcodec}、连带所有链接它的代码
一起推进 GPL 空间。
对照组是 \Opt{enable-{\Bk}libopus} 与 \Opt{enable-{\Bk}libvpx}：同为 BSD-3 的库，
一样需要显式开启，但不要求 \Opt{enable-{\Bk}gpl}，加上之后产物许可不变。
第三种麻烦是 \Opt{enable-{\Bk}nonfree}：它放行 \Opt{enable-{\Bk}libfdk-{\Bk}aac} 这类带费条款的实现，
结果是产物在法律上不可自由再分发——不是``贵''，是``不能发''。
至于 AV1 侧的 \Lib{dav1d}（BSD-2）与 \Lib{libaom}（BSD-3），代码许可宽松不等于专利授权到位，
那是与构建开关无关的另一本账。
结论很硬：许可不是法务在发布前才读的一段 README，而是构建矩阵的一个维度，
应该在 CI 里以``构建产物、许可标签、允许的发布通道''三张表的映射形式固化。

\subsection*{静态、动态与符号版本}
动态链接是 LGPL 义务最容易满足的形态（给用户替换 \Tp{.{\Bk}so} 的可能性），也带来版本漂移风险：
libav* 的每个库都独立走 \Tp{major} 的 ABI 版本，头文件里的 \Tp{version.{\Bk}h} 与运行期
\Fn{avcodec\_{\Bk}version} 可以不一致，用 \Fn{avcodec\_{\Bk}parame{\Bk}ters\_{\Bk}alloc} 一族的代码
必须按编译期与运行期都校验来写。
静态链接把依赖焊死在二进制里，换来可复现，代价是 GPL 义务直接落到你的可执行文件上，
而且必须先链接 \Lib{libavutil}（其他库都依赖它），并显式写出 \Opt{enable-{\Bk}libx264} 一类的开关。

\begin{lstlisting}[style=shellstyle,caption={用 pkg-config 校验一个 GPL 构建是否真的可用},label={lst:ch13-pkgcheck}]
pkg-config --exists libavcodec libavformat libavutil && \
  echo "headers ok: $(pkg-config --modversion libavcodec)"
# 只看能力不看承诺：这一步决定后面能不能用 libx264 编解码
ffmpeg -hide_banner -buildconf | grep -E 'enable-(gpl|nonfree|libx264|libx265)'
ffmpeg -hide_banner -encoders | grep -E 'libx264|libx265|libaom|libopus'
\end{lstlisting}

\section{子进程边界：喂 stdin、读 stdout}

把 ffmpeg 当管道用，是所有集成方式里工程性价比最高的一档，前提是你把三条纪律刻进代码。
第一，视频侧必须显式声明像素几何：容器信息丢了就不可能从字节流反推，所以
`-f rawvideo -pix\_fmt yuv420p -s WxH -framerate 30` 这一串是合同的一部分。
第二，stdin 是管道时永远加 \Tp{-{\Bk}nostdin}（对 ffmpeg 自身）或走 \Tp{pipe:0} 显式地址，
否则子进程会吃掉你留给自己的终端输入。
第三，进度与统计不要抓 stderr 的文本，用 \Tp{-{\Bk}progre{\Bk}ss pipe:2} 拿到 `key=value` 行，
或在 C++ 侧用 \Tp{AVIO} 自定义回调把字节送到自己手上。

\begin{lstlisting}[style=shellstyle,caption={子进程管道：stdin 收 rawvideo，stdout 出 fragmented MP4},label={lst:ch13-pipe}]
# 渲染器每帧写 1920x1080 的 I420 到子进程 stdin；stdout 收 fmp4 字节
ffmpeg -hide_banner -loglevel error -nostdin -y \
  -f rawvideo -pix_fmt yuv420p -s 1920x1080 -framerate 30 -i pipe:0 \
  -f lavfi -i anullsrc=channel_layout=stereo:sample_rate=48000 \
  -c:v libx264 -preset veryfast -tune fastdecode -crf 23 \
  -maxrate 6000k -bufsize 12000k -g 60 -pix_fmt yuv420p \
  -c:a aac -b:a 128k -shortest \
  -movflags frag_keyframe+empty_moov+default_base_moof \
  -f mp4 pipe:1 > out.fragmented.mp4

# 反向：把一个来源的裸帧抽出来给自己的 GPU 上传，避免磁盘往返
ffmpeg -hide_banner -loglevel error -nostdin \
  -hwaccel auto -i in.mp4 \
  -vf scale=1280:720:format=nv12 -pix_fmt nv12 -f rawvideo pipe:1 | \
  ./my_uploader --width 1280 --height 720
\end{lstlisting}

子进程形态的延迟下限由三次调度决定：写满一次管道缓冲、子进程内部一帧的处理、读回一次。
管道是阻塞的，写侧不控制节奏就会在 64 KiB 一类的小缓冲上卡住整个渲染循环；
反过来读侧不排空 stdout，子进程会在写满时收到 SIGPIPE 而退出——这是集成里最常见的``莫名消失''。

\section{游戏引擎侧：视频播放的原生后端与纹理交付}

\subsection*{Unity：VideoPlayer 到 RenderTexture}
\Tp{VideoPla{\Bk}yer} 是引擎自带的播放器节点：源可以是 `VideoClip`（编辑器预转码的资源）或 URL，
输出形态由 `renderMode` 决定——相机平面、屏幕空间、`RenderTexture`、或远端纹理。
拿到纹理就走标准的 `Texture2D`/`RenderTexture` 路径，帧到达用 `frameReceived` 事件驱动，
时间轴由 `time`/`platformTime` 读写，`stepForward` 供剪辑器式逐帧控制，`prepare` 异步预热资源。
音轨用 `SetDirectAudioVolume`/`SetAudioVolume` 一族 API 调。
真正属于本章的结论是：Windows 上这个组件是媒体基础（Media Foundation）的包装，
移动端是各系统播放器包装；它给你的是纹理与时间轴，不是逐帧 CPU 可访问的像素，
也不是解码器线程控制权。需要 YUV 平面、需要硬解失败时可回退、需要外挂自己的解码器时，
正确做法是写一个 native plugin 把帧上传成纹理，或走引擎的自定义抛送视频路径。

\subsection*{Unreal：Media 框架的播放器抽象}
引擎侧的视频是一组接口：\Tp{UMediaPl{\Bk}ayer} 提供播放、暂停、跳转、准备与时间询问；
\Tp{UMediaTe{\Bk}xture} 把视频输出到一张 GPU 纹理；\Tp{UMediaSo{\Bk}undCompo{\Bk}nent} 把媒体里的音频
分派成引擎的音频组件；具体实现由平台插件提供——Windows 的媒体基础播放器、macOS/iOS 的
AVAsset 播放器、Android 播放器、媒体文件与图片序列播放器、以及 Blackmagic DeckLink 一类的采集播放器。
UI 侧通过 UMG 的图像控件绑定这张纹理，或经材质参数进任何网格。
要新增一种后端，就实现这套接口并注册进媒体框架，而不是在渲染代码里塞条件编译。
需要逐帧像素做视觉算法时，纹理的读回（CPU 可访问的映射）是昂贵的同步点；
更好的路径是让算法进 GPU（见引擎的着色器侧），或者用 `IMediaTexture` 的更新事件在同一帧内消费纹理。

\section{浏览器、引擎与许可：三条边界}

\subsection*{浏览器侧：MSE、WebCodecs、EME 的分工}
MSE（\Tp{MediaSou{\Bk}rce} 加 \Tp{SourceBu{\Bk}ffer}）把``谁来解析容器''交给你：页面 append 分片字节，
由浏览器的解码栈决定哪些类型能放——判断手段是 \Fn{MediaSou{\Bk}rce.{\Bk}isType{\Bk}Supporte{\Bk}d}，传
`video/mp4;codecs=avc1.640028,mp4a.40.2` 一类字符串。
它给你的是同步可控的缓冲与播放起点，不是像素；对第 14 章的 ABR 决策，这是唯一现实的接口层。
WebCodecs 才把解码器直接暴露给页面：\Tp{VideoDec{\Bk}oder} 收 \Tp{EncodedV{\Bk}ideoChunk}、吐
\Tp{VideoFra{\Bk}me}，可用 \Tp{Hardware{\Bk}Accelera{\Bk}tion} 偏好请求硬解，用 `latencyMode` 声明实时性；
\Tp{VideoFra{\Bk}me} 可拷贝到 CPU（\Tp{copyTo}），也可直接交给 WebGPU 或 Canvas（\Tp{importExternalTexture} 与 \Tp{createImageBitmap}），跨线程用 \Tp{postMessage} 转移所有权。
这是 2026 年做端到端低延迟 Web 管线的正确基座，但内存与线程成本仍在 JS 边界之外，
需要把帧尽快转成 GPU 资源，别在 JS 堆里持有。
EME（\Tp{navigator.{\Bk}reques{\Bk}tMediaKe{\Bk}ySystemA{\Bk}ccess}）解决的是``密钥与解码在受保护上下文里完成''，
密钥系统串形如 `com.widevine.alpha`、`com.microsoft.playready`，明文侧是 CENC 的 `cenc`/`cbcs` 模式。
它不改变延迟与码率决策，但会把可观测的度量遮住——度量口径必须区分可解与受保护内容。

\subsection*{WASM 的性能边界}
纯 \CPP{} 移植的解码器进 WASM 后，性能由四件事封顶：线性内存以 32 位寻址为主（4 GiB 上界），
SIMD 是固定 128 位宽（与 \Tp{AVX2} 的 256 位、\Tp{NEON} 的用法不能一一对应，需要重排循环），
没有可靠的逐帧硬件解码通道，多线程要 \Tp{SharedAr{\Bk}rayBuffe{\Bk}r} 与跨域隔离头（COOP、COEP）。
工程结论：把 WASM 用在容器解析、字幕、滤镜链与轻量编解码上，不要指望它在 4K 高帧率下
替代原生解码器；把重活留给 WebCodecs，或用 Web Worker 分帧，并接受由此而来的调度延迟。

\subsection*{引擎与浏览器之间：谁来持有解码器}
两条路线的取舍可以统一成一句话：如果视频是``场景里的内容''（过场、屏幕、实拍背景），
用引擎自带的播放器，接受后端与许可；如果视频是``要参与算法的数据流''（绿幕、串流进纹理、
逐帧分析），就在原生层持有解码器，把帧交付点做成引擎的自定义纹理源。
混用会引出两套时钟（引擎时间轴与媒体时间轴），这是音画同步问题最隐蔽的来源之一。

\begin{guideline}[title={形态判据}]
\begin{itemize}
\item 需要逐帧像素、亚帧延迟、解码器线程控制：进进程内链接，同时接受许可、崩溃半径与版本矩阵三项成本。
\item 只需要``文件进、文件出''，且要求故障可隔离、依赖可独立替换：用子进程管道，合同写在命令行上。
\item 需要容量与灰度、可容忍秒级延迟：独立服务，别把转码放进客户端进程。
\item 面向浏览器分发：MSE 管缓冲、WebCodecs 管解码与帧、EME 管受保护内容、WASM 只承担边界内的轻活。
\item 许可是构建维度，不是发布检查项：\Opt{enable-{\Bk}gpl} 与 \Opt{enable-{\Bk}nonfree} 一旦打开，整份产物跟着变。
\end{itemize}
\end{guideline}

% ============================================================
% ---- frag_ch14.tex ----
\chapter{传输质量与自适应码率}
\label{ch:ch14-congestion-abr}

第 11、12 章把流送出去、把抖动压住；本章解决剩下的两个问题：网络到底能给多少，
以及该用多少。带宽估计是传感器，拥塞控制是执行器，ABR 是在分片粒度上的决策器，
度量口径是判据。四件事都在 C/C++ 侧落地——编码器码率热更新、pacer 队列、
分片时长与关键帧间距，全是可以直接写进代码的旋钮。

\section{带宽估计的两条腿与 pacing 缓冲}

\subsection*{速率模型：把速率当成对偶变量}
Reno/CUBIC 一族的 TCP 把可用带宽看成丢包的对偶：不丢包就加大，丢了就乘性退让，
稳态是一条锯齿线。锯齿意味着两件事——带宽被系统性地多占了一个缓冲的深度，
以及``等待丢包''本身就是延迟。
BBR 走的是另一条路：用最大带宽滤波器（\Tp{max\Bk-{\Bk}filter}，窗约 10 个 RTT）取 \Tp{Bt}，
用最小 RTT 滤波器（窗约 15 秒）取 \Tp{Rt}，然后把发送速率钉在 \Tp{Bt} 上、
把拥塞窗当成 \Tp{Bt} 乘 \Tp{Rt} 的 BDP 加一个增益，探测阶段按周期轮转增益，
\Tp{ProbeRTT} 阶段定期把队列抽干以校准 \Tp{Rt}。
它不需要丢包作为信号，因此对浅缓冲网络更友好；但对``谁在同一瓶颈上''的公平性
与对突发的响应速度是另一类问题——v2 起把丢包重新引入为辅助信号。

\subsection*{延迟梯度：在丢包之前看见排队}
实时链路的瓶颈通常在接入网，可用缓冲很浅：一旦排队，代价立刻体现在端到端延迟上，
而不是体现在丢包上。
延迟梯度估计量是``分组到达时延差''（grouped inter-arrival delta，IGI）。
把包按发送时刻成组（突发裁剪：组内发送时刻差小于阈值就并成一组，典型阈值 5 毫秒），
每组的梯度是 ``到达时刻差减去发送时刻差''，再扣掉这组包本身需要的传输时间
（按当前目标速率折算）。
一条斜率为零的梯度说明没有新排队；斜率持续为正说明队列在长。
判决用指数平滑后的趋势线与一个自适应阈值比较，输出三态：过使用、正常、欠使用。
阈值随样本增长上抬、随负梯度回落，量级上限比下限高两个数量级，这个自适应是
``不把正常波动当拥塞''的关键。

\subsection*{到达速率、窗与突发}
纯到达速率（窗内收到字节除以窗时长）是最稳但最迟钝的信号，它会把突发平均掉。
分组对（packet-pair）与分组链（packet-train）试图从两个背靠背包的到达间隔读出
瓶颈容量，代价是极易被交叉流量与网卡聚合（LRO、GRO）破坏——这也是很多
自研估计器在实验室准、在现场不准的原因。
工程上通常三者并用：延迟梯度定方向，到达速率定范围，探测包定上限。

\begin{lstlisting}[style=cppstyle,caption={接收侧延迟梯度与到达速率估计的最小骨架},label={lst:ch14-trendline}]
struct Group {
  uint64_t send_us;      // 组内首包的发送时刻（由 transport-wide 时戳或 SR 反射）
  uint64_t arrival_us;   // 组内首包的本地单调到达时刻
  int64_t  bytes;        // 组内累计 RTP 载荷字节
  int      packets;
};

class DelayGradient {
 public:
  static constexpr int    kMinSamples      = 20;     // 样本不足不判决
  static constexpr int    kMaxWindowPackets = 200;   // 梯度窗上限
  static constexpr double kGroupThresholdMs = 5.0;   // 突发裁剪阈值
  static constexpr double kMinThresholdMs   = 6.0;   // 判决阈值上下限
  static constexpr double kMaxThresholdMs   = 600.0;
  static constexpr double kSmoothing        = 0.99;  // 趋势线平滑系数

  // 返回三态：>0 过使用（队列在长），<0 欠使用，0 正常
  int Update(uint64_t send_us, uint64_t arrival_us, uint32_t payload_bytes,
             double target_rate_bps) {
    // 1) 成组：与当前组太近就并入，否则滚动
    if (cur_.packets == 0) {
      cur_ = {send_us, arrival_us, (int64_t)payload_bytes, 1};
      return 0;
    }
    const double gap_ms = double(send_us - cur_.send_us) / 1000.0;
    if (gap_ms < kGroupThresholdMs) {
      cur_.bytes += payload_bytes;
      cur_.packets++;
      return 0;
    }
    const Group g = cur_;
    cur_ = {send_us, arrival_us, (int64_t)payload_bytes, 1};

    // 2) IGI：到达时延差，扣掉这组字节本应占用的发送时间
    const double igi_ms =
        (double(g.arrival_us) - double(g.send_us)) / 1000.0 -
        (double(prev_.arrival_us) - double(prev_.send_us)) / 1000.0 -
        double(g.bytes - prev_.bytes) * 8e6 / target_rate_bps;  // 字节转毫秒

    // 3) 回归斜率：对窗内 (t_ms, igi_ms) 样本做最小二乘，斜率单位是毫秒每毫秒
    win_.push_back({igi_ms, double(g.arrival_us) / 1000.0});
    if (win_.size() > size_t(kMaxWindowPackets))
      win_.pop_front();
    const double slope = LeastSquaresSlope(win_);
    smoothed_ = kSmoothing * smoothed_ + (1.0 - kSmoothing) * slope;

    // 4) 到达速率：组字节累进速率窗，窗满结算一次，单位比特每秒
    if (rate_start_us_ == 0)
      rate_start_us_ = g.arrival_us;
    rate_bytes_ += g.bytes;
    const int64_t span_us = int64_t(g.arrival_us - rate_start_us_);
    if (span_us >= kRateWindowUs) {
      arrival_rate_bps_ = double(rate_bytes_) * 8e6 / double(span_us);
      rate_bytes_ = 0;
      rate_start_us_ = g.arrival_us;
    }

    const double dt_ms = last_update_us_ == 0
        ? 0.0 : double(g.arrival_us - last_update_us_) / 1000.0;
    prev_ = g;
    last_update_us_ = g.arrival_us;
    if (int(win_.size()) < kMinSamples || dt_ms <= 0.0)
      return 0;

    // 5) 三态判决与阈值自适应：阈值只在判决有效时移动
    if (smoothed_ > threshold_ && slope > 0.0) {
      threshold_ = std::min(threshold_ + 0.01 * dt_ms, kMaxThresholdMs);
      return +1;
    }
    if (smoothed_ < -threshold_) {
      threshold_ = std::max(threshold_ - 0.001 * dt_ms, kMinThresholdMs);
      return -1;
    }
    return 0;
  }

  double arrival_rate_bps() const { return arrival_rate_bps_; }

 private:
  static constexpr int64_t kRateWindowUs = 100000;  // 速率窗 100 毫秒

  std::deque<std::pair<double, double>> win_;  // (igi_ms, t_ms)
  Group cur_{}, prev_{};
  double smoothed_ = 0.0;
  double threshold_ = kMinThresholdMs;
  uint64_t last_update_us_ = 0, rate_start_us_ = 0;
  int64_t rate_bytes_ = 0;
  double arrival_rate_bps_ = 0.0;
};
\end{lstlisting}

常量与阈值在真实实现（WebRTC 的 \Tp{Trendlin{\Bk}eEstimat{\Bk}or}、Kalman 变体、\Tp{DelayBas{\Bk}edBwe}）
里各有自适应回路，上面的数值是量级参考，不是规范。
两个单位陷阱值得单独警告一次，见框。

\begin{warning}[title={时戳与单位的两个经典错误}]
\begin{itemize}
\item RTP 时间戳是媒体时钟（视频 90 kHz、Opus 48 kHz）且会 16/32 位回绕；
      做梯度必须用单调时钟或扩展序列号，把 90 kHz 的差值当毫秒用会放大 90 倍。
\item 突发聚合（网卡的 LRO、GRO、Wi-Fi 的重传合并）会把多个包的到达时刻压成同一时刻，
      梯度瞬间``看起来没有排队''。测延迟梯度时先在网卡侧关掉聚合，
      或在估计时对同一到达时刻的包做去突发处理。
\end{itemize}
\end{warning}

\section{从估计到发送：AIMD、pacer 与关键帧请求}

估计值不等于发送速率。发送侧还要三件事：一个把三态判决变成速率调整的 AIMD，
一个把速率摊成包列的 pacer，一组在丢了之后决定``要什么回来''的请求器。

\begin{lstlisting}[style=cppstyle,caption={估计融合、AIMD 与 pacer 队列的最小骨架},label={lst:ch14-aimd-pacer}]
struct RateState {
  double estimate_bps = 300e3;   // 初始值来自信令或历史，不要从 0 起
  double floor_bps    = 30e3;
  double ceiling_bps  = 20e6;
  double rtt_ms       = 50.0;    // 由 RTCP SR 与 RR 的 DLSR 段算出
  double loss_frac    = 0.0;     // 统计窗（约 20 到 100 个包）内的丢包比例
  double probe_factor = 1.5;     // 探测上限 = 目标速率乘该系数，典型 1.5 到 2.5

  // 三态判决与到达速率共同决定新估计；系数是典型量级
  double Update(int grad_state, double arrival_rate_bps) {
    double est = estimate_bps;
    if (grad_state > 0) {                       // 过使用：乘性退让
      est *= 0.95;
    } else if (grad_state < 0 && loss_frac < 0.08) {
      est *= 1.08;                              // 欠使用：每 RTT 一次加性到乘性爬升
    } else if (loss_frac > 0.10) {
      est *= 0.85;                              // 丢包主导：更狠的退让
    }
    if (arrival_rate_bps > 0)                   // 到达速率给一个可验证的下界
      est = std::max(est, 0.5 * arrival_rate_bps);
    estimate_bps = std::clamp(est, floor_bps, ceiling_bps);
    return estimate_bps;
  }

  double paced_rate_bps() const { return estimate_bps * probe_factor; }
};

struct Pacer {
  double rate_bps;              // 通常 = 1.5 倍估计，留突发余量
  int64_t interval_budget_bits; // 本 interval 可用预算（欠账可为负）
  int64_t media_budget_bits;    // 媒体预算，只允许借还款、不允许凭空造
  int padding_bitrate_bps;      // 维持恒定速率的填充目标
  uint32_t interval_ms = 5;     // 发包节拍，典型 5 毫秒或按 RTT 的分数

  bool CanSend(uint32_t packet_bits) const {
    return interval_budget_bits >= -int64_t(packet_bits) &&
           media_budget_bits >= 0;              // 关键：不看包长即可为负预算放行
  }
  void OnSend(uint32_t packet_bits) { interval_budget_bits -= packet_bits;
                                      media_budget_bits    -= packet_bits; }
  void Tick() {                                  // 每 interval 结算一次
    const int64_t add = int64_t(rate_bps * interval_ms / 1000);
    interval_budget_bits = std::min(interval_budget_bits + add, add);
    media_budget_bits    = std::min(media_budget_bits + add, add / 2);
  }
};
\end{lstlisting}

关键帧请求走 RTCP：PLI（通用 RTCP 反馈，PSFB 且 FMT 为 1）说``给我一帧可解的''，
FIR（RFC 5104 一类，全量刷新）在加入流、解码器复位这类``整帧都不可信''的时刻用。
两者都必须限频——多个订阅者各自请求会造成关键帧风暴，而关键帧是整条流里最贵的包。
接收侧还应配合 SLI、RPSI 之类（RFC 4585 定义的条带与参考点指示）把请求粒度缩小，
以及用 ``先请求一次，若仍不可解再请求一次'' 的退避而不是每秒重复请求整帧。

丢包恢复三件套的分工如下。
NACK（RTPFB、FMT 为 1，携带 PID 加位掩码；重传走独立 SSRC 的 RTX 负载类型，
原序列号以 16 位头前置）是按需、精确、低开销的，代价是恢复耗时一个 RTT。
FEC 分两层：媒体内 RED 与 RTP 级 ULPFEC（RFC 6342，保护相邻包）、FlexFEC
（把保护关系写进扩展头，可跨帧、按时间窗布局），代价是恒定带宽税（典型 20\% 到 50\% 的冗余）。
关键帧请求恢复最彻底，也最贵，且会打乱码率。
三者组合的选择取决于 RTT 与抖动预算：RTT 小的时候优先 NACK，因为它便宜；
RTT 大（例如 150 毫秒以上）时，重传要额外花一个 RTT，只有把 jitter buffer 加深到
``能容纳一次往返''才能用，而深缓冲本身毁掉交互性——这时 FEC 的代价是固定的、
与 RTT 无关的包内延迟，于是更划算。
反过来说，在无损或轻度随机丢包且延迟预算宽裕的地方，FEC 就是纯粹的浪费。
FlexFEC 的时间跨度式布局（一行保护跨若干时间槽）正是为高 RTT 加突发丢包准备的，
而 ULPFEC 的一前一后布局在这种场景里几乎无效。

\section{ABR 决策：分片粒度上的速率选择}

\subsection*{throughput-based 与 buffer-based 的行为差异}
吞吐型估计分片下载耗时算速率，通常对每个分片取若干样本做加权（快样本保守、慢样本敏感，
避免把缓存命中算成带宽），再选下一个分片层级。它响应快、在静止网络上稳，
但两个结构性缺陷：速率噪声直接传到画质，且它只看``刚下载的这段''，不看下游有多少帧。
缓冲型决策（BOLA 一族及其变体）把目标写成缓冲水平的效用函数，与瞬时速率几乎无关，
因此在多流竞争时能收敛到一个稳态分配；代价是对网络变化响应慢，
在网络突然好转时``不敢上''，以及需要按分片时长调参。
混合策略（先由缓冲定范围，再由吞吐定具体档）是目前最实用的做法。

\subsection*{切换滞回与抖动}
滞回有三条正交的旋钮：时间（两次切换至少隔 N 个分片）、电平（缓冲高于某阈值才允许升级、
低于某阈值才强制降级）与代价（升级要求预测收益超过一个百分点门槛）。
只有一条时抖动仍在。LL-HLS 场景更麻烦：分片短到 1 秒、部分分片（partial segment）
短到 200 毫秒，采样变密，单位时间的切换次数自然上升。
对策是把决策频率与分片时长解耦：决策周期至少覆盖 2 到 3 个部分分片，
并在决策里带上``一次只走一档''的约束，配合 side-switching 一类的做法在关键帧边界切换。

\subsection*{启动、seek 与受限设备}
启动阶段用低层级换首帧时间（``快进式''逐级上升），seek 之后重挂 SourceBuffer
会让缓冲短暂归零，此时若按缓冲型决策会强制降到最低档——需要一个``seek 后的宽限期''。
设备解码能力（分辨率上限、档位数、CPU 占用）应作为档位过滤条件，
不要交给运行时估计；把 4K 档留在候选里但实际解不动，是最常见的一类``看起来像网络问题''的故障。

\begin{longtable}{>{\raggedright\arraybackslash}p{0.16\textwidth}
                  >{\raggedright\arraybackslash}p{0.15\textwidth}
                  >{\raggedright\arraybackslash}p{0.15\textwidth}
                  >{\raggedright\arraybackslash}p{0.15\textwidth}
                  >{\raggedright\arraybackslash}p{0.18\textwidth}}
\caption{算法族对比：收敛速度、稳定度、适用协议与已知副作用}\label{tab:ch14-algo-family} \\
\toprule
算法族 & 收敛速度 & 稳定度 & 适用协议 & 主要副作用 \\
\midrule
\endfirsthead
\toprule
算法族 & 收敛速度 & 稳定度 & 适用协议 & 主要副作用 \\
\midrule
\endhead
\midrule
\endfoot
\bottomrule
\endlastfoot
AIMD 加丢包信号（Reno 类） & 慢，锯齿周期正比 RTT & 高，行为可预测 &
TCP 与 QUIC & 把缓冲填满才知边界，排队延迟高 \\
CUBIC & 快（尤其远离拥塞点时） & 中，多流公平但短时波动大 & TCP 与 QUIC &
浅缓冲下丢包频发，与 ABR 叠加易振荡 \\
BBR v1、v2 & 很快，启动即靠最大带宽滤波收敛 & 中高，对浅缓冲友好 & TCP 与 QUIC &
对同瓶颈的 CUBIC 流不友好，v1 有公平性问题 \\
延迟梯度加 AIMD（GCC 一族） & 几秒级，先建立梯度再定标 & 高 &
UDP 加 RTP、RTCP 反馈、DataChannel & 探测流量与填充包占预算，参数依赖现场 \\
分组对与分组链 & 单程即可给容量、收敛最快 & 低，交叉流量与网卡聚合破坏 &
UDP 探测流、DASH 或 QUIC 的显式探测 & 误判突发为容量 \\
吞吐型 ABR & 一个到两个分片 & 低，噪声直接进画质 & HTTP（HLS、DASH） &
切换抖动、升级过冲引起 rebuffer \\
缓冲型 ABR（BOLA 一族） & 慢到收敛 & 高，与速率近乎无关 &
HTTP、LL-HLS、DASH & 网络转好时保守，需按分片时长重调参 \\
\end{longtable}

\section{度量：延迟、rebuffer、画质与它们的可复现性}

端到端延迟必须锚定在一个物理时刻上。采集链路用``曝光或采样时刻''打戳（相机时间戳或
单调钟减去帧龄），播放侧在.presented frames 时刻记录；两者之差才是可比较的量。
用 ``推流端发送时刻减拉流端收到时刻'' 只能测网络段，会漏掉编码与渲染两侧的排队。
rebuffer 比 = 卡顿时长除以（媒体时长加卡顿时长），卡顿次数是每次停顿事件计数
（阈值通常取 200 毫秒，低于它的人眼几乎不可察），两个指标都要报分位数而不是均值。
画质侧，VMAF 要按播放时间加权（``看了多久''而不是``编了多久''），
并用各档位的 VMAF 加权平均来比较 ABR 策略；裸平均会被极少播放的高档样本污染。
可复现性靠三件事：固定随机种子与网络损伤配置、保存原始逐分片日志、
把测试用的参考源与位深/色度格式钉死（VMAF 对 8 位与 10 位、对色度子采样都敏感）。

\begin{lstlisting}[style=shellstyle,caption={读流、读分片与读 RTP 头：一条排障组合},label={lst:ch14-probe}]
# 1) 流本身：分辨率、帧率、码率、关键帧间距（第 7 章的 GOP 旋钮在这里核对）
ffprobe -hide_banner -v error -select_streams v:0 \
  -show_entries stream=codec_name,profile,width,height,r_frame_rate,bit_rate \
  -show_entries frame=pict_type,key_frame -of json in.mp4 | jq '.frames | length'
mediainfo --Inform="Video;%BitRate% kbps %FrameRate% fps %CodecID%" in.mp4

# 2) 抓 RTP 头：序列号间隔、时间戳差、去掉 8 字节 UDP 头后的载荷长
tshark -i any -f "udp port 40000" -Y rtp -T fields \
  -e frame.time_epoch -e rtp.seq -e rtp.timestamp -e rtp.marker -e udp.length \
  | awk 'NR==1{t0=$1} {printf "%9.2f ms seq=%6d ts=%11d m=%d size=%d\n", \
         ($1-t0)*1000, $2, $3, $4, $5-8}'

# 3) RTCP 侧：反馈包类型与 FMT 计数，确认 NACK(205) 与 PLI(206) 的真实频率
tshark -r cap.pcapng -Y rtcp -T fields \
  -e frame.time_relative -e rtcp.type -e rtcp.fmt -e rtcp.ssrc \
  | sort | uniq -c | sort -rn | head
\end{lstlisting}

\section{回到 C/C++ 侧要动的旋钮}

估计值要变成发送行为，中间是四个可动的旋钮。

第一个是编码器码率热更新。最干净的是内建实时栈的显式接口（例如编码器抽象上的
速率设置入口，参数带目标 bps 与网络反馈时间戳）。
走 FFmpeg 链路时，\Tp{AVCodecC{\Bk}ontext} 的 \Tp{bit\_{\Bk}rate}、\Tp{rc\_{\Bk}max\_{\Bk}rate}、
\Tp{rc\_{\Bk}buffer\_{\Bk}size} 在打开编码器时定标，运行中改这些字段的可见性依封装实现而定：
\Lib{x264} 提供原地重配的通道（其头文件里的重配置入口，配合新的参数集），
而多数 FFmpeg 的外壳只在你重新打开编码器后生效。
工程上稳妥的做法是：把码率变化点钉在 GOP 边界，先关闭当前编码上下文再按新参数重开，
或者预留多个档位实例切换，代价是一帧的抖动与一次 IDR。
不要把 ``改了字段但没重开'' 和 ``生效了'' 当成同一件事。

第二个是分片粒度。分片时长决定 ABR 的反应粒度与首帧延迟，也决定 HTTP 侧的请求数量、
LL-HLS 的 PART 数量与播放器端的清单体积。
低延迟目标下用 1 到 2 秒整片加 200 到 600 毫秒部分分片，配合 chunked transfer
（HTTP 分块传输编码，服务端侧语义见第 11 章）；
纯点播则用 4 到 6 秒减少请求数。分片必须落在可解码边界（CMAF 的 fragment 首帧为
独立解码帧），否则切换档位会撕裂。

第三个是关键帧间距与可切换性。间距越短，档位切换越及时、码率越不稳；
典型实时链路取 1 到 2 秒，配合``按需 IDR''通道在请求时提前产生关键帧。
第四个是 pacer 与包大小：MTU 探测、单包不超过路径 MTU 的余量、
把 FEC 包与媒体包按同一速率预算排队（否则 FEC 会在拥塞时火上浇油）。

\begin{guideline}[title={传输质量的四个不变式}]
\begin{itemize}
\item 估计器只能看见你喂给它的时钟：单位与单调性错一次，后面所有参数都是在拟合噪声。
\item 信号要有两条腿：只看丢包会等出缓冲，只看延迟会被突发骗；两者不一致时的取舍要写死在代码里。
\item ABR 的稳定性来自滞回与分片时长的匹配，不来自更聪明的模型。
\item 恢复手段的成本函数不同：NACK 花一个 RTT，FEC 花恒定带宽，IDR 花一次码率尖峰——
      按 RTT 与延迟预算选，不要三个同时全开。
\end{itemize}
\end{guideline}

% ============================================================
% ---- frag_ch15-delivery.tex ----
\chapter{字幕、图文轨与加密投递：清单、表示与密钥}
\label{ch:ch15-delivery}

前十几章把像素与采样送上了线；本章处理三条常被当成附属品的轨道：字幕、图文数据与
加密后的可投递形态。它们的共同点是``不参与解码主循环，却决定用户能不能读懂、能不能选、
能不能在授权内看到画面''。三条谱系的差别不在编码复杂度，而在时间轴由谁给、呈现责任落在
解码器还是渲染器；清单与表示则决定客户端在哪里换轨、在哪里取密钥。分片粒度与档位选择的
算法在第 \ref{ch:ch14-congestion-abr} 章，时间基与错序的账在第 \ref{ch:ch10-mux-sync} 章，
本章只把结构与字段列清楚。

\section{三条轨道谱系：时间轴与呈现责任的归属}
\label{sec:ch15-lineage}

文本式字幕（\Tp{SRT}、\Tp{WebVTT}）把一条提示写成``起止时刻加一段字符''，时间轴与呈现责任
都在终端：字号、换行、字体回退由播放器自己定。标记式（\Tp{TTML}、\Tp{IMSC} 1.1、\Tp{EBU-TT-D}）
仍是文本，但把行框、对齐、字体族与透明度这类呈现意图写进属性与样式，终端只负责执行。图像式
（\Tp{VobSub}、DVD 的 SPU 与蓝光 PGS）没有文本语义：字形已由内容侧栅格化，终端只贴位图。于是
文本式小、可搜索、可被辅助技术朗读但跨终端换行不一致；图像式在哪都长一样，换来体积与不可访问。

FFmpeg 侧三条谱系共用媒体类型 \Tp{AVMEDIA\_{\Bk}TYPE\_{\Bk}SUBTITLE}，靠编码 id 区分：
\Tp{AV\_{\Bk}CODEC\_{\Bk}ID\_{\Bk}SUBRIP}、\Tp{AV\_{\Bk}CODEC\_{\Bk}ID\_{\Bk}WEBVTT}、\Tp{AV\_{\Bk}CODEC\_{\Bk}ID\_{\Bk}MOV\_{\Bk}TEXT}
（3GPP 计时文本，MP4 侧的 \Tp{tx3g}）、\Tp{AV\_{\Bk}CODEC\_{\Bk}ID\_{\Bk}DVB\_{\Bk}SUBTITLE}。解码出口也统一：
\Fn{avcodec\_{\Bk}decode\_{\Bk}subtitle2} 交出 \Tp{AVSubtitle}，可靠字段是 \Tp{num\_rects}、\Tp{rects}、
以毫秒计的 \Tp{pts} 与相对它的 \Tp{start\_{\Bk}display\_{\Bk}time}、\Tp{end\_{\Bk}display\_{\Bk}time}。矩形自带
分支类型与位置尺寸：文本分支给字符或 SSA/ASS 标记串，位图分支给游程数据加调色板。清单
\ref{lst:ch15-subtitle-rect} 是遍历骨架，四条谱系对照见表 \ref{tab:ch15-lineage}。

\begin{lstlisting}[style=cstyle,caption={遍历 \Tp{AVSubtitle}：一条提示的时间窗与三种呈现分支},label={lst:ch15-subtitle-rect}]
static int dump_sub(AVCodecContext *c, AVPacket *p, AVSubtitle *s) {
  int got = 0, ret = avcodec_decode_subtitle2(c, s, &got, p);
  if (ret < 0 || !got) return ret;      /* 未出完整提示时继续喂包 */
  /* s->pts 已是毫秒：容器时基的换算在解封装侧做过，不再乘 time_base */
  int64_t t0 = s->pts + s->start_display_time;
  int64_t t1 = s->pts + s->end_display_time;    /* 窗宽为 0 表示由下一条截断 */
  for (unsigned i = 0; i < s->num_rects; i++) {
    AVSubtitleRect *r = s->rects[i];   /* r->type 分位图 / 纯文本 / SSA-ASS 三种分支 */
    printf("cue %" PRId64 "..%" PRId64 " rect%u type=%d %dx%d@%d,%d\n",
           t0, t1, i, r->type, r->w, r->h, r->x, r->y);
  }
  avsubtitle_free(s);                  /* 按分支释放后置零，可复用 */
  return ret;
}
\end{lstlisting}

文本输入的字符集问题在 \Tp{AVCodecContext} 上：\Tp{sub\_charenc} 指定源字符集（如
\Tp{ISO-8859-1}、\Tp{GB18030}），\Tp{sub\_charenc\_mode} 决定真转换还是原样交出。UTF-8 来源
不需要设置；老式广播字幕与磁盘 \Tp{SRT} 常是单字节码页，不设就出现``字节对但字错''的乱码。

\Lib{libass} 是标记式与 SSA/ASS 的事实路径，C 侧组合三件事：初始化 \Fn{ass\_{\Bk}library\_{\Bk}init}
与 \Fn{ass\_{\Bk}renderer\_{\Bk}init}；用 \Fn{ass\_{\Bk}set\_{\Bk}fonts} 挂字体映射，较新版本另提供
\Fn{ass\_{\Bk}process\_{\Bk}css} 把外部样式折进内部样式对象；喂内容用 \Fn{ass\_{\Bk}process\_{\Bk}codec\_{\Bk}private}
与 \Fn{ass\_{\Bk}process\_{\Bk}data}。渲染侧 \Fn{ass\_{\Bk}render\_{\Bk}frame} 按毫秒时刻返回图像链表，每项带位图
指针、步幅、目标坐标与颜色，终端只做带透明度的叠加；图像式不走这条路，位图从解码器直接来。

\begin{warning}[title={字幕的时间基来自容器，不来自帧}]
\Tp{AVSubtitle.pts} 与 \Tp{start\_display\_time} 来源不同：前者由解封装按流的时间基换算
（TS 里是 90 kHz 的分组基本流时戳，MP4 里是轨道时基，Matroska 里是段时基），后者是解码器
给出的毫秒偏移。帧率、丢帧与解码延迟只影响``这条提示何时被取出''，不影响它该显示在什么时刻。
用帧序号或渲染帧计数当字幕时钟，会在慢放、倍速与可变帧率源上稳定积累偏移。
\end{warning}

\begin{longtable}{@{}>{\raggedright\arraybackslash\footnotesize}p{0.15\textwidth} >{\raggedright\arraybackslash\footnotesize}p{0.20\textwidth} >{\raggedright\arraybackslash\footnotesize}p{0.21\textwidth} >{\raggedright\arraybackslash\footnotesize}p{0.22\textwidth}@{}}
\caption{三条字幕谱系的时间轴与呈现责任}\label{tab:ch15-lineage} \\
\toprule
谱系 & 时间轴来源 & 呈现责任 & 典型标识 \\
\midrule
\endfirsthead
\toprule
谱系 & 时间轴来源 & 呈现责任 & 典型标识 \\
\midrule
\endhead
\bottomrule
\endfoot
\endlastfoot
文本式 & 提示自带起止时刻 & 终端全权：字号、换行、字体回退 & \Tp{AV{\Bk}\_CODEC{\Bk}\_ID{\Bk}\_SUBRIP}、\Tp{AV{\Bk}\_CODEC{\Bk}\_ID{\Bk}\_WEBVTT} \\
标记式 & 提示起止加样式作用域 & 作者定呈现，终端执行 & \Tp{TTML}、\Tp{IMSC} 1.1、\Tp{AV{\Bk}\_CODEC{\Bk}\_ID{\Bk}\_ASS} \\
图像式 & 包时戳，无文本时长语义 & 作者已栅格化，终端只贴 & DVD SPU、\Tp{AV{\Bk}\_CODEC{\Bk}\_ID{\Bk}\_DVB{\Bk}\_SUBTITLE} \\
闭环捕获 & 视频消隐期内到达 & 抽取侧负责与视频帧配对 & \Tp{CEA-608}/708，随包 \Tp{AV{\Bk}\_PKT{\Bk}\_DATA{\Bk}\_A53{\Bk}\_CC} \\
\end{longtable}

\section{图文数据与遥测通道：包级时间戳与可过期性}
\label{sec:ch15-data-sections}

第二路数据轨不发文字，发结构。MPEG-2 Systems 的 section 机制是代表：节目分配表与节目映射表
给出节目到基本流的映射，业务描述表给业务名与运营者名，事件信息表给节目时间表，再往上是私有
section——表号留出一段给运营方自定义，携带版本、区段与校验。FFmpeg 的 TS 解封装把前两者映射
成 \Tp{AVProgram} 与 \Tp{AVStream}；后几类与私有 section 要按原文自行解析，交出的是字节而不是语义。

另一条路线是广播与工业遥测：\Tp{KLV}（键长值）三元组承载 \Tp{MISB} ST 0601 系列的无人机
通用元数据，键是固定长度标识加一项清单编号，值是平台姿态、航向、目标标识这类字段；它在 TS
里通常走私有基本流，并用 SMPTE 时间码或相对时戳把自己钉到视频帧上。CEA-608/708 的用户数据
（含增强型通用容器里的注册用户数据字节）也属此类：数据寄生在视频基本流里而不是独立基本流，
容器侧必须用 side data 保留配对关系，见清单 \ref{lst:ch15-sidedata}。

\begin{lstlisting}[style=cstyle,caption={从视频包与字幕包里取出随附 side data：图文与捕获字幕的配对},label={lst:ch15-sidedata}]
static void scan_side_data(AVPacket *p, int64_t video_pts_ms) {
  const uint8_t *d = NULL;
  size_t sz = 0;
  /* 闭环捕获：608 字节对成对携带，随视频帧到达，自身不带绝对时戳 */
  d = av_packet_get_side_data(p, AV_PKT_DATA_A53_CC, &sz);
  if (d) feed_cc_decoder(d, sz, video_pts_ms);
  /* WebVTT 的标识行与设置行分别作为独立 side data 附着在提示包上 */
  d = av_packet_get_side_data(p, AV_PKT_DATA_WEBVTT_IDENTIFIER, &sz);
  if (d) note_identifier((const char *)d);
  d = av_packet_get_side_data(p, AV_PKT_DATA_WEBVTT_SETTINGS, &sz);
  if (d) note_settings((const char *)d);
}
\end{lstlisting}

为什么这类通道要按``包级时间戳加可过期''设计？因为它没有内在时长。字幕提示天然带起止
（或由下一条截断），过期时刻是内容的一部分；遥测 section 表达的是``采样那一刻状态如此''，
有效期由消费者决定：版本不变就重复携带同一事实，接收端要么按轮询间隔判失效，要么靠
\Tp{current\_next\_indicator} 与版本切换丢弃旧值。工程后果有三条。冗余是设计目标：节目分配表
与节目映射表以固定周期反复插入，接收端可在任意时刻入网拿到完整拓扑，代价是调谐后的信息
真空期。绑定粒度不同：遥测常以一次突发或访问单元为键。丢弃策略不同：字幕落后按时间轴跳过。

\begin{note}[title={图文数据不要伪装成字幕轨}]
把事件信息表风格的私有 section 或 \Tp{KLV} 清单塞进字幕轨，能借到字幕的时间戳与容器支持，
但会污染语义：播放器会把它按字幕排队、按字幕时长显示，辅助技术还会朗读出来。需要与画面绑定
的遥测应走独立基本流加 side data，或事件消息盒（\Tp{emsg}）这类事件通道，并在清单里显式声明
它不是可读字幕。与第 \ref{sec:ch15-lineage} 节对照：差别不是字节形状，而是谁判定有效期。
\end{note}

\section{多轨道、语言与呈现取舍}
\label{sec:ch15-multitrack}

轨道身份有两层：容器内的标识与用途的标识。前者是 TS 的基本流号、MP4 的轨道号、Matroska 的
轨道序号与全局唯一标识；FFmpeg 侧流对象上有一个整型标识字段承载格式专用号，另有一个字符串
形式的标识符字段。后者是一组位标志 \Tp{AV\_{\Bk}DISPOSITION\_{\Bk}*}，常用取值与各自含义如下：
\begin{itemize}
  \item \Tp{AV\_{\Bk}DISPOSITION\_{\Bk}DEFAULT}：入网时的缺省选择。
  \item \Tp{AV\_{\Bk}DISPOSITION\_{\Bk}ORIGINAL}：标原始配音。
  \item \Tp{AV\_{\Bk}DISPOSITION\_{\Bk}DUBBING}：标配音版本。
  \item \Tp{AV\_{\Bk}DISPOSITION\_{\Bk}COMMENT}：给评论轨。
  \item \Tp{AV\_{\Bk}DISPOSITION\_{\Bk}DESCRIPTION}：给音频描述，即为视障者补述画面的轨。
  \item \Tp{AV\_{\Bk}DISPOSITION\_{\Bk}FORCED}：只在画面里有不可省的文字（外语标题、铭牌）时置位。
  \item \Tp{AV\_{\Bk}DISPOSITION\_{\Bk}HEARING\_{\Bk}IMPAIRED}：听障辅助轨。
  \item \Tp{AV\_{\Bk}DISPOSITION\_{\Bk}VISUALLY\_{\Bk}IMPAIRED}：视障辅助轨，与上一条合称辅助技术旗。
  \item \Tp{AV\_{\Bk}DISPOSITION\_{\Bk}DEPENDENT}：这条轨必须与另一条同用才有意义。
  \item \Tp{AV\_{\Bk}DISPOSITION\_{\Bk}ATTACHED\_{\Bk}POK}：把封面图当附件挂进容器。
\end{itemize}
放置提示 \Tp{AV\_{\Bk}TRACK\_{\Bk}PLACING} 表达``贴顶还是贴底''这类留位需求。

语言与声道布局是另外两个正交维度。MP4、Matroska 与 TS 的语言码都是三字母（ISO 639-2 家族），
FFmpeg 侧通过流元数据键 \Tp{language} 写入。音频的声道配置用 \Tp{AVChannelLayout} 显式表达
（默认布局由 \Fn{av\_channel\_layout\_default} 生成，序列化后交给容器），不要从采样数反推布局。
视频轨的色彩解释靠 \Tp{ICC} 配置文件的 side data 附着；同一档码率的两个语言版本必须有相同的
色彩与尺寸描述，否则换音轨会连带触发呈现重配置。色域细节见第 \ref{ch:ch05-pixel-format} 章。

\begin{compare}[title={多语言字幕与同语言多呈现}]
\textbf{一语言一轨}：每条字幕一个轨，声明语言与用途，客户端按用户语言选一条。优点是清单短、
缓存友好、可单独加密或单独撤下；代价是同一时刻只渲染一条，混合需求（外语片段要字幕、旁白要
音频描述）要靠用途标志解决。\\
\textbf{同语言多呈现}：同一语言同时给一条标记式与一条图像式，或给一条全文轨加一条仅含强制段
的 \Tp{FORCED} 轨。优点是新旧终端各取所需、强制段能单独烧录；代价是同一份内容两个时间轴来源，
任何一次编辑都要同步两份，且两条字幕同时被选中会重叠显示。取舍点在终端能力差异有多大：最低端
只支持图像式时就值得付出双份维护；终端矩阵收敛到支持标记式时，多语言单谱系更省。
\end{compare}

\begin{guideline}[title={选轨逻辑的四条硬规则}]
\begin{itemize}
\item 缺省判定要覆盖``没有可用语言''：按用户偏好、与音频语言一致、\Tp{DEFAULT} 标志、
      轨道顺序四级回退，最后一级必须给出非空结果。
\item \Tp{FORCED} 轨在用户关闭字幕时仍应显示，但设置项要能区分``关全文''与``全关''。
\item 辅助技术轨（音频描述、手语、可朗读字幕）不要复用主轨：一条轨一个用途，
      混在一条里就无法在清单里声明。
\item 轨数不是越多越好：每条轨都要过清单、每条轨的密钥都要能取，轨数直接放大首屏清单
      体积与密钥请求次数。
\end{itemize}
\end{guideline}

\section{容器侧承载：盒、轨与时间基对齐}
\label{sec:ch15-containers}

MP4 系用样本入口区分字幕谱系。QuickTime 时代的 \Tp{text} 与 ISO 派生的 \Tp{tx3g} 承载
3GPP 计时文本，样本文本加样式表；\Tp{wvtt} 按通用计时文本规范把提示装进 \Tp{vttc}
（内含 \Tp{cue} 文本）与 \Tp{vtte}（空提示）；\Tp{ttml} 样本入口装整份标记文档或其片段。
共同麻烦是容器时基与提示时长分属两个体系，负偏移与编辑列表会把换算变成坑。

分片化 MP4 的承载由 \Tp{moof} 与 \Tp{mdat} 成对完成：\Tp{traf} 里的 \Tp{tfhd} 给默认标志与
默认时长，\Tp{tfdt} 给片段起始解码时戳，\Tp{trun} 逐样本给时长、标志、偏移与组成时戳差。
字幕轨通常每条提示一个样本，时长写进 \Tp{trun}，所以``提示时长''是容器字段而不是文本内容。
索引盒 \Tp{sidx} 与 \Tp{moof} 的关系是引用而非包含：它按字节或样本列出后续片段的偏移、时长与
最早呈现时戳，让客户端不下载媒体就能算出可寻址区间；它不影响片段的解码独立性，所以片段边界
仍要落在独立解码帧上，这条约束在第 \ref{sec:ch15-manifest} 节展开。

Matroska 把字幕类型直接写进编码标识：\Tp{S\_TEXT/ASS}、\Tp{S\_TEXT/UTF8}、\Tp{S\_HDMV/PGS}、
\Tp{S\_VOBSUB}、\Tp{S\_DVBSUB}，轨头带缺省、强制与辅助技术标志，字体可作为附件进容器
（对应 \Tp{AV\_DISPOSITION\_ATTACHED\_POK}）。TS 侧字幕是私有基本流，由字幕流描述符声明
组成语言、类型与三字母语言码，另有字幕描述符给出格式标识、语言与页号；这层信息在 FFmpeg
侧多半只体现在流表与元数据里，不做语义展开。加扰后的这些轨还要与密钥声明对齐，见第
\ref{sec:ch15-encryption} 节。

\begin{guideline}[title={封装字幕与图文轨时的四个开关}]
\begin{itemize}
\item 分片化输出用 \Opt{movflags} 的组合（按关键帧切片段、空 \Tp{moov} 起始、片段基地址
      按 \Tp{moof} 计），字幕轨同样逐段成形，便于清单侧统一寻址。
\item 每轨语言与标题在封装侧用逐流元数据写（键 \Tp{language}、\Tp{title}），
      不要指望从文件名推断。
\item 附件（字体、封面）走容器自带的附件机制并置 \Tp{ATTACHED\_POK} 标志。
\item 写入一律用交错排队接口，字幕落后时靠权重与统一时间轴处理，不手动补时戳；
      口径见第 \ref{ch:ch10-mux-sync} 章。
\end{itemize}
\end{guideline}

\begin{warning}[title={负偏移与 start\_time 的错配会整体推移字幕}]
\Tp{AVStream.start\_time} 是容器里这条轨第一个可见时戳换算后的值。字幕轨常见两种情况：
提示起始早于视频轨（片头前的对白），或容器为字幕写了负的组成时戳差。若音视频按混流规则
对齐、字幕却直接写原始时戳，就会出现``画面到 00:12、字幕还在 00:09''的恒定偏差，而且它不随
播放进度变化，最难被测试发现。做法是所有轨统一到同一条参考时间轴，写包前统一换算。
\end{warning}

\section{加密投递：通用加密、子样本与密钥生命周期}
\label{sec:ch15-encryption}

通用加密规范（ISO/IEC 23001-7）把``加密什么''与``怎么拿密钥''分开。方案标识有 \Tp{cenc}
（计数器模式加样本内抽样）与 \Tp{cbcs}（密码块链模式加同样抽样），共用同一组盒：\Tp{sinf}
里的 \Tp{schm} 声明方案与版本，轨道加密盒 \Tp{tenc} 给是否加密的默认值、初始化向量长度与
默认密钥标识；每个片段的实际参数在 \Tp{senc}（初始化向量加子样本长度表），样本级偏移由
\Tp{saio} 指、大小由 \Tp{saiz} 给。抽样比例叫 pattern，由 \Tp{crypt\_byte\_block} 与
\Tp{skip\_byte\_block} 两个 4 位字段按密码块交替加密与跳过，两者都为 0 表示整段加密。

子样本抽样是投递侧的关键设计：视频不需要全加密。以 \Tp{H.264}/{\Bk}\Tp{H.265} 为例，每个样本是
若干网络抽象层单元，只加密单元头若干字节与按块抽样出的一部分载荷，其余明文继续参与熵编码。
好处是密文仍可解析结构、开销可控、解密与解码边界清楚；代价是明文字节构成额外攻击面，所以
pattern 要按内容分级确定，不能把``只加密头''当成强保护。清单 \ref{lst:ch15-senc} 是读取骨架。

\begin{lstlisting}[style=cstyle,caption={\Tp{senc} 内逐样本初始化向量与子样本表的读取骨架（只解析，不解密）},label={lst:ch15-senc}]
struct Subsample { uint16_t clear_bytes; uint32_t enc_bytes; };  /* 明文长, 密文长 */
static int read_senc_entry(const uint8_t *p, size_t n, int has_iv, size_t iv_len,
                           uint8_t *iv, struct Subsample *ss, int max_ss) {
  size_t off = 0;
  if (has_iv && (iv_len == 0 || iv_len > n || iv_len > 16)) return -1; /* IV 长 8 或 16 */
  if (has_iv) { memcpy(iv, p, iv_len); off += iv_len; }
  uint16_t cnt = (uint16_t)((p[off] << 8) | p[off + 1]); off += 2;
  for (uint16_t i = 0; i < cnt && off + 5 <= n; i++) {
    ss[i].clear_bytes = (uint16_t)((p[off] << 8) | p[off + 1]); off += 2;
    ss[i].enc_bytes = ((uint32_t)p[off] << 16) | ((uint32_t)p[off + 1] << 8)
                      | (uint32_t)p[off + 2];                   off += 3;
  }
  return cnt;                             /* 0 表示整样本加密 */
}
\end{lstlisting}

清单侧的密钥声明是两套语言。HLS 用 \Tp{EXT-X-KEY}（方法有 \Tp{NONE}、\Tp{AES-128}、
\Tp{SAMPLE-AES} 及其通用加密变体）加密钥 URI、可选显式初始向量与密钥格式，主清单里预声明用
\Tp{EXT-X-SESSION-KEY}；\Tp{AES-128} 是整个分片一次对称加密，粒度粗、实现简单，轮换靠在新
分片处再出现一条声明。\Tp{SAMPLE-AES} 是样本级，按容器与媒体类型各有抽样规则。DASH 用
\Tp{ContentProtection}：通用加密靠方案值 \Tp{cenc} 或 \Tp{cbcs} 加属性里的默认密钥标识，
厂商方案用``\Tp{urn:uuid:} 加系统标识''的形制，引导数据放在 \Tp{cenc:pssh} 属性或引导信息盒
（\Tp{pssh}）里。轮换发生在周期边界或多密钥标识并行处，所以周期、适配集与表示三级都要能各自
携带密钥声明；几种粒度对照见表 \ref{tab:ch15-enc-schemes}。

设备侧的接口不论厂商，形制都是三件套：客户端把内容里的引导数据与设备能力组成一次挑战，
许可服务返回带密钥标识的响应，客户端交给会话，会话交出可用的密钥参数；之后密文与初始化向量
按 \Tp{senc} 交给硬件或经认证的媒体路径解密，区别只在引导数据格式与许可通道协议。开源面能拿
到的是结构与打包：\Lib{libdovi}、\Lib{Shaka Packager} 这类库处理清单与盒解析、加密参数计算。

\begin{longtable}{@{}>{\raggedright\arraybackslash\footnotesize}p{0.14\textwidth} >{\raggedright\arraybackslash\footnotesize}p{0.15\textwidth} >{\raggedright\arraybackslash\footnotesize}p{0.19\textwidth} >{\raggedright\arraybackslash\footnotesize}p{0.18\textwidth} >{\raggedright\arraybackslash\footnotesize}p{0.15\textwidth}@{}}
\caption{加密投递的几种粒度与声明位置}\label{tab:ch15-enc-schemes} \\
\toprule
形态 & 粒度 & 抽样方式 & 密钥标识来源 & 清单侧声明 \\
\midrule
\endfirsthead
\toprule
形态 & 粒度 & 抽样方式 & 密钥标识来源 & 清单侧声明 \\
\midrule
\endhead
\bottomrule
\endfoot
\endlastfoot
整片对称加密 & 一片一把 & 无抽样，整片密文 & 密钥 URI 拉取，可带显式初始向量 & \Tp{EXT{\Bk}-X{\Bk}-KEY} 方法 \Tp{AES-128} \\
通用加密 & 样本级或会话 & 子样本表，按块交替 & 引导信息盒、轨道加密盒或许可下发 & \Tp{Content{\Bk}Protection} 方案值加默认密钥标识 \\
\end{longtable}

\begin{note}[title={本节的边界}]
本章只讲密文的结构、字段与投递语义：为什么这样抽样、密钥标识如何在盒与清单之间对应、轮换点
在哪里。任何规避内容保护、提取受保护密钥或绕过会话的内容都不在范围内，也不提供可运行的解密
路径。命令行侧为已知密钥的自有内容提供一把裸密钥入口（\Opt{decryption\_key}），用途限于自测与
归档；硬件解码与内容保护缓冲的实现见第 \ref{ch:ch08-decode-hw} 章。
\end{note}

\section{清单、表示与切换点：HLS 与 DASH 的对偶}
\label{sec:ch15-manifest}

同一套内容在两种清单里是同构的：HLS 主清单的变体对应 DASH 的表示，HLS 的媒体分组对应
\Tp{AdaptationSet}，时段与事件周期对应 \Tp{Period}。区别在表达重心：HLS 是行式文本清单，
每个标签既是数据也是指令；DASH 是文档树，属性可继承，扩展靠统一资源名。

\begin{lstlisting}[style=shellstyle,caption={主清单节选：档位、rendition、加密声明与事件区间},label={lst:ch15-master-m3u8}]
#EXTM3U
#EXT-X-MEDIA:TYPE=SUBTITLES,GROUP-ID="subs",LANGUAGE="zh",FORCED=NO,URI="zh.m3u8"
#EXT-X-MEDIA:TYPE=CLOSED-CAPTIONS,GROUP-ID="cc",INSTREAM-ID="SERVICE1"
#EXT-X-KEY:METHOD=AES-128,URI="k/1",IV=0x000102030405060708090a0b0c0d0e0f
#EXT-X-STREAM-INF:BANDWIDTH=4800000,CODECS="avc1.640028",SUBTITLES="subs"
720p/playlist.m3u8
#EXT-X-DATERANGE:ID="cue-1",START-DATE="2026-04-01T08:00:00.000Z",DURATION=15
\end{lstlisting}

清单 \ref{lst:ch15-master-m3u8} 里的 \Tp{CODECS} 参数串最容易写错：它是媒体编码参数列表，
每条由配置标识加级别与可选的轮廓、兼容性约束组成，多条用逗号分隔；声明的参数必须与分片实际
内容一致——写宽了让客户端承诺它做不到的解码，写窄了让本可播放的设备滤掉档位。字幕与辅助音轨
靠 \Tp{EXT-X-MEDIA}（类型、分组、语言、名称、缺省与自动选择、强制标志、可访问性特征串）挂到
\Tp{EXT-X-STREAM-INF} 的分组属性上，闭环字幕用同一标签但带通道内标识；初始化段由
\Tp{EXT-X-MAP} 指定，只含独立解码帧分片的档位由 \Tp{EXT-X-I-FRAMES-ONLY} 声明。

DASH 侧的对偶写法是：\Tp{SegmentTemplate} 用初始化与媒体两条模板串（带 \Tp{\$Number\$}、
\Tp{\$Time\$}、\Tp{\$RepresentationID\$} 这类替换项），\Tp{SegmentTimeline} 用 \Tp{S} 项写
时戳、时长与重复次数；\Tp{Role} 用统一资源名加取值（主、配音、字幕、强制字幕、描述）声明用途，
密钥用 \Tp{ContentProtection}，事件用 \Tp{EventStream}。切换点约束两边一致：同一适配集内的
表示必须分片边界对齐、独立解码帧落在同一呈现时戳上，否则切换点撕裂或出黑场。CMAF 把这条约束
标准化为分片起点即独立解码帧、初始化段与媒体段分离。对照见表 \ref{tab:ch15-hls-dash}。

\begin{longtable}{@{}>{\raggedright\arraybackslash\footnotesize}p{0.17\textwidth} >{\raggedright\arraybackslash\footnotesize}p{0.30\textwidth} >{\raggedright\arraybackslash\footnotesize}p{0.31\textwidth}@{}}
\caption{清单要素的 HLS 与 DASH 对偶}\label{tab:ch15-hls-dash} \\
\toprule
要素 & HLS & DASH \\
\midrule
\endfirsthead
\toprule
要素 & HLS & DASH \\
\midrule
\endhead
\bottomrule
\endfoot
\endlastfoot
一组可选码率 & \Tp{EXT-X-STREAM-INF} 及其分组属性 & \Tp{AdaptationSet} 下的 \Tp{Representation} \\
轨用途与语言 & \Tp{EXT-X-MEDIA} 的类型、语言与特征属性 & \Tp{Role}、\Tp{Language} 与 \Tp{Label} \\
事件与广告边界 & \Tp{EXT-X-DATERANGE}（区间加自定义前缀） & \Tp{EventStream} 与周期边界 \\
\end{longtable}

时钟是这一层的隐性要求。\Tp{EXT-X-PROGRAM-DATE-TIME} 把某个分片起点钉到世界时，客户端据此把
呈现时戳换成墙上时钟，广告替换、跨设备续播与多路同步都依赖这条映射；DASH 侧对应的量是可获取
起点加片段可用时间线。两边共同的硬要求是生成端与服务器时钟误差小于一个分片时长，否则表现为
``已列出的分片返回 404''或``直播边界反复回退''。绝对时刻换算见第 \ref{ch:ch02-clock-buffer} 章。

FFmpeg 侧要清楚解封装器的职责边界：\Tp{hls}（历史名 \Tp{applehttp}）读清单，把候选流与其
带宽、分辨率映射成程序级对象；\Tp{dash} 读文档树后交出可解码的流。字幕与音频 rendition、事件
区间、密钥声明这些清单语义并不都会展开成流，要可靠拿到全部轨就得自己解析清单。引擎侧的封装与
选流见第 \ref{ch:ch13-framework-integration} 章，术语对照见附录 \ref{ap:apA-glossary}。

\begin{bestpractice}[title={投递层落地的五条检查}]
\begin{itemize}
\item 字幕轨的时间轴只做一次容器时基换算并断言单调，渲染侧只用毫秒，禁止再乘时基。
\item 图文与遥测一律带``最后已知时刻加可信标志''，不要用显示时长当有效期。
\item 每条轨的用途标志、语言、声道布局与色彩解释都要能独立读出，任一缺失即视为清单缺陷。
\item 加密声明三处对齐：引导信息盒、样本加密参数、清单里的密钥标识指向同一把密钥；
      轮换点写在周期或分片边界上，不在片段中间换。
\item 切换点用独立解码帧与分片边界对齐保证，档位参数串与实际内容自动比对，不一致就阻断
      发布而不是等客户端报错。
\end{itemize}
\end{bestpractice}

% ============================================================
% ---- frag_apA.tex ----
\appendix
\chapter{术语与缩写}
\label{ap:apA-glossary}

本表收录正文用到的术语与缩写，按四条主线分组：音频与设备、视频与像素、
编解码与码控、封装与流媒体。每条给出一句可以独立读懂的定义；``首见'' 列
指向该词获得详细解释的章，正文在其他章复用时以该章定义为准。

四点读表约定。其一，术语列优先给中文通行译名，括注或缩写列给原文短形；
没有公认译名的（如 UVC、v210）直接沿用原文。其二，一个缩写在不同社区可能
含义不同（``plane'' 在显示栈指合成层、在像素格式指分平面），表中注明本章
语境，跨组查阅时注意区分。其三，首见列只标概念章，部件级函数与结构体的
完整签名集中在附录 \ref{ap:apB-api-quickref}，选型对照在附录
\ref{ap:apC-selection}。
其四，列的约定：缩写列的 ``—'' 表示该术语没有通行短形；同一行只收入
定义信息量最大的一处，另一处的讨论顺着首见列回正文查找即可。

\section{音频与设备}

音频链路的概念大多绕不开 ``时钟'' 与 ``缓冲'' 两件事：设备侧以帧为单位
推进，应用侧以周期为单位追赶，两侧时钟永不严格同步（第
\ref{ch:ch02-clock-buffer} 章与第 \ref{ch:ch03-audio-device} 章）。
本组词条因此在 ``量的表示''（采样、位深、电平）与 ``节奏的控制''
（缓冲、周期、漂移）两个层面展开。

\begin{longtable}{>{\raggedright\arraybackslash}p{0.20\textwidth} >{\raggedright\arraybackslash}p{0.12\textwidth} >{\raggedright\arraybackslash}p{0.36\textwidth} >{\raggedright\arraybackslash}p{0.12\textwidth}}
\caption{音频与设备术语}\label{tab:apa-audio} \\
\toprule
术语 & 缩写 & 定义 & 首见 \\
\midrule
\endfirsthead
\toprule
术语 & 缩写 & 定义 & 首见 \\
\midrule
\endhead
\bottomrule
\endfoot
脉冲编码调制 & PCM & 以固定周期采样并对幅度量化成整数的表示法，设备接口的通用形态 & \ref{ch:ch02-clock-buffer} \\
采样点 & — & 量化后的单个幅度数值，一个声道一个 & \ref{ch:ch02-clock-buffer} \\
采样率 & fs & 每秒产生的采样点个数；影视工业以 48 kHz 为主 & \ref{ch:ch02-clock-buffer} \\
样本点 & sample & 单声道视角下的一个采样值，帧按声道数聚合它 & \ref{ch:ch02-clock-buffer} \\
音频帧 & — & 同一时刻全部声道的采样点集合，设备读写的最小单位 & \ref{ch:ch02-clock-buffer} \\
位深 & — & 一个采样点占用的量化位数，常见 16 与 24 位 & \ref{ch:ch02-clock-buffer} \\
声道布局 & — & 声道数量与次序的成对表示（左、右、中置等） & \ref{ch:ch02-clock-buffer} \\
环形缓冲 & — & 写指针与读指针首尾相接构成的循环队列，设备缓冲的标准结构 & \ref{ch:ch02-clock-buffer} \\
缓冲水位 & — & 环形缓冲内当前可读或可写的样本点差额 & \ref{ch:ch02-clock-buffer} \\
缓冲延迟 & — & 设备侧为容纳抖动而驻留数据造成的时延，以周期数折算 & \ref{ch:ch02-clock-buffer} \\
往返延迟 & RTL & 从写入设备到能听到该样本的总回路时长 & \ref{ch:ch02-clock-buffer} \\
时钟漂移 & — & 收发两端晶振的频率差，导致需要周期性补点或重采样修正 & \ref{ch:ch02-clock-buffer} \\
重定时 & resync & 按主时钟滑移对从侧缓冲做拉伸或丢弃的校正 & \ref{ch:ch02-clock-buffer} \\
dBFS & — & 以数字满量程为 0 分贝的电平标尺，常态信号在其下方 & \ref{ch:ch02-clock-buffer} \\
削波 & — & 采样值越过量化器上下界被削平产生的失真 & \ref{ch:ch02-clock-buffer} \\
混音 & — & 多路信号按给定电平叠加成一路，由软件或硬件完成 & \ref{ch:ch03-audio-device} \\
重采样 & — & 把一种采样率换算成另一种；算法见滤镜章 & \ref{ch:ch09-filter-resample} \\
采样格式 & S16LE & 端序与定点或浮点的存储编码名，如 16 位小端 & \ref{ch:ch03-audio-device} \\
ALSA & — & Linux 内核音频栈及其用户态库 & \ref{ch:ch03-audio-device} \\
周期 & period & ALSA 的推进与中断单位，环形缓冲按周期划分 & \ref{ch:ch03-audio-device} \\
xrun & — & ALSA 对欠载或过载的统称，消费或生产失去节奏的症状 & \ref{ch:ch03-audio-device} \\
HW params & — & ALSA 协商采样率、格式与缓冲参数的入口对象 & \ref{ch:ch03-audio-device} \\
直接内存访问 & DMA & 外设不经 CPU 逐字节搬运、直接读写内存的机制 & \ref{ch:ch03-audio-device} \\
WASAPI & — & Windows 从 Vista 起的音频会话接口 & \ref{ch:ch03-audio-device} \\
独占模式 & — & WASAPI 下绕过系统混音、应用直控设备格式的模式 & \ref{ch:ch03-audio-device} \\
事件驱动 & — & WASAPI 用事件句柄唤醒读写而非定时轮询的方式 & \ref{ch:ch03-audio-device} \\
AudioUnit & AU & CoreAudio 在 Apple 平台的处理单元与回调模型 & \ref{ch:ch03-audio-device} \\
音频会话 & AS & iOS 侧管理路由、中断与类别的会话对象 & \ref{ch:ch03-audio-device} \\
JACK & — & 专业音频的路由服务及其低延迟客户端协议 & \ref{ch:ch03-audio-device} \\
回声消除 & AEC & 从采集信号中减去已知的远端播放参考信号 & \ref{ch:ch03-audio-device} \\
热插拔 & — & 设备增删事件及其引发的管线重建流程 & \ref{ch:ch03-audio-device} \\
回环 & loopback & 系统把播放流原样送回采集侧供录制 & \ref{ch:ch03-audio-device} \\
路由变更 & — & 耳机插拔等输入输出端点变化通知 & \ref{ch:ch03-audio-device} \\
\end{longtable}

\section{视频与像素}

视频侧的术语分属两条边界：采集与显示的路径词汇（第
\ref{ch:ch04-capture-display} 章），与帧内容本身的表示词汇（第
\ref{ch:ch05-pixel-format} 章）。同一个词在两侧可能各司其职，
``plane'' 即是一例：显示栈里是合成层，像素格式里是分平面。

\begin{longtable}{>{\raggedright\arraybackslash}p{0.20\textwidth} >{\raggedright\arraybackslash}p{0.12\textwidth} >{\raggedright\arraybackslash}p{0.36\textwidth} >{\raggedright\arraybackslash}p{0.12\textwidth}}
\caption{视频与像素术语}\label{tab:apa-video} \\
\toprule
术语 & 缩写 & 定义 & 首见 \\
\midrule
\endfirsthead
\toprule
术语 & 缩写 & 定义 & 首见 \\
\midrule
\endhead
\bottomrule
\endfoot
Video4Linux2 & V4L2 & Linux 视频采集与输出设备的用户态接口 & \ref{ch:ch04-capture-display} \\
USB Video Class & UVC & USB 摄像头的传输协议与格式协商规范 & \ref{ch:ch04-capture-display} \\
零拷贝 & — & 缓冲在内核、GPU 与编解码器间共享而不重复搬运 & \ref{ch:ch04-capture-display} \\
Media Foundation & MF & Windows 的媒体框架，相机枚举与采集图建立在其上 & \ref{ch:ch04-capture-display} \\
DMA 缓冲 & DMABUF & 内核缓冲共享机制，fd 可跨驱动零拷贝传递 & \ref{ch:ch04-capture-display} \\
直接渲染管理 & DRM & Linux 内核显示栈，含模式设置与平面合成 & \ref{ch:ch04-capture-display} \\
平面（显示） & plane & KMS 中可独立扫描合成的图层，分主显、叠加、光标 & \ref{ch:ch04-capture-display} \\
垂直同步 & vsync & 扫描回到顶端的时刻，帧提交与亮屏的分界 & \ref{ch:ch04-capture-display} \\
消隐期 & vblank & 两帧之间的回扫窗口，翻转事件在此投递 & \ref{ch:ch04-capture-display} \\
原子提交 & — & KMS 一次性校验并生效全部显示状态的接口 & \ref{ch:ch04-capture-display} \\
上屏反馈 & — & Wayland presentation 扩展回报帧真正亮起的时刻 & \ref{ch:ch04-capture-display} \\
OML 同步 & — & X11/GLX 扩展，给出 ust、msc、sbc 三块同步表 & \ref{ch:ch04-capture-display} \\
串行数字接口 & SDI & 广电点对点串行视频电缆及其信号标准 & \ref{ch:ch04-capture-display} \\
帧同步锁定 & genlock & 多台设备借参考信号锁定在同一出帧相位 & \ref{ch:ch04-capture-display} \\
时间码 & TC & 与真实时刻绑定的 SMPTE 帧位置标签 & \ref{ch:ch04-capture-display} \\
四字符码 & fourcc & 以 4 个 ASCII 字符为格式标记的编码习惯 & \ref{ch:ch05-pixel-format} \\
像素格式 & — & 采样、布局与位深约定的合称，管线接口的握手项 & \ref{ch:ch05-pixel-format} \\
亮度色差表示 & YCbCr & 一路亮度加两路色差分量的传输域表示 & \ref{ch:ch05-pixel-format} \\
色差信号 & Cb/Cr & 蓝、红与亮度之差按比例缩放的两个色度分量 & \ref{ch:ch05-pixel-format} \\
高动态范围 & HDR & 用 PQ 或 HLG 曲线压缩更宽亮度范围的传递特性 & \ref{ch:ch05-pixel-format} \\
亮度信号 & luma & 对已做传递函数编码的信号加权之和，非物理亮度 & \ref{ch:ch05-pixel-format} \\
色度次采样 & 4:2:0 & 色道采样密度低于亮度道的比例约定 & \ref{ch:ch05-pixel-format} \\
采样点位置 & siting & 色度样点相对亮度样点钉在居中还是同位 & \ref{ch:ch05-pixel-format} \\
限定范围 & — & 留出上下边界的量化区间，8-bit 亮度 16--235 & \ref{ch:ch05-pixel-format} \\
全范围 & — & 占满 0--255 的量化区间，桌面内容默认 & \ref{ch:ch05-pixel-format} \\
平面格式 & planar & 各分量各占独立数组的布局 & \ref{ch:ch05-pixel-format} \\
半平面格式 & NV12 & 亮度平面加一个两色度交错的平面 & \ref{ch:ch05-pixel-format} \\
打包格式 & UYVY & 各分量沿扫描线按像素交错的单平面布局 & \ref{ch:ch05-pixel-format} \\
行步长 & stride & 平面中一行的实际字节数，常大于可见宽度 & \ref{ch:ch05-pixel-format} \\
对齐 & — & 硬件与向量化对行步长取整的要求，错位则画面倾斜 & \ref{ch:ch05-pixel-format} \\
v210 & — & 广播中间格式：4 个 10-bit 分量打包进 3 个 32-bit 字 & \ref{ch:ch05-pixel-format} \\
\end{longtable}

\section{编解码与码控}

压缩侧的词汇围绕三件事：帧间与帧内的参考结构（第
\ref{ch:ch06-codec-api} 部），比特如何在质量与额度之间分配（第
\ref{ch:ch07-rate-control} 章），以及编码器与解码器作为组件的生命周期
（第 \ref{ch:ch08-decode-hw} 章）。三层词汇常在同一串调用里交织出现，
首见列把每个词拆回它的主场。

\begin{longtable}{>{\raggedright\arraybackslash}p{0.20\textwidth} >{\raggedright\arraybackslash}p{0.12\textwidth} >{\raggedright\arraybackslash}p{0.36\textwidth} >{\raggedright\arraybackslash}p{0.12\textwidth}}
\caption{编解码与码控术语}\label{tab:apa-codec} \\
\toprule
术语 & 缩写 & 定义 & 首见 \\
\midrule
\endfirsthead
\toprule
术语 & 缩写 & 定义 & 首见 \\
\midrule
\endhead
\bottomrule
\endfoot
帧内编码 & I & 只自参考的帧，体积大但不依赖邻居 & \ref{ch:ch06-codec-api} \\
前向预测帧 & P & 参考更早重建帧的帧 & \ref{ch:ch06-codec-api} \\
双向预测帧 & B & 还参考更晚的帧，带来显示序与解码序的错位 & \ref{ch:ch06-codec-api} \\
运动估计 & — & 在参考帧内搜索最匹配块以压缩残差 & \ref{ch:ch06-codec-api} \\
转码 & — & 解码再编码的全流程，区别于只换容器的再封装 & \ref{ch:ch06-codec-api} \\
即时解码刷新 & IDR & 清空全部参考链的最强帧内位置，随机访问的分界 & \ref{ch:ch06-codec-api} \\
图像组 & GOP & 一个随机访问点到下一个之间的帧序列 & \ref{ch:ch06-codec-api} \\
档位 & profile & 编码特性子集的承诺，决定解码器兼容范围 & \ref{ch:ch06-codec-api} \\
级别 & level & 分辨率、帧率与码率等资源上限的承诺 & \ref{ch:ch06-codec-api} \\
NAL 单元 & NAL & H.26x 系码流中带类型头的语法单元 & \ref{ch:ch06-codec-api} \\
参数集 & SPS/PPS & H.264 序列级与图像级的解码初始化数据 & \ref{ch:ch06-codec-api} \\
外部数据 & — & 打开解码器所需的参数集等初始比特流 & \ref{ch:ch06-codec-api} \\
环路滤波 & — & 作用于重建参考帧内部的去块滤波等处理 & \ref{ch:ch06-codec-api} \\
变换 & — & 把残差转到频域集中能量的步骤，如 DCT & \ref{ch:ch06-codec-api} \\
熵编码 & — & 符号级无损压缩环节，如 CABAC & \ref{ch:ch06-codec-api} \\
量化参数 & QP & 控制变换系数取舍粒度的旋钮，越大越粗 & \ref{ch:ch07-rate-control} \\
码率控制 & — & 在质量与比特额度之间做帧级与块级分配 & \ref{ch:ch07-rate-control} \\
码率验证器 & VBV & 解码端缓冲充实度模型，CBR 约束的来源 & \ref{ch:ch07-rate-control} \\
码率 & — & 每秒消耗的比特数，码控的直接额度 & \ref{ch:ch07-rate-control} \\
恒定质量 & CRF & 按固定量化质量压缩、允许码率浮动的模式 & \ref{ch:ch07-rate-control} \\
前瞻 & lookahead & 编码前分析未来若干帧的统计再决定分配 & \ref{ch:ch07-rate-control} \\
两遍编码 & — & 第一遍产复杂度统计，第二遍按统计分配 & \ref{ch:ch07-rate-control} \\
场景切换 & — & 帧间统计突变处强制插入帧内编码 & \ref{ch:ch07-rate-control} \\
延迟档 & — & 编码器低延迟设置与压缩效率的折中选项 & \ref{ch:ch07-rate-control} \\
帧长 & — & 一帧音频覆盖的毫秒数，Opus 可变 2.5--120 & \ref{ch:ch06-codec-api} \\
不连续传输 & DTX & 静音段改发舒适噪声包的压缩 & \ref{ch:ch06-codec-api} \\
前向纠错 & FEC & 在码流内携带冗余以硬抗丢包 & \ref{ch:ch06-codec-api} \\
语音激活检测 & VAD & 判定当前是否有人声，供 DTX 等使用 & \ref{ch:ch06-codec-api} \\
冲刷 & flush & 以空输入榨干编码器或解码器的内部驻留帧 & \ref{ch:ch08-decode-hw} \\
硬件加速 & — & 由 ASIC 或 GPU 固定功能单元承担编解码 & \ref{ch:ch08-decode-hw} \\
\end{longtable}

\section{封装与流媒体}

封装与传输侧的词汇分三层：时间标签与容器（第
\ref{ch:ch10-mux-sync} 章）、协议与组件装配（第
\ref{ch:ch11-streaming} 部、第 \ref{ch:ch13-framework-integration}
章）、以及实时链路的缓冲与带宽应对（第 \ref{ch:ch12-realtime-thread}
章与第 \ref{ch:ch14-congestion-abr} 章）。跨层词（分片、同步点）先在
容器层理解其定义，再到协议层结合承载方式读第二遍。

\begin{longtable}{>{\raggedright\arraybackslash}p{0.20\textwidth} >{\raggedright\arraybackslash}p{0.12\textwidth} >{\raggedright\arraybackslash}p{0.36\textwidth} >{\raggedright\arraybackslash}p{0.12\textwidth}}
\caption{封装与流媒体术语}\label{tab:apa-stream} \\
\toprule
术语 & 缩写 & 定义 & 首见 \\
\midrule
\endfirsthead
\toprule
术语 & 缩写 & 定义 & 首见 \\
\midrule
\endhead
\bottomrule
\endfoot
封装与解封装 & mux & 把多路比特流与其时间标签按版式合写与拆读 & \ref{ch:ch10-mux-sync} \\
容器格式 & — & 复用后比特流在文件或流中的版式规范 & \ref{ch:ch10-mux-sync} \\
时间基 & — & 时间戳每格代表的秒数，逐流各有一份 & \ref{ch:ch10-mux-sync} \\
显示时间戳 & PTS & 帧应当上屏或出声的时刻 & \ref{ch:ch10-mux-sync} \\
解码时间戳 & DTS & 帧进入解码器的时刻，含 B 帧流早于 PTS & \ref{ch:ch10-mux-sync} \\
MP4 盒 & box & ISO 基础媒体格式的嵌套块结构 & \ref{ch:ch10-mux-sync} \\
影片头盒 & moov & MP4 的轨道与索引信息所在盒 & \ref{ch:ch10-mux-sync} \\
快启 & faststart & 把索引挪到文件开头以便边下边播 & \ref{ch:ch10-mux-sync} \\
同步样本表 & stss & 轨道随机访问点的索引表 & \ref{ch:ch10-mux-sync} \\
传送流 & TS & 188 字节定长包结构的广播容器 & \ref{ch:ch10-mux-sync} \\
再封装 & remux & 不重新编码、只更换容器或切点的操作 & \ref{ch:ch10-mux-sync} \\
音画同步 & — & 按时间标签把两流对齐到同一播放轴 & \ref{ch:ch10-mux-sync} \\
分片 MP4 & fMP4 & 初始化段加媒体小段拼成的 MP4 & \ref{ch:ch11-streaming} \\
通用媒体应用格式 & CMAF & HLS 与 DASH 可共用的标准化分片规范 & \ref{ch:ch11-streaming} \\
实时传输协议 & RTP & UDP 上逐帧封装、带序号与时间戳的协议 & \ref{ch:ch11-streaming} \\
RTP 控制 & RTCP & RTP 的伴随反馈通道：统计、重传与关键帧请求 & \ref{ch:ch11-streaming} \\
安全 RTP & SRTP & RTP 载荷的加密变换套件 & \ref{ch:ch11-streaming} \\
SRT 协议 & SRT & 面向 UDP、可配延迟预算的可靠重传直播协议 & \ref{ch:ch11-streaming} \\
分片 & — & HTTP 自适应流按时间切出的小文件单位 & \ref{ch:ch11-streaming} \\
清单文件 & — & 枚举各档码率与分片地址的索引文件 & \ref{ch:ch14-congestion-abr} \\
包时长 & ptime & 每个 RTP 包承载的音频毫秒数 & \ref{ch:ch11-streaming} \\
关键帧请求 & PLI/FIR & RTCP 反馈消息，要求发送端立刻产出一帧帧内编码 & \ref{ch:ch11-streaming} \\
抖动缓冲 & — & 接收端吸收网络忽快忽慢的自适应缓冲 & \ref{ch:ch12-realtime-thread} \\
网络抖动 & — & 分组到达间隔相对发送间隔的波动 & \ref{ch:ch12-realtime-thread} \\
发送整形 & pacing & 把成簇帧摊匀到时间轴上再发出的限速 & \ref{ch:ch12-realtime-thread} \\
端到端延迟 & E2E & 从采集曝光到远端亮屏的总时长口径 & \ref{ch:ch12-realtime-thread} \\
拥塞控制 & GCC & 基于时延梯度估带宽并调发送速率 & \ref{ch:ch14-congestion-abr} \\
码率阶梯 & — & 供自适应切换挑选的多档编码集合 & \ref{ch:ch14-congestion-abr} \\
自适应码率 & ABR & 客户端按可用带宽在阶梯间选档与切换 & \ref{ch:ch14-congestion-abr} \\
\end{longtable}

查询任何术语，先在本附录定位分组与首见章，再回正文该章找带上下文的
解释与代码。若某词不在表中，多半是正文的一次性用法，回目录按首见章定位；
若反复出现仍语义含糊，则应视为勘误候选补进本表。清单 \ref{lst:apa-tags}
给出核对 ``标签类'' 术语（时间戳、色彩三元组、元数据）最常用的读取命令。

\begin{lstlisting}[style=shellstyle,caption={读取文件的时间戳与元数据标签，与术语表逐条对照},label={lst:apa-tags}]
# 打印封装层全部标签与逐流时间基
ffprobe -hide_banner -show_entries format_tags:stream=time_base \
       -of ini sample.mkv

# 看一个 TS 的 PCR 间隔与比特率：实时流侧的健康指标
ffprobe -hide_banner -show_packets -select_streams v \
        -print_format csv sample.ts | head
\end{lstlisting}

\begin{guideline}[title={本册用词约定}]
正文中 \DomAud{}、\DomVid{} 标记音频与视频各自的链路侧，\DomRT{} 与
\DomBatch{} 标记实时与离线批处理场景，四个标签跨章同义、不混用。函数与
方法名一律 \Fn{}，类型与结构名 \Tp{}，命令行开关 \Opt{}，库名 \Lib{}。
中文译名优先采用社区通行写法；首次出现给原文括注，此后直接复用本表缩写。
同一概念存在两种通行叫法时（如上屏反馈的 presentation、行步长的 stride
与 linesize），正文两者并陈并以本表词条为主名。
\end{guideline}

% ============================================================
% ---- frag_apB.tex ----
\chapter{API 与参数速查}
\label{ap:apB-api-quickref}

本附录是随手翻的表，不是教程：每节一张，表头统一为名称、语义、取值或默认、注意四栏。
列出的都是各库在主线版本里长期稳定的调用与开关，用于把正文的讨论落到具体名字上。
凡标注``随版本''的条目，请在动手前核对该库的头文件与手册页。

\section{设备与时钟侧调用}

\begin{longtable}{>{\raggedright\arraybackslash}p{0.17\textwidth}
                  >{\raggedright\arraybackslash}p{0.22\textwidth}
                  >{\raggedright\arraybackslash}p{0.19\textwidth}
                  >{\raggedright\arraybackslash}p{0.20\textwidth}}
\caption{设备与采集：ALSA、WASAPI、CoreAudio、V4L2 关键调用}\label{tab:apB-device} \\
\toprule
名称 & 语义 & 取值或默认 & 注意 \\
\midrule
\endfirsthead
\toprule
名称 & 语义 & 取值或默认 & 注意 \\
\midrule
\endhead
\midrule
\endfoot
\bottomrule
\endlastfoot
\Fn{snd\_{\Bk}pcm\_{\Bk}open}（ALSA） & 打开播放或采集流，拿到 \Tp{snd\_{\Bk}pcm\_{\Bk}t} & \Tp{hw:0,0}、\Tp{plughw:0,0}、\Tp{default}；\Tp{SND\_{\Bk}PCM\_{\Bk}STREAM\_{\Bk}PLAYBA{\Bk}CK} 与 \Tp{SND\_{\Bk}PCM\_{\Bk}STREAM\_{\Bk}CAPTURE}；非阻塞用 \Tp{SND\_{\Bk}PCM\_{\Bk}NONBLO{\Bk}CK} & 只有 \Tp{plug} 与 \Tp{dmix} 层才做格式与速率转换，\Tp{hw} 直接反映硬件能力 \\
\Fn{snd\_{\Bk}pcm\_{\Bk}hw\_{\Bk}params\_{\Bk}set\_{\Bk}rate\_{\Bk}near} 一族 & 在硬件参数窗内协商格式、速率、通道与周期 & \Tp{SND\_{\Bk}PCM\_{\Bk}FORMAT\_{\Bk}S16\_{\Bk}LE}、\Tp{SND\_{\Bk}PCM\_{\Bk}ACCESS\_{\Bk}RW\_{\Bk}INTERL{\Bk}EAVED}；周期常取 5 到 10 毫秒 & 必须先 \Fn{snd\_{\Bk}pcm\_{\Bk}hw\_{\Bk}params\_{\Bk}any} 起窗；``near'' 就近取值而不报错，要回读确认 \\
\Fn{snd\_{\Bk}pcm\_{\Bk}sw\_{\Bk}params} & 设置启动阈值、可用阈值与时间戳模式 & \Tp{start\_{\Bk}thresh{\Bk}old}、\Tp{avail\_{\Bk}min}、\Tp{tstamp\_{\Bk}mode} & 低延迟配置全靠这一层，默认值偏保守 \\
\Fn{snd\_{\Bk}pcm\_{\Bk}readi}、\Fn{snd\_{\Bk}pcm\_{\Bk}writei} & 交织方式读写若干帧 & 返回实际帧数，可能小于请求数 & 短读短写是正常路径，必须循环；\Tp{-{\Bk}EPIPE} 表示 underrun 或 overrun \\
\Fn{snd\_{\Bk}pcm\_{\Bk}avail\_{\Bk}update}、\Fn{snd\_{\Bk}pcm\_{\Bk}delay} & 刷新可用量；取队列内未播帧数 & 单位都是帧 & \Fn{avail\_{\Bk}update} 不是时间源；端到端延迟要再叠设备延迟 \\
\Fn{snd\_{\Bk}pcm\_{\Bk}htimes{\Bk}tamp} & 取帧位置与对应时钟 & \Tp{snd\_{\Bk}timest{\Bk}amp}，精度取决于 \Tp{tstamp\_{\Bk}mode} & 时钟源可能是单调钟或实时钟，混用会让第 2 章的平滑器发散 \\
\Fn{snd\_{\Bk}pcm\_{\Bk}recover}、\Fn{snd\_{\Bk}pcm\_{\Bk}resume} & 一次处理 \Tp{-{\Bk}EPIPE} 与 \Tp{-{\Bk}ESTRPI{\Bk}PE} & 第二参数非零则静默 & 电源暂停要轮询 \Fn{resume} 直到成功或明确失败，不能只 prepare \\
\Fn{snd\_{\Bk}pcm\_{\Bk}drain}、\Fn{snd\_{\Bk}pcm\_{\Bk}drop} & 排空或立即丢弃缓冲 & 无参数 & 切设备与正常关闭用 drain，急停与录音截断用 drop \\
\Fn{IAudioCl{\Bk}ient\_{\Bk}Initia{\Bk}lize}（WASAPI） & 建立共享或独占缓冲 & \Tp{AUDCLNT\_{\Bk}SHAREM{\Bk}ODE\_{\Bk}SHARED} 与 \Tp{AUDCLNT\_{\Bk}SHAREM{\Bk}ODE\_{\Bk}EXCLUS{\Bk}IVE}；\Tp{hnsBuffe{\Bk}rDuratio{\Bk}n} 以 100 纳秒计；\Tp{AUDCLNT\_{\Bk}STREAM{\Bk}FLAGS\_{\Bk}EVENTC{\Bk}ALLBACK} & 事件回调要先 \Fn{SetEvent{\Bk}Handle}；独占模式的格式必须与设备当前格式一致 \\
\Fn{IAudioCl{\Bk}ient\_{\Bk}GetMix{\Bk}Format}、\Fn{GetDevic{\Bk}ePeriod} & 取混音格式与宿主周期 & 混音常为 float32；默认周期约 10 毫秒，最小约 3 毫秒 & 传入更小的缓冲时长不报错，会被抬到最小周期 \\
\Fn{IAudioRe{\Bk}nderClie{\Bk}nt\_{\Bk}GetBuf{\Bk}fer} 与 \Fn{ReleaseB{\Bk}uffer} & 取写入指针并提交上去 & 可写帧数由 \Fn{GetCurre{\Bk}ntPaddin{\Bk}g} 反推 & 两个调用必须成对，漏掉提交时缓冲不动 \\
\Fn{IAudioCa{\Bk}ptureCli{\Bk}ent\_{\Bk}GetNex{\Bk}tPacketS{\Bk}ize} 与 \Fn{GetBuffer} & 按数据包读采集 & \Tp{S\_{\Bk}FALSE} 表示流已停 & 包边界不保证对齐你的回调大小，需自行拼接 \\
\Fn{IAudioCl{\Bk}ient\_{\Bk}Reset} 与设备失效 & 出错后回到可再用状态 & \Tp{AUDCLNT\_{\Bk}E\_{\Bk}DEVICE\_{\Bk}INVALI{\Bk}DATED} & 该错误要重建整套对象并回退默认设备（第 3 章的路径） \\
\Fn{AudioDev{\Bk}iceCreat{\Bk}eIOProcI{\Bk}D}、\Fn{AudioDev{\Bk}iceStart}（CoreAudio） & 注册回调并启动设备 & 回调形参含帧数、输入与输出的 \Tp{AudioBuf{\Bk}ferList}、帧时间 & IOProc 跑在实时线程：禁止分配、加锁、系统调用与 OC 对象 \\
\Fn{AudioStr{\Bk}eamBasic{\Bk}Descript{\Bk}ion} & 描述一种音频格式 & \Tp{mFormatID} 取 \Tp{kAudioFo{\Bk}rmatLine{\Bk}arPCM}，\Tp{mFramesP{\Bk}erPacket} 常为 1 & 各字段必须自洽（每帧字节等于位数乘通道数），不一致会静默失败 \\
\Fn{AudioObj{\Bk}ectGetPr{\Bk}opertyDa{\Bk}ta} & 查询流配置、延迟、缓冲与标称速率 & \Tp{kAudioDe{\Bk}viceProp{\Bk}ertyStre{\Bk}amConfig{\Bk}uration}、\Tp{.\Bk.\Bk.{\Bk}Safety{\Bk}Offset}、\Tp{.\Bk.\Bk.{\Bk}Latency}、\Tp{.\Bk.\Bk.{\Bk}Buffer{\Bk}FrameSiz{\Bk}e} & 先问大小再取值；buffer、latency、safety offset 是三层，相加才可比 \\
\Fn{AudioObj{\Bk}ectAddPr{\Bk}opertyLi{\Bk}stenerBl{\Bk}ock} & 监听设备与速率变化 & \Tp{kAudioDe{\Bk}viceProp{\Bk}ertyDevi{\Bk}ceIsRunn{\Bk}ingSomew{\Bk}here}、标称采样率 & 回调落在你给的队列上，不是实时上下文 \\
\Fn{AVCaptur{\Bk}eSession} 与 \Fn{AVCaptur{\Bk}eVideoDa{\Bk}taOutput} & 相机或采集卡的输入图 & \Tp{activeFo{\Bk}rmat} 决定分辨率与帧率的组合 & 请求会向下取可达模式，必须回读；丢帧策略在输出端的设置里 \\
\Fn{CVPixelB{\Bk}ufferGet{\Bk}BaseAddr{\Bk}essOfPla{\Bk}ne} & 取 YUV 双平面地址 & \Tp{kCVPixel{\Bk}FormatTy{\Bk}pe\_{\Bk}420YpC{\Bk}bCr8BiPl{\Bk}anarVide{\Bk}oRange} 与对应的 FullRange 变体 & 先 \Fn{CVPixelB{\Bk}ufferLock{\Bk}BaseAddr{\Bk}ess}；\Tp{bytesPer{\Bk}Row} 与宽高都不等于你的假设值 \\
\Fn{VIDIOC\_{\Bk}QUERYC{\Bk}AP}（V4L2） & 确认设备能力 & \Tp{V4L2\_{\Bk}CAP\_{\Bk}VIDEO\_{\Bk}CAPTURE}、\Tp{.\Bk.\Bk.{\Bk}STREAM{\Bk}ING}、\Tp{.\Bk.\Bk.{\Bk}READWR{\Bk}ITE} & 只有支持 \Tp{READWRITE} 才能用 \Fn{read}，否则必须走缓冲区队列 \\
\Fn{VIDIOC\_{\Bk}ENUM\_{\Bk}FMT} 与 \Fn{ENUM\_{\Bk}FRAMES{\Bk}IZES}、\Fn{ENUM\_{\Bk}FRAMEI{\Bk}NTERVALS} & 枚举像素格式、尺寸与帧率 & \Tp{V4L2\_{\Bk}PIX\_{\Bk}FMT\_{\Bk}YUYV}、\Tp{.\Bk.\Bk.{\Bk}MJPEG}、\Tp{.\Bk.\Bk.{\Bk}NV12} & 三个枚举的笛卡尔积不等于全部可用组合，需要逐个试 \\
\Fn{VIDIOC\_{\Bk}S\_{\Bk}FMT} & 提交宽高、像素格式与场序 & \Tp{v4l2\_{\Bk}format} 的 \Tp{pix} 子结构：\Tp{width}、\Tp{height}、\Tp{pixelfor{\Bk}mat}、\Tp{sizeimage} & 驱动可以改你的请求，之后一定要回读 \Fn{VIDIOC\_{\Bk}G\_{\Bk}FMT} \\
\Fn{VIDIOC\_{\Bk}REQBUFS}、\Fn{QUERYBUF}、\Fn{QBUF}、\Fn{DQBUF} & 缓冲协商、建立映射、入队与取帧 & \Tp{V4L2\_{\Bk}MEMORY\_{\Bk}MMAP} 常用，跨设备共享走 \Tp{DMABUF}；计数常取 4 到 8 & \Tp{bytesused} 可小于 \Tp{length}（变长压缩）；重设数量会释放既有映射 \\
\Fn{VIDIOC\_{\Bk}STREAM{\Bk}ON}、\Fn{STREAMOFF} & 启停数据流 & 参数是 \Tp{v4l2\_{\Bk}buf\_{\Bk}type} & 停止后队列归驱动，重启前需重新入队或重协商 \\
\Fn{select} 或 \Fn{poll} 配 \Fn{DQBUF} & 等待下一帧可读 & \Tp{POLLRDNO{\Bk}RM} & 忙等抬高 CPU 与延迟抖动；超时通常指向 USB 带宽或驱动问题 \\
\Fn{VIDIOC\_{\Bk}G\_{\Bk}CTRL}、\Fn{S\_{\Bk}CTRL}、\Fn{QUERYCTRL} & 曝光、增益、对焦等控制 & \Tp{V4L2\_{\Bk}CID\_{\Bk}EXPOSU{\Bk}RE\_{\Bk}AUTO}、\Tp{.\Bk.\Bk.{\Bk}FOCUS\_{\Bk}AUTO} & 先看 flag 位判断只读与惰性；手动模式会推翻你的帧率协商 \\
\Fn{VIDIOC\_{\Bk}TRY\_{\Bk}EXT\_{\Bk}CTRLS} 与 \Fn{G\_{\Bk}EXT\_{\Bk}CTRLS} & 批量读写扩展控制 & 一次带多个 \Tp{v4l2\_{\Bk}ext\_{\Bk}control} 结构 & 批量的每一项都要正确初始化保留字段，否则整批失败 \\
\Fn{jack\_{\Bk}port\_{\Bk}regist{\Bk}er}、\Fn{jack\_{\Bk}connect}（JACK） & 注册端口、将输入接到输出 & 端口名与 \Tp{JackRead}、\Tp{JackWrite} 标志；帧数由服务器定 & 采样率与周期都不是你能改的，回调大小变化是常态（第 3 章） \\
\Fn{jack\_{\Bk}transp{\Bk}ort\_{\Bk}query} & 取外部时钟与位置 & \Tp{jack\_{\Bk}position\_{\Bk}t} 加 \Tp{JackTran{\Bk}sportSto{\Bk}pped} 等状态 & 与设备时钟是两套时间源，混用即失步 \\
\Fn{snd\_{\Bk}mixer\_{\Bk}selem\_{\Bk}set\_{\Bk}playba{\Bk}ck\_{\Bk}volume} 与 dB 变体 & 硬件音量与静音 & 量程由对应的 \Fn{snd\_{\Bk}mixer\_{\Bk}selem\_{\Bk}get\_{\Bk}playback\_{\Bk}vol{\Bk}ume\_{\Bk}range} 给出，也可走 dB 接口 & 软件音量与硬件增益不是同一层，叠两次会提前削顶 \\
\Fn{IMMDevic{\Bk}eEnumera{\Bk}tor\_{\Bk}Regist{\Bk}erEndpoi{\Bk}ntNotifi{\Bk}cationCa{\Bk}llback} & 热插拔与默认设备变化 & 回调接口是 \Tp{IMMNotif{\Bk}icationC{\Bk}lient} & 回调里只置标志，重建整套对象放到工作线程上做 \\
\Fn{IAudioCl{\Bk}ient\_{\Bk}GetSer{\Bk}vice} & 从客户端取渲染或捕获接口 & \Tp{IID\_{\Bk}IAudio{\Bk}RenderCl{\Bk}ient}、\Tp{IID\_{\Bk}IAudio{\Bk}CaptureC{\Bk}lient} & 共享与独占都要走这步，只测独占容易漏掉回退路径 \\
\Fn{CMSample{\Bk}BufferGe{\Bk}tImageBu{\Bk}ffer} & 从采集回调里取像素缓冲 & 返回 \Tp{CVPixelB{\Bk}ufferRef} & 回调返回后对象可能回收，跨线程用前要持引用（第 4 章） \\
\end{longtable}

\section{libav* 关键类型与函数}

\begin{longtable}{>{\raggedright\arraybackslash}p{0.17\textwidth}
                  >{\raggedright\arraybackslash}p{0.22\textwidth}
                  >{\raggedright\arraybackslash}p{0.19\textwidth}
                  >{\raggedright\arraybackslash}p{0.20\textwidth}}
\caption{libav* 类型与函数速查}\label{tab:apB-libav} \\
\toprule
名称 & 语义 & 取值或默认 & 注意 \\
\midrule
\endfirsthead
\toprule
名称 & 语义 & 取值或默认 & 注意 \\
\midrule
\endhead
\midrule
\endfoot
\bottomrule
\endlastfoot
\Tp{AVFrame} & 未压缩样点的容器（视频帧或音频帧） & \Tp{data[]} 配 \Tp{linesize[]}、\Tp{format}、\Tp{pts}；音频侧 \Tp{nb\_{\Bk}samples}、\Tp{sample\_{\Bk}rate}、\Tp{ch\_{\Bk}layout} & 平面寻址按 \Tp{linesize} 走；引用计数由 \Fn{av\_{\Bk}frame\_{\Bk}unref} 与 \Fn{av\_{\Bk}frame\_{\Bk}free} 管 \\
\Tp{AVPacket} & 压缩码流的一段 & \Tp{data}、\Tp{size}、\Tp{pts}、\Tp{dts}、\Tp{duration}、\Tp{stream\_{\Bk}index}；\Tp{flags} 含 \Tp{AV\_{\Bk}PKT\_{\Bk}FLAG\_{\Bk}KEY} & \Tp{size} 之后还有强制填充区，不要按内容长度重切缓冲区 \\
\Tp{AVCodecC{\Bk}ontext} & 编解码器实例与其全部参数 & \Tp{time\_{\Bk}base}、\Tp{framerate}、\Tp{pix\_{\Bk}fmt}、\Tp{bit\_{\Bk}rate}、\Tp{rc\_{\Bk}max\_{\Bk}rate}、\Tp{rc\_{\Bk}buffer\_{\Bk}size}、\Tp{gop\_{\Bk}size}、\Tp{thread\_{\Bk}count} & 多数码控参数在 \Fn{avcodec\_{\Bk}open2} 时定标；\Tp{extradata} 的所有权归上下文 \\
\Tp{AVFormat{\Bk}Context} & 容器层的流表与 I/O & \Tp{pb}、\Tp{nb\_{\Bk}streams}、\Tp{streams[]}、\Tp{probesize}、\Tp{max\_{\Bk}analyze\_{\Bk}durati{\Bk}on}、\Tp{flags} & 输出侧必须先 \Fn{avformat\_{\Bk}write\_{\Bk}header}；\Tp{AVFMT\_{\Bk}FLAG\_{\Bk}GENPTS} 影响时间戳补齐 \\
\Tp{AVStream} 与 \Tp{AVCodecP{\Bk}arameter{\Bk}s} & 一条流的时间信息编解码信息 & \Tp{time\_{\Bk}base}、\Tp{start\_{\Bk}time}、\Tp{duration}、\Tp{avg\_{\Bk}frame\_{\Bk}rate}、\Tp{codecpar} & \Tp{AVStream.{\Bk}time\_{\Bk}base} 与 \Tp{AVCodecC{\Bk}ontext.{\Bk}time\_{\Bk}base} 可以不同，跨两者一律换基 \\
\Fn{av\_{\Bk}rescale\_{\Bk}q} & 时戳在两个时间基之间换算 & 入参是值与两个 \Tp{AVRation{\Bk}al}；内部走 64 位并四舍五入 & 不要手写乘积；要指定舍入方向用 \Fn{av\_{\Bk}rescale\_{\Bk}q\_{\Bk}rnd} \\
\Fn{av\_{\Bk}packet\_{\Bk}rescale\_{\Bk}ts} & 把包内 \Tp{pts}、\Tp{dts}、\Tp{duration} 一起换基 & 第二个参数是旧时间基，第三个是目标流的时间基 & 自己换基最常漏掉 \Tp{duration}，导致时长与偏移都偏 \\
\Fn{av\_{\Bk}frame\_{\Bk}alloc}、\Fn{av\_{\Bk}frame\_{\Bk}get\_{\Bk}buffer}、\Fn{av\_{\Bk}frame\_{\Bk}make\_{\Bk}writab{\Bk}le} & 分配、挂缓冲、解除写保护 & 对齐参数常传 32 或 64 & 共享缓冲上 \Fn{make\_{\Bk}writab{\Bk}le} 触发深拷贝，逐帧调用等于放弃引用计数的收益 \\
\Fn{avcodec\_{\Bk}send\_{\Bk}packet}、\Fn{avcodec\_{\Bk}receive\_{\Bk}frame} & 半双工队列的两个端口 & 返回 \Tp{AVERROR{\Bk}(EAGAIN)}、\Tp{AVERROR\_{\Bk}INVALI{\Bk}DDATA}、\Tp{AVERROR\_{\Bk}EOF} & 送 \Tp{NULL} 包进入排空；每次都要按返回值决定下一步，不存在一进一出 \\
\Fn{avcodec\_{\Bk}open2} 配 \Tp{AVDictio{\Bk}nary} & 应用公共与私有选项 & 键如 \Tp{threads}、\Tp{preset}、\Tp{crf}、\Tp{x264-{\Bk}params} & 打开后未消费的键留在出参字典里，这是拼错参数唯一的信号 \\
\Fn{avformat\_{\Bk}open\_{\Bk}input}、\Fn{avformat\_{\Bk}find\_{\Bk}stream\_{\Bk}info}、\Fn{av\_{\Bk}read\_{\Bk}frame} & 打开、嗅探、读下一个包 & 第三参可传自定义 \Tp{AVIOCont{\Bk}ext}；\Fn{av\_{\Bk}read\_{\Bk}frame} 返回 \Tp{AVERROR\_{\Bk}EOF} 结束 & 不保证按 \Tp{dts} 有序交错；嗅探预算不足会漏流（第 10 章） \\
\Fn{avio\_{\Bk}open2} 与自定义 \Tp{AVIO} 回调 & 替换 I/O 层 & \Tp{AVIO\_{\Bk}FLAG\_{\Bk}READ}、\Tp{AVIO\_{\Bk}FLAG\_{\Bk}WRITE}、\Tp{AVIO\_{\Bk}FLAG\_{\Bk}DIRECT} & 回调缓冲区末尾仍需保留输入填充余量，否则解析器会越界读 \\
\Fn{sws\_{\Bk}get{\Bk}Context}、\Fn{sws\_{\Bk}scale} & 缩放与像素格式转换 & \Tp{SWS\_{\Bk}BILINE{\Bk}AR}、\Tp{SWS\_{\Bk}LANCZOS}、\Tp{SWS\_{\Bk}BICUBIC} & 上下文可复用，尺寸或格式变了要重建；矩阵与量程另有一层（第 5、9 章） \\
\Fn{swr\_{\Bk}alloc\_{\Bk}set\_{\Bk}opts2}、\Fn{swr\_{\Bk}convert\_{\Bk}frame} & 重采样、声道与格式转换 & 输出侧参数在 \Fn{swr\_{\Bk}init} 之前定 & 结束时用空输入冲残余样；声道用 \Tp{AVChanne{\Bk}lLayout}，旧位掩码接口已废弃 \\
\Fn{av\_{\Bk}buffer{\Bk}src\_{\Bk}add\_{\Bk}frame\_{\Bk}flags}、\Fn{av\_{\Bk}buffer{\Bk}sink\_{\Bk}get\_{\Bk}frame} & 滤镜图的进与出 & \Tp{AV\_{\Bk}BUFFER{\Bk}SRC\_{\Bk}FLAG\_{\Bk}PUSH}、\Tp{AV\_{\Bk}BUFFER{\Bk}SRC\_{\Bk}FLAG\_{\Bk}KEEP\_{\Bk}REF} & 图由 \Fn{avfilter\_{\Bk}graph\_{\Bk}parse\_{\Bk}ptr} 与 \Fn{avfilter\_{\Bk}graph\_{\Bk}config} 建成，帧所有权仍在调用方 \\
\Fn{av\_{\Bk}channel\_{\Bk}layout\_{\Bk}default}、\Fn{av\_{\Bk}channel\_{\Bk}layout\_{\Bk}descri{\Bk}be} & 新声道布局模型（FFmpeg 5.1 起） & \Tp{order} 配 \Tp{u.{\Bk}mask} 或 \Tp{nb\_{\Bk}channe{\Bk}ls} & 与旧接口混用会静默改变声道顺序，迁移时按第 3 章的对照表核 \\
\Fn{av\_{\Bk}err2str}、\Fn{av\_{\Bk}log\_{\Bk}set\_{\Bk}level}、\Fn{av\_{\Bk}log\_{\Bk}set\_{\Bk}cal{\Bk}lback} & 错误与日志 & \Tp{AV\_{\Bk}LOG\_{\Bk}DEBUG} 到 \Tp{AV\_{\Bk}LOG\_{\Bk}QUIET}；文本缓冲建议 1024 字节 & 日志可能来自编码器内部线程，回调里不要分配也不要加锁 \\
\Fn{avformat\_{\Bk}write\_{\Bk}header}、\Fn{av\_{\Bk}interl{\Bk}eaved\_{\Bk}write\_{\Bk}frame}、\Fn{av\_{\Bk}write\_{\Bk}trailer} & 输出侧的三步 & 交错版会在内存里排序，非交错版要求你自己有序 & 队列上限由 \Tp{max\_{\Bk}muxing\_{\Bk}queue\_{\Bk}size} 决定，超了就整条失败（第 10 章） \\
\Fn{av\_{\Bk}bsf\_{\Bk}get\_{\Bk}by\_{\Bk}name} 与 \Fn{av\_{\Bk}bsf\_{\Bk}send\_{\Bk}packet}、\Fn{av\_{\Bk}bsf\_{\Bk}receive\_{\Bk}packet} & 比特流过滤器 & 例如 \Tp{h264\_{\Bk}mp4toa{\Bk}nnexb}、\Tp{extract\_{\Bk}extrad{\Bk}ata}、\Tp{trace\_{\Bk}headers} & 与解码无关的语法修补；换封装时 extradata 的方向最容易搞反 \\
\Fn{avcodec\_{\Bk}find\_{\Bk}decoder\_{\Bk}by\_{\Bk}name}、\Fn{avcodec\_{\Bk}find\_{\Bk}encoder} & 按名字或 ID 查实现 & 同一 ID 常有多个实现（软件、硬件、厂商） & 选错实现就丢掉硬件路径；生产代码按名字查并保留回退表（第 8 章） \\
\Tp{sample\_{\Bk}aspect\_{\Bk}ratio}（\Tp{AV{\Bk}Stream} 与 \Tp{AVCodec{\Bk}Context} 同名）与 \Tp{AV\_{\Bk}FRAME\_{\Bk}DATA\_{\Bk}DISPL{\Bk}AYM{\Bk}ATRIX} & 采样宽高比与显示旋转 & 前者是 \Tp{AVRation{\Bk}al}，分子为 0 表示未知；后者是 3$\times$3 的 \Tp{int32\_{\Bk}t} 矩阵 & SAR 未知就按 1:1 处理，缩放目标算错会拉扁画面；旋转要从 side data 用 \Fn{av\_{\Bk}display\_{\Bk}rotation\_{\Bk}get} 读角度（第 5 章） \\
\Fn{av\_{\Bk}gettime\_{\Bk}relati{\Bk}ve}、\Fn{av\_{\Bk}usleep} & 单调时钟与短睡 & 单位微秒 & 内部同步一律用单调钟，墙上时钟只用于对外报告（第 2 章） \\
\Tp{AVDictio{\Bk}nary} 与 \Fn{av\_{\Bk}dict\_{\Bk}set}、\Fn{av\_{\Bk}dict\_{\Bk}get} & 选项串到结构体的中转 & 键值都以字符串传；标志位控制是否复制、是否忽略后缀 & 打开后剩余项就是你写错的项，务必检查 \\
\end{longtable}

\section{编码器参数对照}

概念是同一批，名字各不相通。下表横向对照，落到具体版本时以各自的帮助页与头文件为准。

\begin{longtable}{>{\raggedright\arraybackslash}p{0.11\textwidth}
                  >{\raggedright\arraybackslash}p{0.13\textwidth}
                  >{\raggedright\arraybackslash}p{0.13\textwidth}
                  >{\raggedright\arraybackslash}p{0.13\textwidth}
                  >{\raggedright\arraybackslash}p{0.13\textwidth}
                  >{\raggedright\arraybackslash}p{0.12\textwidth}}
\caption{等价概念在 x264、x265、aomenc、SVT-AV1 与 Opus 中的名字}\label{tab:apB-encoders} \\
\toprule
概念 & x264 & x265 & aomenc & SVT-AV1 & Opus \\
\midrule
\endfirsthead
\toprule
概念 & x264 & x265 & aomenc & SVT-AV1 & Opus \\
\midrule
\endhead
\midrule
\endfoot
\bottomrule
\endlastfoot
码控模式 & \Opt{crf}（默认 23）、\Opt{bitrate}、\Opt{qp}、\Opt{qpmin} 与 \Opt{qpmax}、\Opt{ratetol} & \Opt{crf}（默认 28）、\Opt{bitrate}、\Opt{qp}、\Opt{ctu} & \Opt{end-{\Bk}usage} 取 q、vbr、cbr，配 \Opt{cq} & \Opt{rc} 选模式（0 到 4），\Opt{crf} 或 \Opt{bitrate} 给目标 & 默认 VBR，关掉即 CBR，无质量模式 \\
峰值与缓冲 & \Opt{vbv-{\Bk}maxrate}、\Opt{vbv-{\Bk}bufsize}（单位 kbit） & 同名两个开关 & \Opt{bitrate} 配 \Opt{buf-{\Bk}sz}，两遍编码另有初始与最优缓冲 & \Opt{rc} 的 CBR 分支带峰值与缓冲约束 & 不适用 \\
关键帧 & \Opt{keyint}（默认 250）、\Opt{min-{\Bk}keyint}、\Opt{intra-{\Bk}refresh}、\Opt{scenecut} & \Opt{keyint}、\Opt{min-{\Bk}keyint}、\Opt{open-{\Bk}gop}、\Opt{no-{\Bk}scenec{\Bk}ut} & \Opt{kf-{\Bk}min-{\Bk}dist}、\Opt{kf-{\Bk}max-{\Bk}dist}、\Opt{static-{\Bk}thresh} & \Opt{g} 给 GOP 长度，另有场景切换检测开关 & 无对应概念，用包长表达 \\
速度与质量档 & \Opt{preset} 从 ultrafast 到 placebo（默认 medium），配 \Opt{tune} & 同名 \Opt{preset} 一族 & \Opt{cpu-{\Bk}used} 配 \Opt{deadline} 取 good、best、realtime & \Opt{preset}（数字越大越快，默认值随版本） & \Tp{OPUS\_{\Bk}SET\_{\Bk}COMPLE{\Bk}XITY} 取 0 到 10 \\
线程与并行 & \Opt{threads}，帧级或切片级由 \Opt{preset} 决定 & \Opt{pools}、\Opt{frame-{\Bk}threads}、\Opt{wpp} & \Opt{threads}、\Opt{tile-{\Bk}columns}、\Opt{enable-{\Bk}row-{\Bk}mt} & 以内部线程模型为主，核数开关名称随版本 & 单流内无并行；多流走 \Tp{OpusMuxer} \\
前瞻与延迟 & \Opt{rc-{\Bk}lookah{\Bk}ead}（置 0 即关）、\Opt{bframes}、\Opt{tune} 的 zerolatency & \Opt{look{\Bk}ahead}（默认约 20）、\Opt{bframes}、\Opt{lowpass-{\Bk}dct} & \Opt{lag-{\Bk}in-{\Bk}frames} 直接决定缓冲深度与首帧延迟 & 实时分支的前瞻相关开关，名称随版本 & \Tp{OPUS\_{\Bk}GET\_{\Bk}LOOKAH{\Bk}EAD} 报告算法延迟 \\
感知与细节 & \Opt{aq-{\Bk}mode}、\Opt{aq-{\Bk}streng{\Bk}th}、\Opt{deblock} & \Opt{aq-{\Bk}mode}、\Opt{psy-{\Bk}rd}、\Opt{psy-{\Bk}rdoq}、\Opt{min-{\Bk}cu-{\Bk}size} & \Opt{aq-{\Bk}mode}、\Opt{arnr-{\Bk}streng{\Bk}th}、\Opt{enable-{\Bk}cdef} & \Opt{aq-{\Bk}mode} 与环路滤波相关开关 & \Tp{OPUS\_{\Bk}SET\_{\Bk}SIGNAL} 取 voice 或 music \\
码流形态 & \Opt{nalu-{\Bk}file} 便于逐包检查；Annex B 还是索引式由封装决定 & \Opt{nalu-{\Bk}file}、\Opt{aud} 访问单元分隔符 & \Opt{ivf} 输出便于逐帧检查 & 默认输出 IVF，裸 OBU 形态取决于版本开关 & 每包自洽，长度与带宽在包内表达（第 8 章） \\
音频专属 & 不适用 & 不适用 & 不适用 & 不适用 & 帧长 2.5 到 60 毫秒、\Tp{OPUS\_{\Bk}SET\_{\Bk}INBAND\_{\Bk}FEC}、\Tp{OPUS\_{\Bk}SET\_{\Bk}DTX}、\Tp{OPUS\_{\Bk}SET\_{\Bk}BANDWI{\Bk}DTH} \\
多遍编码 & \Opt{stats} 配 \Opt{pass} 1 与 2 & 同名两个开关，统计文件可自定义名 & \Opt{passes} 取 2，第一遍的统计文件由 \Opt{fpf} 指定，第二遍同名传入 & \Opt{pass} 配合统计文件，两遍之间参数必须一致 & 不适用 \\
\end{longtable}

\section{命令行：ffmpeg 与 ffprobe}

\begin{longtable}{>{\raggedright\arraybackslash}p{0.17\textwidth}
                  >{\raggedright\arraybackslash}p{0.22\textwidth}
                  >{\raggedright\arraybackslash}p{0.19\textwidth}
                  >{\raggedright\arraybackslash}p{0.20\textwidth}}
\caption{ffmpeg 与 ffprobe 常用开关}\label{tab:apB-cli} \\
\toprule
名称 & 语义 & 取值或默认 & 注意 \\
\midrule
\endfirsthead
\toprule
名称 & 语义 & 取值或默认 & 注意 \\
\midrule
\endhead
\midrule
\endfoot
\bottomrule
\endlastfoot
\Tp{-{\Bk}c}（\Tp{-{\Bk}codec}） & 指定编解码或直接搬运 & \Tp{copy}、\Tp{libx264}、\Tp{aac}；带流选择写作 \Tp{-{\Bk}c:v}、\Tp{-{\Bk}c:a}，丢弃写作 \Tp{-{\Bk}an}、\Tp{-{\Bk}vn} & \Tp{copy} 不重编码，第 10 章的时基与全局头问题会原样存在 \\
\Tp{-{\Bk}map} & 选流与排列 & \Tp{-{\Bk}map 0:v:0 -{\Bk}map 0:a:1}；\Tp{-{\Bk}map\_{\Bk}metada{\Bk}ta -{\Bk}1} 清元数据 & 不写 \Tp{-{\Bk}map} 时每类只自动挑一条，多轨很容易丢 \\
\Tp{-{\Bk}ss}、\Tp{-{\Bk}to}、\Tp{-{\Bk}t} & 定位与截取 & \Tp{-{\Bk}ss} 放 \Tp{-{\Bk}i} 前是容器级粗定位，放后是解码到该点 & 要求帧精确就用后置形式；\Tp{-{\Bk}accura{\Bk}te\_{\Bk}seek} 默认开 \\
\Tp{-{\Bk}crf} & 恒定质量 & x264 默认 23，x265 默认 28，VP9 与 AV1 走 0 到 63 & 只在对码率有控制力的模式里生效，与 \Tp{-{\Bk}b:v} 同给时语义依编码器 \\
\Tp{-{\Bk}b:v}、\Tp{-{\Bk}maxrate}、\Tp{-{\Bk}bufsize} & 目标码率、峰值与校验缓冲 & 比特每秒，可写 \Tp{6000}、\Tp{4M} & 后两个是约束不是目标；直播链路三个要一起给（第 7 章） \\
\Tp{-{\Bk}g}、\Tp{-{\Bk}keyint\_{\Bk}min}、\Tp{-{\Bk}force\_{\Bk}key\_{\Bk}frames} & 关键帧间距与强制插入点 & \Tp{-{\Bk}g 60}；\Tp{-{\Bk}force\_{\Bk}key\_{\Bk}frames} 可写表达式 \Tp{expr:gte(t,n\_{\Bk}forced*2)} & 间距决定档位切换粒度与抗误码能力，应与分片时长成整数倍 \\
\Tp{-{\Bk}preset}、\Tp{-{\Bk}tune} & 速度档与场景倾向 & x264 从 \Tp{ultrafast} 到 \Tp{placebo}（默认 \Tp{medium}）；\Tp{-{\Bk}tune} 取 zerolatency、fastdecode、film & zerolatency 会关前瞻与 B 帧，代价是同码率下画质下降 \\
\Tp{-{\Bk}pix\_{\Bk}fmt} & 强制输出像素格式 & \Tp{yuv420p}、\Tp{nv12}、\Tp{yuv422p1{\Bk}0le} & 编码器不支持时是硬失败，不要指望自动降级（第 5 章） \\
\Tp{-{\Bk}vf}、\Tp{-{\Bk}filter\_{\Bk}complex} & 单链滤镜与图级滤镜 & \Tp{-{\Bk}vf} 常用 scale、fps、format、overlay、libvmaf & 链的顺序决定精度路径；硬件帧要配套的滤镜与设备参数 \\
\Tp{-{\Bk}movfla{\Bk}gs} & MP4 的盒布局 & \Tp{+faststart} 把索引前置；\Tp{+frag\_{\Bk}keyfra{\Bk}me+empty\_{\Bk}moov} 让字节流可边写边读 & fragmented 输出没有总时长字段，播放器要另给时长或读清单 \\
\Tp{-{\Bk}fflags} & 通用格式标志 & \Tp{+genpts} 生成缺失时间戳；\Tp{+discardc{\Bk}orrupt}、\Tp{+igndts} & 补 PTS 不修内容损坏；反复 seek 时先确认它再查别处 \\
\Tp{-{\Bk}probes{\Bk}ize}、\Tp{-{\Bk}analyz{\Bk}eduratio{\Bk}n} & 打开阶段的嗅探预算 & 默认量级约 5 MiB 与数秒 & 太小会漏流或误报帧率；实时源常要调小以加快启动 \\
\Tp{-{\Bk}hwaccel} & 解码侧硬件加速 & \Tp{auto}、\Tp{cuda}、\Tp{vaapi}、\Tp{videotoo{\Bk}lbox}、\Tp{d3d11va}、\Tp{qsv}；配 \Tp{-{\Bk}hwaccel\_{\Bk}output\_{\Bk}format} 与 \Tp{-{\Bk}extra\_{\Bk}hw\_{\Bk}frames} & 只覆盖解码，缩放与格式转换仍要显式安排；不支持时可能静默回软解 \\
\Tp{-{\Bk}threads}、\Tp{-{\Bk}thread\_{\Bk}type} & 编解码并行度与并行方式 & \Tp{0} 表示自动；类型取 \Tp{frame} 或 \Tp{slice} & 帧级并行引入输出延迟，实时链路要压成切片级 \\
\Tp{-{\Bk}re} & 按源速率读取 & 无参数 & 不加它做推流测试会全速灌包，把拥塞控制测成假象 \\
\Tp{-{\Bk}f} & 强制输入或输出格式 & \Tp{rawvideo}、\Tp{mp4}、\Tp{matroska}、\Tp{hls}、\Tp{dash}、\Tp{srt}、\Tp{udp} & 裸流必须同时给几何参数；HLS 侧看 \Tp{-{\Bk}hls\_{\Bk}time}、\Tp{-{\Bk}hls\_{\Bk}list\_{\Bk}size}、\Tp{-{\Bk}hls\_{\Bk}flags} delete\_segments \\
\Tp{-{\Bk}timeout}、\Tp{-{\Bk}rw\_{\Bk}timeout} & 网络与协议超时 & 单位与含义依协议而定，常见为微秒 & 网络阻塞只表现为返回的超时错误码，不会留下异常文本 \\
\Tp{-{\Bk}progre{\Bk}ss}、\Tp{-{\Bk}nostats}、\Tp{-{\Bk}loglev{\Bk}el} & 机器可读进度与噪声控制 & \Tp{-{\Bk}progre{\Bk}ss} 可指向 \Tp{pipe:1} 或 \Tp{pipe:2}；\Tp{-{\Bk}loglev{\Bk}el} 可叠 \Tp{repeat+{\Bk}level+{\Bk}debug} & \Tp{-{\Bk}progre{\Bk}ss} 输出的是键值行，是子进程集成的正确接口 \\
\Tp{ffprobe -{\Bk}show\_{\Bk}streams}、\Tp{-{\Bk}show\_{\Bk}packets}、\Tp{-{\Bk}show\_{\Bk}frames} & 结构化读流、读包、读帧 & 用 \Tp{-{\Bk}show\_{\Bk}entries} 限定字段；\Tp{-{\Bk}count\_{\Bk}frames} 才会给出逐帧计数 & \Tp{-{\Bk}show\_{\Bk}frames} 需要真解码，长文件上是慢操作 \\
\Tp{ffprobe -{\Bk}print\_{\Bk}format}（别名 \Tp{-{\Bk}of}） & 输出格式 & \Tp{json}、\Tp{csv}、\Tp{xml}、\Tp{ini}；\Tp{csv=p=0} 去掉表头 & 脚本里读 JSON，不要解析人读表格；配 \Tp{-{\Bk}v} error 与 \Tp{-{\Bk}hide\_{\Bk}banner} \\
\Tp{ffprobe -{\Bk}select\_{\Bk}streams} & 挑选要输出的流 & \Tp{v:0}、\Tp{a:1}，序号从 0 起 & 用程序内的流序号，不是传输层 PID；多节目源容易挑错 \\
\Tp{ffprobe -{\Bk}show\_{\Bk}format} & 容器级字段 & \Tp{format={\Bk}duration,{\Bk}bit\_{\Bk}rate,{\Bk}nb\_{\Bk}streams} & 直播或无索引文件可能缺时长，估算策略随版本变化 \\
\Tp{mediainfo -\Bk-{\Bk}Inform} & 精简的属性读取 & 段名 \Tp{General}、\Tp{Video}、\Tp{Audio}，模板写作百分号占位符 & 占位符区分大小写，字段名与 ffprobe 不通用 \\
\Tp{-{\Bk}frag\_{\Bk}durati{\Bk}on}、\Tp{-{\Bk}frag\_{\Bk}type} & 分片输出的切分依据 & 时长按微秒给，或按帧类型切 & 与 \Tp{-{\Bk}movfla{\Bk}gs} 里的分片标志配套使用，单独一个不生效 \\
\Tp{-{\Bk}err\_{\Bk}detect} & 解码侧的检错强度 & 取 \Tp{ignore\_{\Bk}err}、\Tp{compliant}、\Tp{aggressi{\Bk}ve}、\Tp{bitstream} & 宽松策略会把损坏静默变成花屏，排障时先收紧（第 8 章） \\
\Tp{-{\Bk}metada{\Bk}ta}、\Tp{-{\Bk}map\_{\Bk}metada{\Bk}ta} & 写或搬元数据 & \Tp{-{\Bk}metada{\Bk}ta title=.\Bk.\Bk.}；\Tp{-{\Bk}map\_{\Bk}metada{\Bk}ta -{\Bk}1} 清空 & 容器与封装是否支持该键决定它能否落盘 \\
\Tp{-{\Bk}frames:v}、\Tp{-{\Bk}frames:a} & 限制输出的帧数 & 正整数，常用于切片与测试 & 与 \Tp{-{\Bk}t} 的区别是计数而非时长，变速源上结果不同 \\
\Tp{-{\Bk}x264-{\Bk}params}、\Tp{-{\Bk}x265-{\Bk}params} & 把私有参数直传编码器 & 形如 \Tp{key={\Bk}value:{\Bk}value={\Bk}value} 的串 & 与同名公共开关同时给时以私有参数为准，容易掩盖笔误 \\
\end{longtable}

\begin{lstlisting}[style=jsonstyle,caption={ffprobe 的 JSON 输出：读流信息与时间基},label={lst:apB-ffprobe}]
{
  "streams": [
    { "index": 0, "codec_name": "h264", "codec_type": "video",
      "profile": "High", "level": 40, "pix_fmt": "yuv420p",
      "width": 1920, "height": 1080, "has_b_frames": 2,
      "r_frame_rate": "30000/1001", "avg_frame_rate": "30000/1001",
      "time_base": "1/90000", "start_time": "0.000000",
      "duration": "10.010000", "bit_rate": "4210000", "nb_frames": "300" },
    { "index": 1, "codec_name": "aac", "codec_type": "audio",
      "sample_rate": "48000", "channels": 2, "channel_layout": "stereo",
      "time_base": "1/48000", "bit_rate": "128000" }
  ],
  "format": { "nb_streams": 2, "duration": "10.010000", "bit_rate": "4352000" }
}
\end{lstlisting}

\section{GStreamer 常用元素与属性}

\begin{longtable}{>{\raggedright\arraybackslash}p{0.17\textwidth}
                  >{\raggedright\arraybackslash}p{0.22\textwidth}
                  >{\raggedright\arraybackslash}p{0.19\textwidth}
                  >{\raggedright\arraybackslash}p{0.20\textwidth}}
\caption{GStreamer 元素与属性速查}\label{tab:apB-gst} \\
\toprule
名称 & 语义 & 取值或默认 & 注意 \\
\midrule
\endfirsthead
\toprule
名称 & 语义 & 取值或默认 & 注意 \\
\midrule
\endhead
\midrule
\endfoot
\bottomrule
\endlastfoot
\Fn{gst\_{\Bk}element\_{\Bk}factory\_{\Bk}make}、\Fn{gst\_{\Bk}parse\_{\Bk}launch} & 由类型名造元素，或由描述串建管线 & 串里 ``!'' 连静态 pad，动态 pad 要自己接 \Tp{pad-{\Bk}added} & 解析版只在构造时协商，运行期改结构得手搭；解析错误走 \Tp{GError} \\
\Tp{appsrc} & 把应用产生的缓冲推进链 & \Tp{caps}、\Tp{format} 取 \Tp{GST\_{\Bk}FORMAT\_{\Bk}TIME}、\Tp{max-{\Bk}bytes}、\Tp{do-{\Bk}timest{\Bk}amp}、\Tp{is-{\Bk}live} & 用 \Tp{need-{\Bk}data} 与 \Tp{enough-{\Bk}data} 信号做背压，不要自驱循环塞满 \\
\Tp{appsink} & 从链尾取回缓冲 & \Tp{emit-{\Bk}signals}、\Tp{max-{\Bk}buffers}、\Tp{drop}、\Tp{sync}、\Tp{caps} & \Tp{drop} 为真时超量丢最旧帧；\Tp{new-{\Bk}sample} 回调在流线程上，重活要转手 \\
\Tp{decodebin}、\Tp{decodebi{\Bk}n3} & 自动挑解封装与解码元素 & 输出 pad 由信号给出，可再套 \Tp{capsfilt{\Bk}er} 收窄 & \Tp{autoplug} 一族的属性在两代之间含义不同，改结构前先按 \Fn{gst-{\Bk}inspect-{\Bk}1.{\Bk}0} 核对 \\
\Tp{playbin3}、\Tp{uridecod{\Bk}ebin} & 高层播放器与按 URI 解码 & \Tp{uri}、\Tp{video-{\Bk}si{\Bk}nk}、\Tp{audio-{\Bk}si{\Bk}nk}、\Tp{flags} & 想控线程与缓冲就换手搭，这两个会把结构藏起来 \\
\Tp{queue}、\Tp{tee} & 线程边界、缓冲解耦与分流 & \Tp{max-{\Bk}size-{\Bk}buffers}、\Tp{max-{\Bk}size-{\Bk}bytes}、\Tp{max-{\Bk}size-{\Bk}time}、\Tp{leaky} & 队列就是第 12 章线程边界的实现；漏队列的方向决定丢最旧还是最新 \\
\Tp{capsfilt{\Bk}er}、\Tp{videocon{\Bk}vert}、\Tp{videosca{\Bk}le}、\Tp{videorate} & 固定 caps，自动补齐转换 & caps 串写作 \Tp{video/{\Bk}x-{\Bk}raw} 加 \Tp{format=NV12} 一类字段 & 只约束一个字段时其余交给协商，跨模块的问题常出在这 \\
\Tp{v4l2src}、\Tp{avfvideo{\Bk}src}、\Tp{nvargusc{\Bk}amerasrc} 等平台采集元素 & 从设备取帧 & \Tp{device}、\Tp{io-{\Bk}mode}，尺寸与帧率走 caps & 各平台元素支持范围差别大，先看 \Fn{gst-{\Bk}inspect-{\Bk}1.{\Bk}0} 给的 caps \\
\Tp{alsa{\Bk}sink}、\Tp{pulse{\Bk}sink}、\Tp{coreaudio{\Bk}sink}、\Tp{wasapi{\Bk}sink} & 音频输出 & \Tp{device}；ALSA 的 \Tp{buffer-{\Bk}time} 与 \Tp{latency-{\Bk}time} 单位微秒 & 设得比设备能力小会被就近抬高，回读才可信（第 3 章） \\
\Tp{audiocon{\Bk}vert}、\Tp{audiores{\Bk}ample}、\Tp{autoaudio{\Bk}sink} & 音频侧格式与速率补齐 & 由 caps 协商触发 & 与视频侧同理：显式插一个胜过让协商猜 \\
\Tp{h264parse}、\Tp{h265parse}、\Tp{avdec\_{\Bk}h264}、\Tp{vaapih26{\Bk}4dec}、\Tp{nvh264dec} & 参数集整理，软件或硬件解码 & \Tp{config-{\Bk}interv{\Bk}al} 控制参数集重复间隔（秒） & 进 RTP 载荷前几乎必过 parse，否则缺参数集的首帧解不开 \\
\Tp{x264enc}、\Tp{vp8enc}、\Tp{vp9enc}、\Tp{vaapih26{\Bk}4enc}、\Tp{nvh264enc} & 编码元素 & \Tp{bitrate}（kbit）、\Tp{crf}、\Tp{speed-{\Bk}preset}、\Tp{tune}、\Tp{keyint-{\Bk}max}、\Tp{threads}；硬件一族多为码控模式与 GOP 开关 & 与 FFmpeg 同名不同义，属性名以 \Fn{gst-{\Bk}inspect-{\Bk}1.{\Bk}0} 为准；\Tp{byte-{\Bk}stream} 决定是否 Annex B \\
\Tp{qtdemux}、\Tp{matroska{\Bk}demux}、\Tp{mpegtsde{\Bk}mux}、\Tp{mpegtsmux} & 解封装与封装 & pad 随流动态增删 & TS 侧的节目表与参考时间间隔决定加入流的首帧等待，名称随元素版本 \\
\Tp{rtph264p{\Bk}ay}、\Tp{rtph264d{\Bk}epay}、\Tp{rtpjitte{\Bk}rbuffer}、\Tp{rtpbin} & RTP 载荷化、去载荷、抖动缓冲与会话关联 & \Tp{pt}、\Tp{mtu}、\Tp{config-{\Bk}interv{\Bk}al}；缓冲侧 \Tp{latency}、\Tp{do-{\Bk}lost}、\Tp{drop-{\Bk}on-{\Bk}latency} & \Tp{drop-{\Bk}on-{\Bk}latency} 为真时用丢帧换延迟，与第 12 章的取舍一致 \\
\Tp{udp{\Bk}sink}、\Tp{udpsrc}、\Tp{srtpenc}、\Tp{srtpdec} & 传输与加密边界 & \Tp{host}、\Tp{port}、\Tp{buffer-{\Bk}size}、\Tp{reuse}、\Tp{socket} & 端口写 0 表示让内核选；多路复用要自己建套接字对象共享 \\
\Tp{webrtcbin} & 连接、ICE、DTLS、SRTP 与统计 & \Tp{bundle-{\Bk}policy}、\Tp{stun-{\Bk}server}，传输走 \Tp{transport} 的 pad & 它把第 11 到 14 章的环节打包成元素，改行为靠属性与信令层 \\
\Tp{identity}、\Tp{fake{\Bk}sink}、\Tp{fpsdispl{\Bk}aysi{\Bk}nk} & 打点、丢弃与统计显示 & \Tp{sync}、\Tp{silent}；后者直接报帧率 & \Tp{identity} 的回调是最小侵入的每帧时延测法（第 2 章） \\
\Fn{gst\_{\Bk}pad\_{\Bk}add\_{\Bk}probe}、\Fn{gst\_{\Bk}buffer\_{\Bk}map}、\Fn{gst\_{\Bk}pad\_{\Bk}query\_{\Bk}caps} & 挂探针、映射字节、问 caps & 探针类型取 BUFFER、EVENT、QUERY，方向与阻塞性另设 & 探针里不可重入管线锁；映射必须有配对的 unmap，否则缓冲池不回收 \\
环境变量 \Tp{GST\_{\Bk}DEBUG}、\Tp{GST\_{\Bk}PLUGIN\_{\Bk}PATH}、\Tp{GST\_{\Bk}REGIST{\Bk}RY}、\Tp{GST\_{\Bk}TRACERS} & 调试输出、插件搜索、注册表缓存与跟踪器 & \Tp{GST\_{\Bk}DEBUG=3}，或 \Tp{GST\_{\Bk}DEBUG=v4l2src:5} 按类别调；\Tp{GST\_{\Bk}TRACERS} 取 latency 一类 & 嵌入形态下注册表决定协商结果，容器里要确认插件目录与缓存策略 \\
\Fn{gst\_{\Bk}bin\_{\Bk}add\_{\Bk}many}、\Fn{gst\_{\Bk}element\_{\Bk}link\_{\Bk}many}、\Fn{gst\_{\Bk}element\_{\Bk}get\_{\Bk}request\_{\Bk}pad} & 手搭容器、静态连接与索取动态 pad & 每个连接都要检查返回值；模板 pad 名字形如 \Tp{src\%u} & 动态 pad 要自己接 \Tp{pad-{\Bk}added} 并持引用，忘记 unref 是常见的泄漏源 \\
\Fn{gst\_{\Bk}element\_{\Bk}get\_{\Bk}state}、\Fn{gst\_{\Bk}bus\_{\Bk}timed\_{\Bk}pop\_{\Bk}filter{\Bk}ed} & 查询状态、从总线取消息 & \Tp{GST\_{\Bk}STATE\_{\Bk}PAUSED} 用于预滚动；筛 \Tp{GST\_{\Bk}MESSAGE\_{\Bk}ERROR} 与 \Tp{EOS} & 返回 \Tp{NO\_{\Bk}PAID} 只是异步没走完，别当成失败（第 13 章的启动顺序） \\
\Tp{fdsrc}、\Tp{fd{\Bk}sink}、\Tp{shmsrc}、\Tp{shmsink} & 跨边界的字节与缓冲传输 & \Tp{fd}、\Tp{socket} 与 \Tp{shm-{\Bk}size}、\Tp{wait-{\Bk}for-{\Bk}connect} & 进程间传整帧时比管道少一次拷贝，但缓冲所有权要自己约定 \\
\Tp{rtprtxse{\Bk}nd}、\Tp{rtprtxre{\Bk}ceive}、\Tp{ulpfecenc}、\Tp{flexfece{\Bk}nc} & RTX 重传与两种前向纠错 & 载荷类型通过 \Tp{apt} 关联到主载荷；冗余度是属性参数 & 与 \Tp{rtpjitte{\Bk}rbuffer} 的 \Tp{do-{\Bk}lost} 配套才有效，单独开一个看不到收益 \\
\Tp{svtav1enc}、\Tp{openh264{\Bk}enc}、\Tp{vtenc\_{\Bk}h264} 一族 & 另一批编码器元素 & 属性名沿用各自库的习惯（码率、GOP、preset） & 与 FFmpeg 的同名概念不同键，仍以 \Fn{gst-{\Bk}inspect-{\Bk}1.{\Bk}0} 输出为准 \\
\end{longtable}

\begin{guideline}[title={速查表的边界}]
\begin{itemize}
\item 本表只给名字、语义与坑，不含完整签名：参数个数、返回值集合与所有权约定，
      一律以你链接的那个版本的头文件与官方文档为准。
\item 带``随版本''字样的条目（尤其码控与硬件加速的开关名）会在小版本之间改名或合并，
      发布前用 \Fn{gst-{\Bk}inspect-{\Bk}1.{\Bk}0}、\Tp{ffmpeg -{\Bk}h full}、\Tp{ffprobe -{\Bk}help} 现场核对一遍。
\item 表里的默认值是常见主线构建的取值，发行版与自研构建可能重编过选项：
      先读 \Tp{-{\Bk}buildc{\Bk}onf} 与 \Tp{-{\Bk}version}，再引用默认值。
\end{itemize}
\end{guideline}

% ============================================================
% ---- frag_apC.tex ----
\chapter{选型矩阵与排障手册}
\label{ap:apC-selection}

本附录回答一个动作：要做决定时翻哪一页。附录 \ref{ap:apA-glossary} 解决
``这个词指什么''，附录 \ref{ap:apB-api-quickref} 解决``这个调用叫什么''，
本附录解决剩下两类问题——``这套需求该选哪套组合''与``现在这个症状先查什么''。
前三节是选型矩阵：引擎与框架、编解码器、许可。第四节把延迟预算拆成一条
可执行的决策链。第五节是按症状索引的排障表，每条给首查项、常见根因与
对应章。

六张表的读法一致。``首选''列只收不带 copyleft 义务的组件（BSD、
Apache-2.0、MIT、ISC 条款）与平台自带接口；``替代''列把 \Lib{FFmpeg}、
\Lib{GStreamer} 这类 LGPL 或 GPL 的实现与专有运行库放在那里，只作客观陈述
并说明分发影响，不作为推荐项。``章''列给出正文里展开讨论的章号，需要论证
过程时回正文；需要具体开关名时回附录 \ref{ap:apB-api-quickref}。

\section{引擎与框架的选型矩阵}

框架选型实际选的是四件事：链接形态、许可义务、崩溃半径与升级成本（第
\ref{ch:ch13-framework-integration} 章把它们列成嵌入形态对比表）。下表按
需求形态排，第一列是``谁在用这套组合''，不是``谁功能多''。

\small
\begin{longtable}{>{\raggedright\arraybackslash}p{0.16\textwidth}
                  >{\raggedright\arraybackslash}p{0.19\textwidth}
                  >{\raggedright\arraybackslash}p{0.30\textwidth}
                  >{\raggedright\arraybackslash}p{0.11\textwidth}}
\caption{按需求形态排的引擎与框架选择}\label{tab:apC-engine} \\
\toprule
需求形态 & 首选（宽松条款或平台自带） & 常见替代与客观差异 & 章 \\
\midrule
\endfirsthead
\toprule
需求形态 & 首选（宽松条款或平台自带） & 常见替代与客观差异 & 章 \\
\midrule
\endhead
\midrule
\endfoot
\bottomrule
\endlastfoot
只做播放与解码 & \Lib{Media Foundation}、\Lib{Video{\Bk}Toolbox} 等平台自带接口，不引入第三方开源条款 & \Lib{libav*} 默认按 LGPL-2.1+ 构建，\Lib{GStreamer} 核心 LGPL 而插件域混杂；两者的格式覆盖与逐帧控制都更强，代价是产物许可由链入项决定 & \ref{ch:ch08-decode-hw}、\ref{ch:ch13-framework-integration} \\
需要逐帧像素与自定义处理 & 平台接口交付帧，转换核自己写 & \Lib{libav*} 的缩放、重采样与滤镜图省事，但同一批开关里 \Opt{enable-\Bk gpl} 会把整份产物推进 GPL；滤镜图的线程与帧引用模型要读两遍 & \ref{ch:ch05-pixel-format}、\ref{ch:ch09-filter-resample} \\
实时互动（通话、连麦） & \Lib{libwebrtc}（BSD-3）：回声消除、降噪、带宽估计与抖动缓冲是一整套 & \Tp{webrtcbin} 把同样的环节打包成元素，改行为靠属性与信令层；自研 RTP 栈自由度最高，但第 12 章列出的每一项都要自己补齐 & \ref{ch:ch11-streaming}、\ref{ch:ch12-realtime-thread} \\
服务端多路转码与录制 & 平台自带的硬件编码入口，配自己的调度层 & \Lib{GStreamer} 的元素粒度适合长管线；把 \Tp{ffmpeg} 当子进程用能把许可与崩溃都关在进程外，只留下管道协议这一层合同 & \ref{ch:ch13-framework-integration} \\
游戏引擎或客户端内播放 & 引擎自带后端，多数落到平台接口 & 进程内链接换到逐帧控制与亚帧延迟，同时接受许可、崩溃半径与版本矩阵三项成本 & \ref{ch:ch13-framework-integration} \\
设备与采集层 & 平台原生：\DomAud{} 用 ALSA、WASAPI、CoreAudio，\DomVid{} 用 V4L2、\Lib{Media Foundation} 与 \Lib{AV{\Bk}Foundation}；Linux 侧 \Lib{PipeWire}（MIT）同时管两类 & 在设备层再叠一层框架抽象会先吃掉延迟预算；热插拔、独占与回环的语义各家不同，抽象容易把差异抹平成不可控 & \ref{ch:ch03-audio-device}、\ref{ch:ch04-capture-display} \\
无真实设备的仿真与合成 & 自己按周期造帧，用管道喂下游，零第三方条款 & 现成的合成源省事，但时间轴由库推进，与第 2 章的时钟模型必须显式对齐，否则测出来的抖动是假的 & \ref{ch:ch02-clock-buffer} \\
只做再封装与切片 & 容器版式自己写，盒与索引自己维护 & \Lib{libav\Bk format} 久经考验，动态链接是 LGPL 义务最省事的满足方式；时间基换算与带外参数集这两笔账无论走哪条路都要自己算 & \ref{ch:ch10-mux-sync} \\
\end{longtable}
\normalsize

\section{编解码器的选型矩阵}

编解码器选型有两条互相拉扯的轴：一条是质量、延迟与复杂度之间的技术折中
（第 \ref{ch:ch07-rate-control} 章），另一条是条款与专利的账本（下一节的
许可矩阵）。下表按``要压什么内容''排，``首选''列一律给条款宽松或格式免版税
的实现；同一行右两列的差别只写决策时要看的那一点。

\small
\begin{longtable}{>{\raggedright\arraybackslash}p{0.16\textwidth}
                  >{\raggedright\arraybackslash}p{0.19\textwidth}
                  >{\raggedright\arraybackslash}p{0.30\textwidth}
                  >{\raggedright\arraybackslash}p{0.11\textwidth}}
\caption{按内容类型排的编解码器选择}\label{tab:apC-codec} \\
\toprule
内容类型与约束 & 首选（宽松条款或免版税格式） & 常见替代与客观差异 & 章 \\
\midrule
\endfirsthead
\toprule
内容类型与约束 & 首选（宽松条款或免版税格式） & 常见替代与客观差异 & 章 \\
\midrule
\endhead
\midrule
\endfoot
\bottomrule
\endlastfoot
\DomAud{}：通话与对讲 & \Lib{Opus}（BSD-3，格式免版税）：帧长、带宽与复杂度都能在运行期调 & 专用语音实现（低延迟 AAC 与 3GPP 一族的语音编解码）需要单独授权；\Lib{Opus} 的 DTX、带内 FEC 与静音检测在包内自洽，不需要额外的载荷格式 & \ref{ch:ch06-codec-api} \\
\DomAud{}：语音与音乐同一条流 & \Lib{Opus} 用 \Tp{OPUS\Bk\_SET\Bk\_SIGNAL} 切倾向 & 拆成两路编码更省，但两路的时间轴与缓冲要各自对齐；混合内容下单纯压码率会先把高频抹平 & \ref{ch:ch06-codec-api}、\ref{ch:ch07-rate-control} \\
\DomAud{}：极窄带宽回传 & \Lib{Opus} 的窄带宽档与语音倾向 & 更省码率的专用编解码几乎都带费条款；先确认远端解什么，再决定压多狠 & \ref{ch:ch06-codec-api} \\
\DomVid{}：实时上行（会议与连麦） & \Lib{OpenH264}（Cisco 维护，条款宽松）或平台自带编码器 & \Lib{x264} 是画质参照系但为 GPL-2+，另有商用授权通道；H.264 的专利池授权与代码许可无关，仍需自己确认 & \ref{ch:ch06-codec-api}、\ref{ch:ch07-rate-control} \\
\DomVid{}：离线高质量归档 & 平台自带编码器，或换到免版税格式绕开费条款 & \Lib{x264}、\Lib{x265}（均 GPL-2+）在两遍编码与慢档上仍是基线；HEVC 侧的专利池不止一个，覆盖的地域也不同 & \ref{ch:ch06-codec-api}、\ref{ch:ch07-rate-control} \\
\DomVid{}：开放生态的分发 & \Lib{libvpx}（VP8/VP9，BSD-3 加专利授权） & 硬件解码覆盖面不如 H.26x 广，出货前按目标设备查；浏览器与开放播放器生态对这两个格式支持稳定 & \ref{ch:ch06-codec-api}、\ref{ch:ch08-decode-hw} \\
\DomVid{}：AV1 编码 & \Lib{libaom}（BSD 条款加专利授权）或 \Lib{SVT-AV1}（多线程扩展好，宽松条款加专利授权） & \Lib{rav1e} 是另一条宽松路线；AV1 的编码速度与质量档位耦合很紧，两遍之间参数必须完全一致 & \ref{ch:ch06-codec-api}、\ref{ch:ch07-rate-control} \\
\DomVid{}：AV1 解码 & \Lib{dav1d}（BSD-2）作为软件解码的基线 & \Lib{libaom} 自带解码参数全但慢；硬件解码条目要看 GPU 代际与驱动，第 8 章按后端列了支持矩阵 & \ref{ch:ch08-decode-hw} \\
\DomVid{}：硬件 ASIC、并发路数优先 & 通过平台接口取用硬件编码器，把授权与驱动交给系统 & 厂商闭源运行库的再分发受其许可文本约束；Linux 侧的开源驱动可随发行版分发，但能力与位深随代际变化；硬件的码控开关少，同码率画质对软件有代差 & \ref{ch:ch08-decode-hw} \\
\DomVid{}：母版与中间格式 & 未压缩平面配免版税的中间表示，链路内不重编码 & 专有中间格式画质好、生态成熟，但绑定剪辑工具链与授权方式；跨平台交付时优先第 5 章的通用平面约定 & \ref{ch:ch05-pixel-format} \\
\end{longtable}
\normalsize

\section{许可矩阵}

许可不是发布前才读的一段说明文字，而是构建矩阵的一个维度（第
\ref{ch:ch13-framework-integration} 章）。对 \Lib{libav*} 而言，同一份源码树
加上不同的 configure 开关会产出许可完全不同的二进制：\Opt{enable-\Bk gpl}
把产物从 LGPL 翻成 GPL，\Opt{enable-\Bk nonfree} 直接让产物失去自由再分发的
资格。因此出货前把 \Fn{avcodec\_configuration} 与 \Fn{avcodec\_license} 打进
日志与``关于''面板，再逐行核对本表，是第 6 章已经点名的动作。

\begin{note}[title={本手册``首选''列的口径}]
本附录所有矩阵的``首选''列只收 permissive 条款（BSD 系、Apache-2.0、MIT、
ISC）的组件，以及不引入第三方开源义务的平台自带接口。LGPL、GPL 与专有实现
一律只进``替代''列作客观陈述，并说明其对分发的影响，不作为推荐项。这条口径
不是价值判断：\Lib{libav*} 与 \Lib{GStreamer} 在功能广度、文档与久经考验的
程度上常常是更工程化的答案，但``能不能静态链接进闭源二进制再发出去''是一个
二元问题，本表把它单独列出来供核对。宽松条款也不是零义务：BSD 系再分发要
保留版权声明与许可文本，Apache-2.0 还要附 NOTICE 并说明改动。
\end{note}

\small
\begin{longtable}{>{\raggedright\arraybackslash}p{0.15\textwidth}
                  >{\raggedright\arraybackslash}p{0.16\textwidth}
                  >{\raggedright\arraybackslash}p{0.20\textwidth}
                  >{\raggedright\arraybackslash}p{0.28\textwidth}}
\caption{组件的代码许可、专利条款与分发影响}\label{tab:apC-license} \\
\toprule
组件 & 代码许可 & 专利与其他条款 & 对可分发产物的影响 \\
\midrule
\endfirsthead
\toprule
组件 & 代码许可 & 专利与其他条款 & 对可分发产物的影响 \\
\midrule
\endhead
\midrule
\endfoot
\bottomrule
\endlastfoot
\Lib{libav\Bk util}、\Lib{libav\Bk codec}、\Lib{libav\Bk format}、\Lib{libsw\Bk scale}、\Lib{libsw\Bk resample} & 默认 LGPL-2.1+；树内散见 BSD、MIT、ISC 与 Apache-2.0 文本 & 库本身不给格式专利授权，用了哪个编码器就面对哪个池 & 动态链接并让用户能替换库是最省事的合规方式；静态链接要留出重链接的可能；先读 \Tp{-{\Bk}buildc\Bk onf} 再谈义务 \\
\Opt{enable-\Bk gpl}（configure 开关） & 把产物许可整体翻成 GPL-2+ & 开关本身不改变专利账本 & 只在同时编入 GPL-only 组件时才需要，也只在允许 GPL 分发的通道上才可接受 \\
\Opt{enable-\Bk nonfree}（configure 开关） & 产物既不是 LGPL 也不是 GPL & 放行的正是带费条款的实现 & 这类构建通常不可自由再分发：不是``贵''，是``不能发'' \\
\Lib{x264} & GPL-2.0+，另有商用授权通道 & H.264 的专利池授权与代码许可是两本账 & 链入即把产物拖进 GPL；把 \Tp{ffmpeg} 当子进程调用不改变你自己代码的许可，但那个二进制仍受 GPL 约束 \\
\Lib{x265} & GPL-2.0+，可选组件的条款以仓库文本为准 & HEVC 的池不止一个，费用与地域范围各自不同 & 同上；它是离线归档链路的画质基线，出货前把池的清单与产品线过一遍 \\
\Lib{libaom}（AV1） & BSD 条款加专利授权与终止条件 & 联盟声明 AV1 免版税，个别专利池是否主张权利要单独确认 & 代码层面可静态链接进闭源产物；终止条款意味着涉诉即收回授权，法务要独立过一遍 \\
\Lib{SVT-AV1} & 宽松条款（BSD 系）加专利授权，条款数以仓库 \Tp{LICENSE} 为准 & 同上 & 多线程扩展适合服务器批量；产物许可不被它拖累 \\
\Lib{libvpx}（VP8/VP9） & BSD-3 条款加专利授权 & 免版税声明与专利终止条款并存 & 可静态链接分发；硬件解码覆盖是它相对 H.26x 的主要短板 \\
\Lib{dav1d} & BSD-2 条款 & 解码不引入新的代码许可问题，内容侧专利另计 & 最适合放进静态产品二进制；注意它只是解码器 \\
\Lib{rav1e} & MIT 与 Apache-2.0 双许可（以仓库为准） & 同 AV1 的专利账本 & 双许可可择一满足；选 Apache-2.0 时要保留 NOTICE 并说明改动 \\
\Lib{Opus}（\Lib{libopus}） & BSD-3 条款 & 格式本身免版税 & 可静态链接进闭源产物；这是本表里最省心的一行 \\
\Lib{fdk-aac} & Fraunhofer IIS 的专有许可 & AAC 格式专利费另计，池随 profile 不同 & 需 \Opt{enable-\Bk nonfree} 才可编入；编进可再分发的二进制需单独授权 \\
\Lib{openh264} & BSD 风格条款（Cisco 名义），条款数以仓库为准 & Cisco 对其提供的预编译二进制另有再分发与赔偿承诺，属单方政策 & 代码层面可静态；不要把 Cisco 的二进制政策默认继承给你自己编译的那份 \\
\Lib{GStreamer} 核心与库 & LGPL-2.1+ & 框架不给格式专利授权 & 动态链接是常规满足方式；真正决定产物许可的是链入的插件集合 \\
\Lib{GStreamer} 插件域 & 按域分散，从宽松到 GPL 乃至非 free 都有 & 同上 & 用 \Fn{gst-inspect-\Bk 1.0} 逐个元素读许可字段，把``元素名到许可标签''固化进 CI \\
\Lib{WebRTC} 参考实现 & BSD-3 条款 & 自带的音频与视频路线偏免版税格式；H.264 要另接 & 可静态链接；信令与 STUN、TURN 服务不在许可范围内，是另一层成本 \\
内核侧：V4L2 与 ALSA & 内核 GPL-2.0-only（UAPI 头文件带内核说明） & 不适用 & 系统调用与 ioctl 边界不把用户态程序变成 GPL；这条区分很值钱 \\
\Lib{alsa-lib} & LGPL-2.1+ & 不适用 & 动态链接省事；静态链接要给替换可能 \\
\Lib{PipeWire}（含 \Lib{Wire{\Bk}Plumber}） & MIT & 不适用 & 与产物许可无冲突；它以独立服务进程运行，集成成本在协议侧而非许可侧 \\
\Lib{JACK}（jack2） & 以发行版携带的许可文本为准，服务端与客户端库历史上条款不一致 & 不适用 & 链接客户端库前把头文件与许可文件一起读，别按``音频开源库''想当然 \\
NVIDIA 编解码运行库（NVENC\Bk /NVDEC） & 闭源二进制；\Lib{FFmpeg} 侧的接口头文件是开源发布 & 无额外条款，但受其 SDK 许可约束 & 再分发规则按 SDK 文本与你的打包方式核对；第 8 章的后端表按平台列了可达性 \\
Intel 侧驱动与运行时 & Linux 侧的媒体驱动与运行时以 MIT 发布；Windows 侧运行库专有 & 不适用 & 两条路径的再分发结论不同，跨平台产品要把它们分开写进构建矩阵 \\
\Lib{Media Foundation}、\Lib{Video{\Bk}Toolbox}、\Lib{AV{\Bk}Foundation} & 操作系统自带接口，无第三方开源条款 & 平台替你处理了部分格式的授权状态，不保证全部 & 不引入开源义务，但内容侧的专利账本仍要自己面对；能力与位深随系统版本变化 \\
\end{longtable}
\normalsize

\begin{warning}[title={代码许可不等于专利授权}]
上表第二列与第三列是两本独立的账。BSD-2、BSD-3、MIT 与 Apache-2.0 只约束
``这份源码能不能用、怎么带着版权声明用''；H.264、H.265、VC-1 这些格式的池
主张的是``实现这个标准要付什么''。AV1 与 Opus 一侧的免版税声明来自标准组织
与贡献者协议，个别池在哪些地区主张权利会随时间变化，这不是本手册能替你回答
的问题。落到工程上的做法只有两条：把代码许可固化成 CI 里的一张映射表（产物、
许可标签、允许的发布通道），把专利问题作为独立的法务检查项。
\end{warning}

\section{从延迟预算到参数与协议的决策链}

选型不是一次性投票，而是一条有方向的链：延迟预算决定编码参数，编码参数
决定可用的传输协议，协议与并发规模再决定引擎形态。顺序反过来写就会得到
一份``哪一环都在猜''的架构。下表按这一顺序自上而下列出每一环要钉住的量。

\begin{longtable}{>{\raggedright\arraybackslash}p{0.14\textwidth}
                  >{\raggedright\arraybackslash}p{0.26\textwidth}
                  >{\raggedright\arraybackslash}p{0.30\textwidth}
                  >{\raggedright\arraybackslash}p{0.11\textwidth}}
\caption{决策链：每一环要钉住的量与钉法}\label{tab:apC-chain} \\
\toprule
环节 & 这一环要钉住的量 & 常见钉法与代价 & 章 \\
\midrule
\endfirsthead
\toprule
环节 & 这一环要钉住的量 & 常见钉法与代价 & 章 \\
\midrule
\endhead
\midrule
\endfoot
\bottomrule
\endlastfoot
端到端预算 & 目标分位、允许的上限、超预算时的降级动作 & 先把这三个数写进文档再碰热路径代码；没有它们，``实时''只是形容词，后面每一环都没有判据 & \ref{ch:ch01-boundary} \\
采集与设备缓冲 & 周期与缓冲水位、共享或独占、系统报回的延迟与安全余量 & 以设备回读值为准而不是以你请求的值为准；把周期、\Tp{avail\Bk\_min} 与单次写入量对齐到同一个毫秒数 & \ref{ch:ch03-audio-device}、\ref{ch:ch04-capture-display} \\
编码参数 & 关键帧间距、双向预测帧数、前瞻深度、速度档与线程方式 & 压延迟就关前瞻与双向预测、把并行压成切片级；代价是同样码率下画质下降，这笔账要在验收表上写清 & \ref{ch:ch07-rate-control} \\
封装与切点 & 时间基、随机访问点位置与分片边界的对齐关系 & 关键帧间距与分片时长取整数倍；帧级并行与前瞻会让解码序与显示序错位，切之前先看两侧各是什么序 & \ref{ch:ch10-mux-sync} \\
传输协议 & 允许的缓冲时长，与用这个缓冲换来的重传或纠错机会 & 协议排在延迟轴上：越往下越靠``等得起''；把同一笔预算同时给编码与网络是重复计费 & \ref{ch:ch11-streaming} \\
线程与丢弃 & 队列个数、每段队列的容量、满了丢哪一头 & 每一处队列都是一段有名字的延迟；漏队列的丢弃方向决定丢最旧还是最新，两者症状完全不同 & \ref{ch:ch12-realtime-thread} \\
\end{longtable}

\small
\begin{longtable}{>{\raggedright\arraybackslash}p{0.12\textwidth}
                  >{\raggedright\arraybackslash}p{0.17\textwidth}
                  >{\raggedright\arraybackslash}p{0.18\textwidth}
                  >{\raggedright\arraybackslash}p{0.17\textwidth}
                  >{\raggedright\arraybackslash}p{0.16\textwidth}}
\caption{四档延迟预算对应的典型配置组合}\label{tab:apC-latency} \\
\toprule
档位 & 采集与缓冲 & 编码侧 & 传输侧 & 典型装配 \\
\midrule
\endfirsthead
\toprule
档位 & 采集与缓冲 & 编码侧 & 传输侧 & 典型装配 \\
\midrule
\endhead
\midrule
\endfoot
\bottomrule
\endlastfoot
互动级：最低 & 短周期、独占或直连，把系统延迟与安全余量都算进总额 & \DomAud{} 用短包长且不叠多余缓冲；\DomVid{} 关前瞻与双向预测 & 实时传输协议族，抖动缓冲按下限，宁丢不叠 & 自研实时线程加平台设备接口 \\
会议级：百毫秒量级 & 共享模式配系统混音，接受宿主周期 & 恒定或带峰值约束的码控，切片级并行 & 带重传与前向纠错的实时会话，拥塞控制闭环 & \Lib{libwebrtc} 或其元素化等价物 \\
直播级：秒级 & 无硬约束，稳定优先 & 慢档与更大前瞻，码率阶梯按档位铺 & 用延迟买重传机会的协议，或按分片走自适应 & 编码器加切片器与清单，进程边界清晰 \\
归档级：不限 & 无 & 两遍编码或恒定质量，最慢的速度档 & 落文件即可，索引前置换快启 & 批处理编排，逐帧校验与统计 \\
\end{longtable}
\normalsize

\begin{bestpractice}[title={决策链的四条检查项}]
\begin{itemize}
\item 每一环的延迟都要有一个\textbf{名字}和一个上限数字，能在架构图上被单独指出来；
      指不出来的环节就是后来``多出来的那两百毫秒''。
\item 编码参数与传输协议在同一笔预算下互相竞争：前瞻吃掉的是留给重传窗口的额度，
      抖动缓冲吃掉的是本可以让编码更早结束的那部分。
\item 降级动作要提前写成代码路径（丢帧、降档、关掉某一路叠加），而不是超预算时临场加。
\item 用真实量级检验：周期时长、一帧字节数、每秒帧数写在注释里，链路预算表跟着
      度量走（第 \ref{ch:ch01-boundary} 章的验收标准一节）。
\end{itemize}
\end{bestpractice}

\section{按症状索引的排障手册}

排障表按症状索引，不按模块索引：出问题时手里只有症状。每行的``首查项''是
十分钟内做得完、并且能定性的动作；``常见根因''是这些动作排完仍然不对时的
优先猜测。顺序是先定性（哪一侧、哪一层）再定位（哪个调用、哪个参数）。

\small
\begin{longtable}{>{\raggedright\arraybackslash}p{0.17\textwidth}
                  >{\raggedright\arraybackslash}p{0.22\textwidth}
                  >{\raggedright\arraybackslash}p{0.28\textwidth}
                  >{\raggedright\arraybackslash}p{0.12\textwidth}}
\caption{按症状索引的排障表}\label{tab:apC-symptom} \\
\toprule
症状 & 首查项 & 常见根因 & 对应章 \\
\midrule
\endfirsthead
\toprule
症状 & 首查项 & 常见根因 & 对应章 \\
\midrule
\endhead
\midrule
\endfoot
\bottomrule
\endlastfoot
\DomAud{} 断续、周期性咔哒 & 回读设备实际生效的周期与缓冲参数，与你的读写周期比对；读一次 xrun 计数 & 欠载或过载：写入节奏跟不住；把短读短写当失败而不是循环；软件音量与硬件增益叠了两次；实时上下文里发生了分配或加锁 & \ref{ch:ch03-audio-device}、\ref{ch:ch12-realtime-thread} \\
开播黑屏，过一会儿才有画面 & 第一个随机访问点之前有没有带外参数集；打开阶段的探测预算与首包时间戳 & 带内与带外只写了其中一份；双向预测让显示序与解码序错位，首帧时间戳为负；切片没切在刷新点上 & \ref{ch:ch06-codec-api}、\ref{ch:ch10-mux-sync} \\
音画漂移随分钟增大 & 两条流各自的时间基与采样时钟来源；内部同步用的是哪一种时钟 & 两端是两台独立晶振，本来就会漂，没做重定时；把墙上时钟当内部时间源；一侧按帧计数另一侧按字节计数 & \ref{ch:ch02-clock-buffer}、\ref{ch:ch10-mux-sync} \\
\Tp{AVERROR\Bk(EAGAIN)} 风暴、延迟一路累积 & 把 send 与 receive 的交替按返回值逐次打点，看是哪一侧先堵 & 只在输入侧循环；正常途中没取空输出就再送帧；把``再多喂一帧''读成``结束了''；要中止本轮应调 flush 而不是继续送 & \ref{ch:ch06-codec-api} \\
花屏、局部马赛克并沿参考链扩散 & 先收紧检错强度，再看损坏包计数与关键帧请求是否真到达发送端 & 参考帧被污染后没有刷新点可退；请求到真正插帧之间往返太长；只重传了载荷没重传参数集 & \ref{ch:ch08-decode-hw}、\ref{ch:ch14-congestion-abr} \\
硬解没生效、CPU 跟着帧率涨 & 打印实际选中的解码器名与设备上下文；确认输出格式与帧池深度参数 & 按 ID 查实现查到了软件那一份；位深、profile 或像素格式不被驱动接受而静默回落；帧又额外下载回主内存，多付一次拷贝 & \ref{ch:ch08-decode-hw} \\
低码率下色带、渐变处断层 & 查量化区间与矩阵标记（限定还是全范围、哪一代标准），再查目标码率与峰值约束 & 峰值与缓冲约束只给了一个，另一个还是默认值；范围标记错位把暗部削掉；同一链路上矩阵被替换两次，越换越糊 & \ref{ch:ch05-pixel-format}、\ref{ch:ch07-rate-control} \\
WebRTC 上行码率压不下去 & 看带宽估计值、发送整形队列长度，以及反馈是否真回到发送端 & 反馈通道被中间设备挡住；估计器只升不降；关键帧请求成风暴，每轮都从最贵的一帧起；上限与缓冲配置互相矛盾 & \ref{ch:ch12-realtime-thread}、\ref{ch:ch14-congestion-abr} \\
HLS 分片时长不齐、清单反复重写 & 逐片看首帧是不是随机访问点，比对关键帧间距与目标分片时长 & 两者不成整数倍；场景切换检测强制插帧把间距打散；前瞻使帧的到达序与时间序不一致；重写清单的判据用了墙上钟 & \ref{ch:ch07-rate-control}、\ref{ch:ch11-streaming} \\
时间戳回绕、时长忽然变负或暴涨 & 看容器的计数字段位宽与跨流换算；确认不连续标志是否被处理 & 定宽计数器本来就会回绕；换基时漏了时长字段；用 32 位存微秒；把不连续标志当错误而不是裁决点 & \ref{ch:ch10-mux-sync} \\
本地测试一切正常、真实网络卡顿 & 测试命令里是否按源速率读取输入 & 全速灌包把拥塞控制与抖动缓冲测成假象；测试环境的往返延迟远小于现场 & \ref{ch:ch11-streaming}、\ref{ch:ch14-congestion-abr} \\
配置``没生效''但没有报错 & 打印打开上下文之后字典里剩下的键，逐个当错误抛出 & 参数名拼错或该版本已改名而被静默忽略；公共开关与私有参数串给了同名项，后者覆盖前者 & \ref{ch:ch06-codec-api}、\ref{ch:ch13-framework-integration} \\
\end{longtable}
\normalsize

\begin{lstlisting}[style=shellstyle,caption={排障起手式：三条只读命令先定性与定界},label={lst:apC-triage}]
# 一、时间基与首帧偏移：黑屏、时长变负都先落到这一行
ffprobe -hide_banner -v error \
  -show_entries stream=index,codec_type,time_base,start_time,duration \
  -of csv=p=0 in.mp4

# 二、实际选中的解码实现：这里没有硬件条目就说明已回落
ffmpeg -hide_banner -hwaccel auto -i in.mp4 -f null - 2>&1 | grep 'Stream #'

# 三、逐包看关键帧位置：切片不齐与花屏恢复慢都查这张表
ffprobe -hide_banner -v error -select_streams v:0 \
  -show_entries packet=pts,dts,flags -of csv=p=0 in.mp4 | grep ',K'
\end{lstlisting}

\begin{note}[title={读表的两个前提}]
本表的命令与字段名只用于定性，不能替代逐帧核对：\Tp{ffprobe} 的字段集合与
默认构建的可用能力都随版本变化，先看 \Tp{-{\Bk}buildc\Bk onf} 再引用默认值。
排障时改的是观测点，不是产品配置——把参数改回来之前，先记下改动前后的两组
度量（延迟分位、码率、丢帧计数），否则``修好了''无法复现。
\end{note}

\begin{bestpractice}[title={选型与排障共享的三条纪律}]
\begin{enumerate}
\item 每次只动一层：设备、编码、封装、传输、线程，五层里同一次只改一层，
      并留下改前的度量快照。跨层同时调优等于把根因搅进噪声。
\item 许可与能力分开查：构建开关决定``能不能发''，能力表决定``能不能跑''，
      两者用同一份日志取证（\Fn{avcodec\_configuration}、\Fn{gst-inspect-\Bk 1.0}）。
\item 把本附录的表当成检查清单而不是答案：每个``常见根因''背后是一章论证，
      现场能省时间，但不能替代对第 \ref{ch:ch02-clock-buffer} 章到第
      \ref{ch:ch14-congestion-abr} 章的理解。
\end{enumerate}
\end{bestpractice}

\end{document}
